% !TeX root = ../main.tex
\clearpage
\section*{附录3\quad 核心代码（Python详编）}
\phantomsection\label{appendix:code}
\addcontentsline{toc}{section}{附录3：核心代码（Python详编）}

本附录沿用原有12段代码的顺序，补入共享热资源、逐楼热状态、配电边界、硬约束权重及紧凑启停实现，共17段、1066行原代码和224行补充注释。代码均从随包R8O固定源码连续节选，按输入输出、实现思路、关键注释和结果核对展开；正文模型与算法结构保持原样。

\textbf{阅读顺序。}第1段说明入口；第2--4段连接建筑、储热与节点热平衡；第5--6段完成稀疏求解及数值复核；第7--8段将规划资产交给连续供暖季运行；第9--12段说明多目标、权重诊断和管网反馈。新增第13--15段补齐跨模块的物理衔接，第16--17段展示权重优化与整数规模压缩的具体实现。

显示行号从每段1开始并包含新增注释，索引给出原文件准确行号。\path{code_annotated_r8p_a4}内保存原样节选、注释版及逐行映射，\path{code_annotation_manifest_a4.json}记录来源和指纹。删除新增注释后，原代码行及可执行token一致；可以独立解析的片段另核对AST。代码排版沿用原有配色、字体与框线，上下文片段仍须在完整工程中调用。

源码说明各接口的实现，实验是否执行及其结果以对应运行记录为准。代表日规划、固定投资复算、两阶段端点和独立碳帽扫描分别保留执行身份；下文在相关代码处说明开关和衔接条件。

\begin{longtable}{@{}cL{0.34\linewidth}L{0.54\linewidth}@{}}
\caption{详编代码来源与正文方法对应}\label{tab:code-index-a4}\\
\toprule
序号 & 文件与原行号 & 对应内容\\\midrule
\endfirsthead
\toprule 序号 & 文件与原行号 & 对应内容\\\midrule
\endhead
\bottomrule\endfoot
1 & \path{unified_cli.py}，93--118 & 单次求解入口与导出\\
2 & \path{building_thermal.py}，31--46 & 逐楼1R1C精确离散系数\\
3 & \path{model.py}，1391--1471 & 储热容量、运行与状态边界\\
4 & \path{model.py}，1603--1697 & 节点温位逐时热平衡\\
5 & \path{linear.py}，177--205 & SciPy稀疏MILP求解接口\\
6 & \path{linear.py}，40--105 & 独立数值与守恒残差复核\\
7 & \path{experiments.py}，13--147 & 完整投资冻结与费用账本\\
8 & \path{heating_season.py}，183--246 & 固定投资的连续供暖季复算\\
9 & \path{model.py}，2829--2888 & 带主目标容差的两阶段端点\\
10 & \path{experiments.py}，351--411 & 独立碳帽采样与非支配筛选\\
11 & \path{representative_adoption.py}，48--94 & 代表日硬能量约束可行性诊断\\
12 & \path{network_feedback.py}，440--481 & 有限模式管网反馈与失败处理\\
13 & \path{physical_extensions.py}，25--79 & 热源出力的上游余热共用约束\\
14 & \path{building_thermal.py}，463--525 & 逐楼RC状态、舒适边界与节点交付量\\
15 & \path{model.py}，2030--2089 & 逐时配电接入容量及可靠容量方法入口\\
16 & \path{hard_weight_calibration.py}，13--97 & 硬能量守恒下的最小一维W1权重线性规划\\
17 & \path{model.py}，938--1018 & 热源启停三行紧凑凸包与容量关联爬坡\\
\end{longtable}

正文算法\ref{alg:season-workflow}对应规划与固定投资复算；算法\ref{alg:epsilon-workflow}对应独立碳帽采样与复核。第16段补充的是代表日权重优化本体，第11段则是等式与下界能否同时成立的诊断，两者作用分开。

\begin{paperalgorithm}{独立费用与碳排放采样}{alg:epsilon-workflow}
\noindent\textbf{输入：}同一配置与时间序列、碳帽档数 $K\geq2$。\par
\noindent\textbf{输出：}通过复核的非支配方案、边际减排费用与折中建议。
\begin{enumerate}[label=\arabic*.,leftmargin=1.8em,itemsep=2pt,topsep=3pt,parsep=0pt]
\item 清除原碳帽并令碳价为零，保留设备、运行与可靠性约束。
\item 分别以费用和排放为主目标求端点，据两端排放值等距设置 $K$ 档碳帽。
\item 对每档碳帽重新构模求解，独立检查数值残差与碳帽是否成立。
\item 对通过的样本去重并剔除受支配点，计算相邻方案的边际减排费用，用归一化理想点距离给出折中建议。
\end{enumerate}
\end{paperalgorithm}


\Needspace{12\baselineskip}
\subsection*{1．单次求解入口与导出}
主入口中没有政策路径时的单次求解分支；配置可以按需展开相关天气场景，随后统一求解和导出。

\noindent\textbf{输入：}当前profile的配置路径、逐时输入路径、可选场景说明及求解预算。\par
\noindent\textbf{输出：}投资、逐时运行、汇总指标与场景结果文件。该节选依赖外层profile循环。\par

\noindent\textbf{实现思路。}该入口负责把配置、逐时输入和求解预算连接起来。需求列从节点配置取得，避免在场景展开时漏掉某一节点；随后将场景风险设置写入配置副本，再统一校验。\path{HeatPlanningModel}接收校验后的同一组输入，导出也使用模型实际采用的时序，因而输入、求解与结果文件可以一一对应。

{\small 来源：\path{unified_cli.py}，原第93--118行；26行原代码、10行补充注释。}\par

\begingroup
\setstretch{1.0}
\begin{Verbatim}[commandchars=\\\{\},numbers=left,firstnumber=1,stepnumber=1,fontsize=\fontsize{9.5}{11.5}\selectfont,breaklines=true,breakanywhere=true,frame=single,framesep=3mm,numbersep=5pt,rulecolor=\color{black},xleftmargin=18pt]
\PYG{k}{else}\PYG{p}{:}
    \PYG{c+c1}{\PYGZsh{} \CJKCode{配置与时序共同定义算例身份；本段不自动启动碳帽扫描。}}
    \PYG{n}{config}\PYG{p}{,} \PYG{n}{frame} \PYG{o}{=} \PYG{n}{load\PYGZus{}inputs}\PYG{p}{(}\PYG{n}{config\PYGZus{}path}\PYG{p}{,} \PYG{n}{timeseries\PYGZus{}path}\PYG{p}{)}
    \PYG{c+c1}{\PYGZsh{} \CJKCode{仅在提供场景文件时展开不确定场景；无文件则沿用读入时序。}}
    \PYG{k}{if} \PYG{n}{scenario\PYGZus{}specs\PYGZus{}path}\PYG{p}{:}
        \PYG{n}{specs} \PYG{o}{=} \PYG{n}{yaml}\PYG{o}{.}\PYG{n}{safe\PYGZus{}load}\PYG{p}{(}\PYG{n}{scenario\PYGZus{}specs\PYGZus{}path}\PYG{o}{.}\PYG{n}{read\PYGZus{}text}\PYG{p}{(}\PYG{n}{encoding}\PYG{o}{=}\PYG{l+s+s2}{\PYGZdq{}}\PYG{l+s+s2}{utf\PYGZhy{}8}\PYG{l+s+s2}{\PYGZdq{}}\PYG{p}{)}\PYG{p}{)}
        \PYG{c+c1}{\PYGZsh{} \CJKCode{按节点配置收集需求列，保持场景扩展与需求接口一致。}}
        \PYG{n}{demand\PYGZus{}columns} \PYG{o}{=} \PYG{p}{[}
            \PYG{n}{column} \PYG{k}{for} \PYG{n}{node} \PYG{o+ow}{in} \PYG{n}{config}\PYG{p}{[}\PYG{l+s+s2}{\PYGZdq{}}\PYG{l+s+s2}{nodes}\PYG{l+s+s2}{\PYGZdq{}}\PYG{p}{]}
            \PYG{k}{for} \PYG{n}{column} \PYG{o+ow}{in} \PYG{n}{node}\PYG{o}{.}\PYG{n}{get}\PYG{p}{(}\PYG{l+s+s2}{\PYGZdq{}}\PYG{l+s+s2}{demand\PYGZus{}columns}\PYG{l+s+s2}{\PYGZdq{}}\PYG{p}{,} \PYG{p}{\PYGZob{}}\PYG{p}{\PYGZcb{}}\PYG{p}{)}\PYG{o}{.}\PYG{n}{values}\PYG{p}{(}\PYG{p}{)}
        \PYG{p}{]}
        \PYG{n}{frame} \PYG{o}{=} \PYG{n}{build\PYGZus{}correlated\PYGZus{}northern\PYGZus{}scenario\PYGZus{}bundle}\PYG{p}{(}
            \PYG{n}{frame}\PYG{p}{,} \PYG{n}{demand\PYGZus{}columns}\PYG{p}{,} \PYG{n}{specs}\PYG{p}{[}\PYG{l+s+s2}{\PYGZdq{}}\PYG{l+s+s2}{scenarios}\PYG{l+s+s2}{\PYGZdq{}}\PYG{p}{]}\PYG{p}{,}
            \PYG{c+c1}{\PYGZsh{} period\PYGZus{}weight\CJKCode{是代表时段权重，不等同于额外的物理小时步长。}}
            \PYG{n}{period\PYGZus{}weight}\PYG{o}{=}\PYG{n+nb}{float}\PYG{p}{(}\PYG{n}{specs}\PYG{o}{.}\PYG{n}{get}\PYG{p}{(}\PYG{l+s+s2}{\PYGZdq{}}\PYG{l+s+s2}{period\PYGZus{}weight}\PYG{l+s+s2}{\PYGZdq{}}\PYG{p}{,} \PYG{n}{config}\PYG{p}{[}\PYG{l+s+s2}{\PYGZdq{}}\PYG{l+s+s2}{planning}\PYG{l+s+s2}{\PYGZdq{}}\PYG{p}{]}\PYG{o}{.}\PYG{n}{get}\PYG{p}{(}\PYG{l+s+s2}{\PYGZdq{}}\PYG{l+s+s2}{representative\PYGZus{}day\PYGZus{}weight}\PYG{l+s+s2}{\PYGZdq{}}\PYG{p}{,} \PYG{l+m+mf}{1.0}\PYG{p}{)}\PYG{p}{)}\PYG{p}{)}\PYG{p}{,}
        \PYG{p}{)}
        \PYG{c+c1}{\PYGZsh{} \CJKCode{复制后修改风险设置，保留此前读入的配置对象。}}
        \PYG{n}{config} \PYG{o}{=} \PYG{n}{copy}\PYG{o}{.}\PYG{n}{deepcopy}\PYG{p}{(}\PYG{n}{config}\PYG{p}{)}
        \PYG{n}{config}\PYG{p}{[}\PYG{l+s+s2}{\PYGZdq{}}\PYG{l+s+s2}{planning}\PYG{l+s+s2}{\PYGZdq{}}\PYG{p}{]}\PYG{p}{[}\PYG{l+s+s2}{\PYGZdq{}}\PYG{l+s+s2}{risk}\PYG{l+s+s2}{\PYGZdq{}}\PYG{p}{]} \PYG{o}{=} \PYG{n+nb}{dict}\PYG{p}{(}\PYG{n}{specs}\PYG{o}{.}\PYG{n}{get}\PYG{p}{(}\PYG{l+s+s2}{\PYGZdq{}}\PYG{l+s+s2}{risk}\PYG{l+s+s2}{\PYGZdq{}}\PYG{p}{,} \PYG{p}{\PYGZob{}}\PYG{p}{\PYGZcb{}}\PYG{p}{)}\PYG{p}{)}
        \PYG{c+c1}{\PYGZsh{} \CJKCode{场景展开后再次校验列名、参数和时序约定。}}
        \PYG{n}{validate\PYGZus{}inputs}\PYG{p}{(}\PYG{n}{config}\PYG{p}{,} \PYG{n}{frame}\PYG{p}{)}
    \PYG{c+c1}{\PYGZsh{} \CJKCode{构造公用模型；源、网、荷、储与配电模块由配置装配。}}
    \PYG{n}{model} \PYG{o}{=} \PYG{n}{HeatPlanningModel}\PYG{p}{(}\PYG{n}{config}\PYG{p}{,} \PYG{n}{frame}\PYG{p}{)}
    \PYG{c+c1}{\PYGZsh{} mip\PYGZus{}gap\CJKCode{和}time\PYGZus{}limit\CJKCode{限制求解预算，不是物理可行性容差。}}
    \PYG{n}{artifacts} \PYG{o}{=} \PYG{n}{model}\PYG{o}{.}\PYG{n}{solve}\PYG{p}{(}
        \PYG{n}{mip\PYGZus{}gap}\PYG{o}{=}\PYG{n}{args}\PYG{o}{.}\PYG{n}{mip\PYGZus{}gap}\PYG{p}{,}
        \PYG{n}{time\PYGZus{}limit}\PYG{o}{=}\PYG{n}{args}\PYG{o}{.}\PYG{n}{time\PYGZus{}limit}\PYG{p}{,}
    \PYG{p}{)}
    \PYG{c+c1}{\PYGZsh{} \CJKCode{导出所得解及其实际输入；无可行解会由上层模型抛出异常。}}
    \PYG{n}{export\PYGZus{}results}\PYG{p}{(}\PYG{n}{profile\PYGZus{}output}\PYG{p}{,} \PYG{n}{artifacts}\PYG{p}{,} \PYG{n}{config}\PYG{p}{,} \PYG{n}{model}\PYG{o}{.}\PYG{n}{ts}\PYG{p}{)}
    \PYG{n}{summary} \PYG{o}{=} \PYG{p}{\PYGZob{}}\PYG{n}{key}\PYG{p}{:} \PYG{n}{value} \PYG{k}{for} \PYG{n}{key}\PYG{p}{,} \PYG{n}{value} \PYG{o+ow}{in} \PYG{n}{artifacts}\PYG{o}{.}\PYG{n}{summary}\PYG{o}{.}\PYG{n}{items}\PYG{p}{(}\PYG{p}{)} \PYG{k}{if} \PYG{o+ow}{not} \PYG{n+nb}{isinstance}\PYG{p}{(}\PYG{n}{value}\PYG{p}{,} \PYG{n+nb}{dict}\PYG{p}{)}\PYG{p}{\PYGZcb{}}
    \PYG{c+c1}{\PYGZsh{} \CJKCode{场景统计另存，避免以全系统平均掩盖某个场景的失供。}}
    \PYG{p}{(}\PYG{n}{profile\PYGZus{}output} \PYG{o}{/} \PYG{l+s+s2}{\PYGZdq{}}\PYG{l+s+s2}{scenario\PYGZus{}results.json}\PYG{l+s+s2}{\PYGZdq{}}\PYG{p}{)}\PYG{o}{.}\PYG{n}{write\PYGZus{}text}\PYG{p}{(}
        \PYG{n}{json}\PYG{o}{.}\PYG{n}{dumps}\PYG{p}{(}\PYG{n}{artifacts}\PYG{o}{.}\PYG{n}{summary}\PYG{o}{.}\PYG{n}{get}\PYG{p}{(}\PYG{l+s+s2}{\PYGZdq{}}\PYG{l+s+s2}{scenario\PYGZus{}results}\PYG{l+s+s2}{\PYGZdq{}}\PYG{p}{,} \PYG{p}{\PYGZob{}}\PYG{p}{\PYGZcb{}}\PYG{p}{)}\PYG{p}{,} \PYG{n}{ensure\PYGZus{}ascii}\PYG{o}{=}\PYG{k+kc}{False}\PYG{p}{,} \PYG{n}{indent}\PYG{o}{=}\PYG{l+m+mi}{2}\PYG{p}{)}\PYG{p}{,}
        \PYG{n}{encoding}\PYG{o}{=}\PYG{l+s+s2}{\PYGZdq{}}\PYG{l+s+s2}{utf\PYGZhy{}8}\PYG{l+s+s2}{\PYGZdq{}}\PYG{p}{,}
    \PYG{p}{)}
\end{Verbatim}

\endgroup

\noindent\textbf{结果核对。}阅读结果时先检查场景文件是否提供，再核对导出的时序、求解状态及逐场景结果。单次入口的结果不能仅凭文件存在就视为已执行碳帽扫描。

\Needspace{12\baselineskip}
\subsection*{2．逐楼1R1C精确离散系数}
独立逐楼热模型的完整系数函数。逐小时输入保持常值时，系数连接小时初末室温与交付冷热量。

\noindent\textbf{输入：}热损失系数H（kW/K）、有效热容C（kWh/K）和步长（h）。\par
\noindent\textbf{输出：}无量纲衰减系数a和温升系数b（K/kW）。函数不进行参数辨识。\par

\noindent\textbf{实现思路。}热平衡微分方程在一个时段内采用分段常值外温与净热输入。代码计算$a=\exp(-H\Delta t/C)$和$b=(1-a)/H$，使下一时段室温可直接写成当前室温、外温和净热输入的线性组合。系数由输入参数预先计算，优化时只含状态和供热变量，既保留热惯性，又不引入决策变量之间的乘积。

{\small 来源：\path{building_thermal.py}，原第31--46行；16行原代码、5行补充注释。}\par

\begingroup
\setstretch{1.0}
\begin{Verbatim}[commandchars=\\\{\},numbers=left,firstnumber=1,stepnumber=1,fontsize=\fontsize{9.5}{11.5}\selectfont,breaklines=true,breakanywhere=true,frame=single,framesep=3mm,numbersep=5pt,rulecolor=\color{black},xleftmargin=18pt]
\PYG{c+c1}{\PYGZsh{} \CJKCode{本函数属于独立}1R1C\CJKCode{分支，不替换历史聚合需求响应状态。}}
\PYG{k}{def}\PYG{+w}{ }\PYG{n+nf}{exact\PYGZus{}rc\PYGZus{}coefficients}\PYG{p}{(}
    \PYG{n}{heat\PYGZus{}loss\PYGZus{}kw\PYGZus{}per\PYGZus{}c}\PYG{p}{:} \PYG{n+nb}{float}\PYG{p}{,} \PYG{n}{thermal\PYGZus{}capacity\PYGZus{}kwh\PYGZus{}per\PYGZus{}c}\PYG{p}{:} \PYG{n+nb}{float}\PYG{p}{,} \PYG{n}{dt\PYGZus{}hours}\PYG{p}{:} \PYG{n+nb}{float}\PYG{p}{,}
\PYG{p}{)} \PYG{o}{\PYGZhy{}}\PYG{o}{\PYGZgt{}} \PYG{n+nb}{tuple}\PYG{p}{[}\PYG{n+nb}{float}\PYG{p}{,} \PYG{n+nb}{float}\PYG{p}{]}\PYG{p}{:}
\PYG{+w}{    }\PYG{l+s+sd}{\PYGZdq{}\PYGZdq{}\PYGZdq{}Return a, b for T\PYGZus{}end = a*T\PYGZus{}start + (1\PYGZhy{}a)*Tout + b*(heat\PYGZhy{}cool+gains).}

\PYG{l+s+sd}{    expm1 avoids cancellation at small H*dt/C. H=0 is deliberately rejected:}
\PYG{l+s+sd}{    this module requires a positive envelope/ventilation heat loss coefficient.}
\PYG{l+s+sd}{    \PYGZdq{}\PYGZdq{}\PYGZdq{}}
    \PYG{c+c1}{\PYGZsh{} H\CJKCode{、}C\CJKCode{与步长必须为正且有限，单位应先统一。}}
    \PYG{n}{h} \PYG{o}{=} \PYG{n}{\PYGZus{}number}\PYG{p}{(}\PYG{n}{heat\PYGZus{}loss\PYGZus{}kw\PYGZus{}per\PYGZus{}c}\PYG{p}{,} \PYG{l+s+s2}{\PYGZdq{}}\PYG{l+s+s2}{heat\PYGZus{}loss\PYGZus{}kw\PYGZus{}per\PYGZus{}c}\PYG{l+s+s2}{\PYGZdq{}}\PYG{p}{,} \PYG{n}{positive}\PYG{o}{=}\PYG{k+kc}{True}\PYG{p}{)}
    \PYG{n}{c} \PYG{o}{=} \PYG{n}{\PYGZus{}number}\PYG{p}{(}\PYG{n}{thermal\PYGZus{}capacity\PYGZus{}kwh\PYGZus{}per\PYGZus{}c}\PYG{p}{,} \PYG{l+s+s2}{\PYGZdq{}}\PYG{l+s+s2}{thermal\PYGZus{}capacity\PYGZus{}kwh\PYGZus{}per\PYGZus{}c}\PYG{l+s+s2}{\PYGZdq{}}\PYG{p}{,} \PYG{n}{positive}\PYG{o}{=}\PYG{k+kc}{True}\PYG{p}{)}
    \PYG{n}{dt} \PYG{o}{=} \PYG{n}{\PYGZus{}number}\PYG{p}{(}\PYG{n}{dt\PYGZus{}hours}\PYG{p}{,} \PYG{l+s+s2}{\PYGZdq{}}\PYG{l+s+s2}{dt\PYGZus{}hours}\PYG{l+s+s2}{\PYGZdq{}}\PYG{p}{,} \PYG{n}{positive}\PYG{o}{=}\PYG{k+kc}{True}\PYG{p}{)}
    \PYG{c+c1}{\PYGZsh{} H\CJKCode{乘步长除以}C\CJKCode{为无量纲指数，避免将秒与小时混算。}}
    \PYG{n}{decay} \PYG{o}{=} \PYG{n}{h} \PYG{o}{*} \PYG{p}{(}\PYG{n}{dt} \PYG{o}{/} \PYG{n}{c}\PYG{p}{)}
    \PYG{c+c1}{\PYGZsh{} expm1\CJKCode{减小小指数下的数值相消；}b\CJKCode{乘净输入功率得到温升。}}
    \PYG{n}{a}\PYG{p}{,} \PYG{n}{b} \PYG{o}{=} \PYG{n}{exp}\PYG{p}{(}\PYG{o}{\PYGZhy{}}\PYG{n}{decay}\PYG{p}{)}\PYG{p}{,} \PYG{o}{\PYGZhy{}}\PYG{n}{expm1}\PYG{p}{(}\PYG{o}{\PYGZhy{}}\PYG{n}{decay}\PYG{p}{)} \PYG{o}{/} \PYG{n}{h}
    \PYG{c+c1}{\PYGZsh{} \CJKCode{系数低于矩阵数值分辨能力时明确报错，不返回无效热状态。}}
    \PYG{k}{if} \PYG{o+ow}{not} \PYG{n}{isfinite}\PYG{p}{(}\PYG{n}{b}\PYG{p}{)} \PYG{o+ow}{or} \PYG{n}{b} \PYG{o}{\PYGZlt{}}\PYG{o}{=} \PYG{l+m+mf}{1e\PYGZhy{}12}\PYG{p}{:}
        \PYG{k}{raise} \PYG{n+ne}{ValueError}\PYG{p}{(}\PYG{l+s+s2}{\PYGZdq{}}\PYG{l+s+s2}{RC coefficient b is outside LinearMILP numerical resolution; rescale units}\PYG{l+s+s2}{\PYGZdq{}}\PYG{p}{)}
    \PYG{k}{return} \PYG{n}{a}\PYG{p}{,} \PYG{n}{b}
\end{Verbatim}

\endgroup

\noindent\textbf{结果核对。}核对单位时，$H$为kW/K，$C$为kWh/K，$\Delta t$为h，故$b$为K/kW。逐楼状态如何接入模型及舒适边界，见后面的RC装配代码。

\Needspace{12\baselineskip}
\subsection*{3．储热容量、运行与状态边界}
完整展示储热投资与逐时运行的这一连续代码块。既有与新增能量容量共同限制SOC和充放功率。

\noindent\textbf{输入：}储热位置、容量（kWh）、功率能量比（1/h）、效率、自损失及边界模式。\par
\noindent\textbf{输出：}容量、充热、放热、SOC变量和相应线性约束；严格互斥模式另含二元变量。\par

\noindent\textbf{实现思路。}代码先建立新增储热容量，再用同一容量同时约束SOC与充放功率，避免仅满足能量容量却忽略设备功率。严格模式下，充、放热分别受一个二元状态控制；线性化常数取最大允许容量对应的功率。状态递推采用物理步长和充放效率，代表时段权重只参与相应统计，不在此重复放大能量转移。

{\small 来源：\path{model.py}，原第1391--1471行；81行原代码、19行补充注释。}\par

\begingroup
\setstretch{1.0}
\begin{Verbatim}[commandchars=\\\{\},numbers=left,firstnumber=1,stepnumber=1,fontsize=\fontsize{9.5}{11.5}\selectfont,breaklines=true,breakanywhere=true,frame=single,framesep=3mm,numbersep=5pt,rulecolor=\color{black},xleftmargin=18pt]
\PYG{c+c1}{\PYGZsh{} \CJKCode{以下为}build\CJKCode{方法内部片段，依赖既有模型、时序与变量登记表。}}
\PYG{c+c1}{\PYGZsh{} Storage investments and operation. Short\PYGZhy{}term storage can be cyclic;}
\PYG{c+c1}{\PYGZsh{} seasonal storage can use fixed initial/final boundary conditions.}
\PYG{k}{for} \PYG{n}{st\PYGZus{}name}\PYG{p}{,} \PYG{n}{storage} \PYG{o+ow}{in} \PYG{n+nb+bp}{self}\PYG{o}{.}\PYG{n}{storages}\PYG{o}{.}\PYG{n}{items}\PYG{p}{(}\PYG{p}{)}\PYG{p}{:}
    \PYG{c+c1}{\PYGZsh{} existing\CJKCode{和}max\PYGZus{}new\CJKCode{均为能量容量}kWh\CJKCode{，不是充放功率}kW\CJKCode{。}}
    \PYG{n}{existing} \PYG{o}{=} \PYG{n+nb}{float}\PYG{p}{(}\PYG{n}{storage}\PYG{o}{.}\PYG{n}{get}\PYG{p}{(}\PYG{l+s+s2}{\PYGZdq{}}\PYG{l+s+s2}{existing\PYGZus{}capacity\PYGZus{}kwh}\PYG{l+s+s2}{\PYGZdq{}}\PYG{p}{,} \PYG{l+m+mf}{0.0}\PYG{p}{)}\PYG{p}{)}
    \PYG{n}{max\PYGZus{}new} \PYG{o}{=} \PYG{n+nb}{float}\PYG{p}{(}\PYG{n}{storage}\PYG{o}{.}\PYG{n}{get}\PYG{p}{(}\PYG{l+s+s2}{\PYGZdq{}}\PYG{l+s+s2}{max\PYGZus{}new\PYGZus{}capacity\PYGZus{}kwh}\PYG{l+s+s2}{\PYGZdq{}}\PYG{p}{,} \PYG{l+m+mf}{0.0}\PYG{p}{)}\PYG{p}{)}
    \PYG{n}{cap\PYGZus{}idx} \PYG{o}{=} \PYG{k+kc}{None}
    \PYG{c+c1}{\PYGZsh{} \CJKCode{只有允许新增容量时才建立投资变量；既有容量作为常量。}}
    \PYG{k}{if} \PYG{n}{max\PYGZus{}new} \PYG{o}{\PYGZgt{}} \PYG{l+m+mi}{0}\PYG{p}{:}
        \PYG{n}{cap\PYGZus{}idx} \PYG{o}{=} \PYG{n}{p}\PYG{o}{.}\PYG{n}{add\PYGZus{}var}\PYG{p}{(}\PYG{l+s+sa}{f}\PYG{l+s+s2}{\PYGZdq{}}\PYG{l+s+s2}{storage\PYGZus{}capacity[}\PYG{l+s+si}{\PYGZob{}}\PYG{n}{st\PYGZus{}name}\PYG{l+s+si}{\PYGZcb{}}\PYG{l+s+s2}{]}\PYG{l+s+s2}{\PYGZdq{}}\PYG{p}{,} \PYG{n}{ub}\PYG{o}{=}\PYG{n}{max\PYGZus{}new}\PYG{p}{)}
        \PYG{n+nb+bp}{self}\PYG{o}{.}\PYG{n}{variables}\PYG{p}{[}\PYG{l+s+s2}{\PYGZdq{}}\PYG{l+s+s2}{storage\PYGZus{}capacity}\PYG{l+s+s2}{\PYGZdq{}}\PYG{p}{]}\PYG{p}{[}\PYG{n}{st\PYGZus{}name}\PYG{p}{]} \PYG{o}{=} \PYG{n}{cap\PYGZus{}idx}
        \PYG{n+nb+bp}{self}\PYG{o}{.}\PYG{n}{\PYGZus{}record\PYGZus{}investment\PYGZus{}cost}\PYG{p}{(}
            \PYG{n}{cap\PYGZus{}idx}\PYG{p}{,} \PYG{l+s+s2}{\PYGZdq{}}\PYG{l+s+s2}{storage\PYGZus{}capacity}\PYG{l+s+s2}{\PYGZdq{}}\PYG{p}{,} \PYG{n}{st\PYGZus{}name}\PYG{p}{,} \PYG{n+nb}{float}\PYG{p}{(}\PYG{n}{storage}\PYG{o}{.}\PYG{n}{get}\PYG{p}{(}\PYG{l+s+s2}{\PYGZdq{}}\PYG{l+s+s2}{capex\PYGZus{}yuan\PYGZus{}per\PYGZus{}kwh}\PYG{l+s+s2}{\PYGZdq{}}\PYG{p}{,} \PYG{l+m+mf}{0.0}\PYG{p}{)}\PYG{p}{)}
        \PYG{p}{)}

    \PYG{c+c1}{\PYGZsh{} \CJKCode{功率能量比乘装机能量容量得到逐时充放功率上限。}}
    \PYG{n}{power\PYGZus{}ratio} \PYG{o}{=} \PYG{n+nb}{float}\PYG{p}{(}\PYG{n}{storage}\PYG{o}{.}\PYG{n}{get}\PYG{p}{(}\PYG{l+s+s2}{\PYGZdq{}}\PYG{l+s+s2}{power\PYGZus{}to\PYGZus{}energy\PYGZus{}ratio\PYGZus{}per\PYGZus{}hour}\PYG{l+s+s2}{\PYGZdq{}}\PYG{p}{,} \PYG{l+m+mf}{0.25}\PYG{p}{)}\PYG{p}{)}
    \PYG{n}{max\PYGZus{}power} \PYG{o}{=} \PYG{n}{power\PYGZus{}ratio} \PYG{o}{*} \PYG{p}{(}\PYG{n}{existing} \PYG{o}{+} \PYG{n}{max\PYGZus{}new}\PYG{p}{)}
    \PYG{c+c1}{\PYGZsh{} \CJKCode{严格模式按配置启用；关闭时不能宣称二元互斥已实施。}}
    \PYG{n}{strict\PYGZus{}storage\PYGZus{}mode} \PYG{o}{=} \PYG{n+nb}{bool}\PYG{p}{(}\PYG{n}{storage}\PYG{o}{.}\PYG{n}{get}\PYG{p}{(}\PYG{l+s+s2}{\PYGZdq{}}\PYG{l+s+s2}{strict\PYGZus{}charge\PYGZus{}discharge\PYGZus{}mode}\PYG{l+s+s2}{\PYGZdq{}}\PYG{p}{,} \PYG{k+kc}{True}\PYG{p}{)}\PYG{p}{)}
    \PYG{n}{throughput\PYGZus{}cost} \PYG{o}{=} \PYG{n+nb}{float}\PYG{p}{(}
        \PYG{n}{storage}\PYG{o}{.}\PYG{n}{get}\PYG{p}{(}\PYG{l+s+s2}{\PYGZdq{}}\PYG{l+s+s2}{throughput\PYGZus{}cost\PYGZus{}yuan\PYGZus{}per\PYGZus{}kwh}\PYG{l+s+s2}{\PYGZdq{}}\PYG{p}{,} \PYG{l+m+mf}{0.0}\PYG{p}{)}
    \PYG{p}{)}
    \PYG{c+c1}{\PYGZsh{} \CJKCode{每一真实时段建立充热、放热和}SOC\CJKCode{；状态连续性在后文连接。}}
    \PYG{k}{for} \PYG{n}{t} \PYG{o+ow}{in} \PYG{n+nb+bp}{self}\PYG{o}{.}\PYG{n}{times}\PYG{p}{:}
        \PYG{n}{ch} \PYG{o}{=} \PYG{n}{p}\PYG{o}{.}\PYG{n}{add\PYGZus{}var}\PYG{p}{(}\PYG{l+s+sa}{f}\PYG{l+s+s2}{\PYGZdq{}}\PYG{l+s+s2}{storage\PYGZus{}charge[}\PYG{l+s+si}{\PYGZob{}}\PYG{n}{st\PYGZus{}name}\PYG{l+s+si}{\PYGZcb{}}\PYG{l+s+s2}{,}\PYG{l+s+si}{\PYGZob{}}\PYG{n}{t}\PYG{l+s+si}{\PYGZcb{}}\PYG{l+s+s2}{]}\PYG{l+s+s2}{\PYGZdq{}}\PYG{p}{)}
        \PYG{n}{dis} \PYG{o}{=} \PYG{n}{p}\PYG{o}{.}\PYG{n}{add\PYGZus{}var}\PYG{p}{(}\PYG{l+s+sa}{f}\PYG{l+s+s2}{\PYGZdq{}}\PYG{l+s+s2}{storage\PYGZus{}discharge[}\PYG{l+s+si}{\PYGZob{}}\PYG{n}{st\PYGZus{}name}\PYG{l+s+si}{\PYGZcb{}}\PYG{l+s+s2}{,}\PYG{l+s+si}{\PYGZob{}}\PYG{n}{t}\PYG{l+s+si}{\PYGZcb{}}\PYG{l+s+s2}{]}\PYG{l+s+s2}{\PYGZdq{}}\PYG{p}{)}
        \PYG{n}{soc} \PYG{o}{=} \PYG{n}{p}\PYG{o}{.}\PYG{n}{add\PYGZus{}var}\PYG{p}{(}\PYG{l+s+sa}{f}\PYG{l+s+s2}{\PYGZdq{}}\PYG{l+s+s2}{storage\PYGZus{}soc[}\PYG{l+s+si}{\PYGZob{}}\PYG{n}{st\PYGZus{}name}\PYG{l+s+si}{\PYGZcb{}}\PYG{l+s+s2}{,}\PYG{l+s+si}{\PYGZob{}}\PYG{n}{t}\PYG{l+s+si}{\PYGZcb{}}\PYG{l+s+s2}{]}\PYG{l+s+s2}{\PYGZdq{}}\PYG{p}{)}
        \PYG{n+nb+bp}{self}\PYG{o}{.}\PYG{n}{variables}\PYG{p}{[}\PYG{l+s+s2}{\PYGZdq{}}\PYG{l+s+s2}{storage\PYGZus{}charge}\PYG{l+s+s2}{\PYGZdq{}}\PYG{p}{]}\PYG{p}{[}\PYG{p}{(}\PYG{n}{st\PYGZus{}name}\PYG{p}{,} \PYG{n}{t}\PYG{p}{)}\PYG{p}{]} \PYG{o}{=} \PYG{n}{ch}
        \PYG{n+nb+bp}{self}\PYG{o}{.}\PYG{n}{variables}\PYG{p}{[}\PYG{l+s+s2}{\PYGZdq{}}\PYG{l+s+s2}{storage\PYGZus{}discharge}\PYG{l+s+s2}{\PYGZdq{}}\PYG{p}{]}\PYG{p}{[}\PYG{p}{(}\PYG{n}{st\PYGZus{}name}\PYG{p}{,} \PYG{n}{t}\PYG{p}{)}\PYG{p}{]} \PYG{o}{=} \PYG{n}{dis}
        \PYG{n+nb+bp}{self}\PYG{o}{.}\PYG{n}{variables}\PYG{p}{[}\PYG{l+s+s2}{\PYGZdq{}}\PYG{l+s+s2}{storage\PYGZus{}soc}\PYG{l+s+s2}{\PYGZdq{}}\PYG{p}{]}\PYG{p}{[}\PYG{p}{(}\PYG{n}{st\PYGZus{}name}\PYG{p}{,} \PYG{n}{t}\PYG{p}{)}\PYG{p}{]} \PYG{o}{=} \PYG{n}{soc}
        \PYG{c+c1}{\PYGZsh{} \CJKCode{可选吞吐费用作用于充放热量，不改变能量守恒方程。}}
        \PYG{k}{if} \PYG{n}{throughput\PYGZus{}cost} \PYG{o}{\PYGZgt{}} \PYG{l+m+mf}{0.0}\PYG{p}{:}
            \PYG{n+nb+bp}{self}\PYG{o}{.}\PYG{n}{\PYGZus{}record\PYGZus{}time\PYGZus{}component}\PYG{p}{(}
                \PYG{l+s+s2}{\PYGZdq{}}\PYG{l+s+s2}{operating\PYGZus{}cost}\PYG{l+s+s2}{\PYGZdq{}}\PYG{p}{,} \PYG{n}{ch}\PYG{p}{,} \PYG{n}{throughput\PYGZus{}cost}\PYG{p}{,} \PYG{n}{t}
            \PYG{p}{)}
            \PYG{n+nb+bp}{self}\PYG{o}{.}\PYG{n}{\PYGZus{}record\PYGZus{}time\PYGZus{}component}\PYG{p}{(}
                \PYG{l+s+s2}{\PYGZdq{}}\PYG{l+s+s2}{operating\PYGZus{}cost}\PYG{l+s+s2}{\PYGZdq{}}\PYG{p}{,} \PYG{n}{dis}\PYG{p}{,} \PYG{n}{throughput\PYGZus{}cost}\PYG{p}{,} \PYG{n}{t}
            \PYG{p}{)}

        \PYG{c+c1}{\PYGZsh{} \CJKCode{新增容量以负系数进入上界，形成}SOC\CJKCode{与装机容量的线性关系。}}
        \PYG{c+c1}{\PYGZsh{} \CJKCode{该行左端为}SOC\CJKCode{减新增容量，右端保留既有容量。}}
        \PYG{n}{soc\PYGZus{}row} \PYG{o}{=} \PYG{p}{\PYGZob{}}\PYG{n}{soc}\PYG{p}{:} \PYG{l+m+mf}{1.0}\PYG{p}{\PYGZcb{}}
        \PYG{n}{ch\PYGZus{}row} \PYG{o}{=} \PYG{p}{\PYGZob{}}\PYG{n}{ch}\PYG{p}{:} \PYG{l+m+mf}{1.0}\PYG{p}{\PYGZcb{}}
        \PYG{n}{dis\PYGZus{}row} \PYG{o}{=} \PYG{p}{\PYGZob{}}\PYG{n}{dis}\PYG{p}{:} \PYG{l+m+mf}{1.0}\PYG{p}{\PYGZcb{}}
        \PYG{k}{if} \PYG{n}{cap\PYGZus{}idx} \PYG{o+ow}{is} \PYG{o+ow}{not} \PYG{k+kc}{None}\PYG{p}{:}
            \PYG{n}{soc\PYGZus{}row}\PYG{p}{[}\PYG{n}{cap\PYGZus{}idx}\PYG{p}{]} \PYG{o}{=} \PYG{o}{\PYGZhy{}}\PYG{l+m+mf}{1.0}
            \PYG{n}{ch\PYGZus{}row}\PYG{p}{[}\PYG{n}{cap\PYGZus{}idx}\PYG{p}{]} \PYG{o}{=} \PYG{o}{\PYGZhy{}}\PYG{n}{power\PYGZus{}ratio}
            \PYG{n}{dis\PYGZus{}row}\PYG{p}{[}\PYG{n}{cap\PYGZus{}idx}\PYG{p}{]} \PYG{o}{=} \PYG{o}{\PYGZhy{}}\PYG{n}{power\PYGZus{}ratio}
        \PYG{n}{p}\PYG{o}{.}\PYG{n}{add\PYGZus{}constraint}\PYG{p}{(}\PYG{n}{soc\PYGZus{}row}\PYG{p}{,} \PYG{n}{ub}\PYG{o}{=}\PYG{n}{existing}\PYG{p}{)}
        \PYG{n}{p}\PYG{o}{.}\PYG{n}{add\PYGZus{}constraint}\PYG{p}{(}\PYG{n}{ch\PYGZus{}row}\PYG{p}{,} \PYG{n}{ub}\PYG{o}{=}\PYG{n}{power\PYGZus{}ratio} \PYG{o}{*} \PYG{n}{existing}\PYG{p}{)}
        \PYG{n}{p}\PYG{o}{.}\PYG{n}{add\PYGZus{}constraint}\PYG{p}{(}\PYG{n}{dis\PYGZus{}row}\PYG{p}{,} \PYG{n}{ub}\PYG{o}{=}\PYG{n}{power\PYGZus{}ratio} \PYG{o}{*} \PYG{n}{existing}\PYG{p}{)}
        \PYG{c+c1}{\PYGZsh{} M\CJKCode{取候选最大功率常数；不使用投资变量乘二元变量。}}
        \PYG{k}{if} \PYG{n}{strict\PYGZus{}storage\PYGZus{}mode}\PYG{p}{:}
            \PYG{n}{mode} \PYG{o}{=} \PYG{n}{p}\PYG{o}{.}\PYG{n}{add\PYGZus{}var}\PYG{p}{(}\PYG{l+s+sa}{f}\PYG{l+s+s2}{\PYGZdq{}}\PYG{l+s+s2}{storage\PYGZus{}mode[}\PYG{l+s+si}{\PYGZob{}}\PYG{n}{st\PYGZus{}name}\PYG{l+s+si}{\PYGZcb{}}\PYG{l+s+s2}{,}\PYG{l+s+si}{\PYGZob{}}\PYG{n}{t}\PYG{l+s+si}{\PYGZcb{}}\PYG{l+s+s2}{]}\PYG{l+s+s2}{\PYGZdq{}}\PYG{p}{,} \PYG{n}{ub}\PYG{o}{=}\PYG{l+m+mf}{1.0}\PYG{p}{,} \PYG{n}{integer}\PYG{o}{=}\PYG{k+kc}{True}\PYG{p}{)}
            \PYG{n+nb+bp}{self}\PYG{o}{.}\PYG{n}{variables}\PYG{p}{[}\PYG{l+s+s2}{\PYGZdq{}}\PYG{l+s+s2}{storage\PYGZus{}mode}\PYG{l+s+s2}{\PYGZdq{}}\PYG{p}{]}\PYG{p}{[}\PYG{p}{(}\PYG{n}{st\PYGZus{}name}\PYG{p}{,} \PYG{n}{t}\PYG{p}{)}\PYG{p}{]} \PYG{o}{=} \PYG{n}{mode}
            \PYG{n}{p}\PYG{o}{.}\PYG{n}{add\PYGZus{}constraint}\PYG{p}{(}\PYG{p}{\PYGZob{}}\PYG{n}{ch}\PYG{p}{:} \PYG{l+m+mf}{1.0}\PYG{p}{,} \PYG{n}{mode}\PYG{p}{:} \PYG{o}{\PYGZhy{}}\PYG{n}{max\PYGZus{}power}\PYG{p}{\PYGZcb{}}\PYG{p}{,} \PYG{n}{ub}\PYG{o}{=}\PYG{l+m+mf}{0.0}\PYG{p}{)}
            \PYG{n}{p}\PYG{o}{.}\PYG{n}{add\PYGZus{}constraint}\PYG{p}{(}\PYG{p}{\PYGZob{}}\PYG{n}{dis}\PYG{p}{:} \PYG{l+m+mf}{1.0}\PYG{p}{,} \PYG{n}{mode}\PYG{p}{:} \PYG{n}{max\PYGZus{}power}\PYG{p}{\PYGZcb{}}\PYG{p}{,} \PYG{n}{ub}\PYG{o}{=}\PYG{n}{max\PYGZus{}power}\PYG{p}{)}

    \PYG{c+c1}{\PYGZsh{} cyclic\CJKCode{与}fixed\CJKCode{是两种边界身份，不能一律要求初末}SOC\CJKCode{相等。}}
    \PYG{n}{boundary\PYGZus{}mode} \PYG{o}{=} \PYG{n+nb}{str}\PYG{p}{(}\PYG{n}{storage}\PYG{o}{.}\PYG{n}{get}\PYG{p}{(}\PYG{l+s+s2}{\PYGZdq{}}\PYG{l+s+s2}{boundary\PYGZus{}mode}\PYG{l+s+s2}{\PYGZdq{}}\PYG{p}{,} \PYG{l+s+s2}{\PYGZdq{}}\PYG{l+s+s2}{cyclic}\PYG{l+s+s2}{\PYGZdq{}}\PYG{p}{)}\PYG{p}{)}
    \PYG{c+c1}{\PYGZsh{} \CJKCode{自损失按小时给定，再按实际步长取幂换算。}}
    \PYG{n}{loss\PYGZus{}factor} \PYG{o}{=} \PYG{p}{(}\PYG{l+m+mf}{1.0} \PYG{o}{\PYGZhy{}} \PYG{n+nb}{float}\PYG{p}{(}\PYG{n}{storage}\PYG{o}{.}\PYG{n}{get}\PYG{p}{(}\PYG{l+s+s2}{\PYGZdq{}}\PYG{l+s+s2}{self\PYGZus{}loss\PYGZus{}per\PYGZus{}hour}\PYG{l+s+s2}{\PYGZdq{}}\PYG{p}{,} \PYG{l+m+mf}{0.0}\PYG{p}{)}\PYG{p}{)}\PYG{p}{)} \PYG{o}{*}\PYG{o}{*} \PYG{n+nb+bp}{self}\PYG{o}{.}\PYG{n}{dt}
    \PYG{n}{eta\PYGZus{}ch} \PYG{o}{=} \PYG{n+nb}{float}\PYG{p}{(}\PYG{n}{storage}\PYG{o}{.}\PYG{n}{get}\PYG{p}{(}\PYG{l+s+s2}{\PYGZdq{}}\PYG{l+s+s2}{charge\PYGZus{}efficiency}\PYG{l+s+s2}{\PYGZdq{}}\PYG{p}{,} \PYG{l+m+mf}{0.95}\PYG{p}{)}\PYG{p}{)}
    \PYG{n}{eta\PYGZus{}dis} \PYG{o}{=} \PYG{n+nb}{float}\PYG{p}{(}\PYG{n}{storage}\PYG{o}{.}\PYG{n}{get}\PYG{p}{(}\PYG{l+s+s2}{\PYGZdq{}}\PYG{l+s+s2}{discharge\PYGZus{}efficiency}\PYG{l+s+s2}{\PYGZdq{}}\PYG{p}{,} \PYG{l+m+mf}{0.95}\PYG{p}{)}\PYG{p}{)}
    \PYG{k}{if} \PYG{n}{boundary\PYGZus{}mode} \PYG{o+ow}{not} \PYG{o+ow}{in} \PYG{p}{\PYGZob{}}\PYG{l+s+s2}{\PYGZdq{}}\PYG{l+s+s2}{cyclic}\PYG{l+s+s2}{\PYGZdq{}}\PYG{p}{,} \PYG{l+s+s2}{\PYGZdq{}}\PYG{l+s+s2}{fixed}\PYG{l+s+s2}{\PYGZdq{}}\PYG{p}{\PYGZcb{}}\PYG{p}{:}
        \PYG{k}{raise} \PYG{n+ne}{ValueError}\PYG{p}{(}\PYG{l+s+sa}{f}\PYG{l+s+s2}{\PYGZdq{}}\PYG{l+s+s2}{\CJKCode{储热} }\PYG{l+s+si}{\PYGZob{}}\PYG{n}{st\PYGZus{}name}\PYG{l+s+si}{\PYGZcb{}}\PYG{l+s+s2}{ \CJKCode{的} boundary\PYGZus{}mode \CJKCode{仅支持} cyclic \CJKCode{或} fixed}\PYG{l+s+s2}{\PYGZdq{}}\PYG{p}{)}
    \PYG{c+c1}{\PYGZsh{} \CJKCode{充放功率乘步长后才是能量；放热需除以放热效率。}}
    \PYG{k}{for} \PYG{n}{t} \PYG{o+ow}{in} \PYG{n+nb+bp}{self}\PYG{o}{.}\PYG{n}{times}\PYG{p}{:}
        \PYG{n}{row} \PYG{o}{=} \PYG{p}{\PYGZob{}}
            \PYG{n+nb+bp}{self}\PYG{o}{.}\PYG{n}{variables}\PYG{p}{[}\PYG{l+s+s2}{\PYGZdq{}}\PYG{l+s+s2}{storage\PYGZus{}soc}\PYG{l+s+s2}{\PYGZdq{}}\PYG{p}{]}\PYG{p}{[}\PYG{p}{(}\PYG{n}{st\PYGZus{}name}\PYG{p}{,} \PYG{n}{t}\PYG{p}{)}\PYG{p}{]}\PYG{p}{:} \PYG{l+m+mf}{1.0}\PYG{p}{,}
            \PYG{n+nb+bp}{self}\PYG{o}{.}\PYG{n}{variables}\PYG{p}{[}\PYG{l+s+s2}{\PYGZdq{}}\PYG{l+s+s2}{storage\PYGZus{}charge}\PYG{l+s+s2}{\PYGZdq{}}\PYG{p}{]}\PYG{p}{[}\PYG{p}{(}\PYG{n}{st\PYGZus{}name}\PYG{p}{,} \PYG{n}{t}\PYG{p}{)}\PYG{p}{]}\PYG{p}{:} \PYG{o}{\PYGZhy{}}\PYG{n}{eta\PYGZus{}ch} \PYG{o}{*} \PYG{n+nb+bp}{self}\PYG{o}{.}\PYG{n}{dt}\PYG{p}{,}
            \PYG{n+nb+bp}{self}\PYG{o}{.}\PYG{n}{variables}\PYG{p}{[}\PYG{l+s+s2}{\PYGZdq{}}\PYG{l+s+s2}{storage\PYGZus{}discharge}\PYG{l+s+s2}{\PYGZdq{}}\PYG{p}{]}\PYG{p}{[}\PYG{p}{(}\PYG{n}{st\PYGZus{}name}\PYG{p}{,} \PYG{n}{t}\PYG{p}{)}\PYG{p}{]}\PYG{p}{:} \PYG{n+nb+bp}{self}\PYG{o}{.}\PYG{n}{dt} \PYG{o}{/} \PYG{n}{eta\PYGZus{}dis}\PYG{p}{,}
        \PYG{p}{\PYGZcb{}}
        \PYG{c+c1}{\PYGZsh{} \CJKCode{状态仅在所属边界组内衔接，禁止跨无关场景串联。}}
        \PYG{n}{group\PYGZus{}times} \PYG{o}{=} \PYG{n+nb+bp}{self}\PYG{o}{.}\PYG{n}{boundary\PYGZus{}groups}\PYG{p}{[}\PYG{n+nb+bp}{self}\PYG{o}{.}\PYG{n}{time\PYGZus{}boundary\PYGZus{}group}\PYG{p}{[}\PYG{n}{t}\PYG{p}{]}\PYG{p}{]}
        \PYG{n}{is\PYGZus{}group\PYGZus{}start} \PYG{o}{=} \PYG{n}{t} \PYG{o}{==} \PYG{n}{group\PYGZus{}times}\PYG{p}{[}\PYG{l+m+mi}{0}\PYG{p}{]}
        \PYG{c+c1}{\PYGZsh{} \CJKCode{循环模式使用组内前驱；非首时段同样递推。}}
        \PYG{k}{if} \PYG{n}{boundary\PYGZus{}mode} \PYG{o}{==} \PYG{l+s+s2}{\PYGZdq{}}\PYG{l+s+s2}{cyclic}\PYG{l+s+s2}{\PYGZdq{}} \PYG{o+ow}{or} \PYG{o+ow}{not} \PYG{n}{is\PYGZus{}group\PYGZus{}start}\PYG{p}{:}
            \PYG{n}{prev} \PYG{o}{=} \PYG{n+nb+bp}{self}\PYG{o}{.}\PYG{n}{time\PYGZus{}predecessor}\PYG{p}{[}\PYG{n}{t}\PYG{p}{]}
            \PYG{n+nb+bp}{self}\PYG{o}{.}\PYG{n}{\PYGZus{}add}\PYG{p}{(}\PYG{n}{row}\PYG{p}{,} \PYG{n+nb+bp}{self}\PYG{o}{.}\PYG{n}{variables}\PYG{p}{[}\PYG{l+s+s2}{\PYGZdq{}}\PYG{l+s+s2}{storage\PYGZus{}soc}\PYG{l+s+s2}{\PYGZdq{}}\PYG{p}{]}\PYG{p}{[}\PYG{p}{(}\PYG{n}{st\PYGZus{}name}\PYG{p}{,} \PYG{n}{prev}\PYG{p}{)}\PYG{p}{]}\PYG{p}{,} \PYG{o}{\PYGZhy{}}\PYG{n}{loss\PYGZus{}factor}\PYG{p}{)}
            \PYG{n}{p}\PYG{o}{.}\PYG{n}{add\PYGZus{}constraint}\PYG{p}{(}\PYG{n}{row}\PYG{p}{,} \PYG{n}{lb}\PYG{o}{=}\PYG{l+m+mf}{0.0}\PYG{p}{,} \PYG{n}{ub}\PYG{o}{=}\PYG{l+m+mf}{0.0}\PYG{p}{)}
        \PYG{k}{else}\PYG{p}{:}
            \PYG{c+c1}{\PYGZsh{} \CJKCode{固定边界以装机能量容量的声明比例给定初态。}}
            \PYG{n}{initial\PYGZus{}fraction} \PYG{o}{=} \PYG{n+nb}{float}\PYG{p}{(}\PYG{n}{storage}\PYG{o}{.}\PYG{n}{get}\PYG{p}{(}\PYG{l+s+s2}{\PYGZdq{}}\PYG{l+s+s2}{initial\PYGZus{}soc\PYGZus{}fraction}\PYG{l+s+s2}{\PYGZdq{}}\PYG{p}{,} \PYG{l+m+mf}{0.5}\PYG{p}{)}\PYG{p}{)}
            \PYG{k}{if} \PYG{n}{cap\PYGZus{}idx} \PYG{o+ow}{is} \PYG{o+ow}{not} \PYG{k+kc}{None}\PYG{p}{:}
                \PYG{n}{row}\PYG{p}{[}\PYG{n}{cap\PYGZus{}idx}\PYG{p}{]} \PYG{o}{=} \PYG{o}{\PYGZhy{}}\PYG{n}{loss\PYGZus{}factor} \PYG{o}{*} \PYG{n}{initial\PYGZus{}fraction}
            \PYG{c+c1}{\PYGZsh{} \CJKCode{首小时先计初始储热的自损，再计本小时充放。}}
            \PYG{n}{initial\PYGZus{}energy} \PYG{o}{=} \PYG{n}{loss\PYGZus{}factor} \PYG{o}{*} \PYG{n}{initial\PYGZus{}fraction} \PYG{o}{*} \PYG{n}{existing}
            \PYG{n}{p}\PYG{o}{.}\PYG{n}{add\PYGZus{}constraint}\PYG{p}{(}\PYG{n}{row}\PYG{p}{,} \PYG{n}{lb}\PYG{o}{=}\PYG{n}{initial\PYGZus{}energy}\PYG{p}{,} \PYG{n}{ub}\PYG{o}{=}\PYG{n}{initial\PYGZus{}energy}\PYG{p}{)}
    \PYG{c+c1}{\PYGZsh{} \CJKCode{固定模式期末是最低}SOC\CJKCode{比例约束，不强制等于初态。}}
    \PYG{k}{if} \PYG{n}{boundary\PYGZus{}mode} \PYG{o}{==} \PYG{l+s+s2}{\PYGZdq{}}\PYG{l+s+s2}{fixed}\PYG{l+s+s2}{\PYGZdq{}}\PYG{p}{:}
        \PYG{n}{final\PYGZus{}fraction} \PYG{o}{=} \PYG{n+nb}{float}\PYG{p}{(}\PYG{n}{storage}\PYG{o}{.}\PYG{n}{get}\PYG{p}{(}\PYG{l+s+s2}{\PYGZdq{}}\PYG{l+s+s2}{final\PYGZus{}soc\PYGZus{}min\PYGZus{}fraction}\PYG{l+s+s2}{\PYGZdq{}}\PYG{p}{,} \PYG{l+m+mf}{0.0}\PYG{p}{)}\PYG{p}{)}
        \PYG{k}{for} \PYG{n}{group\PYGZus{}times} \PYG{o+ow}{in} \PYG{n+nb+bp}{self}\PYG{o}{.}\PYG{n}{boundary\PYGZus{}groups}\PYG{o}{.}\PYG{n}{values}\PYG{p}{(}\PYG{p}{)}\PYG{p}{:}
            \PYG{n}{final\PYGZus{}row} \PYG{o}{=} \PYG{p}{\PYGZob{}}\PYG{n+nb+bp}{self}\PYG{o}{.}\PYG{n}{variables}\PYG{p}{[}\PYG{l+s+s2}{\PYGZdq{}}\PYG{l+s+s2}{storage\PYGZus{}soc}\PYG{l+s+s2}{\PYGZdq{}}\PYG{p}{]}\PYG{p}{[}\PYG{p}{(}\PYG{n}{st\PYGZus{}name}\PYG{p}{,} \PYG{n}{group\PYGZus{}times}\PYG{p}{[}\PYG{o}{\PYGZhy{}}\PYG{l+m+mi}{1}\PYG{p}{]}\PYG{p}{)}\PYG{p}{]}\PYG{p}{:} \PYG{l+m+mf}{1.0}\PYG{p}{\PYGZcb{}}
            \PYG{k}{if} \PYG{n}{cap\PYGZus{}idx} \PYG{o+ow}{is} \PYG{o+ow}{not} \PYG{k+kc}{None}\PYG{p}{:}
                \PYG{n}{final\PYGZus{}row}\PYG{p}{[}\PYG{n}{cap\PYGZus{}idx}\PYG{p}{]} \PYG{o}{=} \PYG{o}{\PYGZhy{}}\PYG{n}{final\PYGZus{}fraction}
            \PYG{c+c1}{\PYGZsh{} \CJKCode{期末库存下界与总装机容量保持同一比例。}}
            \PYG{n}{p}\PYG{o}{.}\PYG{n}{add\PYGZus{}constraint}\PYG{p}{(}\PYG{n}{final\PYGZus{}row}\PYG{p}{,} \PYG{n}{lb}\PYG{o}{=}\PYG{n}{final\PYGZus{}fraction} \PYG{o}{*} \PYG{n}{existing}\PYG{p}{)}
\end{Verbatim}

\endgroup

\noindent\textbf{结果核对。}复核顺序为容量上界、充放互斥、逐时SOC递推、初末边界。循环边界在组内首尾相接，固定边界按初态比例起算并检查期末最低SOC。

\Needspace{12\baselineskip}
\subsection*{4．节点温位逐时热平衡}
将热源、双向管流、储热、需求响应和交付负荷接入同一节点温位等式；不以总量平衡代替局部平衡。

\noindent\textbf{输入：}已注册的出力与输配变量（kW）、管段效率、需求列及建筑模型覆盖范围。\par
\noindent\textbf{输出：}每个节点、温位、时段的热量等式。其他变量注册及指标统计见完整文件。\par

\noindent\textbf{实现思路。}节点平衡把每个变量在该节点的净作用写成稀疏系数行：本地产热及入流为正，出流与充热为负，到达端热流乘管段效率。储热、需求响应、计量和末端温位改造都在同一等式中登记。对启用逐楼RC的节点，需求由决策中的实际交付热量承接，避免再次加上同一建筑的固定负荷。

{\small 来源：\path{model.py}，原第1603--1697行；95行原代码、14行补充注释。}\par

\begingroup
\setstretch{1.0}
\begin{Verbatim}[commandchars=\\\{\},numbers=left,firstnumber=1,stepnumber=1,fontsize=\fontsize{9.5}{11.5}\selectfont,breaklines=true,breakanywhere=true,frame=single,framesep=3mm,numbersep=5pt,rulecolor=\color{black},xleftmargin=18pt]
\PYG{c+c1}{\PYGZsh{} \CJKCode{每一行对应节点}n\CJKCode{、温位}g\CJKCode{、时段}t\CJKCode{，热流系数均采用功率口径。}}
\PYG{c+c1}{\PYGZsh{} Node\PYGZhy{}grade\PYGZhy{}time heat balances.}
\PYG{k}{for} \PYG{n}{node\PYGZus{}name}\PYG{p}{,} \PYG{n}{node} \PYG{o+ow}{in} \PYG{n+nb+bp}{self}\PYG{o}{.}\PYG{n}{nodes}\PYG{o}{.}\PYG{n}{items}\PYG{p}{(}\PYG{p}{)}\PYG{p}{:}
    \PYG{k}{for} \PYG{n}{grade\PYGZus{}name}\PYG{p}{,} \PYG{n}{grade} \PYG{o+ow}{in} \PYG{n+nb+bp}{self}\PYG{o}{.}\PYG{n}{grades}\PYG{o}{.}\PYG{n}{items}\PYG{p}{(}\PYG{p}{)}\PYG{p}{:}
        \PYG{n}{demand\PYGZus{}column} \PYG{o}{=} \PYG{n}{node}\PYG{o}{.}\PYG{n}{get}\PYG{p}{(}\PYG{l+s+s2}{\PYGZdq{}}\PYG{l+s+s2}{demand\PYGZus{}columns}\PYG{l+s+s2}{\PYGZdq{}}\PYG{p}{,} \PYG{p}{\PYGZob{}}\PYG{p}{\PYGZcb{}}\PYG{p}{)}\PYG{o}{.}\PYG{n}{get}\PYG{p}{(}\PYG{n}{grade\PYGZus{}name}\PYG{p}{)}
        \PYG{k}{for} \PYG{n}{t} \PYG{o+ow}{in} \PYG{n+nb+bp}{self}\PYG{o}{.}\PYG{n}{times}\PYG{p}{:}
            \PYG{n}{row}\PYG{p}{:} \PYG{n+nb}{dict}\PYG{p}{[}\PYG{n+nb}{int}\PYG{p}{,} \PYG{n+nb}{float}\PYG{p}{]} \PYG{o}{=} \PYG{p}{\PYGZob{}}\PYG{p}{\PYGZcb{}}
            \PYG{c+c1}{\PYGZsh{} \CJKCode{仅汇入安装在本节点且允许服务该温位的热源出力。}}
            \PYG{k}{for} \PYG{n}{s\PYGZus{}name}\PYG{p}{,} \PYG{n}{source} \PYG{o+ow}{in} \PYG{n+nb+bp}{self}\PYG{o}{.}\PYG{n}{sources}\PYG{o}{.}\PYG{n}{items}\PYG{p}{(}\PYG{p}{)}\PYG{p}{:}
                \PYG{n}{idx} \PYG{o}{=} \PYG{n+nb+bp}{self}\PYG{o}{.}\PYG{n}{variables}\PYG{p}{[}\PYG{l+s+s2}{\PYGZdq{}}\PYG{l+s+s2}{source\PYGZus{}heat}\PYG{l+s+s2}{\PYGZdq{}}\PYG{p}{]}\PYG{o}{.}\PYG{n}{get}\PYG{p}{(}\PYG{p}{(}\PYG{n}{s\PYGZus{}name}\PYG{p}{,} \PYG{n}{grade\PYGZus{}name}\PYG{p}{,} \PYG{n}{t}\PYG{p}{)}\PYG{p}{)}
                \PYG{k}{if} \PYG{n}{idx} \PYG{o+ow}{is} \PYG{o+ow}{not} \PYG{k+kc}{None} \PYG{o+ow}{and} \PYG{n}{source}\PYG{p}{[}\PYG{l+s+s2}{\PYGZdq{}}\PYG{l+s+s2}{node}\PYG{l+s+s2}{\PYGZdq{}}\PYG{p}{]} \PYG{o}{==} \PYG{n}{node\PYGZus{}name}\PYG{p}{:}
                    \PYG{n+nb+bp}{self}\PYG{o}{.}\PYG{n}{\PYGZus{}add}\PYG{p}{(}\PYG{n}{row}\PYG{p}{,} \PYG{n}{idx}\PYG{p}{,} \PYG{l+m+mf}{1.0}\PYG{p}{)}
            \PYG{c+c1}{\PYGZsh{} \CJKCode{正反方向分开登记；管容量与方向互斥约束在此前装配。}}
            \PYG{k}{for} \PYG{n}{e\PYGZus{}name}\PYG{p}{,} \PYG{n}{pipe} \PYG{o+ow}{in} \PYG{n+nb+bp}{self}\PYG{o}{.}\PYG{n}{pipes}\PYG{o}{.}\PYG{n}{items}\PYG{p}{(}\PYG{p}{)}\PYG{p}{:}
                \PYG{n}{forward} \PYG{o}{=} \PYG{n+nb+bp}{self}\PYG{o}{.}\PYG{n}{variables}\PYG{p}{[}\PYG{l+s+s2}{\PYGZdq{}}\PYG{l+s+s2}{pipe\PYGZus{}flow\PYGZus{}forward}\PYG{l+s+s2}{\PYGZdq{}}\PYG{p}{]}\PYG{p}{[}\PYG{p}{(}\PYG{n}{e\PYGZus{}name}\PYG{p}{,} \PYG{n}{grade\PYGZus{}name}\PYG{p}{,} \PYG{n}{t}\PYG{p}{)}\PYG{p}{]}
                \PYG{c+c1}{\PYGZsh{} \CJKCode{离开本节点的热流以负号计入，尚未扣除下游管损。}}
                \PYG{k}{if} \PYG{n}{pipe}\PYG{p}{[}\PYG{l+s+s2}{\PYGZdq{}}\PYG{l+s+s2}{from}\PYG{l+s+s2}{\PYGZdq{}}\PYG{p}{]} \PYG{o}{==} \PYG{n}{node\PYGZus{}name}\PYG{p}{:}
                    \PYG{n+nb+bp}{self}\PYG{o}{.}\PYG{n}{\PYGZus{}add}\PYG{p}{(}\PYG{n}{row}\PYG{p}{,} \PYG{n}{forward}\PYG{p}{,} \PYG{o}{\PYGZhy{}}\PYG{l+m+mf}{1.0}\PYG{p}{)}
                \PYG{c+c1}{\PYGZsh{} \CJKCode{到达节点的热流乘效率；接收端承担输配损失。}}
                \PYG{k}{if} \PYG{n}{pipe}\PYG{p}{[}\PYG{l+s+s2}{\PYGZdq{}}\PYG{l+s+s2}{to}\PYG{l+s+s2}{\PYGZdq{}}\PYG{p}{]} \PYG{o}{==} \PYG{n}{node\PYGZus{}name}\PYG{p}{:}
                    \PYG{n+nb+bp}{self}\PYG{o}{.}\PYG{n}{\PYGZus{}add}\PYG{p}{(}\PYG{n}{row}\PYG{p}{,} \PYG{n}{forward}\PYG{p}{,} \PYG{n+nb+bp}{self}\PYG{o}{.}\PYG{n}{\PYGZus{}pipe\PYGZus{}efficiency}\PYG{p}{(}\PYG{n}{pipe}\PYG{p}{,} \PYG{n}{t}\PYG{p}{)}\PYG{p}{)}
                    \PYG{n}{activated} \PYG{o}{=} \PYG{n+nb+bp}{self}\PYG{o}{.}\PYG{n}{variables}\PYG{p}{[}\PYG{l+s+s2}{\PYGZdq{}}\PYG{l+s+s2}{pipe\PYGZus{}retrofit\PYGZus{}flow\PYGZus{}forward}\PYG{l+s+s2}{\PYGZdq{}}\PYG{p}{]}\PYG{o}{.}\PYG{n}{get}\PYG{p}{(}
                        \PYG{p}{(}\PYG{n}{e\PYGZus{}name}\PYG{p}{,} \PYG{n}{grade\PYGZus{}name}\PYG{p}{,} \PYG{n}{t}\PYG{p}{)}
                    \PYG{p}{)}
                    \PYG{k}{if} \PYG{n}{activated} \PYG{o+ow}{is} \PYG{o+ow}{not} \PYG{k+kc}{None}\PYG{p}{:}
                        \PYG{n}{gain} \PYG{o}{=} \PYG{n+nb+bp}{self}\PYG{o}{.}\PYG{n}{\PYGZus{}pipe\PYGZus{}retrofit\PYGZus{}efficiency}\PYG{p}{(}\PYG{n}{pipe}\PYG{p}{,} \PYG{n}{t}\PYG{p}{)} \PYG{o}{\PYGZhy{}} \PYG{n+nb+bp}{self}\PYG{o}{.}\PYG{n}{\PYGZus{}pipe\PYGZus{}efficiency}\PYG{p}{(}\PYG{n}{pipe}\PYG{p}{,} \PYG{n}{t}\PYG{p}{)}
                        \PYG{n+nb+bp}{self}\PYG{o}{.}\PYG{n}{\PYGZus{}add}\PYG{p}{(}\PYG{n}{row}\PYG{p}{,} \PYG{n}{activated}\PYG{p}{,} \PYG{n+nb}{max}\PYG{p}{(}\PYG{l+m+mf}{0.0}\PYG{p}{,} \PYG{n}{gain}\PYG{p}{)}\PYG{p}{)}
                \PYG{c+c1}{\PYGZsh{} \CJKCode{允许反向输送时使用独立变量，符号按端点反转。}}
                \PYG{n}{reverse} \PYG{o}{=} \PYG{n+nb+bp}{self}\PYG{o}{.}\PYG{n}{variables}\PYG{p}{[}\PYG{l+s+s2}{\PYGZdq{}}\PYG{l+s+s2}{pipe\PYGZus{}flow\PYGZus{}reverse}\PYG{l+s+s2}{\PYGZdq{}}\PYG{p}{]}\PYG{o}{.}\PYG{n}{get}\PYG{p}{(}\PYG{p}{(}\PYG{n}{e\PYGZus{}name}\PYG{p}{,} \PYG{n}{grade\PYGZus{}name}\PYG{p}{,} \PYG{n}{t}\PYG{p}{)}\PYG{p}{)}
                \PYG{k}{if} \PYG{n}{reverse} \PYG{o+ow}{is} \PYG{o+ow}{not} \PYG{k+kc}{None}\PYG{p}{:}
                    \PYG{k}{if} \PYG{n}{pipe}\PYG{p}{[}\PYG{l+s+s2}{\PYGZdq{}}\PYG{l+s+s2}{to}\PYG{l+s+s2}{\PYGZdq{}}\PYG{p}{]} \PYG{o}{==} \PYG{n}{node\PYGZus{}name}\PYG{p}{:}
                        \PYG{n+nb+bp}{self}\PYG{o}{.}\PYG{n}{\PYGZus{}add}\PYG{p}{(}\PYG{n}{row}\PYG{p}{,} \PYG{n}{reverse}\PYG{p}{,} \PYG{o}{\PYGZhy{}}\PYG{l+m+mf}{1.0}\PYG{p}{)}
                    \PYG{k}{if} \PYG{n}{pipe}\PYG{p}{[}\PYG{l+s+s2}{\PYGZdq{}}\PYG{l+s+s2}{from}\PYG{l+s+s2}{\PYGZdq{}}\PYG{p}{]} \PYG{o}{==} \PYG{n}{node\PYGZus{}name}\PYG{p}{:}
                        \PYG{n+nb+bp}{self}\PYG{o}{.}\PYG{n}{\PYGZus{}add}\PYG{p}{(}\PYG{n}{row}\PYG{p}{,} \PYG{n}{reverse}\PYG{p}{,} \PYG{n+nb+bp}{self}\PYG{o}{.}\PYG{n}{\PYGZus{}pipe\PYGZus{}efficiency}\PYG{p}{(}\PYG{n}{pipe}\PYG{p}{,} \PYG{n}{t}\PYG{p}{)}\PYG{p}{)}
                        \PYG{n}{activated\PYGZus{}reverse} \PYG{o}{=} \PYG{n+nb+bp}{self}\PYG{o}{.}\PYG{n}{variables}\PYG{p}{[}\PYG{l+s+s2}{\PYGZdq{}}\PYG{l+s+s2}{pipe\PYGZus{}retrofit\PYGZus{}flow\PYGZus{}reverse}\PYG{l+s+s2}{\PYGZdq{}}\PYG{p}{]}\PYG{o}{.}\PYG{n}{get}\PYG{p}{(}
                            \PYG{p}{(}\PYG{n}{e\PYGZus{}name}\PYG{p}{,} \PYG{n}{grade\PYGZus{}name}\PYG{p}{,} \PYG{n}{t}\PYG{p}{)}
                        \PYG{p}{)}
                        \PYG{k}{if} \PYG{n}{activated\PYGZus{}reverse} \PYG{o+ow}{is} \PYG{o+ow}{not} \PYG{k+kc}{None}\PYG{p}{:}
                            \PYG{n}{gain} \PYG{o}{=} \PYG{n+nb+bp}{self}\PYG{o}{.}\PYG{n}{\PYGZus{}pipe\PYGZus{}retrofit\PYGZus{}efficiency}\PYG{p}{(}\PYG{n}{pipe}\PYG{p}{,} \PYG{n}{t}\PYG{p}{)} \PYG{o}{\PYGZhy{}} \PYG{n+nb+bp}{self}\PYG{o}{.}\PYG{n}{\PYGZus{}pipe\PYGZus{}efficiency}\PYG{p}{(}\PYG{n}{pipe}\PYG{p}{,} \PYG{n}{t}\PYG{p}{)}
                            \PYG{n+nb+bp}{self}\PYG{o}{.}\PYG{n}{\PYGZus{}add}\PYG{p}{(}\PYG{n}{row}\PYG{p}{,} \PYG{n}{activated\PYGZus{}reverse}\PYG{p}{,} \PYG{n+nb}{max}\PYG{p}{(}\PYG{l+m+mf}{0.0}\PYG{p}{,} \PYG{n}{gain}\PYG{p}{)}\PYG{p}{)}
            \PYG{c+c1}{\PYGZsh{} \CJKCode{放热增加节点供给，充热占用节点供给。}}
            \PYG{k}{for} \PYG{n}{st\PYGZus{}name}\PYG{p}{,} \PYG{n}{storage} \PYG{o+ow}{in} \PYG{n+nb+bp}{self}\PYG{o}{.}\PYG{n}{storages}\PYG{o}{.}\PYG{n}{items}\PYG{p}{(}\PYG{p}{)}\PYG{p}{:}
                \PYG{k}{if} \PYG{n}{storage}\PYG{p}{[}\PYG{l+s+s2}{\PYGZdq{}}\PYG{l+s+s2}{node}\PYG{l+s+s2}{\PYGZdq{}}\PYG{p}{]} \PYG{o}{==} \PYG{n}{node\PYGZus{}name} \PYG{o+ow}{and} \PYG{n}{storage}\PYG{p}{[}\PYG{l+s+s2}{\PYGZdq{}}\PYG{l+s+s2}{grade}\PYG{l+s+s2}{\PYGZdq{}}\PYG{p}{]} \PYG{o}{==} \PYG{n}{grade\PYGZus{}name}\PYG{p}{:}
                    \PYG{n+nb+bp}{self}\PYG{o}{.}\PYG{n}{\PYGZus{}add}\PYG{p}{(}\PYG{n}{row}\PYG{p}{,} \PYG{n+nb+bp}{self}\PYG{o}{.}\PYG{n}{variables}\PYG{p}{[}\PYG{l+s+s2}{\PYGZdq{}}\PYG{l+s+s2}{storage\PYGZus{}discharge}\PYG{l+s+s2}{\PYGZdq{}}\PYG{p}{]}\PYG{p}{[}\PYG{p}{(}\PYG{n}{st\PYGZus{}name}\PYG{p}{,} \PYG{n}{t}\PYG{p}{)}\PYG{p}{]}\PYG{p}{,} \PYG{l+m+mf}{1.0}\PYG{p}{)}
                    \PYG{n+nb+bp}{self}\PYG{o}{.}\PYG{n}{\PYGZus{}add}\PYG{p}{(}\PYG{n}{row}\PYG{p}{,} \PYG{n+nb+bp}{self}\PYG{o}{.}\PYG{n}{variables}\PYG{p}{[}\PYG{l+s+s2}{\PYGZdq{}}\PYG{l+s+s2}{storage\PYGZus{}charge}\PYG{l+s+s2}{\PYGZdq{}}\PYG{p}{]}\PYG{p}{[}\PYG{p}{(}\PYG{n}{st\PYGZus{}name}\PYG{p}{,} \PYG{n}{t}\PYG{p}{)}\PYG{p}{]}\PYG{p}{,} \PYG{o}{\PYGZhy{}}\PYG{l+m+mf}{1.0}\PYG{p}{)}
            \PYG{c+c1}{\PYGZsh{} \CJKCode{只有配置允许并注册失供变量时才进入等式；严格主模型可无此项。}}
            \PYG{n}{unmet} \PYG{o}{=} \PYG{n+nb+bp}{self}\PYG{o}{.}\PYG{n}{variables}\PYG{p}{[}\PYG{l+s+s2}{\PYGZdq{}}\PYG{l+s+s2}{unserved}\PYG{l+s+s2}{\PYGZdq{}}\PYG{p}{]}\PYG{o}{.}\PYG{n}{get}\PYG{p}{(}\PYG{p}{(}\PYG{n}{node\PYGZus{}name}\PYG{p}{,} \PYG{n}{grade\PYGZus{}name}\PYG{p}{,} \PYG{n}{t}\PYG{p}{)}\PYG{p}{)}
            \PYG{k}{if} \PYG{n}{unmet} \PYG{o+ow}{is} \PYG{o+ow}{not} \PYG{k+kc}{None}\PYG{p}{:}
                \PYG{n+nb+bp}{self}\PYG{o}{.}\PYG{n}{\PYGZus{}add}\PYG{p}{(}\PYG{n}{row}\PYG{p}{,} \PYG{n}{unmet}\PYG{p}{,} \PYG{l+m+mf}{1.0}\PYG{p}{)}
            \PYG{c+c1}{\PYGZsh{} \CJKCode{上调增加需求、下调减少需求，分别按相反符号登记。}}
            \PYG{n}{up} \PYG{o}{=} \PYG{n+nb+bp}{self}\PYG{o}{.}\PYG{n}{variables}\PYG{p}{[}\PYG{l+s+s2}{\PYGZdq{}}\PYG{l+s+s2}{demand\PYGZus{}up}\PYG{l+s+s2}{\PYGZdq{}}\PYG{p}{]}\PYG{o}{.}\PYG{n}{get}\PYG{p}{(}\PYG{p}{(}\PYG{n}{node\PYGZus{}name}\PYG{p}{,} \PYG{n}{grade\PYGZus{}name}\PYG{p}{,} \PYG{n}{t}\PYG{p}{)}\PYG{p}{)}
            \PYG{n}{down} \PYG{o}{=} \PYG{n+nb+bp}{self}\PYG{o}{.}\PYG{n}{variables}\PYG{p}{[}\PYG{l+s+s2}{\PYGZdq{}}\PYG{l+s+s2}{demand\PYGZus{}down}\PYG{l+s+s2}{\PYGZdq{}}\PYG{p}{]}\PYG{o}{.}\PYG{n}{get}\PYG{p}{(}\PYG{p}{(}\PYG{n}{node\PYGZus{}name}\PYG{p}{,} \PYG{n}{grade\PYGZus{}name}\PYG{p}{,} \PYG{n}{t}\PYG{p}{)}\PYG{p}{)}
            \PYG{k}{if} \PYG{n}{up} \PYG{o+ow}{is} \PYG{o+ow}{not} \PYG{k+kc}{None}\PYG{p}{:}
                \PYG{n+nb+bp}{self}\PYG{o}{.}\PYG{n}{\PYGZus{}add}\PYG{p}{(}\PYG{n}{row}\PYG{p}{,} \PYG{n}{up}\PYG{p}{,} \PYG{o}{\PYGZhy{}}\PYG{l+m+mf}{1.0}\PYG{p}{)}
            \PYG{k}{if} \PYG{n}{down} \PYG{o+ow}{is} \PYG{o+ow}{not} \PYG{k+kc}{None}\PYG{p}{:}
                \PYG{n+nb+bp}{self}\PYG{o}{.}\PYG{n}{\PYGZus{}add}\PYG{p}{(}\PYG{n}{row}\PYG{p}{,} \PYG{n}{down}\PYG{p}{,} \PYG{l+m+mf}{1.0}\PYG{p}{)}
            \PYG{c+c1}{\PYGZsh{} RC\CJKCode{覆盖节点以实际交付量闭合，不能再叠加同一建筑固定负荷。}}
            \PYG{k}{if} \PYG{n+nb+bp}{self}\PYG{o}{.}\PYG{n}{building\PYGZus{}thermal} \PYG{o+ow}{is} \PYG{o+ow}{not} \PYG{k+kc}{None} \PYG{o+ow}{and} \PYG{p}{(}\PYG{n}{node\PYGZus{}name}\PYG{p}{,}\PYG{n}{grade\PYGZus{}name}\PYG{p}{)} \PYG{o+ow}{in} \PYG{n+nb+bp}{self}\PYG{o}{.}\PYG{n}{building\PYGZus{}thermal}\PYG{o}{.}\PYG{n}{covered\PYGZus{}node\PYGZus{}grades}\PYG{p}{:}
                \PYG{n}{delivered} \PYG{o}{=} \PYG{n+nb+bp}{self}\PYG{o}{.}\PYG{n}{building\PYGZus{}thermal}\PYG{o}{.}\PYG{n}{node\PYGZus{}heat}\PYG{p}{[}\PYG{n}{node\PYGZus{}name}\PYG{p}{,}\PYG{n}{grade\PYGZus{}name}\PYG{p}{,}\PYG{n}{t}\PYG{p}{]}
                \PYG{n+nb+bp}{self}\PYG{o}{.}\PYG{n}{\PYGZus{}add}\PYG{p}{(}\PYG{n}{row}\PYG{p}{,}\PYG{n}{delivered}\PYG{p}{,}\PYG{o}{\PYGZhy{}}\PYG{l+m+mf}{1.0}\PYG{p}{)}
                \PYG{n}{p}\PYG{o}{.}\PYG{n}{add\PYGZus{}constraint}\PYG{p}{(}\PYG{n}{row}\PYG{p}{,}\PYG{n}{lb}\PYG{o}{=}\PYG{l+m+mf}{0.0}\PYG{p}{,}\PYG{n}{ub}\PYG{o}{=}\PYG{l+m+mf}{0.0}\PYG{p}{)}
                \PYG{n+nb+bp}{self}\PYG{o}{.}\PYG{n}{\PYGZus{}record\PYGZus{}time\PYGZus{}metric}\PYG{p}{(}\PYG{l+s+s2}{\PYGZdq{}}\PYG{l+s+s2}{demand\PYGZus{}kwh}\PYG{l+s+s2}{\PYGZdq{}}\PYG{p}{,}\PYG{n}{delivered}\PYG{p}{,}\PYG{l+m+mf}{1.0}\PYG{p}{,}\PYG{n}{t}\PYG{p}{)}
                \PYG{n+nb+bp}{self}\PYG{o}{.}\PYG{n}{\PYGZus{}record\PYGZus{}time\PYGZus{}metric}\PYG{p}{(}\PYG{l+s+s2}{\PYGZdq{}}\PYG{l+s+s2}{building\PYGZus{}heat\PYGZus{}delivered\PYGZus{}kwh}\PYG{l+s+s2}{\PYGZdq{}}\PYG{p}{,}\PYG{n}{delivered}\PYG{p}{,}\PYG{l+m+mf}{1.0}\PYG{p}{,}\PYG{n}{t}\PYG{p}{)}
                \PYG{n+nb+bp}{self}\PYG{o}{.}\PYG{n}{\PYGZus{}record\PYGZus{}time\PYGZus{}metric}\PYG{p}{(}\PYG{l+s+s2}{\PYGZdq{}}\PYG{l+s+s2}{demand\PYGZus{}exergy\PYGZus{}kwh}\PYG{l+s+s2}{\PYGZdq{}}\PYG{p}{,}\PYG{n}{delivered}\PYG{p}{,}\PYG{n+nb+bp}{self}\PYG{o}{.}\PYG{n}{\PYGZus{}exergy\PYGZus{}factor}\PYG{p}{(}\PYG{n+nb}{float}\PYG{p}{(}\PYG{n}{grade}\PYG{p}{[}\PYG{l+s+s2}{\PYGZdq{}}\PYG{l+s+s2}{temperature\PYGZus{}c}\PYG{l+s+s2}{\PYGZdq{}}\PYG{p}{]}\PYG{p}{)}\PYG{p}{)}\PYG{p}{,}\PYG{n}{t}\PYG{p}{)}
                \PYG{n}{reference} \PYG{o}{=} \PYG{n+nb+bp}{self}\PYG{o}{.}\PYG{n}{\PYGZus{}demand\PYGZus{}value}\PYG{p}{(}\PYG{n}{node}\PYG{p}{,}\PYG{n}{demand\PYGZus{}column}\PYG{p}{,}\PYG{n}{t}\PYG{p}{)} \PYG{k}{if} \PYG{n}{demand\PYGZus{}column} \PYG{k}{else} \PYG{l+m+mf}{0.0}
                \PYG{n+nb+bp}{self}\PYG{o}{.}\PYG{n}{metric\PYGZus{}constants}\PYG{p}{[}\PYG{l+s+s2}{\PYGZdq{}}\PYG{l+s+s2}{building\PYGZus{}reference\PYGZus{}heat\PYGZus{}kwh}\PYG{l+s+s2}{\PYGZdq{}}\PYG{p}{]} \PYG{o}{+}\PYG{o}{=} \PYG{n}{reference}\PYG{o}{*}\PYG{n+nb+bp}{self}\PYG{o}{.}\PYG{n}{\PYGZus{}op\PYGZus{}weight}\PYG{p}{(}\PYG{n}{t}\PYG{p}{)}
                \PYG{n+nb+bp}{self}\PYG{o}{.}\PYG{n}{scenario\PYGZus{}metric\PYGZus{}constants}\PYG{p}{[}\PYG{n+nb+bp}{self}\PYG{o}{.}\PYG{n}{\PYGZus{}scenario}\PYG{p}{(}\PYG{n}{t}\PYG{p}{)}\PYG{p}{]}\PYG{p}{[}\PYG{l+s+s2}{\PYGZdq{}}\PYG{l+s+s2}{building\PYGZus{}reference\PYGZus{}heat\PYGZus{}kwh}\PYG{l+s+s2}{\PYGZdq{}}\PYG{p}{]} \PYG{o}{+}\PYG{o}{=} \PYG{n}{reference}\PYG{o}{*}\PYG{n+nb+bp}{self}\PYG{o}{.}\PYG{n}{\PYGZus{}physical\PYGZus{}weight}\PYG{p}{(}\PYG{n}{t}\PYG{p}{)}
                \PYG{n+nb+bp}{self}\PYG{o}{.}\PYG{n}{\PYGZus{}record\PYGZus{}metering\PYGZus{}reference}\PYG{p}{(}\PYG{n}{node\PYGZus{}name}\PYG{p}{,} \PYG{n}{grade}\PYG{p}{,} \PYG{n}{reference}\PYG{p}{,} \PYG{n}{t}\PYG{p}{)}
                \PYG{k}{continue}
            \PYG{c+c1}{\PYGZsh{} \CJKCode{未被}RC\CJKCode{覆盖的温位沿用登记需求列及其调节接口。}}
            \PYG{n}{demand} \PYG{o}{=} \PYG{n+nb+bp}{self}\PYG{o}{.}\PYG{n}{\PYGZus{}demand\PYGZus{}value}\PYG{p}{(}\PYG{n}{node}\PYG{p}{,} \PYG{n}{demand\PYGZus{}column}\PYG{p}{,} \PYG{n}{t}\PYG{p}{)} \PYG{k}{if} \PYG{n}{demand\PYGZus{}column} \PYG{k}{else} \PYG{l+m+mf}{0.0}
            \PYG{n}{metering} \PYG{o}{=} \PYG{n}{node}\PYG{o}{.}\PYG{n}{get}\PYG{p}{(}\PYG{l+s+s2}{\PYGZdq{}}\PYG{l+s+s2}{heat\PYGZus{}metering}\PYG{l+s+s2}{\PYGZdq{}}\PYG{p}{,} \PYG{p}{\PYGZob{}}\PYG{p}{\PYGZcb{}}\PYG{p}{)}
            \PYG{n}{existing\PYGZus{}metering} \PYG{o}{=} \PYG{n+nb}{float}\PYG{p}{(}
                \PYG{n}{metering}\PYG{o}{.}\PYG{n}{get}\PYG{p}{(}\PYG{l+s+s2}{\PYGZdq{}}\PYG{l+s+s2}{existing\PYGZus{}coverage\PYGZus{}fraction}\PYG{l+s+s2}{\PYGZdq{}}\PYG{p}{,} \PYG{l+m+mf}{0.0}\PYG{p}{)}
            \PYG{p}{)}
            \PYG{n}{saving\PYGZus{}fraction} \PYG{o}{=} \PYG{n+nb}{float}\PYG{p}{(}
                \PYG{n}{metering}\PYG{o}{.}\PYG{n}{get}\PYG{p}{(}\PYG{l+s+s2}{\PYGZdq{}}\PYG{l+s+s2}{saving\PYGZus{}fraction\PYGZus{}at\PYGZus{}full\PYGZus{}coverage}\PYG{l+s+s2}{\PYGZdq{}}\PYG{p}{,} \PYG{l+m+mf}{0.0}\PYG{p}{)}
            \PYG{p}{)}
            \PYG{n}{existing\PYGZus{}metering\PYGZus{}saving} \PYG{o}{=} \PYG{p}{(}
                \PYG{n}{demand} \PYG{o}{*} \PYG{n}{existing\PYGZus{}metering} \PYG{o}{*} \PYG{n}{saving\PYGZus{}fraction}
            \PYG{p}{)}
            \PYG{n}{metering\PYGZus{}idx} \PYG{o}{=} \PYG{n+nb+bp}{self}\PYG{o}{.}\PYG{n}{variables}\PYG{p}{[}\PYG{l+s+s2}{\PYGZdq{}}\PYG{l+s+s2}{heat\PYGZus{}metering\PYGZus{}coverage}\PYG{l+s+s2}{\PYGZdq{}}\PYG{p}{]}\PYG{o}{.}\PYG{n}{get}\PYG{p}{(}
                \PYG{n}{node\PYGZus{}name}
            \PYG{p}{)}
            \PYG{c+c1}{\PYGZsh{} \CJKCode{计量覆盖投资按登记节约系数调整需求，须与原需求保持同一单位。}}
            \PYG{k}{if} \PYG{n}{metering\PYGZus{}idx} \PYG{o+ow}{is} \PYG{o+ow}{not} \PYG{k+kc}{None} \PYG{o+ow}{and} \PYG{n}{demand} \PYG{o}{\PYGZgt{}} \PYG{l+m+mf}{0.0} \PYG{o+ow}{and} \PYG{n}{saving\PYGZus{}fraction} \PYG{o}{\PYGZgt{}} \PYG{l+m+mf}{0.0}\PYG{p}{:}
                \PYG{n+nb+bp}{self}\PYG{o}{.}\PYG{n}{\PYGZus{}add}\PYG{p}{(}\PYG{n}{row}\PYG{p}{,} \PYG{n}{metering\PYGZus{}idx}\PYG{p}{,} \PYG{n}{demand} \PYG{o}{*} \PYG{n}{saving\PYGZus{}fraction}\PYG{p}{)}
            \PYG{c+c1}{\PYGZsh{} \CJKCode{末端改造在两个温位间转移符合条件的需求，不凭空产生热量。}}
            \PYG{n}{retrofit} \PYG{o}{=} \PYG{n}{node}\PYG{o}{.}\PYG{n}{get}\PYG{p}{(}\PYG{l+s+s2}{\PYGZdq{}}\PYG{l+s+s2}{temperature\PYGZus{}retrofit}\PYG{l+s+s2}{\PYGZdq{}}\PYG{p}{)}
            \PYG{n}{retrofit\PYGZus{}idx} \PYG{o}{=} \PYG{n+nb+bp}{self}\PYG{o}{.}\PYG{n}{variables}\PYG{p}{[}\PYG{l+s+s2}{\PYGZdq{}}\PYG{l+s+s2}{terminal\PYGZus{}temperature\PYGZus{}retrofit}\PYG{l+s+s2}{\PYGZdq{}}\PYG{p}{]}\PYG{o}{.}\PYG{n}{get}\PYG{p}{(}\PYG{n}{node\PYGZus{}name}\PYG{p}{)}
            \PYG{k}{if} \PYG{n}{retrofit} \PYG{o+ow}{and} \PYG{n}{retrofit\PYGZus{}idx} \PYG{o+ow}{is} \PYG{o+ow}{not} \PYG{k+kc}{None}\PYG{p}{:}
                \PYG{n}{source\PYGZus{}grade} \PYG{o}{=} \PYG{n+nb}{str}\PYG{p}{(}\PYG{n}{retrofit}\PYG{p}{[}\PYG{l+s+s2}{\PYGZdq{}}\PYG{l+s+s2}{source\PYGZus{}grade}\PYG{l+s+s2}{\PYGZdq{}}\PYG{p}{]}\PYG{p}{)}
                \PYG{n}{target\PYGZus{}grade} \PYG{o}{=} \PYG{n+nb}{str}\PYG{p}{(}\PYG{n}{retrofit}\PYG{p}{[}\PYG{l+s+s2}{\PYGZdq{}}\PYG{l+s+s2}{target\PYGZus{}grade}\PYG{l+s+s2}{\PYGZdq{}}\PYG{p}{]}\PYG{p}{)}
                \PYG{n}{source\PYGZus{}column} \PYG{o}{=} \PYG{n}{node}\PYG{p}{[}\PYG{l+s+s2}{\PYGZdq{}}\PYG{l+s+s2}{demand\PYGZus{}columns}\PYG{l+s+s2}{\PYGZdq{}}\PYG{p}{]}\PYG{p}{[}\PYG{n}{source\PYGZus{}grade}\PYG{p}{]}
                \PYG{n}{shift} \PYG{o}{=} \PYG{p}{(}
                    \PYG{n+nb}{float}\PYG{p}{(}\PYG{n}{retrofit}\PYG{o}{.}\PYG{n}{get}\PYG{p}{(}\PYG{l+s+s2}{\PYGZdq{}}\PYG{l+s+s2}{eligible\PYGZus{}fraction}\PYG{l+s+s2}{\PYGZdq{}}\PYG{p}{,} \PYG{l+m+mf}{1.0}\PYG{p}{)}\PYG{p}{)}
                    \PYG{o}{*} \PYG{n+nb+bp}{self}\PYG{o}{.}\PYG{n}{\PYGZus{}demand\PYGZus{}value}\PYG{p}{(}\PYG{n}{node}\PYG{p}{,} \PYG{n}{source\PYGZus{}column}\PYG{p}{,} \PYG{n}{t}\PYG{p}{)}
                \PYG{p}{)}
                \PYG{k}{if} \PYG{n}{grade\PYGZus{}name} \PYG{o}{==} \PYG{n}{source\PYGZus{}grade}\PYG{p}{:}
                    \PYG{n+nb+bp}{self}\PYG{o}{.}\PYG{n}{\PYGZus{}add}\PYG{p}{(}\PYG{n}{row}\PYG{p}{,} \PYG{n}{retrofit\PYGZus{}idx}\PYG{p}{,} \PYG{n}{shift}\PYG{p}{)}
                \PYG{k}{elif} \PYG{n}{grade\PYGZus{}name} \PYG{o}{==} \PYG{n}{target\PYGZus{}grade}\PYG{p}{:}
                    \PYG{n+nb+bp}{self}\PYG{o}{.}\PYG{n}{\PYGZus{}add}\PYG{p}{(}\PYG{n}{row}\PYG{p}{,} \PYG{n}{retrofit\PYGZus{}idx}\PYG{p}{,} \PYG{o}{\PYGZhy{}}\PYG{n}{shift}\PYG{p}{)}
            \PYG{c+c1}{\PYGZsh{} \CJKCode{固定需求分支以双侧相等界闭合；此前调节项已在系数行内。}}
            \PYG{n}{effective\PYGZus{}constant\PYGZus{}demand} \PYG{o}{=} \PYG{n}{demand} \PYG{o}{\PYGZhy{}} \PYG{n}{existing\PYGZus{}metering\PYGZus{}saving}
            \PYG{n}{p}\PYG{o}{.}\PYG{n}{add\PYGZus{}constraint}\PYG{p}{(}
                \PYG{n}{row}\PYG{p}{,}
                \PYG{n}{lb}\PYG{o}{=}\PYG{n}{effective\PYGZus{}constant\PYGZus{}demand}\PYG{p}{,}
                \PYG{n}{ub}\PYG{o}{=}\PYG{n}{effective\PYGZus{}constant\PYGZus{}demand}\PYG{p}{,}
            \PYG{p}{)}
\end{Verbatim}

\endgroup

\noindent\textbf{结果核对。}检查时按节点、温位、小时逐行重算，重点核对管段首尾方向、入流效率及RC覆盖集合。全系统供需总量相等不能替代这些局部等式。

\Needspace{12\baselineskip}
\subsection*{5．SciPy稀疏MILP求解接口}
完整求解方法，将已装配的稀疏约束、边界及整数类型传给SciPy的milp接口。

\noindent\textbf{输入：}目标系数、约束矩阵、变量上下界、整数类型、间隙和时间预算。\par
\noindent\textbf{输出：}求解器原始结果。是否存在可行解及独立审计由上层模型处理。\par

\noindent\textbf{实现思路。}上层模型以变量索引装配目标系数、约束行和边界，本接口把它们整理为求解器接受的稀疏矩阵。整数类型与连续变量边界分别传入，因此容量、逐时运行和启停状态可以在同一个MILP中处理。矩阵缓存减少相同结构的重复整理，但费用或碳排目标切换时仍保存本次实际目标，供返回后的复核使用。

{\small 来源：\path{linear.py}，原第177--205行；29行原代码、6行补充注释。}\par

\begingroup
\setstretch{1.0}
\begin{Verbatim}[commandchars=\\\{\},numbers=left,firstnumber=1,stepnumber=1,fontsize=\fontsize{9.5}{11.5}\selectfont,breaklines=true,breakanywhere=true,frame=single,framesep=3mm,numbersep=5pt,rulecolor=\color{black},xleftmargin=18pt]
\PYG{c+c1}{\PYGZsh{} \CJKCode{本方法不构建热物理方程，只求解上层已装配的}MILP\CJKCode{。}}
\PYG{k}{def}\PYG{+w}{ }\PYG{n+nf}{solve}\PYG{p}{(}\PYG{n+nb+bp}{self}\PYG{p}{,} \PYG{n}{mip\PYGZus{}gap}\PYG{p}{:} \PYG{n+nb}{float} \PYG{o}{=} \PYG{l+m+mf}{0.005}\PYG{p}{,} \PYG{n}{time\PYGZus{}limit}\PYG{p}{:} \PYG{n+nb}{float} \PYG{o}{|} \PYG{k+kc}{None} \PYG{o}{=} \PYG{l+m+mf}{120.0}\PYG{p}{)}\PYG{p}{:}
    \PYG{c+c1}{\PYGZsh{} \CJKCode{矩阵和边界可缓存；约束修改后由模型维护缓存失效。}}
    \PYG{n}{matrix} \PYG{o}{=} \PYG{n+nb+bp}{self}\PYG{o}{.}\PYG{n}{constraint\PYGZus{}matrix}\PYG{p}{(}\PYG{p}{)}
    \PYG{k}{if} \PYG{n+nb+bp}{self}\PYG{o}{.}\PYG{n}{\PYGZus{}compiled\PYGZus{}constraints} \PYG{o+ow}{is} \PYG{k+kc}{None}\PYG{p}{:}
        \PYG{n+nb+bp}{self}\PYG{o}{.}\PYG{n}{\PYGZus{}compiled\PYGZus{}constraints} \PYG{o}{=} \PYG{n}{LinearConstraint}\PYG{p}{(}
            \PYG{n}{matrix}\PYG{p}{,}
            \PYG{n}{np}\PYG{o}{.}\PYG{n}{asarray}\PYG{p}{(}\PYG{n+nb+bp}{self}\PYG{o}{.}\PYG{n}{row\PYGZus{}lower}\PYG{p}{)}\PYG{p}{,}
            \PYG{n}{np}\PYG{o}{.}\PYG{n}{asarray}\PYG{p}{(}\PYG{n+nb+bp}{self}\PYG{o}{.}\PYG{n}{row\PYGZus{}upper}\PYG{p}{)}\PYG{p}{,}
        \PYG{p}{)}
    \PYG{k}{if} \PYG{n+nb+bp}{self}\PYG{o}{.}\PYG{n}{\PYGZus{}compiled\PYGZus{}bounds} \PYG{o+ow}{is} \PYG{k+kc}{None}\PYG{p}{:}
        \PYG{n+nb+bp}{self}\PYG{o}{.}\PYG{n}{\PYGZus{}compiled\PYGZus{}bounds} \PYG{o}{=} \PYG{n}{Bounds}\PYG{p}{(}
            \PYG{n}{np}\PYG{o}{.}\PYG{n}{asarray}\PYG{p}{(}\PYG{n+nb+bp}{self}\PYG{o}{.}\PYG{n}{lower}\PYG{p}{)}\PYG{p}{,} \PYG{n}{np}\PYG{o}{.}\PYG{n}{asarray}\PYG{p}{(}\PYG{n+nb+bp}{self}\PYG{o}{.}\PYG{n}{upper}\PYG{p}{)}
        \PYG{p}{)}
    \PYG{k}{if} \PYG{n+nb+bp}{self}\PYG{o}{.}\PYG{n}{\PYGZus{}compiled\PYGZus{}integrality} \PYG{o+ow}{is} \PYG{k+kc}{None}\PYG{p}{:}
        \PYG{n+nb+bp}{self}\PYG{o}{.}\PYG{n}{\PYGZus{}compiled\PYGZus{}integrality} \PYG{o}{=} \PYG{n}{np}\PYG{o}{.}\PYG{n}{asarray}\PYG{p}{(}\PYG{n+nb+bp}{self}\PYG{o}{.}\PYG{n}{integrality}\PYG{p}{)}
    \PYG{c+c1}{\PYGZsh{} \CJKCode{相对间隙控制停止条件，达到时间上限也可能仅返回可行候选。}}
    \PYG{n}{options}\PYG{p}{:} \PYG{n+nb}{dict}\PYG{p}{[}\PYG{n+nb}{str}\PYG{p}{,} \PYG{n+nb}{float} \PYG{o}{|} \PYG{n+nb}{bool}\PYG{p}{]} \PYG{o}{=} \PYG{p}{\PYGZob{}}\PYG{l+s+s2}{\PYGZdq{}}\PYG{l+s+s2}{mip\PYGZus{}rel\PYGZus{}gap}\PYG{l+s+s2}{\PYGZdq{}}\PYG{p}{:} \PYG{n}{mip\PYGZus{}gap}\PYG{p}{,} \PYG{l+s+s2}{\PYGZdq{}}\PYG{l+s+s2}{disp}\PYG{l+s+s2}{\PYGZdq{}}\PYG{p}{:} \PYG{k+kc}{False}\PYG{p}{\PYGZcb{}}
    \PYG{k}{if} \PYG{n}{time\PYGZus{}limit} \PYG{o+ow}{is} \PYG{o+ow}{not} \PYG{k+kc}{None}\PYG{p}{:}
        \PYG{n}{options}\PYG{p}{[}\PYG{l+s+s2}{\PYGZdq{}}\PYG{l+s+s2}{time\PYGZus{}limit}\PYG{l+s+s2}{\PYGZdq{}}\PYG{p}{]} \PYG{o}{=} \PYG{n+nb}{float}\PYG{p}{(}\PYG{n}{time\PYGZus{}limit}\PYG{p}{)}
    \PYG{c+c1}{\PYGZsh{} \CJKCode{记录本次实际目标，便于上层改回目标后仍能独立复算。}}
    \PYG{n+nb+bp}{self}\PYG{o}{.}\PYG{n}{\PYGZus{}last\PYGZus{}solve\PYGZus{}objective} \PYG{o}{=} \PYG{n}{np}\PYG{o}{.}\PYG{n}{asarray}\PYG{p}{(}\PYG{n+nb+bp}{self}\PYG{o}{.}\PYG{n}{objective}\PYG{p}{,} \PYG{n}{dtype}\PYG{o}{=}\PYG{n+nb}{float}\PYG{p}{)}\PYG{o}{.}\PYG{n}{copy}\PYG{p}{(}\PYG{p}{)}
    \PYG{n+nb+bp}{self}\PYG{o}{.}\PYG{n}{\PYGZus{}last\PYGZus{}solve\PYGZus{}fun} \PYG{o}{=} \PYG{k+kc}{None}
    \PYG{c+c1}{\PYGZsh{} \CJKCode{整数类型、边界与矩阵同时传入；不能只检查返回目标值。}}
    \PYG{n}{result} \PYG{o}{=} \PYG{n}{milp}\PYG{p}{(}
        \PYG{n}{c}\PYG{o}{=}\PYG{n+nb+bp}{self}\PYG{o}{.}\PYG{n}{\PYGZus{}last\PYGZus{}solve\PYGZus{}objective}\PYG{p}{,}
        \PYG{n}{integrality}\PYG{o}{=}\PYG{n+nb+bp}{self}\PYG{o}{.}\PYG{n}{\PYGZus{}compiled\PYGZus{}integrality}\PYG{p}{,}
        \PYG{n}{bounds}\PYG{o}{=}\PYG{n+nb+bp}{self}\PYG{o}{.}\PYG{n}{\PYGZus{}compiled\PYGZus{}bounds}\PYG{p}{,}
        \PYG{n}{constraints}\PYG{o}{=}\PYG{n+nb+bp}{self}\PYG{o}{.}\PYG{n}{\PYGZus{}compiled\PYGZus{}constraints}\PYG{p}{,}
        \PYG{n}{options}\PYG{o}{=}\PYG{n}{options}\PYG{p}{,}
    \PYG{p}{)}
    \PYG{c+c1}{\PYGZsh{} \CJKCode{有目标值时保留它；无可行解的异常报告在}HeatPlanningModel.solve\CJKCode{中。}}
    \PYG{k}{if} \PYG{n}{result}\PYG{o}{.}\PYG{n}{fun} \PYG{o+ow}{is} \PYG{o+ow}{not} \PYG{k+kc}{None}\PYG{p}{:}
        \PYG{n+nb+bp}{self}\PYG{o}{.}\PYG{n}{\PYGZus{}last\PYGZus{}solve\PYGZus{}fun} \PYG{o}{=} \PYG{n+nb}{float}\PYG{p}{(}\PYG{n}{result}\PYG{o}{.}\PYG{n}{fun}\PYG{p}{)}
    \PYG{k}{return} \PYG{n}{result}
\end{Verbatim}

\endgroup

\noindent\textbf{结果核对。}原始求解结果还需由上层检查可行向量、状态、间隙及数值审计；达到时间上限与满足全部物理约束是两个不同检查项。

\Needspace{12\baselineskip}
\subsection*{6．独立数值与守恒残差复核}
完整数值审计方法，从全解重新计算每一约束行、变量边界、整数性及本次目标。

\noindent\textbf{输入：}求解向量、绝对容差和相对容差，以及已装配的原始矩阵。\par
\noindent\textbf{输出：}残差、最差约束和变量、整数偏差、目标复算误差及通过标记。\par

\noindent\textbf{实现思路。}审计不直接沿用求解器给出的目标和可行性标签，而是从返回向量重新计算$Ax$、变量界、整数偏差及目标值。约束行与变量边界残差采用绝对及缩放相对容差，整数偏差单独采用绝对容差，目标另按自身量级复算，避免不同检查项混用阈值。记录最差约束行与变量名，可将异常定位回具体模块。

{\small 来源：\path{linear.py}，原第40--105行；66行原代码、12行补充注释。}\par

\begingroup
\setstretch{1.0}
\begin{Verbatim}[commandchars=\\\{\},numbers=left,firstnumber=1,stepnumber=1,fontsize=\fontsize{9.5}{11.5}\selectfont,breaklines=true,breakanywhere=true,frame=single,framesep=3mm,numbersep=5pt,rulecolor=\color{black},xleftmargin=18pt]
\PYG{c+c1}{\PYGZsh{} \CJKCode{审计独立重算矩阵行；它不证明输入来自现场或未装配的工程约束成立。}}
\PYG{k}{def}\PYG{+w}{ }\PYG{n+nf}{audit\PYGZus{}solution}\PYG{p}{(}
    \PYG{n+nb+bp}{self}\PYG{p}{,} \PYG{n}{values}\PYG{p}{:} \PYG{n}{np}\PYG{o}{.}\PYG{n}{ndarray}\PYG{p}{,} \PYG{o}{*}\PYG{p}{,} \PYG{n}{absolute\PYGZus{}tolerance}\PYG{p}{:} \PYG{n+nb}{float} \PYG{o}{=} \PYG{l+m+mf}{1e\PYGZhy{}6}\PYG{p}{,}
    \PYG{n}{relative\PYGZus{}tolerance}\PYG{p}{:} \PYG{n+nb}{float} \PYG{o}{=} \PYG{l+m+mf}{1e\PYGZhy{}8}\PYG{p}{,}
\PYG{p}{)} \PYG{o}{\PYGZhy{}}\PYG{o}{\PYGZgt{}} \PYG{n+nb}{dict}\PYG{p}{:}
\PYG{+w}{    }\PYG{l+s+sd}{\PYGZdq{}\PYGZdq{}\PYGZdq{}Recompute every row, bound and integer restriction from the full solution.}

\PYG{l+s+sd}{    Scaled violations divide by atol + rtol * row magnitude; \PYGZlt{}= 1 passes.}
\PYG{l+s+sd}{    The objective snapshot belongs to the last solve, including emissions}
\PYG{l+s+sd}{    objectives that a caller restores after the solver returns.}
\PYG{l+s+sd}{    \PYGZdq{}\PYGZdq{}\PYGZdq{}}
    \PYG{c+c1}{\PYGZsh{} \CJKCode{拒绝无效容差；不可通过任意放大容差掩盖物理超限。}}
    \PYG{k}{if} \PYG{o+ow}{not} \PYG{n}{np}\PYG{o}{.}\PYG{n}{isfinite}\PYG{p}{(}\PYG{n}{absolute\PYGZus{}tolerance}\PYG{p}{)} \PYG{o+ow}{or} \PYG{o+ow}{not} \PYG{n}{np}\PYG{o}{.}\PYG{n}{isfinite}\PYG{p}{(}\PYG{n}{relative\PYGZus{}tolerance}\PYG{p}{)} \PYG{o+ow}{or} \PYG{n}{absolute\PYGZus{}tolerance} \PYG{o}{\PYGZlt{}}\PYG{o}{=} \PYG{l+m+mi}{0} \PYG{o+ow}{or} \PYG{n}{relative\PYGZus{}tolerance} \PYG{o}{\PYGZlt{}} \PYG{l+m+mi}{0}\PYG{p}{:}
        \PYG{k}{raise} \PYG{n+ne}{ValueError}\PYG{p}{(}\PYG{l+s+s2}{\PYGZdq{}}\PYG{l+s+s2}{audit tolerances must be positive/nonnegative}\PYG{l+s+s2}{\PYGZdq{}}\PYG{p}{)}
    \PYG{n}{x} \PYG{o}{=} \PYG{n}{np}\PYG{o}{.}\PYG{n}{asarray}\PYG{p}{(}\PYG{n}{values}\PYG{p}{,} \PYG{n}{dtype}\PYG{o}{=}\PYG{n+nb}{float}\PYG{p}{)}
    \PYG{n}{report} \PYG{o}{=} \PYG{p}{\PYGZob{}}
        \PYG{l+s+s2}{\PYGZdq{}}\PYG{l+s+s2}{numerical\PYGZus{}audit\PYGZus{}passed}\PYG{l+s+s2}{\PYGZdq{}}\PYG{p}{:} \PYG{k+kc}{False}\PYG{p}{,}
        \PYG{l+s+s2}{\PYGZdq{}}\PYG{l+s+s2}{numerical\PYGZus{}audit\PYGZus{}tolerance}\PYG{l+s+s2}{\PYGZdq{}}\PYG{p}{:} \PYG{p}{\PYGZob{}}
            \PYG{l+s+s2}{\PYGZdq{}}\PYG{l+s+s2}{absolute}\PYG{l+s+s2}{\PYGZdq{}}\PYG{p}{:} \PYG{n}{absolute\PYGZus{}tolerance}\PYG{p}{,} \PYG{l+s+s2}{\PYGZdq{}}\PYG{l+s+s2}{relative}\PYG{l+s+s2}{\PYGZdq{}}\PYG{p}{:} \PYG{n}{relative\PYGZus{}tolerance}\PYG{p}{,}
            \PYG{l+s+s2}{\PYGZdq{}}\PYG{l+s+s2}{scaled\PYGZus{}violation\PYGZus{}limit}\PYG{l+s+s2}{\PYGZdq{}}\PYG{p}{:} \PYG{l+m+mf}{1.0}\PYG{p}{,}
        \PYG{p}{\PYGZcb{}}\PYG{p}{,}
        \PYG{c+c1}{\PYGZsh{} \CJKCode{首先检查向量维数和有限性，缺失或非有限解立即不通过。}}
        \PYG{l+s+s2}{\PYGZdq{}}\PYG{l+s+s2}{numerical\PYGZus{}solution\PYGZus{}finite}\PYG{l+s+s2}{\PYGZdq{}}\PYG{p}{:} \PYG{n+nb}{bool}\PYG{p}{(}\PYG{n}{x}\PYG{o}{.}\PYG{n}{shape} \PYG{o}{==} \PYG{p}{(}\PYG{n+nb}{len}\PYG{p}{(}\PYG{n+nb+bp}{self}\PYG{o}{.}\PYG{n}{names}\PYG{p}{)}\PYG{p}{,}\PYG{p}{)} \PYG{o+ow}{and} \PYG{n}{np}\PYG{o}{.}\PYG{n}{isfinite}\PYG{p}{(}\PYG{n}{x}\PYG{p}{)}\PYG{o}{.}\PYG{n}{all}\PYG{p}{(}\PYG{p}{)}\PYG{p}{)}\PYG{p}{,}
    \PYG{p}{\PYGZcb{}}
    \PYG{k}{if} \PYG{o+ow}{not} \PYG{n}{report}\PYG{p}{[}\PYG{l+s+s2}{\PYGZdq{}}\PYG{l+s+s2}{numerical\PYGZus{}solution\PYGZus{}finite}\PYG{l+s+s2}{\PYGZdq{}}\PYG{p}{]}\PYG{p}{:}
        \PYG{k}{return} \PYG{n}{report}
    \PYG{n}{compile\PYGZus{}seconds}\PYG{p}{,} \PYG{n}{cache\PYGZus{}hit} \PYG{o}{=} \PYG{n+nb+bp}{self}\PYG{o}{.}\PYG{n}{last\PYGZus{}matrix\PYGZus{}compile\PYGZus{}seconds}\PYG{p}{,} \PYG{n+nb+bp}{self}\PYG{o}{.}\PYG{n}{last\PYGZus{}matrix\PYGZus{}cache\PYGZus{}hit}
    \PYG{n}{matrix} \PYG{o}{=} \PYG{n+nb+bp}{self}\PYG{o}{.}\PYG{n}{constraint\PYGZus{}matrix}\PYG{p}{(}\PYG{p}{)}
    \PYG{n+nb+bp}{self}\PYG{o}{.}\PYG{n}{last\PYGZus{}matrix\PYGZus{}compile\PYGZus{}seconds}\PYG{p}{,} \PYG{n+nb+bp}{self}\PYG{o}{.}\PYG{n}{last\PYGZus{}matrix\PYGZus{}cache\PYGZus{}hit} \PYG{o}{=} \PYG{n}{compile\PYGZus{}seconds}\PYG{p}{,} \PYG{n}{cache\PYGZus{}hit}
    \PYG{c+c1}{\PYGZsh{} activity\CJKCode{为}Ax\CJKCode{；热平衡与储热递推等式也包含在这些行内。}}
    \PYG{n}{activity} \PYG{o}{=} \PYG{n}{np}\PYG{o}{.}\PYG{n}{asarray}\PYG{p}{(}\PYG{n}{matrix} \PYG{o}{@} \PYG{n}{x}\PYG{p}{)}
    \PYG{c+c1}{\PYGZsh{} \CJKCode{用各项绝对贡献估计量级，避免正负抵消使容差失真。}}
    \PYG{n}{magnitude} \PYG{o}{=} \PYG{n}{np}\PYG{o}{.}\PYG{n}{asarray}\PYG{p}{(}\PYG{n+nb}{abs}\PYG{p}{(}\PYG{n}{matrix}\PYG{p}{)} \PYG{o}{@} \PYG{n}{np}\PYG{o}{.}\PYG{n}{abs}\PYG{p}{(}\PYG{n}{x}\PYG{p}{)}\PYG{p}{)}
    \PYG{n}{row\PYGZus{}lo}\PYG{p}{,} \PYG{n}{row\PYGZus{}hi} \PYG{o}{=} \PYG{n}{np}\PYG{o}{.}\PYG{n}{asarray}\PYG{p}{(}\PYG{n+nb+bp}{self}\PYG{o}{.}\PYG{n}{row\PYGZus{}lower}\PYG{p}{)}\PYG{p}{,} \PYG{n}{np}\PYG{o}{.}\PYG{n}{asarray}\PYG{p}{(}\PYG{n+nb+bp}{self}\PYG{o}{.}\PYG{n}{row\PYGZus{}upper}\PYG{p}{)}
    \PYG{c+c1}{\PYGZsh{} \CJKCode{同时检查行下界与上界，等式并非仅看某一方向的残差。}}
    \PYG{n}{row\PYGZus{}violation} \PYG{o}{=} \PYG{n}{np}\PYG{o}{.}\PYG{n}{maximum}\PYG{p}{(}\PYG{l+m+mf}{0.0}\PYG{p}{,} \PYG{n}{np}\PYG{o}{.}\PYG{n}{maximum}\PYG{p}{(}\PYG{n}{row\PYGZus{}lo} \PYG{o}{\PYGZhy{}} \PYG{n}{activity}\PYG{p}{,} \PYG{n}{activity} \PYG{o}{\PYGZhy{}} \PYG{n}{row\PYGZus{}hi}\PYG{p}{)}\PYG{p}{)}
    \PYG{n}{row\PYGZus{}scaled} \PYG{o}{=} \PYG{n}{np}\PYG{o}{.}\PYG{n}{zeros\PYGZus{}like}\PYG{p}{(}\PYG{n}{activity}\PYG{p}{)}
    \PYG{k}{for} \PYG{n}{bound}\PYG{p}{,} \PYG{n}{violation} \PYG{o+ow}{in} \PYG{p}{(}\PYG{p}{(}\PYG{n}{row\PYGZus{}lo}\PYG{p}{,} \PYG{n}{np}\PYG{o}{.}\PYG{n}{maximum}\PYG{p}{(}\PYG{l+m+mf}{0.0}\PYG{p}{,} \PYG{n}{row\PYGZus{}lo} \PYG{o}{\PYGZhy{}} \PYG{n}{activity}\PYG{p}{)}\PYG{p}{)}\PYG{p}{,} \PYG{p}{(}\PYG{n}{row\PYGZus{}hi}\PYG{p}{,} \PYG{n}{np}\PYG{o}{.}\PYG{n}{maximum}\PYG{p}{(}\PYG{l+m+mf}{0.0}\PYG{p}{,} \PYG{n}{activity} \PYG{o}{\PYGZhy{}} \PYG{n}{row\PYGZus{}hi}\PYG{p}{)}\PYG{p}{)}\PYG{p}{)}\PYG{p}{:}
        \PYG{n}{scale} \PYG{o}{=} \PYG{n}{np}\PYG{o}{.}\PYG{n}{maximum}\PYG{p}{(}\PYG{n}{np}\PYG{o}{.}\PYG{n}{maximum}\PYG{p}{(}\PYG{l+m+mf}{1.0}\PYG{p}{,} \PYG{n}{magnitude}\PYG{p}{)}\PYG{p}{,} \PYG{n}{np}\PYG{o}{.}\PYG{n}{where}\PYG{p}{(}\PYG{n}{np}\PYG{o}{.}\PYG{n}{isfinite}\PYG{p}{(}\PYG{n}{bound}\PYG{p}{)}\PYG{p}{,} \PYG{n}{np}\PYG{o}{.}\PYG{n}{abs}\PYG{p}{(}\PYG{n}{bound}\PYG{p}{)}\PYG{p}{,} \PYG{l+m+mf}{0.0}\PYG{p}{)}\PYG{p}{)}
        \PYG{n}{row\PYGZus{}scaled} \PYG{o}{=} \PYG{n}{np}\PYG{o}{.}\PYG{n}{maximum}\PYG{p}{(}\PYG{n}{row\PYGZus{}scaled}\PYG{p}{,} \PYG{n}{violation} \PYG{o}{/} \PYG{p}{(}\PYG{n}{absolute\PYGZus{}tolerance} \PYG{o}{+} \PYG{n}{relative\PYGZus{}tolerance} \PYG{o}{*} \PYG{n}{scale}\PYG{p}{)}\PYG{p}{)}
    \PYG{c+c1}{\PYGZsh{} \CJKCode{变量上下界单独复核，防止行约束满足而容量或非负界被违反。}}
    \PYG{n}{lower}\PYG{p}{,} \PYG{n}{upper} \PYG{o}{=} \PYG{n}{np}\PYG{o}{.}\PYG{n}{asarray}\PYG{p}{(}\PYG{n+nb+bp}{self}\PYG{o}{.}\PYG{n}{lower}\PYG{p}{)}\PYG{p}{,} \PYG{n}{np}\PYG{o}{.}\PYG{n}{asarray}\PYG{p}{(}\PYG{n+nb+bp}{self}\PYG{o}{.}\PYG{n}{upper}\PYG{p}{)}
    \PYG{n}{bound\PYGZus{}violation} \PYG{o}{=} \PYG{n}{np}\PYG{o}{.}\PYG{n}{maximum}\PYG{p}{(}\PYG{l+m+mf}{0.0}\PYG{p}{,} \PYG{n}{np}\PYG{o}{.}\PYG{n}{maximum}\PYG{p}{(}\PYG{n}{lower} \PYG{o}{\PYGZhy{}} \PYG{n}{x}\PYG{p}{,} \PYG{n}{x} \PYG{o}{\PYGZhy{}} \PYG{n}{upper}\PYG{p}{)}\PYG{p}{)}
    \PYG{n}{bound\PYGZus{}scaled} \PYG{o}{=} \PYG{n}{np}\PYG{o}{.}\PYG{n}{zeros\PYGZus{}like}\PYG{p}{(}\PYG{n}{x}\PYG{p}{)}
    \PYG{k}{for} \PYG{n}{bound}\PYG{p}{,} \PYG{n}{violation} \PYG{o+ow}{in} \PYG{p}{(}\PYG{p}{(}\PYG{n}{lower}\PYG{p}{,} \PYG{n}{np}\PYG{o}{.}\PYG{n}{maximum}\PYG{p}{(}\PYG{l+m+mf}{0.0}\PYG{p}{,} \PYG{n}{lower} \PYG{o}{\PYGZhy{}} \PYG{n}{x}\PYG{p}{)}\PYG{p}{)}\PYG{p}{,} \PYG{p}{(}\PYG{n}{upper}\PYG{p}{,} \PYG{n}{np}\PYG{o}{.}\PYG{n}{maximum}\PYG{p}{(}\PYG{l+m+mf}{0.0}\PYG{p}{,} \PYG{n}{x} \PYG{o}{\PYGZhy{}} \PYG{n}{upper}\PYG{p}{)}\PYG{p}{)}\PYG{p}{)}\PYG{p}{:}
        \PYG{n}{scale} \PYG{o}{=} \PYG{n}{np}\PYG{o}{.}\PYG{n}{maximum}\PYG{p}{(}\PYG{n}{np}\PYG{o}{.}\PYG{n}{maximum}\PYG{p}{(}\PYG{l+m+mf}{1.0}\PYG{p}{,} \PYG{n}{np}\PYG{o}{.}\PYG{n}{abs}\PYG{p}{(}\PYG{n}{x}\PYG{p}{)}\PYG{p}{)}\PYG{p}{,} \PYG{n}{np}\PYG{o}{.}\PYG{n}{where}\PYG{p}{(}\PYG{n}{np}\PYG{o}{.}\PYG{n}{isfinite}\PYG{p}{(}\PYG{n}{bound}\PYG{p}{)}\PYG{p}{,} \PYG{n}{np}\PYG{o}{.}\PYG{n}{abs}\PYG{p}{(}\PYG{n}{bound}\PYG{p}{)}\PYG{p}{,} \PYG{l+m+mf}{0.0}\PYG{p}{)}\PYG{p}{)}
        \PYG{n}{bound\PYGZus{}scaled} \PYG{o}{=} \PYG{n}{np}\PYG{o}{.}\PYG{n}{maximum}\PYG{p}{(}\PYG{n}{bound\PYGZus{}scaled}\PYG{p}{,} \PYG{n}{violation} \PYG{o}{/} \PYG{p}{(}\PYG{n}{absolute\PYGZus{}tolerance} \PYG{o}{+} \PYG{n}{relative\PYGZus{}tolerance} \PYG{o}{*} \PYG{n}{scale}\PYG{p}{)}\PYG{p}{)}
    \PYG{c+c1}{\PYGZsh{} \CJKCode{整数变量与最近整数比较，不将近似浮点数直接当作合法离散决策。}}
    \PYG{n}{integer\PYGZus{}values} \PYG{o}{=} \PYG{n}{x}\PYG{p}{[}\PYG{n}{np}\PYG{o}{.}\PYG{n}{asarray}\PYG{p}{(}\PYG{n+nb+bp}{self}\PYG{o}{.}\PYG{n}{integrality}\PYG{p}{)} \PYG{o}{!=} \PYG{l+m+mi}{0}\PYG{p}{]}
    \PYG{n}{integer\PYGZus{}violation} \PYG{o}{=} \PYG{n+nb}{float}\PYG{p}{(}\PYG{n}{np}\PYG{o}{.}\PYG{n}{max}\PYG{p}{(}\PYG{n}{np}\PYG{o}{.}\PYG{n}{abs}\PYG{p}{(}\PYG{n}{integer\PYGZus{}values} \PYG{o}{\PYGZhy{}} \PYG{n}{np}\PYG{o}{.}\PYG{n}{rint}\PYG{p}{(}\PYG{n}{integer\PYGZus{}values}\PYG{p}{)}\PYG{p}{)}\PYG{p}{,} \PYG{n}{initial}\PYG{o}{=}\PYG{l+m+mf}{0.0}\PYG{p}{)}\PYG{p}{)}
    \PYG{c+c1}{\PYGZsh{} \CJKCode{尚无本次目标快照时，不伪造目标复算误差。}}
    \PYG{n}{objective\PYGZus{}error} \PYG{o}{=} \PYG{k+kc}{None}
    \PYG{n}{objective\PYGZus{}pass} \PYG{o}{=} \PYG{k+kc}{True}
    \PYG{c+c1}{\PYGZsh{} \CJKCode{按本次实际目标复算，适用于费用与排放目标切换。}}
    \PYG{k}{if} \PYG{n+nb+bp}{self}\PYG{o}{.}\PYG{n}{\PYGZus{}last\PYGZus{}solve\PYGZus{}objective} \PYG{o+ow}{is} \PYG{o+ow}{not} \PYG{k+kc}{None} \PYG{o+ow}{and} \PYG{n+nb+bp}{self}\PYG{o}{.}\PYG{n}{\PYGZus{}last\PYGZus{}solve\PYGZus{}fun} \PYG{o+ow}{is} \PYG{o+ow}{not} \PYG{k+kc}{None}\PYG{p}{:}
        \PYG{n}{recomputed} \PYG{o}{=} \PYG{n+nb}{float}\PYG{p}{(}\PYG{n+nb+bp}{self}\PYG{o}{.}\PYG{n}{\PYGZus{}last\PYGZus{}solve\PYGZus{}objective} \PYG{o}{@} \PYG{n}{x}\PYG{p}{)}
        \PYG{n}{objective\PYGZus{}error} \PYG{o}{=} \PYG{n+nb}{abs}\PYG{p}{(}\PYG{n}{recomputed} \PYG{o}{\PYGZhy{}} \PYG{n+nb+bp}{self}\PYG{o}{.}\PYG{n}{\PYGZus{}last\PYGZus{}solve\PYGZus{}fun}\PYG{p}{)}
        \PYG{c+c1}{\PYGZsh{} \CJKCode{费用与排放目标均按当前目标各项贡献确定比较量级。}}
        \PYG{n}{objective\PYGZus{}scale} \PYG{o}{=} \PYG{n+nb}{max}\PYG{p}{(}\PYG{l+m+mf}{1.0}\PYG{p}{,} \PYG{n+nb}{float}\PYG{p}{(}\PYG{n}{np}\PYG{o}{.}\PYG{n}{abs}\PYG{p}{(}\PYG{n+nb+bp}{self}\PYG{o}{.}\PYG{n}{\PYGZus{}last\PYGZus{}solve\PYGZus{}objective} \PYG{o}{*} \PYG{n}{x}\PYG{p}{)}\PYG{o}{.}\PYG{n}{sum}\PYG{p}{(}\PYG{p}{)}\PYG{p}{)}\PYG{p}{,} \PYG{n+nb}{abs}\PYG{p}{(}\PYG{n+nb+bp}{self}\PYG{o}{.}\PYG{n}{\PYGZus{}last\PYGZus{}solve\PYGZus{}fun}\PYG{p}{)}\PYG{p}{)}
        \PYG{n}{objective\PYGZus{}pass} \PYG{o}{=} \PYG{n+nb}{bool}\PYG{p}{(}\PYG{n}{objective\PYGZus{}error} \PYG{o}{\PYGZlt{}}\PYG{o}{=} \PYG{n}{absolute\PYGZus{}tolerance} \PYG{o}{+} \PYG{n}{relative\PYGZus{}tolerance} \PYG{o}{*} \PYG{n}{objective\PYGZus{}scale}\PYG{p}{)}
    \PYG{n}{report}\PYG{o}{.}\PYG{n}{update}\PYG{p}{(}\PYG{p}{\PYGZob{}}
        \PYG{l+s+s2}{\PYGZdq{}}\PYG{l+s+s2}{numerical\PYGZus{}constraint\PYGZus{}max\PYGZus{}violation}\PYG{l+s+s2}{\PYGZdq{}}\PYG{p}{:} \PYG{n+nb}{float}\PYG{p}{(}\PYG{n}{np}\PYG{o}{.}\PYG{n}{max}\PYG{p}{(}\PYG{n}{row\PYGZus{}violation}\PYG{p}{,} \PYG{n}{initial}\PYG{o}{=}\PYG{l+m+mf}{0.0}\PYG{p}{)}\PYG{p}{)}\PYG{p}{,}
        \PYG{l+s+s2}{\PYGZdq{}}\PYG{l+s+s2}{numerical\PYGZus{}constraint\PYGZus{}max\PYGZus{}scaled\PYGZus{}violation}\PYG{l+s+s2}{\PYGZdq{}}\PYG{p}{:} \PYG{n+nb}{float}\PYG{p}{(}\PYG{n}{np}\PYG{o}{.}\PYG{n}{max}\PYG{p}{(}\PYG{n}{row\PYGZus{}scaled}\PYG{p}{,} \PYG{n}{initial}\PYG{o}{=}\PYG{l+m+mf}{0.0}\PYG{p}{)}\PYG{p}{)}\PYG{p}{,}
        \PYG{l+s+s2}{\PYGZdq{}}\PYG{l+s+s2}{numerical\PYGZus{}constraint\PYGZus{}worst\PYGZus{}row}\PYG{l+s+s2}{\PYGZdq{}}\PYG{p}{:} \PYG{n+nb}{int}\PYG{p}{(}\PYG{n}{np}\PYG{o}{.}\PYG{n}{argmax}\PYG{p}{(}\PYG{n}{row\PYGZus{}scaled}\PYG{p}{)}\PYG{p}{)} \PYG{k}{if} \PYG{n+nb}{len}\PYG{p}{(}\PYG{n}{row\PYGZus{}scaled}\PYG{p}{)} \PYG{k}{else} \PYG{k+kc}{None}\PYG{p}{,}
        \PYG{l+s+s2}{\PYGZdq{}}\PYG{l+s+s2}{numerical\PYGZus{}bound\PYGZus{}max\PYGZus{}violation}\PYG{l+s+s2}{\PYGZdq{}}\PYG{p}{:} \PYG{n+nb}{float}\PYG{p}{(}\PYG{n}{np}\PYG{o}{.}\PYG{n}{max}\PYG{p}{(}\PYG{n}{bound\PYGZus{}violation}\PYG{p}{,} \PYG{n}{initial}\PYG{o}{=}\PYG{l+m+mf}{0.0}\PYG{p}{)}\PYG{p}{)}\PYG{p}{,}
        \PYG{l+s+s2}{\PYGZdq{}}\PYG{l+s+s2}{numerical\PYGZus{}bound\PYGZus{}max\PYGZus{}scaled\PYGZus{}violation}\PYG{l+s+s2}{\PYGZdq{}}\PYG{p}{:} \PYG{n+nb}{float}\PYG{p}{(}\PYG{n}{np}\PYG{o}{.}\PYG{n}{max}\PYG{p}{(}\PYG{n}{bound\PYGZus{}scaled}\PYG{p}{,} \PYG{n}{initial}\PYG{o}{=}\PYG{l+m+mf}{0.0}\PYG{p}{)}\PYG{p}{)}\PYG{p}{,}
        \PYG{l+s+s2}{\PYGZdq{}}\PYG{l+s+s2}{numerical\PYGZus{}bound\PYGZus{}worst\PYGZus{}variable}\PYG{l+s+s2}{\PYGZdq{}}\PYG{p}{:} \PYG{n+nb+bp}{self}\PYG{o}{.}\PYG{n}{names}\PYG{p}{[}\PYG{n+nb}{int}\PYG{p}{(}\PYG{n}{np}\PYG{o}{.}\PYG{n}{argmax}\PYG{p}{(}\PYG{n}{bound\PYGZus{}scaled}\PYG{p}{)}\PYG{p}{)}\PYG{p}{]} \PYG{k}{if} \PYG{n+nb}{len}\PYG{p}{(}\PYG{n}{x}\PYG{p}{)} \PYG{k}{else} \PYG{k+kc}{None}\PYG{p}{,}
        \PYG{l+s+s2}{\PYGZdq{}}\PYG{l+s+s2}{numerical\PYGZus{}integrality\PYGZus{}max\PYGZus{}violation}\PYG{l+s+s2}{\PYGZdq{}}\PYG{p}{:} \PYG{n}{integer\PYGZus{}violation}\PYG{p}{,}
        \PYG{l+s+s2}{\PYGZdq{}}\PYG{l+s+s2}{numerical\PYGZus{}objective\PYGZus{}recompute\PYGZus{}error}\PYG{l+s+s2}{\PYGZdq{}}\PYG{p}{:} \PYG{n}{objective\PYGZus{}error}\PYG{p}{,}
        \PYG{l+s+s2}{\PYGZdq{}}\PYG{l+s+s2}{numerical\PYGZus{}objective\PYGZus{}checked}\PYG{l+s+s2}{\PYGZdq{}}\PYG{p}{:} \PYG{n}{objective\PYGZus{}error} \PYG{o+ow}{is} \PYG{o+ow}{not} \PYG{k+kc}{None}\PYG{p}{,}
        \PYG{c+c1}{\PYGZsh{} \CJKCode{所有数值门槛共同决定通过，}scaled violation\CJKCode{不超过}1\CJKCode{才满足对应容差。}}
        \PYG{l+s+s2}{\PYGZdq{}}\PYG{l+s+s2}{numerical\PYGZus{}audit\PYGZus{}passed}\PYG{l+s+s2}{\PYGZdq{}}\PYG{p}{:} \PYG{n+nb}{bool}\PYG{p}{(}
            \PYG{n}{np}\PYG{o}{.}\PYG{n}{isfinite}\PYG{p}{(}\PYG{n}{activity}\PYG{p}{)}\PYG{o}{.}\PYG{n}{all}\PYG{p}{(}\PYG{p}{)} \PYG{o+ow}{and} \PYG{n}{np}\PYG{o}{.}\PYG{n}{max}\PYG{p}{(}\PYG{n}{row\PYGZus{}scaled}\PYG{p}{,} \PYG{n}{initial}\PYG{o}{=}\PYG{l+m+mf}{0.0}\PYG{p}{)} \PYG{o}{\PYGZlt{}}\PYG{o}{=} \PYG{l+m+mf}{1.0}
            \PYG{o+ow}{and} \PYG{n}{np}\PYG{o}{.}\PYG{n}{max}\PYG{p}{(}\PYG{n}{bound\PYGZus{}scaled}\PYG{p}{,} \PYG{n}{initial}\PYG{o}{=}\PYG{l+m+mf}{0.0}\PYG{p}{)} \PYG{o}{\PYGZlt{}}\PYG{o}{=} \PYG{l+m+mf}{1.0}
            \PYG{o+ow}{and} \PYG{n}{integer\PYGZus{}violation} \PYG{o}{\PYGZlt{}}\PYG{o}{=} \PYG{n}{absolute\PYGZus{}tolerance} \PYG{o+ow}{and} \PYG{n}{objective\PYGZus{}pass}
        \PYG{p}{)}\PYG{p}{,}
    \PYG{p}{\PYGZcb{}}\PYG{p}{)}
    \PYG{k}{return} \PYG{n}{report}
\end{Verbatim}

\endgroup

\noindent\textbf{结果核对。}检查\path{numerical_audit_passed}时同时保留最大缩放残差及最差行号；缩放残差不大于1表示在登记容差内，不能据此替代数据来源检查。

\Needspace{12\baselineskip}
\subsection*{7．完整投资冻结与费用账本}
完整冻结函数。容量投资转入存量，离散改造保存固定取值，并关闭后续新增或退役自由度。

\noindent\textbf{输入：}原配置、规划投资表及规划费用摘要；默认不放宽可靠性。\par
\noindent\textbf{输出：}深复制后的固定配置，包含所选资产年化费用账本及复算政策身份。\par

\noindent\textbf{实现思路。}规划与运行复算通过资产冻结函数衔接。函数复制规划配置，将所选新增容量转为运行期存量，并扣除已退役容量；管径、末端改造等离散选择保留为固定取值，同时关闭再次投资的自由度。冻结覆盖热源、储热、管网、建筑侧改造、配电和燃料库存等登记资产，费用账本随所选资产保留。

{\small 来源：\path{experiments.py}，原第13--147行；135行原代码、16行补充注释。}\par

\begingroup
\setstretch{1.0}
\begin{Verbatim}[commandchars=\\\{\},numbers=left,firstnumber=1,stepnumber=1,fontsize=\fontsize{9.5}{11.5}\selectfont,breaklines=true,breakanywhere=true,frame=single,framesep=3mm,numbersep=5pt,rulecolor=\color{black},xleftmargin=18pt]
\PYG{c+c1}{\PYGZsh{} \CJKCode{必须调用完整函数，不能仅冻结正文中示出的热源和储热两类。}}
\PYG{k}{def}\PYG{+w}{ }\PYG{n+nf}{freeze\PYGZus{}investment\PYGZus{}plan}\PYG{p}{(}
    \PYG{n}{config}\PYG{p}{:} \PYG{n+nb}{dict}\PYG{p}{,}
    \PYG{n}{nominal}\PYG{p}{,}
    \PYG{o}{*}\PYG{p}{,}
    \PYG{n}{relax\PYGZus{}reliability}\PYG{p}{:} \PYG{n+nb}{bool} \PYG{o}{=} \PYG{k+kc}{False}\PYG{p}{,}
\PYG{p}{)} \PYG{o}{\PYGZhy{}}\PYG{o}{\PYGZgt{}} \PYG{n+nb}{dict}\PYG{p}{:}
\PYG{+w}{    }\PYG{l+s+sd}{\PYGZdq{}\PYGZdq{}\PYGZdq{}Convert an optimized plan into fixed existing assets for replay.}

\PYG{l+s+sd}{    Reliability is retained by default. A diagnostic stress test may explicitly}
\PYG{l+s+sd}{    pass ``relax\PYGZus{}reliability=True`` to quantify shortage, but that relaxed run}
\PYG{l+s+sd}{    is not evidence of compliance with the original service requirement.}
\PYG{l+s+sd}{    \PYGZdq{}\PYGZdq{}\PYGZdq{}}
    \PYG{c+c1}{\PYGZsh{} \CJKCode{深复制保护原规划配置，复算修改不回写原数据。}}
    \PYG{n}{fixed} \PYG{o}{=} \PYG{n}{copy}\PYG{o}{.}\PYG{n}{deepcopy}\PYG{p}{(}\PYG{n}{config}\PYG{p}{)}
    \PYG{c+c1}{\PYGZsh{} \CJKCode{默认保留原服务可靠性；诊断性放宽会被单独标记。}}
    \PYG{n}{fixed}\PYG{p}{[}\PYG{l+s+s2}{\PYGZdq{}}\PYG{l+s+s2}{replay\PYGZus{}service\PYGZus{}policy}\PYG{l+s+s2}{\PYGZdq{}}\PYG{p}{]} \PYG{o}{=} \PYG{p}{(}
        \PYG{l+s+s2}{\PYGZdq{}}\PYG{l+s+s2}{diagnostic\PYGZus{}relaxed\PYGZus{}reliability}\PYG{l+s+s2}{\PYGZdq{}} \PYG{k}{if} \PYG{n}{relax\PYGZus{}reliability} \PYG{k}{else} \PYG{l+s+s2}{\PYGZdq{}}\PYG{l+s+s2}{original\PYGZus{}service\PYGZus{}requirements}\PYG{l+s+s2}{\PYGZdq{}}
    \PYG{p}{)}
    \PYG{c+c1}{\PYGZsh{} \CJKCode{冻结不抹去所选资产的年化费用；季节范围转换另由季节接口处理。}}
    \PYG{c+c1}{\PYGZsh{} A frozen asset has no *new* investment variable, but its selected annual}
    \PYG{c+c1}{\PYGZsh{} capital charge still belongs in annual replay and affordability tests.}
    \PYG{n}{planning\PYGZus{}summary} \PYG{o}{=} \PYG{n+nb}{getattr}\PYG{p}{(}\PYG{n}{nominal}\PYG{p}{,} \PYG{l+s+s2}{\PYGZdq{}}\PYG{l+s+s2}{summary}\PYG{l+s+s2}{\PYGZdq{}}\PYG{p}{,} \PYG{p}{\PYGZob{}}\PYG{p}{\PYGZcb{}}\PYG{p}{)}
    \PYG{n}{components} \PYG{o}{=} \PYG{n}{planning\PYGZus{}summary}\PYG{o}{.}\PYG{n}{get}\PYG{p}{(}\PYG{l+s+s2}{\PYGZdq{}}\PYG{l+s+s2}{components\PYGZus{}yuan\PYGZus{}per\PYGZus{}year}\PYG{l+s+s2}{\PYGZdq{}}\PYG{p}{,} \PYG{p}{\PYGZob{}}\PYG{p}{\PYGZcb{}}\PYG{p}{)}
    \PYG{n}{fixed}\PYG{p}{[}\PYG{l+s+s2}{\PYGZdq{}}\PYG{l+s+s2}{fixed\PYGZus{}investment\PYGZus{}accounting}\PYG{l+s+s2}{\PYGZdq{}}\PYG{p}{]} \PYG{o}{=} \PYG{p}{\PYGZob{}}
        \PYG{l+s+s2}{\PYGZdq{}}\PYG{l+s+s2}{scope}\PYG{l+s+s2}{\PYGZdq{}}\PYG{p}{:} \PYG{l+s+s2}{\PYGZdq{}}\PYG{l+s+s2}{annual\PYGZus{}selected\PYGZus{}plan\PYGZus{}charges}\PYG{l+s+s2}{\PYGZdq{}}\PYG{p}{,}
        \PYG{l+s+s2}{\PYGZdq{}}\PYG{l+s+s2}{components}\PYG{l+s+s2}{\PYGZdq{}}\PYG{p}{:} \PYG{p}{\PYGZob{}}\PYG{n}{name}\PYG{p}{:} \PYG{n+nb}{float}\PYG{p}{(}\PYG{n}{components}\PYG{o}{.}\PYG{n}{get}\PYG{p}{(}\PYG{n}{name}\PYG{p}{,} \PYG{l+m+mf}{0.0}\PYG{p}{)}\PYG{p}{)} \PYG{k}{for} \PYG{n}{name} \PYG{o+ow}{in} \PYG{p}{(}\PYG{l+s+s2}{\PYGZdq{}}\PYG{l+s+s2}{annualized\PYGZus{}capex}\PYG{l+s+s2}{\PYGZdq{}}\PYG{p}{,} \PYG{l+s+s2}{\PYGZdq{}}\PYG{l+s+s2}{fixed\PYGZus{}om\PYGZus{}cost}\PYG{l+s+s2}{\PYGZdq{}}\PYG{p}{)}\PYG{p}{\PYGZcb{}}\PYG{p}{,}
        \PYG{l+s+s2}{\PYGZdq{}}\PYG{l+s+s2}{metrics}\PYG{l+s+s2}{\PYGZdq{}}\PYG{p}{:} \PYG{p}{\PYGZob{}}\PYG{n}{name}\PYG{p}{:} \PYG{n+nb}{float}\PYG{p}{(}\PYG{n}{planning\PYGZus{}summary}\PYG{o}{.}\PYG{n}{get}\PYG{p}{(}\PYG{n}{name}\PYG{p}{,} \PYG{l+m+mf}{0.0}\PYG{p}{)}\PYG{p}{)} \PYG{k}{for} \PYG{n}{name} \PYG{o+ow}{in} \PYG{p}{(}
            \PYG{l+s+s2}{\PYGZdq{}}\PYG{l+s+s2}{regulated\PYGZus{}annualized\PYGZus{}capex\PYGZus{}yuan\PYGZus{}per\PYGZus{}year}\PYG{l+s+s2}{\PYGZdq{}}\PYG{p}{,} \PYG{l+s+s2}{\PYGZdq{}}\PYG{l+s+s2}{shared\PYGZus{}non\PYGZus{}heating\PYGZus{}capex\PYGZus{}yuan\PYGZus{}per\PYGZus{}year}\PYG{l+s+s2}{\PYGZdq{}}\PYG{p}{,} \PYG{l+s+s2}{\PYGZdq{}}\PYG{l+s+s2}{public\PYGZus{}subsidy\PYGZus{}yuan\PYGZus{}per\PYGZus{}year}\PYG{l+s+s2}{\PYGZdq{}}\PYG{p}{)}\PYG{p}{\PYGZcb{}}\PYG{p}{,}
    \PYG{p}{\PYGZcb{}}
    \PYG{k}{if} \PYG{o+ow}{not} \PYG{n+nb}{hasattr}\PYG{p}{(}\PYG{n}{nominal}\PYG{p}{,} \PYG{l+s+s2}{\PYGZdq{}}\PYG{l+s+s2}{summary}\PYG{l+s+s2}{\PYGZdq{}}\PYG{p}{)}\PYG{p}{:}
        \PYG{c+c1}{\PYGZsh{} Investment\PYGZhy{}only callers can still inspect installed assets. A model}
        \PYG{c+c1}{\PYGZsh{} replay must refuse these unknown charges instead of assuming zero.}
        \PYG{n}{fixed}\PYG{p}{[}\PYG{l+s+s2}{\PYGZdq{}}\PYG{l+s+s2}{fixed\PYGZus{}investment\PYGZus{}accounting}\PYG{l+s+s2}{\PYGZdq{}}\PYG{p}{]} \PYG{o}{=} \PYG{p}{\PYGZob{}}\PYG{l+s+s2}{\PYGZdq{}}\PYG{l+s+s2}{scope}\PYG{l+s+s2}{\PYGZdq{}}\PYG{p}{:} \PYG{l+s+s2}{\PYGZdq{}}\PYG{l+s+s2}{unavailable\PYGZus{}no\PYGZus{}planning\PYGZus{}summary}\PYG{l+s+s2}{\PYGZdq{}}\PYG{p}{\PYGZcb{}}
    \PYG{c+c1}{\PYGZsh{} \CJKCode{热源容量按既有加新增减退役转为存量，随后锁定两个投资上限。}}
    \PYG{k}{for} \PYG{n}{source} \PYG{o+ow}{in} \PYG{n}{fixed}\PYG{p}{[}\PYG{l+s+s2}{\PYGZdq{}}\PYG{l+s+s2}{sources}\PYG{l+s+s2}{\PYGZdq{}}\PYG{p}{]}\PYG{p}{:}
        \PYG{n}{row} \PYG{o}{=} \PYG{n}{nominal}\PYG{o}{.}\PYG{n}{investments}\PYG{p}{[}
            \PYG{p}{(}\PYG{n}{nominal}\PYG{o}{.}\PYG{n}{investments}\PYG{p}{[}\PYG{l+s+s2}{\PYGZdq{}}\PYG{l+s+s2}{category}\PYG{l+s+s2}{\PYGZdq{}}\PYG{p}{]} \PYG{o}{==} \PYG{l+s+s2}{\PYGZdq{}}\PYG{l+s+s2}{source\PYGZus{}capacity}\PYG{l+s+s2}{\PYGZdq{}}\PYG{p}{)}
            \PYG{o}{\PYGZam{}} \PYG{p}{(}\PYG{n}{nominal}\PYG{o}{.}\PYG{n}{investments}\PYG{p}{[}\PYG{l+s+s2}{\PYGZdq{}}\PYG{l+s+s2}{asset}\PYG{l+s+s2}{\PYGZdq{}}\PYG{p}{]} \PYG{o}{==} \PYG{n}{source}\PYG{p}{[}\PYG{l+s+s2}{\PYGZdq{}}\PYG{l+s+s2}{name}\PYG{l+s+s2}{\PYGZdq{}}\PYG{p}{]}\PYG{p}{)}
        \PYG{p}{]}
        \PYG{n}{addition} \PYG{o}{=} \PYG{n+nb}{float}\PYG{p}{(}\PYG{n}{row}\PYG{p}{[}\PYG{l+s+s2}{\PYGZdq{}}\PYG{l+s+s2}{value}\PYG{l+s+s2}{\PYGZdq{}}\PYG{p}{]}\PYG{o}{.}\PYG{n}{iloc}\PYG{p}{[}\PYG{l+m+mi}{0}\PYG{p}{]}\PYG{p}{)} \PYG{k}{if} \PYG{o+ow}{not} \PYG{n}{row}\PYG{o}{.}\PYG{n}{empty} \PYG{k}{else} \PYG{l+m+mf}{0.0}
        \PYG{c+c1}{\PYGZsh{} \CJKCode{退役量需要扣除，不可把已退役容量继续用于复算。}}
        \PYG{n}{retired} \PYG{o}{=} \PYG{n}{nominal}\PYG{o}{.}\PYG{n}{investments}\PYG{p}{[}
            \PYG{p}{(}\PYG{n}{nominal}\PYG{o}{.}\PYG{n}{investments}\PYG{p}{[}\PYG{l+s+s2}{\PYGZdq{}}\PYG{l+s+s2}{category}\PYG{l+s+s2}{\PYGZdq{}}\PYG{p}{]} \PYG{o}{==} \PYG{l+s+s2}{\PYGZdq{}}\PYG{l+s+s2}{source\PYGZus{}retirement}\PYG{l+s+s2}{\PYGZdq{}}\PYG{p}{)}
            \PYG{o}{\PYGZam{}} \PYG{p}{(}\PYG{n}{nominal}\PYG{o}{.}\PYG{n}{investments}\PYG{p}{[}\PYG{l+s+s2}{\PYGZdq{}}\PYG{l+s+s2}{asset}\PYG{l+s+s2}{\PYGZdq{}}\PYG{p}{]} \PYG{o}{==} \PYG{n}{source}\PYG{p}{[}\PYG{l+s+s2}{\PYGZdq{}}\PYG{l+s+s2}{name}\PYG{l+s+s2}{\PYGZdq{}}\PYG{p}{]}\PYG{p}{)}
        \PYG{p}{]}
        \PYG{n}{retirement} \PYG{o}{=} \PYG{n+nb}{float}\PYG{p}{(}\PYG{n}{retired}\PYG{p}{[}\PYG{l+s+s2}{\PYGZdq{}}\PYG{l+s+s2}{value}\PYG{l+s+s2}{\PYGZdq{}}\PYG{p}{]}\PYG{o}{.}\PYG{n}{iloc}\PYG{p}{[}\PYG{l+m+mi}{0}\PYG{p}{]}\PYG{p}{)} \PYG{k}{if} \PYG{o+ow}{not} \PYG{n}{retired}\PYG{o}{.}\PYG{n}{empty} \PYG{k}{else} \PYG{l+m+mf}{0.0}
        \PYG{n}{source}\PYG{p}{[}\PYG{l+s+s2}{\PYGZdq{}}\PYG{l+s+s2}{existing\PYGZus{}capacity\PYGZus{}kw}\PYG{l+s+s2}{\PYGZdq{}}\PYG{p}{]} \PYG{o}{=} \PYG{n+nb}{float}\PYG{p}{(}\PYG{n}{source}\PYG{o}{.}\PYG{n}{get}\PYG{p}{(}\PYG{l+s+s2}{\PYGZdq{}}\PYG{l+s+s2}{existing\PYGZus{}capacity\PYGZus{}kw}\PYG{l+s+s2}{\PYGZdq{}}\PYG{p}{,} \PYG{l+m+mf}{0.0}\PYG{p}{)}\PYG{p}{)} \PYG{o}{+} \PYG{n}{addition} \PYG{o}{\PYGZhy{}} \PYG{n}{retirement}
        \PYG{n}{source}\PYG{p}{[}\PYG{l+s+s2}{\PYGZdq{}}\PYG{l+s+s2}{max\PYGZus{}new\PYGZus{}capacity\PYGZus{}kw}\PYG{l+s+s2}{\PYGZdq{}}\PYG{p}{]} \PYG{o}{=} \PYG{l+m+mf}{0.0}
        \PYG{n}{source}\PYG{p}{[}\PYG{l+s+s2}{\PYGZdq{}}\PYG{l+s+s2}{max\PYGZus{}retire\PYGZus{}capacity\PYGZus{}kw}\PYG{l+s+s2}{\PYGZdq{}}\PYG{p}{]} \PYG{o}{=} \PYG{l+m+mf}{0.0}

    \PYG{c+c1}{\PYGZsh{} \CJKCode{储热新增能量容量并入}kWh\CJKCode{存量，后续不再增加储热容量。}}
    \PYG{k}{for} \PYG{n}{storage} \PYG{o+ow}{in} \PYG{n}{fixed}\PYG{p}{[}\PYG{l+s+s2}{\PYGZdq{}}\PYG{l+s+s2}{storages}\PYG{l+s+s2}{\PYGZdq{}}\PYG{p}{]}\PYG{p}{:}
        \PYG{n}{row} \PYG{o}{=} \PYG{n}{nominal}\PYG{o}{.}\PYG{n}{investments}\PYG{p}{[}
            \PYG{p}{(}\PYG{n}{nominal}\PYG{o}{.}\PYG{n}{investments}\PYG{p}{[}\PYG{l+s+s2}{\PYGZdq{}}\PYG{l+s+s2}{category}\PYG{l+s+s2}{\PYGZdq{}}\PYG{p}{]} \PYG{o}{==} \PYG{l+s+s2}{\PYGZdq{}}\PYG{l+s+s2}{storage\PYGZus{}capacity}\PYG{l+s+s2}{\PYGZdq{}}\PYG{p}{)}
            \PYG{o}{\PYGZam{}} \PYG{p}{(}\PYG{n}{nominal}\PYG{o}{.}\PYG{n}{investments}\PYG{p}{[}\PYG{l+s+s2}{\PYGZdq{}}\PYG{l+s+s2}{asset}\PYG{l+s+s2}{\PYGZdq{}}\PYG{p}{]} \PYG{o}{==} \PYG{n}{storage}\PYG{p}{[}\PYG{l+s+s2}{\PYGZdq{}}\PYG{l+s+s2}{name}\PYG{l+s+s2}{\PYGZdq{}}\PYG{p}{]}\PYG{p}{)}
        \PYG{p}{]}
        \PYG{n}{addition} \PYG{o}{=} \PYG{n+nb}{float}\PYG{p}{(}\PYG{n}{row}\PYG{p}{[}\PYG{l+s+s2}{\PYGZdq{}}\PYG{l+s+s2}{value}\PYG{l+s+s2}{\PYGZdq{}}\PYG{p}{]}\PYG{o}{.}\PYG{n}{iloc}\PYG{p}{[}\PYG{l+m+mi}{0}\PYG{p}{]}\PYG{p}{)} \PYG{k}{if} \PYG{o+ow}{not} \PYG{n}{row}\PYG{o}{.}\PYG{n}{empty} \PYG{k}{else} \PYG{l+m+mf}{0.0}
        \PYG{n}{storage}\PYG{p}{[}\PYG{l+s+s2}{\PYGZdq{}}\PYG{l+s+s2}{existing\PYGZus{}capacity\PYGZus{}kwh}\PYG{l+s+s2}{\PYGZdq{}}\PYG{p}{]} \PYG{o}{=} \PYG{n+nb}{float}\PYG{p}{(}\PYG{n}{storage}\PYG{o}{.}\PYG{n}{get}\PYG{p}{(}\PYG{l+s+s2}{\PYGZdq{}}\PYG{l+s+s2}{existing\PYGZus{}capacity\PYGZus{}kwh}\PYG{l+s+s2}{\PYGZdq{}}\PYG{p}{,} \PYG{l+m+mf}{0.0}\PYG{p}{)}\PYG{p}{)} \PYG{o}{+} \PYG{n}{addition}
        \PYG{n}{storage}\PYG{p}{[}\PYG{l+s+s2}{\PYGZdq{}}\PYG{l+s+s2}{max\PYGZus{}new\PYGZus{}capacity\PYGZus{}kwh}\PYG{l+s+s2}{\PYGZdq{}}\PYG{p}{]} \PYG{o}{=} \PYG{l+m+mf}{0.0}

    \PYG{c+c1}{\PYGZsh{} \CJKCode{管道改造以固定值保存；离散改造不必伪装成连续新增容量。}}
    \PYG{k}{for} \PYG{n}{pipe} \PYG{o+ow}{in} \PYG{n}{fixed}\PYG{p}{[}\PYG{l+s+s2}{\PYGZdq{}}\PYG{l+s+s2}{pipes}\PYG{l+s+s2}{\PYGZdq{}}\PYG{p}{]}\PYG{p}{:}
        \PYG{n}{retrofit\PYGZus{}row} \PYG{o}{=} \PYG{n}{nominal}\PYG{o}{.}\PYG{n}{investments}\PYG{p}{[}
            \PYG{p}{(}\PYG{n}{nominal}\PYG{o}{.}\PYG{n}{investments}\PYG{p}{[}\PYG{l+s+s2}{\PYGZdq{}}\PYG{l+s+s2}{category}\PYG{l+s+s2}{\PYGZdq{}}\PYG{p}{]} \PYG{o}{==} \PYG{l+s+s2}{\PYGZdq{}}\PYG{l+s+s2}{pipe\PYGZus{}retrofit}\PYG{l+s+s2}{\PYGZdq{}}\PYG{p}{)}
            \PYG{o}{\PYGZam{}} \PYG{p}{(}\PYG{n}{nominal}\PYG{o}{.}\PYG{n}{investments}\PYG{p}{[}\PYG{l+s+s2}{\PYGZdq{}}\PYG{l+s+s2}{asset}\PYG{l+s+s2}{\PYGZdq{}}\PYG{p}{]} \PYG{o}{==} \PYG{n}{pipe}\PYG{p}{[}\PYG{l+s+s2}{\PYGZdq{}}\PYG{l+s+s2}{name}\PYG{l+s+s2}{\PYGZdq{}}\PYG{p}{]}\PYG{p}{)}
        \PYG{p}{]}
        \PYG{k}{if} \PYG{n}{pipe}\PYG{o}{.}\PYG{n}{get}\PYG{p}{(}\PYG{l+s+s2}{\PYGZdq{}}\PYG{l+s+s2}{retrofit}\PYG{l+s+s2}{\PYGZdq{}}\PYG{p}{)}\PYG{p}{:}
            \PYG{n}{pipe}\PYG{p}{[}\PYG{l+s+s2}{\PYGZdq{}}\PYG{l+s+s2}{retrofit}\PYG{l+s+s2}{\PYGZdq{}}\PYG{p}{]}\PYG{p}{[}\PYG{l+s+s2}{\PYGZdq{}}\PYG{l+s+s2}{fixed\PYGZus{}value}\PYG{l+s+s2}{\PYGZdq{}}\PYG{p}{]} \PYG{o}{=} \PYG{p}{(}
                \PYG{n+nb}{float}\PYG{p}{(}\PYG{n}{retrofit\PYGZus{}row}\PYG{p}{[}\PYG{l+s+s2}{\PYGZdq{}}\PYG{l+s+s2}{value}\PYG{l+s+s2}{\PYGZdq{}}\PYG{p}{]}\PYG{o}{.}\PYG{n}{iloc}\PYG{p}{[}\PYG{l+m+mi}{0}\PYG{p}{]}\PYG{p}{)} \PYG{k}{if} \PYG{o+ow}{not} \PYG{n}{retrofit\PYGZus{}row}\PYG{o}{.}\PYG{n}{empty} \PYG{k}{else} \PYG{l+m+mf}{0.0}
            \PYG{p}{)}
            \PYG{n}{pipe}\PYG{p}{[}\PYG{l+s+s2}{\PYGZdq{}}\PYG{l+s+s2}{retrofit}\PYG{l+s+s2}{\PYGZdq{}}\PYG{p}{]}\PYG{p}{[}\PYG{l+s+s2}{\PYGZdq{}}\PYG{l+s+s2}{capex\PYGZus{}total\PYGZus{}yuan}\PYG{l+s+s2}{\PYGZdq{}}\PYG{p}{]} \PYG{o}{=} \PYG{l+m+mf}{0.0}
        \PYG{c+c1}{\PYGZsh{} \CJKCode{被选管径容量并入存量，再清空候选管径，防止复算重新选型。}}
        \PYG{n}{selected} \PYG{o}{=} \PYG{n}{nominal}\PYG{o}{.}\PYG{n}{investments}\PYG{p}{[}
            \PYG{p}{(}\PYG{n}{nominal}\PYG{o}{.}\PYG{n}{investments}\PYG{p}{[}\PYG{l+s+s2}{\PYGZdq{}}\PYG{l+s+s2}{category}\PYG{l+s+s2}{\PYGZdq{}}\PYG{p}{]} \PYG{o}{==} \PYG{l+s+s2}{\PYGZdq{}}\PYG{l+s+s2}{pipe\PYGZus{}choice}\PYG{l+s+s2}{\PYGZdq{}}\PYG{p}{)}
            \PYG{o}{\PYGZam{}} \PYG{p}{(}\PYG{n}{nominal}\PYG{o}{.}\PYG{n}{investments}\PYG{p}{[}\PYG{l+s+s2}{\PYGZdq{}}\PYG{l+s+s2}{asset}\PYG{l+s+s2}{\PYGZdq{}}\PYG{p}{]} \PYG{o}{==} \PYG{n}{pipe}\PYG{p}{[}\PYG{l+s+s2}{\PYGZdq{}}\PYG{l+s+s2}{name}\PYG{l+s+s2}{\PYGZdq{}}\PYG{p}{]}\PYG{p}{)}
            \PYG{o}{\PYGZam{}} \PYG{p}{(}\PYG{n}{nominal}\PYG{o}{.}\PYG{n}{investments}\PYG{p}{[}\PYG{l+s+s2}{\PYGZdq{}}\PYG{l+s+s2}{value}\PYG{l+s+s2}{\PYGZdq{}}\PYG{p}{]} \PYG{o}{\PYGZgt{}} \PYG{l+m+mf}{0.5}\PYG{p}{)}
        \PYG{p}{]}
        \PYG{k}{if} \PYG{o+ow}{not} \PYG{n}{selected}\PYG{o}{.}\PYG{n}{empty}\PYG{p}{:}
            \PYG{n}{option\PYGZus{}name} \PYG{o}{=} \PYG{n+nb}{str}\PYG{p}{(}\PYG{n}{selected}\PYG{p}{[}\PYG{l+s+s2}{\PYGZdq{}}\PYG{l+s+s2}{option}\PYG{l+s+s2}{\PYGZdq{}}\PYG{p}{]}\PYG{o}{.}\PYG{n}{iloc}\PYG{p}{[}\PYG{l+m+mi}{0}\PYG{p}{]}\PYG{p}{)}
            \PYG{n}{option} \PYG{o}{=} \PYG{n+nb}{next}\PYG{p}{(}\PYG{n}{x} \PYG{k}{for} \PYG{n}{x} \PYG{o+ow}{in} \PYG{n}{pipe}\PYG{o}{.}\PYG{n}{get}\PYG{p}{(}\PYG{l+s+s2}{\PYGZdq{}}\PYG{l+s+s2}{diameter\PYGZus{}options}\PYG{l+s+s2}{\PYGZdq{}}\PYG{p}{,} \PYG{p}{[}\PYG{p}{]}\PYG{p}{)} \PYG{k}{if} \PYG{n+nb}{str}\PYG{p}{(}\PYG{n}{x}\PYG{p}{[}\PYG{l+s+s2}{\PYGZdq{}}\PYG{l+s+s2}{name}\PYG{l+s+s2}{\PYGZdq{}}\PYG{p}{]}\PYG{p}{)} \PYG{o}{==} \PYG{n}{option\PYGZus{}name}\PYG{p}{)}
            \PYG{n}{pipe}\PYG{p}{[}\PYG{l+s+s2}{\PYGZdq{}}\PYG{l+s+s2}{existing\PYGZus{}capacity\PYGZus{}kw}\PYG{l+s+s2}{\PYGZdq{}}\PYG{p}{]} \PYG{o}{=} \PYG{n+nb}{float}\PYG{p}{(}\PYG{n}{pipe}\PYG{o}{.}\PYG{n}{get}\PYG{p}{(}\PYG{l+s+s2}{\PYGZdq{}}\PYG{l+s+s2}{existing\PYGZus{}capacity\PYGZus{}kw}\PYG{l+s+s2}{\PYGZdq{}}\PYG{p}{,} \PYG{l+m+mf}{0.0}\PYG{p}{)}\PYG{p}{)} \PYG{o}{+} \PYG{n+nb}{float}\PYG{p}{(}\PYG{n}{option}\PYG{p}{[}\PYG{l+s+s2}{\PYGZdq{}}\PYG{l+s+s2}{capacity\PYGZus{}kw}\PYG{l+s+s2}{\PYGZdq{}}\PYG{p}{]}\PYG{p}{)}
        \PYG{n}{pipe}\PYG{p}{[}\PYG{l+s+s2}{\PYGZdq{}}\PYG{l+s+s2}{diameter\PYGZus{}options}\PYG{l+s+s2}{\PYGZdq{}}\PYG{p}{]} \PYG{o}{=} \PYG{p}{[}\PYG{p}{]}

    \PYG{c+c1}{\PYGZsh{} \CJKCode{计量覆盖率并入已覆盖比例，关闭新增覆盖自由度。}}
    \PYG{k}{for} \PYG{n}{node} \PYG{o+ow}{in} \PYG{n}{fixed}\PYG{p}{[}\PYG{l+s+s2}{\PYGZdq{}}\PYG{l+s+s2}{nodes}\PYG{l+s+s2}{\PYGZdq{}}\PYG{p}{]}\PYG{p}{:}
        \PYG{n}{metering} \PYG{o}{=} \PYG{n}{node}\PYG{o}{.}\PYG{n}{get}\PYG{p}{(}\PYG{l+s+s2}{\PYGZdq{}}\PYG{l+s+s2}{heat\PYGZus{}metering}\PYG{l+s+s2}{\PYGZdq{}}\PYG{p}{)}
        \PYG{k}{if} \PYG{n}{metering}\PYG{p}{:}
            \PYG{n}{row} \PYG{o}{=} \PYG{n}{nominal}\PYG{o}{.}\PYG{n}{investments}\PYG{p}{[}
                \PYG{p}{(}\PYG{n}{nominal}\PYG{o}{.}\PYG{n}{investments}\PYG{p}{[}\PYG{l+s+s2}{\PYGZdq{}}\PYG{l+s+s2}{category}\PYG{l+s+s2}{\PYGZdq{}}\PYG{p}{]} \PYG{o}{==} \PYG{l+s+s2}{\PYGZdq{}}\PYG{l+s+s2}{heat\PYGZus{}metering\PYGZus{}coverage}\PYG{l+s+s2}{\PYGZdq{}}\PYG{p}{)}
                \PYG{o}{\PYGZam{}} \PYG{p}{(}\PYG{n}{nominal}\PYG{o}{.}\PYG{n}{investments}\PYG{p}{[}\PYG{l+s+s2}{\PYGZdq{}}\PYG{l+s+s2}{asset}\PYG{l+s+s2}{\PYGZdq{}}\PYG{p}{]} \PYG{o}{==} \PYG{n}{node}\PYG{p}{[}\PYG{l+s+s2}{\PYGZdq{}}\PYG{l+s+s2}{name}\PYG{l+s+s2}{\PYGZdq{}}\PYG{p}{]}\PYG{p}{)}
            \PYG{p}{]}
            \PYG{n}{addition} \PYG{o}{=} \PYG{n+nb}{float}\PYG{p}{(}\PYG{n}{row}\PYG{p}{[}\PYG{l+s+s2}{\PYGZdq{}}\PYG{l+s+s2}{value}\PYG{l+s+s2}{\PYGZdq{}}\PYG{p}{]}\PYG{o}{.}\PYG{n}{iloc}\PYG{p}{[}\PYG{l+m+mi}{0}\PYG{p}{]}\PYG{p}{)} \PYG{k}{if} \PYG{o+ow}{not} \PYG{n}{row}\PYG{o}{.}\PYG{n}{empty} \PYG{k}{else} \PYG{l+m+mf}{0.0}
            \PYG{n}{metering}\PYG{p}{[}\PYG{l+s+s2}{\PYGZdq{}}\PYG{l+s+s2}{existing\PYGZus{}coverage\PYGZus{}fraction}\PYG{l+s+s2}{\PYGZdq{}}\PYG{p}{]} \PYG{o}{=} \PYG{n+nb}{min}\PYG{p}{(}
                \PYG{l+m+mf}{1.0}\PYG{p}{,}
                \PYG{n+nb}{float}\PYG{p}{(}\PYG{n}{metering}\PYG{o}{.}\PYG{n}{get}\PYG{p}{(}\PYG{l+s+s2}{\PYGZdq{}}\PYG{l+s+s2}{existing\PYGZus{}coverage\PYGZus{}fraction}\PYG{l+s+s2}{\PYGZdq{}}\PYG{p}{,} \PYG{l+m+mf}{0.0}\PYG{p}{)}\PYG{p}{)}
                \PYG{o}{+} \PYG{n}{addition}\PYG{p}{,}
            \PYG{p}{)}
            \PYG{n}{metering}\PYG{p}{[}\PYG{l+s+s2}{\PYGZdq{}}\PYG{l+s+s2}{max\PYGZus{}new\PYGZus{}coverage\PYGZus{}fraction}\PYG{l+s+s2}{\PYGZdq{}}\PYG{p}{]} \PYG{o}{=} \PYG{l+m+mf}{0.0}
            \PYG{n}{metering}\PYG{o}{.}\PYG{n}{pop}\PYG{p}{(}\PYG{l+s+s2}{\PYGZdq{}}\PYG{l+s+s2}{fixed\PYGZus{}new\PYGZus{}coverage\PYGZus{}fraction}\PYG{l+s+s2}{\PYGZdq{}}\PYG{p}{,} \PYG{k+kc}{None}\PYG{p}{)}
            \PYG{n}{metering}\PYG{p}{[}\PYG{l+s+s2}{\PYGZdq{}}\PYG{l+s+s2}{capex\PYGZus{}yuan\PYGZus{}per\PYGZus{}coverage\PYGZus{}fraction}\PYG{l+s+s2}{\PYGZdq{}}\PYG{p}{]} \PYG{o}{=} \PYG{l+m+mf}{0.0}
        \PYG{c+c1}{\PYGZsh{} \CJKCode{末端温位改造保存规划取值，避免逐时复算改变投资模式。}}
        \PYG{n}{retrofit} \PYG{o}{=} \PYG{n}{node}\PYG{o}{.}\PYG{n}{get}\PYG{p}{(}\PYG{l+s+s2}{\PYGZdq{}}\PYG{l+s+s2}{temperature\PYGZus{}retrofit}\PYG{l+s+s2}{\PYGZdq{}}\PYG{p}{)}
        \PYG{k}{if} \PYG{o+ow}{not} \PYG{n}{retrofit}\PYG{p}{:}
            \PYG{k}{continue}
        \PYG{n}{row} \PYG{o}{=} \PYG{n}{nominal}\PYG{o}{.}\PYG{n}{investments}\PYG{p}{[}
            \PYG{p}{(}\PYG{n}{nominal}\PYG{o}{.}\PYG{n}{investments}\PYG{p}{[}\PYG{l+s+s2}{\PYGZdq{}}\PYG{l+s+s2}{category}\PYG{l+s+s2}{\PYGZdq{}}\PYG{p}{]} \PYG{o}{==} \PYG{l+s+s2}{\PYGZdq{}}\PYG{l+s+s2}{terminal\PYGZus{}temperature\PYGZus{}retrofit}\PYG{l+s+s2}{\PYGZdq{}}\PYG{p}{)}
            \PYG{o}{\PYGZam{}} \PYG{p}{(}\PYG{n}{nominal}\PYG{o}{.}\PYG{n}{investments}\PYG{p}{[}\PYG{l+s+s2}{\PYGZdq{}}\PYG{l+s+s2}{asset}\PYG{l+s+s2}{\PYGZdq{}}\PYG{p}{]} \PYG{o}{==} \PYG{n}{node}\PYG{p}{[}\PYG{l+s+s2}{\PYGZdq{}}\PYG{l+s+s2}{name}\PYG{l+s+s2}{\PYGZdq{}}\PYG{p}{]}\PYG{p}{)}
        \PYG{p}{]}
        \PYG{n}{retrofit}\PYG{p}{[}\PYG{l+s+s2}{\PYGZdq{}}\PYG{l+s+s2}{fixed\PYGZus{}value}\PYG{l+s+s2}{\PYGZdq{}}\PYG{p}{]} \PYG{o}{=} \PYG{n+nb}{float}\PYG{p}{(}\PYG{n}{row}\PYG{p}{[}\PYG{l+s+s2}{\PYGZdq{}}\PYG{l+s+s2}{value}\PYG{l+s+s2}{\PYGZdq{}}\PYG{p}{]}\PYG{o}{.}\PYG{n}{iloc}\PYG{p}{[}\PYG{l+m+mi}{0}\PYG{p}{]}\PYG{p}{)} \PYG{k}{if} \PYG{o+ow}{not} \PYG{n}{row}\PYG{o}{.}\PYG{n}{empty} \PYG{k}{else} \PYG{l+m+mf}{0.0}
        \PYG{n}{retrofit}\PYG{p}{[}\PYG{l+s+s2}{\PYGZdq{}}\PYG{l+s+s2}{capex\PYGZus{}total\PYGZus{}yuan}\PYG{l+s+s2}{\PYGZdq{}}\PYG{p}{]} \PYG{o}{=} \PYG{l+m+mf}{0.0}

    \PYG{c+c1}{\PYGZsh{} \CJKCode{电网接入增容同样冻结，不能靠新增配电容量通过复算。}}
    \PYG{n}{grid} \PYG{o}{=} \PYG{n}{fixed}\PYG{o}{.}\PYG{n}{get}\PYG{p}{(}\PYG{l+s+s2}{\PYGZdq{}}\PYG{l+s+s2}{grid\PYGZus{}connection}\PYG{l+s+s2}{\PYGZdq{}}\PYG{p}{)}
    \PYG{k}{if} \PYG{n}{grid}\PYG{p}{:}
        \PYG{n}{row} \PYG{o}{=} \PYG{n}{nominal}\PYG{o}{.}\PYG{n}{investments}\PYG{p}{[}\PYG{n}{nominal}\PYG{o}{.}\PYG{n}{investments}\PYG{p}{[}\PYG{l+s+s2}{\PYGZdq{}}\PYG{l+s+s2}{category}\PYG{l+s+s2}{\PYGZdq{}}\PYG{p}{]} \PYG{o}{==} \PYG{l+s+s2}{\PYGZdq{}}\PYG{l+s+s2}{grid\PYGZus{}connection\PYGZus{}capacity}\PYG{l+s+s2}{\PYGZdq{}}\PYG{p}{]}
        \PYG{n}{addition} \PYG{o}{=} \PYG{n+nb}{float}\PYG{p}{(}\PYG{n}{row}\PYG{p}{[}\PYG{l+s+s2}{\PYGZdq{}}\PYG{l+s+s2}{value}\PYG{l+s+s2}{\PYGZdq{}}\PYG{p}{]}\PYG{o}{.}\PYG{n}{iloc}\PYG{p}{[}\PYG{l+m+mi}{0}\PYG{p}{]}\PYG{p}{)} \PYG{k}{if} \PYG{o+ow}{not} \PYG{n}{row}\PYG{o}{.}\PYG{n}{empty} \PYG{k}{else} \PYG{l+m+mf}{0.0}
        \PYG{n}{grid}\PYG{p}{[}\PYG{l+s+s2}{\PYGZdq{}}\PYG{l+s+s2}{existing\PYGZus{}capacity\PYGZus{}kw}\PYG{l+s+s2}{\PYGZdq{}}\PYG{p}{]} \PYG{o}{=} \PYG{n+nb}{float}\PYG{p}{(}\PYG{n}{grid}\PYG{o}{.}\PYG{n}{get}\PYG{p}{(}\PYG{l+s+s2}{\PYGZdq{}}\PYG{l+s+s2}{existing\PYGZus{}capacity\PYGZus{}kw}\PYG{l+s+s2}{\PYGZdq{}}\PYG{p}{,} \PYG{l+m+mf}{0.0}\PYG{p}{)}\PYG{p}{)} \PYG{o}{+} \PYG{n}{addition}
        \PYG{n}{grid}\PYG{p}{[}\PYG{l+s+s2}{\PYGZdq{}}\PYG{l+s+s2}{max\PYGZus{}new\PYGZus{}capacity\PYGZus{}kw}\PYG{l+s+s2}{\PYGZdq{}}\PYG{p}{]} \PYG{o}{=} \PYG{l+m+mf}{0.0}
    \PYG{c+c1}{\PYGZsh{} \CJKCode{生物质库存容量以}kg\CJKCode{登记；库存容量不同于逐小时到货速率。}}
    \PYG{k}{for} \PYG{n}{chain} \PYG{o+ow}{in} \PYG{n}{fixed}\PYG{o}{.}\PYG{n}{get}\PYG{p}{(}\PYG{l+s+s2}{\PYGZdq{}}\PYG{l+s+s2}{biomass\PYGZus{}supply\PYGZus{}chains}\PYG{l+s+s2}{\PYGZdq{}}\PYG{p}{,} \PYG{p}{[}\PYG{p}{]}\PYG{p}{)}\PYG{p}{:}
        \PYG{n}{row} \PYG{o}{=} \PYG{n}{nominal}\PYG{o}{.}\PYG{n}{investments}\PYG{p}{[}
            \PYG{p}{(}\PYG{n}{nominal}\PYG{o}{.}\PYG{n}{investments}\PYG{p}{[}\PYG{l+s+s2}{\PYGZdq{}}\PYG{l+s+s2}{category}\PYG{l+s+s2}{\PYGZdq{}}\PYG{p}{]} \PYG{o}{==} \PYG{l+s+s2}{\PYGZdq{}}\PYG{l+s+s2}{biomass\PYGZus{}inventory\PYGZus{}capacity}\PYG{l+s+s2}{\PYGZdq{}}\PYG{p}{)}
            \PYG{o}{\PYGZam{}} \PYG{p}{(}\PYG{n}{nominal}\PYG{o}{.}\PYG{n}{investments}\PYG{p}{[}\PYG{l+s+s2}{\PYGZdq{}}\PYG{l+s+s2}{asset}\PYG{l+s+s2}{\PYGZdq{}}\PYG{p}{]} \PYG{o}{==} \PYG{n}{chain}\PYG{p}{[}\PYG{l+s+s2}{\PYGZdq{}}\PYG{l+s+s2}{name}\PYG{l+s+s2}{\PYGZdq{}}\PYG{p}{]}\PYG{p}{)}
        \PYG{p}{]}
        \PYG{n}{addition} \PYG{o}{=} \PYG{n+nb}{float}\PYG{p}{(}\PYG{n}{row}\PYG{p}{[}\PYG{l+s+s2}{\PYGZdq{}}\PYG{l+s+s2}{value}\PYG{l+s+s2}{\PYGZdq{}}\PYG{p}{]}\PYG{o}{.}\PYG{n}{iloc}\PYG{p}{[}\PYG{l+m+mi}{0}\PYG{p}{]}\PYG{p}{)} \PYG{k}{if} \PYG{o+ow}{not} \PYG{n}{row}\PYG{o}{.}\PYG{n}{empty} \PYG{k}{else} \PYG{l+m+mf}{0.0}
        \PYG{n}{chain}\PYG{p}{[}\PYG{l+s+s2}{\PYGZdq{}}\PYG{l+s+s2}{existing\PYGZus{}inventory\PYGZus{}capacity\PYGZus{}kg}\PYG{l+s+s2}{\PYGZdq{}}\PYG{p}{]} \PYG{o}{=} \PYG{n+nb}{float}\PYG{p}{(}
            \PYG{n}{chain}\PYG{o}{.}\PYG{n}{get}\PYG{p}{(}\PYG{l+s+s2}{\PYGZdq{}}\PYG{l+s+s2}{existing\PYGZus{}inventory\PYGZus{}capacity\PYGZus{}kg}\PYG{l+s+s2}{\PYGZdq{}}\PYG{p}{,} \PYG{l+m+mf}{0.0}\PYG{p}{)}
        \PYG{p}{)} \PYG{o}{+} \PYG{n}{addition}
        \PYG{n}{chain}\PYG{p}{[}\PYG{l+s+s2}{\PYGZdq{}}\PYG{l+s+s2}{max\PYGZus{}new\PYGZus{}inventory\PYGZus{}capacity\PYGZus{}kg}\PYG{l+s+s2}{\PYGZdq{}}\PYG{p}{]} \PYG{o}{=} \PYG{l+m+mf}{0.0}
    \PYG{c+c1}{\PYGZsh{} \CJKCode{可再生电源、储电功率与能量、备用电源均纳入冻结。}}
    \PYG{n}{microgrid} \PYG{o}{=} \PYG{n}{fixed}\PYG{o}{.}\PYG{n}{get}\PYG{p}{(}\PYG{l+s+s2}{\PYGZdq{}}\PYG{l+s+s2}{microgrid\PYGZus{}islanding}\PYG{l+s+s2}{\PYGZdq{}}\PYG{p}{)}
    \PYG{k}{if} \PYG{n}{microgrid}\PYG{p}{:}
        \PYG{n}{fields} \PYG{o}{=} \PYG{p}{\PYGZob{}}
            \PYG{l+s+s2}{\PYGZdq{}}\PYG{l+s+s2}{microgrid\PYGZus{}renewable\PYGZus{}capacity}\PYG{l+s+s2}{\PYGZdq{}}\PYG{p}{:} \PYG{p}{(}\PYG{l+s+s2}{\PYGZdq{}}\PYG{l+s+s2}{existing\PYGZus{}renewable\PYGZus{}capacity\PYGZus{}kw}\PYG{l+s+s2}{\PYGZdq{}}\PYG{p}{,} \PYG{l+s+s2}{\PYGZdq{}}\PYG{l+s+s2}{max\PYGZus{}new\PYGZus{}renewable\PYGZus{}capacity\PYGZus{}kw}\PYG{l+s+s2}{\PYGZdq{}}\PYG{p}{)}\PYG{p}{,}
            \PYG{l+s+s2}{\PYGZdq{}}\PYG{l+s+s2}{microgrid\PYGZus{}battery\PYGZus{}power}\PYG{l+s+s2}{\PYGZdq{}}\PYG{p}{:} \PYG{p}{(}\PYG{l+s+s2}{\PYGZdq{}}\PYG{l+s+s2}{existing\PYGZus{}battery\PYGZus{}power\PYGZus{}kw}\PYG{l+s+s2}{\PYGZdq{}}\PYG{p}{,} \PYG{l+s+s2}{\PYGZdq{}}\PYG{l+s+s2}{max\PYGZus{}new\PYGZus{}battery\PYGZus{}power\PYGZus{}kw}\PYG{l+s+s2}{\PYGZdq{}}\PYG{p}{)}\PYG{p}{,}
            \PYG{l+s+s2}{\PYGZdq{}}\PYG{l+s+s2}{microgrid\PYGZus{}battery\PYGZus{}energy}\PYG{l+s+s2}{\PYGZdq{}}\PYG{p}{:} \PYG{p}{(}\PYG{l+s+s2}{\PYGZdq{}}\PYG{l+s+s2}{existing\PYGZus{}battery\PYGZus{}energy\PYGZus{}kwh}\PYG{l+s+s2}{\PYGZdq{}}\PYG{p}{,} \PYG{l+s+s2}{\PYGZdq{}}\PYG{l+s+s2}{max\PYGZus{}new\PYGZus{}battery\PYGZus{}energy\PYGZus{}kwh}\PYG{l+s+s2}{\PYGZdq{}}\PYG{p}{)}\PYG{p}{,}
            \PYG{l+s+s2}{\PYGZdq{}}\PYG{l+s+s2}{microgrid\PYGZus{}backup\PYGZus{}capacity}\PYG{l+s+s2}{\PYGZdq{}}\PYG{p}{:} \PYG{p}{(}\PYG{l+s+s2}{\PYGZdq{}}\PYG{l+s+s2}{existing\PYGZus{}backup\PYGZus{}capacity\PYGZus{}kw}\PYG{l+s+s2}{\PYGZdq{}}\PYG{p}{,} \PYG{l+s+s2}{\PYGZdq{}}\PYG{l+s+s2}{max\PYGZus{}new\PYGZus{}backup\PYGZus{}capacity\PYGZus{}kw}\PYG{l+s+s2}{\PYGZdq{}}\PYG{p}{)}\PYG{p}{,}
        \PYG{p}{\PYGZcb{}}
        \PYG{k}{for} \PYG{n}{category}\PYG{p}{,} \PYG{p}{(}\PYG{n}{existing\PYGZus{}key}\PYG{p}{,} \PYG{n}{max\PYGZus{}key}\PYG{p}{)} \PYG{o+ow}{in} \PYG{n}{fields}\PYG{o}{.}\PYG{n}{items}\PYG{p}{(}\PYG{p}{)}\PYG{p}{:}
            \PYG{n}{row} \PYG{o}{=} \PYG{n}{nominal}\PYG{o}{.}\PYG{n}{investments}\PYG{p}{[}\PYG{n}{nominal}\PYG{o}{.}\PYG{n}{investments}\PYG{p}{[}\PYG{l+s+s2}{\PYGZdq{}}\PYG{l+s+s2}{category}\PYG{l+s+s2}{\PYGZdq{}}\PYG{p}{]} \PYG{o}{==} \PYG{n}{category}\PYG{p}{]}
            \PYG{n}{addition} \PYG{o}{=} \PYG{n+nb}{float}\PYG{p}{(}\PYG{n}{row}\PYG{p}{[}\PYG{l+s+s2}{\PYGZdq{}}\PYG{l+s+s2}{value}\PYG{l+s+s2}{\PYGZdq{}}\PYG{p}{]}\PYG{o}{.}\PYG{n}{iloc}\PYG{p}{[}\PYG{l+m+mi}{0}\PYG{p}{]}\PYG{p}{)} \PYG{k}{if} \PYG{o+ow}{not} \PYG{n}{row}\PYG{o}{.}\PYG{n}{empty} \PYG{k}{else} \PYG{l+m+mf}{0.0}
            \PYG{n}{microgrid}\PYG{p}{[}\PYG{n}{existing\PYGZus{}key}\PYG{p}{]} \PYG{o}{=} \PYG{n+nb}{float}\PYG{p}{(}\PYG{n}{microgrid}\PYG{o}{.}\PYG{n}{get}\PYG{p}{(}\PYG{n}{existing\PYGZus{}key}\PYG{p}{,} \PYG{l+m+mf}{0.0}\PYG{p}{)}\PYG{p}{)} \PYG{o}{+} \PYG{n}{addition}
            \PYG{n}{microgrid}\PYG{p}{[}\PYG{n}{max\PYGZus{}key}\PYG{p}{]} \PYG{o}{=} \PYG{l+m+mf}{0.0}
    \PYG{c+c1}{\PYGZsh{} \CJKCode{仅显式诊断请求可放宽失供率和备用裕度；该分支不构成原要求验收。}}
    \PYG{k}{if} \PYG{n}{relax\PYGZus{}reliability}\PYG{p}{:}
        \PYG{n}{fixed}\PYG{p}{[}\PYG{l+s+s2}{\PYGZdq{}}\PYG{l+s+s2}{planning}\PYG{l+s+s2}{\PYGZdq{}}\PYG{p}{]}\PYG{p}{[}\PYG{l+s+s2}{\PYGZdq{}}\PYG{l+s+s2}{max\PYGZus{}unserved\PYGZus{}rate}\PYG{l+s+s2}{\PYGZdq{}}\PYG{p}{]} \PYG{o}{=} \PYG{l+m+mf}{1.0}
        \PYG{n}{fixed}\PYG{p}{[}\PYG{l+s+s2}{\PYGZdq{}}\PYG{l+s+s2}{planning}\PYG{l+s+s2}{\PYGZdq{}}\PYG{p}{]}\PYG{p}{[}\PYG{l+s+s2}{\PYGZdq{}}\PYG{l+s+s2}{max\PYGZus{}scenario\PYGZus{}unserved\PYGZus{}rate}\PYG{l+s+s2}{\PYGZdq{}}\PYG{p}{]} \PYG{o}{=} \PYG{l+m+mf}{1.0}
        \PYG{n}{fixed}\PYG{p}{[}\PYG{l+s+s2}{\PYGZdq{}}\PYG{l+s+s2}{planning}\PYG{l+s+s2}{\PYGZdq{}}\PYG{p}{]}\PYG{p}{[}\PYG{l+s+s2}{\PYGZdq{}}\PYG{l+s+s2}{firm\PYGZus{}capacity\PYGZus{}reserve\PYGZus{}margin}\PYG{l+s+s2}{\PYGZdq{}}\PYG{p}{]} \PYG{o}{=} \PYG{l+m+mf}{0.0}
    \PYG{c+c1}{\PYGZsh{} \CJKCode{返回完整固定配置；是否冻结成功另由投资冻结审计核对。}}
    \PYG{k}{return} \PYG{n}{fixed}
\end{Verbatim}

\endgroup

\noindent\textbf{结果核对。}复算前同时核对资产清单与新增上限，不能只锁定热源。资本费用是否计入及是否重复计入，应结合保存的费用账本和服务期口径检查。

\Needspace{12\baselineskip}
\subsection*{8．固定投资的连续供暖季复算}
完整服务期复算函数，求规划方案或接收既有固定方案，保留小时顺序，并汇总季节范围检查。

\noindent\textbf{输入：}规划输入、连续时序、明确服务窗口、求解预算和可选已登记固定方案。\par
\noindent\textbf{输出：}季节运行解、储热检查、服务期证书及期间汇总；不生成年度认证。\par

\noindent\textbf{实现思路。}该函数先确定有时区的服务窗口，再准备规划方案或复用已冻结方案。完整时序作为运行输入，建筑与储热状态按其边界组传递。随后把资产冻结、供能能力、数值求解和储热校核合并为服务期证书，并额外加入窗口精确匹配检查，使“求解成功”和“覆盖声明供暖季”分别有可查记录。

{\small 来源：\path{heating_season.py}，原第183--246行；64行原代码、15行补充注释。}\par

\begingroup
\setstretch{1.0}
\begin{Verbatim}[commandchars=\\\{\},numbers=left,firstnumber=1,stepnumber=1,fontsize=\fontsize{9.5}{11.5}\selectfont,breaklines=true,breakanywhere=true,frame=single,framesep=3mm,numbersep=5pt,rulecolor=\color{black},xleftmargin=18pt]
\PYG{c+c1}{\PYGZsh{} \CJKCode{本接口检验声明供暖季，不等同于在全年窗口联合重求投资。}}
\PYG{k}{def}\PYG{+w}{ }\PYG{n+nf}{run\PYGZus{}fixed\PYGZus{}investment\PYGZus{}heating\PYGZus{}season\PYGZus{}replay}\PYG{p}{(}
    \PYG{n}{config}\PYG{p}{:} \PYG{n+nb}{dict}\PYG{p}{[}\PYG{n+nb}{str}\PYG{p}{,} \PYG{n}{Any}\PYG{p}{]}\PYG{p}{,} \PYG{n}{representative\PYGZus{}timeseries}\PYG{p}{:} \PYG{n}{pd}\PYG{o}{.}\PYG{n}{DataFrame}\PYG{p}{,}
    \PYG{n}{chronological\PYGZus{}source}\PYG{p}{:} \PYG{n}{pd}\PYG{o}{.}\PYG{n}{DataFrame}\PYG{p}{,} \PYG{n}{scope}\PYG{p}{:} \PYG{n+nb}{dict}\PYG{p}{[}\PYG{n+nb}{str}\PYG{p}{,} \PYG{n}{Any}\PYG{p}{]}\PYG{p}{,} \PYG{o}{*}\PYG{p}{,}
    \PYG{n}{timestamp\PYGZus{}column}\PYG{p}{:} \PYG{n+nb}{str} \PYG{o}{=} \PYG{l+s+s2}{\PYGZdq{}}\PYG{l+s+s2}{timestamp}\PYG{l+s+s2}{\PYGZdq{}}\PYG{p}{,} \PYG{n}{input\PYGZus{}timezone}\PYG{p}{:} \PYG{n+nb}{str} \PYG{o}{|} \PYG{k+kc}{None} \PYG{o}{=} \PYG{k+kc}{None}\PYG{p}{,}
    \PYG{n}{planning\PYGZus{}time\PYGZus{}limit}\PYG{p}{:} \PYG{n+nb}{float} \PYG{o}{=} \PYG{l+m+mf}{180.0}\PYG{p}{,} \PYG{n}{replay\PYGZus{}time\PYGZus{}limit}\PYG{p}{:} \PYG{n+nb}{float} \PYG{o}{=} \PYG{l+m+mf}{1800.0}\PYG{p}{,}
    \PYG{n}{planning\PYGZus{}mip\PYGZus{}gap}\PYG{p}{:} \PYG{n+nb}{float} \PYG{o}{|} \PYG{k+kc}{None} \PYG{o}{=} \PYG{k+kc}{None}\PYG{p}{,} \PYG{n}{replay\PYGZus{}mip\PYGZus{}gap}\PYG{p}{:} \PYG{n+nb}{float} \PYG{o}{|} \PYG{k+kc}{None} \PYG{o}{=} \PYG{k+kc}{None}\PYG{p}{,}
    \PYG{n}{allow\PYGZus{}cyclic\PYGZus{}state\PYGZus{}assumptions}\PYG{p}{:} \PYG{n+nb}{bool} \PYG{o}{=} \PYG{k+kc}{False}\PYG{p}{,}
    \PYG{n}{fixed\PYGZus{}plan}\PYG{p}{:} \PYG{n}{FixedInvestmentPlanArtifacts} \PYG{o}{|} \PYG{k+kc}{None} \PYG{o}{=} \PYG{k+kc}{None}\PYG{p}{,}
\PYG{p}{)} \PYG{o}{\PYGZhy{}}\PYG{o}{\PYGZgt{}} \PYG{n}{HeatingSeasonReplayArtifacts}\PYG{p}{:}
\PYG{+w}{    }\PYG{l+s+sd}{\PYGZdq{}\PYGZdq{}\PYGZdq{}Plan on the supplied planning horizon and validate only the stated season.\PYGZdq{}\PYGZdq{}\PYGZdq{}}
    \PYG{c+c1}{\PYGZsh{} \CJKCode{按有时区的左闭右开窗口切取时序，检查步长和完整覆盖。}}
    \PYG{n}{chronology}\PYG{p}{,} \PYG{n}{scope\PYGZus{}audit} \PYG{o}{=} \PYG{n}{prepare\PYGZus{}heating\PYGZus{}season\PYGZus{}timeseries}\PYG{p}{(}
        \PYG{n}{config}\PYG{p}{,} \PYG{n}{chronological\PYGZus{}source}\PYG{p}{,} \PYG{n}{scope}\PYG{p}{,} \PYG{n}{timestamp\PYGZus{}column}\PYG{o}{=}\PYG{n}{timestamp\PYGZus{}column}\PYG{p}{,} \PYG{n}{input\PYGZus{}timezone}\PYG{o}{=}\PYG{n}{input\PYGZus{}timezone}\PYG{p}{,}
    \PYG{p}{)}
    \PYG{c+c1}{\PYGZsh{} \CJKCode{不支持的状态交接假设在昂贵求解前拒绝。}}
    \PYG{c+c1}{\PYGZsh{} Reject unsupported state assumptions before a potentially expensive solve.}
    \PYG{n}{\PYGZus{}season\PYGZus{}config}\PYG{p}{(}\PYG{n}{config}\PYG{p}{,} \PYG{n}{allow\PYGZus{}cyclic\PYGZus{}state\PYGZus{}assumptions}\PYG{p}{)}
    \PYG{c+c1}{\PYGZsh{} \CJKCode{无既有固定方案时先规划并完整冻结；有方案时复用该方案。}}
    \PYG{n}{plan} \PYG{o}{=} \PYG{n}{fixed\PYGZus{}plan} \PYG{k}{if} \PYG{n}{fixed\PYGZus{}plan} \PYG{o+ow}{is} \PYG{o+ow}{not} \PYG{k+kc}{None} \PYG{k}{else} \PYG{n}{prepare\PYGZus{}fixed\PYGZus{}investment\PYGZus{}plan}\PYG{p}{(}
        \PYG{n}{config}\PYG{p}{,}\PYG{n}{representative\PYGZus{}timeseries}\PYG{p}{,}\PYG{n}{planning\PYGZus{}time\PYGZus{}limit}\PYG{o}{=}\PYG{n}{planning\PYGZus{}time\PYGZus{}limit}\PYG{p}{,}\PYG{n}{planning\PYGZus{}mip\PYGZus{}gap}\PYG{o}{=}\PYG{n}{planning\PYGZus{}mip\PYGZus{}gap}\PYG{p}{)}
    \PYG{c+c1}{\PYGZsh{} \CJKCode{季节转换保存假设；年度经济与资源约束另需年度复算。}}
    \PYG{n}{fixed\PYGZus{}config}\PYG{p}{,} \PYG{n}{assumptions} \PYG{o}{=} \PYG{n}{\PYGZus{}season\PYGZus{}config}\PYG{p}{(}\PYG{n}{plan}\PYG{o}{.}\PYG{n}{fixed\PYGZus{}config}\PYG{p}{,} \PYG{n}{allow\PYGZus{}cyclic\PYGZus{}state\PYGZus{}assumptions}\PYG{p}{)}
    \PYG{n}{plan} \PYG{o}{=} \PYG{n}{replace}\PYG{p}{(}\PYG{n}{plan}\PYG{p}{,} \PYG{n}{fixed\PYGZus{}config}\PYG{o}{=}\PYG{n}{fixed\PYGZus{}config}\PYG{p}{)}
    \PYG{n}{\PYGZus{}}\PYG{p}{,} \PYG{n}{scenario\PYGZus{}audit} \PYG{o}{=} \PYG{n}{adapt\PYGZus{}config\PYGZus{}to\PYGZus{}chronology\PYGZus{}scenarios}\PYG{p}{(}\PYG{n}{fixed\PYGZus{}config}\PYG{p}{,} \PYG{n}{chronology}\PYG{p}{)}
    \PYG{n}{preflight} \PYG{o}{=} \PYG{n}{AnnualReplayPreflight}\PYG{p}{(}
        \PYG{n}{chronological\PYGZus{}timeseries}\PYG{o}{=}\PYG{n}{chronology}\PYG{p}{,} \PYG{n}{chronology\PYGZus{}audit}\PYG{o}{=}\PYG{n}{scope\PYGZus{}audit}\PYG{p}{,}
        \PYG{n}{representative\PYGZus{}coverage\PYGZus{}audit}\PYG{o}{=}\PYG{n}{pd}\PYG{o}{.}\PYG{n}{DataFrame}\PYG{p}{(}\PYG{n}{columns}\PYG{o}{=}\PYG{p}{[}\PYG{l+s+s2}{\PYGZdq{}}\PYG{l+s+s2}{coverage\PYGZus{}status}\PYG{l+s+s2}{\PYGZdq{}}\PYG{p}{]}\PYG{p}{)}\PYG{p}{,}
        \PYG{n}{representative\PYGZus{}recommendations}\PYG{o}{=}\PYG{n}{pd}\PYG{o}{.}\PYG{n}{DataFrame}\PYG{p}{(}\PYG{p}{)}\PYG{p}{,} \PYG{n}{scenario\PYGZus{}scope\PYGZus{}audit}\PYG{o}{=}\PYG{n}{scenario\PYGZus{}audit}\PYG{p}{,}
    \PYG{p}{)}
    \PYG{c+c1}{\PYGZsh{} \CJKCode{运行模型使用已冻结配置，}require\PYGZus{}full\PYGZus{}year=False\CJKCode{仅标明季节范围。}}
    \PYG{n}{result} \PYG{o}{=} \PYG{n}{replay\PYGZus{}prepared\PYGZus{}fixed\PYGZus{}investment\PYGZus{}plan}\PYG{p}{(}
        \PYG{n}{config}\PYG{p}{,} \PYG{n}{representative\PYGZus{}timeseries}\PYG{p}{,} \PYG{n}{chronology}\PYG{p}{,} \PYG{n}{plan}\PYG{p}{,}
        \PYG{n}{timestamp\PYGZus{}column}\PYG{o}{=}\PYG{n}{timestamp\PYGZus{}column}\PYG{p}{,} \PYG{n}{replay\PYGZus{}time\PYGZus{}limit}\PYG{o}{=}\PYG{n}{replay\PYGZus{}time\PYGZus{}limit}\PYG{p}{,}
        \PYG{n}{replay\PYGZus{}mip\PYGZus{}gap}\PYG{o}{=}\PYG{n}{replay\PYGZus{}mip\PYGZus{}gap}\PYG{p}{,} \PYG{n}{require\PYGZus{}full\PYGZus{}year}\PYG{o}{=}\PYG{k+kc}{False}\PYG{p}{,} \PYG{n}{preflight}\PYG{o}{=}\PYG{n}{preflight}\PYG{p}{,}
    \PYG{p}{)}
    \PYG{c+c1}{\PYGZsh{} \CJKCode{不把年度规划总量与季节总量直接比较，也不按小时比例外推全年费用。}}
    \PYG{c+c1}{\PYGZsh{} Annual planning totals and seasonal totals are intentionally not compared.}
    \PYG{c+c1}{\PYGZsh{} \CJKCode{清空跨年度与季节的直接总量比较，避免口径错配。}}
    \PYG{n}{result}\PYG{o}{.}\PYG{n}{kpi\PYGZus{}comparison} \PYG{o}{=} \PYG{n}{pd}\PYG{o}{.}\PYG{n}{DataFrame}\PYG{p}{(}\PYG{p}{)}
    \PYG{c+c1}{\PYGZsh{} \CJKCode{独立复核储热状态、初末边界和对应时序。}}
    \PYG{n}{storage} \PYG{o}{=} \PYG{n}{seasonal\PYGZus{}storage\PYGZus{}postcheck}\PYG{p}{(}\PYG{n}{result}\PYG{o}{.}\PYG{n}{fixed\PYGZus{}config}\PYG{p}{,} \PYG{n}{result}\PYG{o}{.}\PYG{n}{replay}\PYG{o}{.}\PYG{n}{dispatch}\PYG{p}{,} \PYG{n}{expected\PYGZus{}timeseries}\PYG{o}{=}\PYG{n}{chronology}\PYG{p}{)}
    \PYG{c+c1}{\PYGZsh{} \CJKCode{共同验收检查结合冻结、可靠容量、求解与储热证据。}}
    \PYG{n}{common} \PYG{o}{=} \PYG{n}{annual\PYGZus{}replay\PYGZus{}acceptance\PYGZus{}audit}\PYG{p}{(}
        \PYG{n}{result}\PYG{o}{.}\PYG{n}{fixed\PYGZus{}config}\PYG{p}{,} \PYG{n}{chronology\PYGZus{}audit}\PYG{o}{=}\PYG{n}{scope\PYGZus{}audit}\PYG{p}{,}
        \PYG{n}{representative\PYGZus{}coverage\PYGZus{}audit}\PYG{o}{=}\PYG{n}{preflight}\PYG{o}{.}\PYG{n}{representative\PYGZus{}coverage\PYGZus{}audit}\PYG{p}{,}
        \PYG{n}{investment\PYGZus{}freeze\PYGZus{}audit}\PYG{o}{=}\PYG{n}{result}\PYG{o}{.}\PYG{n}{investment\PYGZus{}freeze\PYGZus{}audit}\PYG{p}{,}
        \PYG{n}{fixed\PYGZus{}investment\PYGZus{}readiness\PYGZus{}audit}\PYG{o}{=}\PYG{n}{result}\PYG{o}{.}\PYG{n}{fixed\PYGZus{}investment\PYGZus{}readiness\PYGZus{}audit}\PYG{p}{,}
        \PYG{n}{planning\PYGZus{}summary}\PYG{o}{=}\PYG{n}{result}\PYG{o}{.}\PYG{n}{planning}\PYG{o}{.}\PYG{n}{summary}\PYG{p}{,} \PYG{n}{replay\PYGZus{}summary}\PYG{o}{=}\PYG{n}{result}\PYG{o}{.}\PYG{n}{replay}\PYG{o}{.}\PYG{n}{summary}\PYG{p}{,}
        \PYG{n}{storage\PYGZus{}boundary\PYGZus{}audit}\PYG{o}{=}\PYG{n}{storage}\PYG{p}{,}
    \PYG{p}{)}
    \PYG{c+c1}{\PYGZsh{} \CJKCode{删除年度专属门槛；季节证书不能因这些门槛被删除而晋级年度证书。}}
    \PYG{n}{gates} \PYG{o}{=} \PYG{p}{[}\PYG{n}{gate} \PYG{k}{for} \PYG{n}{gate} \PYG{o+ow}{in} \PYG{n}{common}\PYG{p}{[}\PYG{l+s+s2}{\PYGZdq{}}\PYG{l+s+s2}{gates}\PYG{l+s+s2}{\PYGZdq{}}\PYG{p}{]} \PYG{k}{if} \PYG{n}{gate}\PYG{p}{[}\PYG{l+s+s2}{\PYGZdq{}}\PYG{l+s+s2}{gate}\PYG{l+s+s2}{\PYGZdq{}}\PYG{p}{]} \PYG{o+ow}{not} \PYG{o+ow}{in} \PYG{p}{\PYGZob{}}\PYG{l+s+s2}{\PYGZdq{}}\PYG{l+s+s2}{continuous\PYGZus{}full\PYGZus{}calendar\PYGZus{}year}\PYG{l+s+s2}{\PYGZdq{}}\PYG{p}{,} \PYG{l+s+s2}{\PYGZdq{}}\PYG{l+s+s2}{continuous\PYGZus{}full\PYGZus{}service\PYGZus{}year}\PYG{l+s+s2}{\PYGZdq{}}\PYG{p}{,} \PYG{l+s+s2}{\PYGZdq{}}\PYG{l+s+s2}{representative\PYGZus{}period\PYGZus{}coverage}\PYG{l+s+s2}{\PYGZdq{}}\PYG{p}{\PYGZcb{}}\PYG{p}{]}
    \PYG{k}{for} \PYG{n}{gate} \PYG{o+ow}{in} \PYG{n}{gates}\PYG{p}{:}
        \PYG{n}{gate}\PYG{p}{[}\PYG{l+s+s2}{\PYGZdq{}}\PYG{l+s+s2}{gate}\PYG{l+s+s2}{\PYGZdq{}}\PYG{p}{]} \PYG{o}{=} \PYG{n}{gate}\PYG{p}{[}\PYG{l+s+s2}{\PYGZdq{}}\PYG{l+s+s2}{gate}\PYG{l+s+s2}{\PYGZdq{}}\PYG{p}{]}\PYG{o}{.}\PYG{n}{replace}\PYG{p}{(}\PYG{l+s+s2}{\PYGZdq{}}\PYG{l+s+s2}{annual\PYGZus{}}\PYG{l+s+s2}{\PYGZdq{}}\PYG{p}{,} \PYG{l+s+s2}{\PYGZdq{}}\PYG{l+s+s2}{period\PYGZus{}}\PYG{l+s+s2}{\PYGZdq{}}\PYG{p}{)}
        \PYG{n}{gate}\PYG{p}{[}\PYG{l+s+s2}{\PYGZdq{}}\PYG{l+s+s2}{requirement}\PYG{l+s+s2}{\PYGZdq{}}\PYG{p}{]} \PYG{o}{=} \PYG{n}{gate}\PYG{p}{[}\PYG{l+s+s2}{\PYGZdq{}}\PYG{l+s+s2}{requirement}\PYG{l+s+s2}{\PYGZdq{}}\PYG{p}{]}\PYG{o}{.}\PYG{n}{replace}\PYG{p}{(}\PYG{l+s+s2}{\PYGZdq{}}\PYG{l+s+s2}{\CJKCode{全年}}\PYG{l+s+s2}{\PYGZdq{}}\PYG{p}{,} \PYG{l+s+s2}{\PYGZdq{}}\PYG{l+s+s2}{\CJKCode{本服务期}}\PYG{l+s+s2}{\PYGZdq{}}\PYG{p}{)}
    \PYG{c+c1}{\PYGZsh{} \CJKCode{另加服务窗口精确匹配门槛，防止短片段冒充完整供暖季。}}
    \PYG{n}{gates}\PYG{o}{.}\PYG{n}{insert}\PYG{p}{(}\PYG{l+m+mi}{0}\PYG{p}{,} \PYG{p}{\PYGZob{}}\PYG{l+s+s2}{\PYGZdq{}}\PYG{l+s+s2}{gate}\PYG{l+s+s2}{\PYGZdq{}}\PYG{p}{:} \PYG{l+s+s2}{\PYGZdq{}}\PYG{l+s+s2}{exact\PYGZus{}service\PYGZus{}period}\PYG{l+s+s2}{\PYGZdq{}}\PYG{p}{,} \PYG{l+s+s2}{\PYGZdq{}}\PYG{l+s+s2}{passed}\PYG{l+s+s2}{\PYGZdq{}}\PYG{p}{:} \PYG{n+nb}{bool}\PYG{p}{(}\PYG{n}{scope\PYGZus{}audit}\PYG{p}{[}\PYG{l+s+s2}{\PYGZdq{}}\PYG{l+s+s2}{exact\PYGZus{}scope\PYGZus{}match}\PYG{l+s+s2}{\PYGZdq{}}\PYG{p}{]}\PYG{p}{)}\PYG{p}{,}
                     \PYG{l+s+s2}{\PYGZdq{}}\PYG{l+s+s2}{observed}\PYG{l+s+s2}{\PYGZdq{}}\PYG{p}{:} \PYG{n}{scope\PYGZus{}audit}\PYG{p}{,} \PYG{l+s+s2}{\PYGZdq{}}\PYG{l+s+s2}{requirement}\PYG{l+s+s2}{\PYGZdq{}}\PYG{p}{:} \PYG{l+s+s2}{\PYGZdq{}}\PYG{l+s+s2}{exact timezone\PYGZhy{}aware half\PYGZhy{}open service window}\PYG{l+s+s2}{\PYGZdq{}}\PYG{p}{\PYGZcb{}}\PYG{p}{)}
    \PYG{c+c1}{\PYGZsh{} \CJKCode{保留全部失败门槛及假设，不把失败结果写成通过。}}
    \PYG{c+c1}{\PYGZsh{} \CJKCode{证书由全部门槛共同决定，保留未通过项目名称。}}
    \PYG{n}{failed} \PYG{o}{=} \PYG{p}{[}\PYG{n}{gate}\PYG{p}{[}\PYG{l+s+s2}{\PYGZdq{}}\PYG{l+s+s2}{gate}\PYG{l+s+s2}{\PYGZdq{}}\PYG{p}{]} \PYG{k}{for} \PYG{n}{gate} \PYG{o+ow}{in} \PYG{n}{gates} \PYG{k}{if} \PYG{o+ow}{not} \PYG{n}{gate}\PYG{p}{[}\PYG{l+s+s2}{\PYGZdq{}}\PYG{l+s+s2}{passed}\PYG{l+s+s2}{\PYGZdq{}}\PYG{p}{]}\PYG{p}{]}
    \PYG{n}{certificate} \PYG{o}{=} \PYG{p}{\PYGZob{}}\PYG{l+s+s2}{\PYGZdq{}}\PYG{l+s+s2}{schema}\PYG{l+s+s2}{\PYGZdq{}}\PYG{p}{:} \PYG{l+s+s2}{\PYGZdq{}}\PYG{l+s+s2}{heatplan.heating\PYGZus{}season\PYGZus{}certificate.v1}\PYG{l+s+s2}{\PYGZdq{}}\PYG{p}{,} \PYG{l+s+s2}{\PYGZdq{}}\PYG{l+s+s2}{scope\PYGZus{}type}\PYG{l+s+s2}{\PYGZdq{}}\PYG{p}{:} \PYG{l+s+s2}{\PYGZdq{}}\PYG{l+s+s2}{heating\PYGZus{}season}\PYG{l+s+s2}{\PYGZdq{}}\PYG{p}{,}
                   \PYG{l+s+s2}{\PYGZdq{}}\PYG{l+s+s2}{certification\PYGZus{}pass}\PYG{l+s+s2}{\PYGZdq{}}\PYG{p}{:} \PYG{o+ow}{not} \PYG{n}{failed}\PYG{p}{,} \PYG{l+s+s2}{\PYGZdq{}}\PYG{l+s+s2}{gate\PYGZus{}count}\PYG{l+s+s2}{\PYGZdq{}}\PYG{p}{:} \PYG{n+nb}{len}\PYG{p}{(}\PYG{n}{gates}\PYG{p}{)}\PYG{p}{,}
                   \PYG{l+s+s2}{\PYGZdq{}}\PYG{l+s+s2}{passed\PYGZus{}gate\PYGZus{}count}\PYG{l+s+s2}{\PYGZdq{}}\PYG{p}{:} \PYG{n+nb}{len}\PYG{p}{(}\PYG{n}{gates}\PYG{p}{)} \PYG{o}{\PYGZhy{}} \PYG{n+nb}{len}\PYG{p}{(}\PYG{n}{failed}\PYG{p}{)}\PYG{p}{,} \PYG{l+s+s2}{\PYGZdq{}}\PYG{l+s+s2}{failed\PYGZus{}gate\PYGZus{}count}\PYG{l+s+s2}{\PYGZdq{}}\PYG{p}{:} \PYG{n+nb}{len}\PYG{p}{(}\PYG{n}{failed}\PYG{p}{)}\PYG{p}{,}
                   \PYG{l+s+s2}{\PYGZdq{}}\PYG{l+s+s2}{failed\PYGZus{}gates}\PYG{l+s+s2}{\PYGZdq{}}\PYG{p}{:} \PYG{n}{failed}\PYG{p}{,} \PYG{l+s+s2}{\PYGZdq{}}\PYG{l+s+s2}{gates}\PYG{l+s+s2}{\PYGZdq{}}\PYG{p}{:} \PYG{n}{gates}\PYG{p}{,} \PYG{l+s+s2}{\PYGZdq{}}\PYG{l+s+s2}{assumptions}\PYG{l+s+s2}{\PYGZdq{}}\PYG{p}{:} \PYG{n}{assumptions}\PYG{p}{,}
                   \PYG{l+s+s2}{\PYGZdq{}}\PYG{l+s+s2}{annual\PYGZus{}replay\PYGZus{}certified}\PYG{l+s+s2}{\PYGZdq{}}\PYG{p}{:} \PYG{k+kc}{False}\PYG{p}{\PYGZcb{}}
    \PYG{n}{summary} \PYG{o}{=} \PYG{p}{\PYGZob{}}\PYG{l+s+s2}{\PYGZdq{}}\PYG{l+s+s2}{scope}\PYG{l+s+s2}{\PYGZdq{}}\PYG{p}{:} \PYG{n}{scope\PYGZus{}audit}\PYG{p}{,} \PYG{l+s+s2}{\PYGZdq{}}\PYG{l+s+s2}{heating\PYGZus{}season\PYGZus{}certified}\PYG{l+s+s2}{\PYGZdq{}}\PYG{p}{:} \PYG{o+ow}{not} \PYG{n}{failed}\PYG{p}{,}
               \PYG{l+s+s2}{\PYGZdq{}}\PYG{l+s+s2}{annual\PYGZus{}replay\PYGZus{}certified}\PYG{l+s+s2}{\PYGZdq{}}\PYG{p}{:} \PYG{k+kc}{False}\PYG{p}{,} \PYG{l+s+s2}{\PYGZdq{}}\PYG{l+s+s2}{assumptions}\PYG{l+s+s2}{\PYGZdq{}}\PYG{p}{:} \PYG{n}{assumptions}\PYG{p}{,}
               \PYG{l+s+s2}{\PYGZdq{}}\PYG{l+s+s2}{solver\PYGZus{}status}\PYG{l+s+s2}{\PYGZdq{}}\PYG{p}{:} \PYG{n}{result}\PYG{o}{.}\PYG{n}{replay}\PYG{o}{.}\PYG{n}{summary}\PYG{p}{[}\PYG{l+s+s2}{\PYGZdq{}}\PYG{l+s+s2}{solver\PYGZus{}status}\PYG{l+s+s2}{\PYGZdq{}}\PYG{p}{]}\PYG{p}{,}
               \PYG{l+s+s2}{\PYGZdq{}}\PYG{l+s+s2}{unserved\PYGZus{}rate}\PYG{l+s+s2}{\PYGZdq{}}\PYG{p}{:} \PYG{n}{result}\PYG{o}{.}\PYG{n}{replay}\PYG{o}{.}\PYG{n}{summary}\PYG{p}{[}\PYG{l+s+s2}{\PYGZdq{}}\PYG{l+s+s2}{unserved\PYGZus{}rate}\PYG{l+s+s2}{\PYGZdq{}}\PYG{p}{]}\PYG{p}{,}
               \PYG{l+s+s2}{\PYGZdq{}}\PYG{l+s+s2}{numerical\PYGZus{}audit\PYGZus{}passed}\PYG{l+s+s2}{\PYGZdq{}}\PYG{p}{:} \PYG{n}{result}\PYG{o}{.}\PYG{n}{replay}\PYG{o}{.}\PYG{n}{summary}\PYG{o}{.}\PYG{n}{get}\PYG{p}{(}\PYG{l+s+s2}{\PYGZdq{}}\PYG{l+s+s2}{numerical\PYGZus{}audit\PYGZus{}passed}\PYG{l+s+s2}{\PYGZdq{}}\PYG{p}{,} \PYG{k+kc}{False}\PYG{p}{)}\PYG{p}{\PYGZcb{}}
    \PYG{c+c1}{\PYGZsh{} \CJKCode{期间指标重新命名为}over\PYGZus{}period\CJKCode{，避免}per\PYGZus{}year\CJKCode{标签误导。}}
    \PYG{k}{for} \PYG{n}{key} \PYG{o+ow}{in} \PYG{p}{(}\PYG{l+s+s2}{\PYGZdq{}}\PYG{l+s+s2}{demand\PYGZus{}kwh\PYGZus{}per\PYGZus{}year}\PYG{l+s+s2}{\PYGZdq{}}\PYG{p}{,} \PYG{l+s+s2}{\PYGZdq{}}\PYG{l+s+s2}{effective\PYGZus{}heat\PYGZus{}demand\PYGZus{}kwh\PYGZus{}per\PYGZus{}year}\PYG{l+s+s2}{\PYGZdq{}}\PYG{p}{,} \PYG{l+s+s2}{\PYGZdq{}}\PYG{l+s+s2}{unserved\PYGZus{}kwh\PYGZus{}per\PYGZus{}year}\PYG{l+s+s2}{\PYGZdq{}}\PYG{p}{,}
                \PYG{l+s+s2}{\PYGZdq{}}\PYG{l+s+s2}{emissions\PYGZus{}kg\PYGZus{}per\PYGZus{}year}\PYG{l+s+s2}{\PYGZdq{}}\PYG{p}{,} \PYG{l+s+s2}{\PYGZdq{}}\PYG{l+s+s2}{total\PYGZus{}service\PYGZus{}emissions\PYGZus{}kg\PYGZus{}per\PYGZus{}year}\PYG{l+s+s2}{\PYGZdq{}}\PYG{p}{,} \PYG{l+s+s2}{\PYGZdq{}}\PYG{l+s+s2}{electricity\PYGZus{}kwh\PYGZus{}per\PYGZus{}year}\PYG{l+s+s2}{\PYGZdq{}}\PYG{p}{,}
                \PYG{l+s+s2}{\PYGZdq{}}\PYG{l+s+s2}{heating\PYGZus{}cash\PYGZus{}operating\PYGZus{}cost\PYGZus{}yuan\PYGZus{}per\PYGZus{}year}\PYG{l+s+s2}{\PYGZdq{}}\PYG{p}{,} \PYG{l+s+s2}{\PYGZdq{}}\PYG{l+s+s2}{cooling\PYGZus{}demand\PYGZus{}kwh\PYGZus{}per\PYGZus{}year}\PYG{l+s+s2}{\PYGZdq{}}\PYG{p}{)}\PYG{p}{:}
        \PYG{k}{if} \PYG{n}{key} \PYG{o+ow}{in} \PYG{n}{result}\PYG{o}{.}\PYG{n}{replay}\PYG{o}{.}\PYG{n}{summary}\PYG{p}{:}
            \PYG{n}{summary}\PYG{p}{[}\PYG{n}{key}\PYG{o}{.}\PYG{n}{replace}\PYG{p}{(}\PYG{l+s+s2}{\PYGZdq{}}\PYG{l+s+s2}{\PYGZus{}per\PYGZus{}year}\PYG{l+s+s2}{\PYGZdq{}}\PYG{p}{,} \PYG{l+s+s2}{\PYGZdq{}}\PYG{l+s+s2}{\PYGZus{}over\PYGZus{}period}\PYG{l+s+s2}{\PYGZdq{}}\PYG{p}{)}\PYG{p}{]} \PYG{o}{=} \PYG{n}{result}\PYG{o}{.}\PYG{n}{replay}\PYG{o}{.}\PYG{n}{summary}\PYG{p}{[}\PYG{n}{key}\PYG{p}{]}
    \PYG{k}{return} \PYG{n}{HeatingSeasonReplayArtifacts}\PYG{p}{(}\PYG{n}{result}\PYG{p}{,} \PYG{n}{scope\PYGZus{}audit}\PYG{p}{,} \PYG{n}{assumptions}\PYG{p}{,} \PYG{n}{storage}\PYG{p}{,} \PYG{n}{certificate}\PYG{p}{,} \PYG{n}{summary}\PYG{p}{)}
\end{Verbatim}

\endgroup

\noindent\textbf{结果核对。}优先核对\path{exact_service_period}、\path{failed_gates}和\path{heating_season_certified}，再读费用与能量。期间指标使用\path{over_period}，按对应窗口解释。

\Needspace{12\baselineskip}
\subsection*{9．带主目标容差的两阶段端点}
完整端点方法，先求主目标，再施加单侧主目标上界，求解次目标并保留两阶段证据。

\noindent\textbf{输入：}主目标类型、绝对和相对容差，以及各阶段求解预算。\par
\noindent\textbf{输出：}次目标阶段的解与两阶段状态。限时可行值不等于全局词典序最优。\par

\noindent\textbf{实现思路。}两阶段端点与“先定投资、再算运行”用途不同。此处先求主目标，再把首阶段所得可行值加绝对、相对容差写成单侧上界，在该约束下优化次目标。排放表达含常量时，加入矩阵行的右端必须扣除常量；否则约束的实际范围会偏移。两阶段状态与预算分别保存，便于判断端点证据强度。

{\small 来源：\path{model.py}，原第2829--2888行；60行原代码、11行补充注释。}\par

\begingroup
\setstretch{1.0}
\begin{Verbatim}[commandchars=\\\{\},numbers=left,firstnumber=1,stepnumber=1,fontsize=\fontsize{9.5}{11.5}\selectfont,breaklines=true,breakanywhere=true,frame=single,framesep=3mm,numbersep=5pt,rulecolor=\color{black},xleftmargin=18pt]
\PYG{c+c1}{\PYGZsh{} \CJKCode{与固定投资两阶段不同：此处两阶段切换目标，不自动冻结资产。}}
\PYG{k}{def}\PYG{+w}{ }\PYG{n+nf}{solve\PYGZus{}lexicographic}\PYG{p}{(}
    \PYG{n+nb+bp}{self}\PYG{p}{,} \PYG{n}{primary}\PYG{p}{:} \PYG{n+nb}{str} \PYG{o}{=} \PYG{l+s+s2}{\PYGZdq{}}\PYG{l+s+s2}{emissions}\PYG{l+s+s2}{\PYGZdq{}}\PYG{p}{,} \PYG{o}{*}\PYG{p}{,} \PYG{n}{absolute\PYGZus{}tolerance}\PYG{p}{:} \PYG{n+nb}{float} \PYG{o}{=} \PYG{l+m+mf}{1e\PYGZhy{}7}\PYG{p}{,}
    \PYG{n}{relative\PYGZus{}tolerance}\PYG{p}{:} \PYG{n+nb}{float} \PYG{o}{=} \PYG{l+m+mf}{1e\PYGZhy{}8}\PYG{p}{,} \PYG{n}{mip\PYGZus{}gap}\PYG{p}{:} \PYG{n+nb}{float} \PYG{o}{|} \PYG{k+kc}{None} \PYG{o}{=} \PYG{k+kc}{None}\PYG{p}{,}
    \PYG{n}{time\PYGZus{}limit}\PYG{p}{:} \PYG{n+nb}{float} \PYG{o}{|} \PYG{k+kc}{None} \PYG{o}{=} \PYG{k+kc}{None}\PYG{p}{,}
\PYG{p}{)} \PYG{o}{\PYGZhy{}}\PYG{o}{\PYGZgt{}} \PYG{n}{SolveArtifacts}\PYG{p}{:}
\PYG{+w}{    }\PYG{l+s+sd}{\PYGZdq{}\PYGZdq{}\PYGZdq{}Two\PYGZhy{}stage endpoint with an explicit primary\PYGZhy{}objective allowance.}

\PYG{l+s+sd}{    Each stage uses the configured solve budget. A time\PYGZhy{}limited first}
\PYG{l+s+sd}{    stage only certifies a secondary solution within its incumbent bound,}
\PYG{l+s+sd}{    not a globally proven lexicographic optimum. The returned summary}
\PYG{l+s+sd}{    retains both stages\PYGZsq{} termination evidence. This instance keeps the}
\PYG{l+s+sd}{    primary bound, so use a fresh model for each independent endpoint.}
\PYG{l+s+sd}{    \PYGZdq{}\PYGZdq{}\PYGZdq{}}
    \PYG{c+c1}{\PYGZsh{} \CJKCode{只允许费用或排放主目标，容差必须有限且非负。}}
    \PYG{k}{if} \PYG{n}{primary} \PYG{o+ow}{not} \PYG{o+ow}{in} \PYG{p}{\PYGZob{}}\PYG{l+s+s2}{\PYGZdq{}}\PYG{l+s+s2}{total\PYGZus{}cost}\PYG{l+s+s2}{\PYGZdq{}}\PYG{p}{,} \PYG{l+s+s2}{\PYGZdq{}}\PYG{l+s+s2}{emissions}\PYG{l+s+s2}{\PYGZdq{}}\PYG{p}{\PYGZcb{}}\PYG{p}{:}
        \PYG{k}{raise} \PYG{n+ne}{ValueError}\PYG{p}{(}\PYG{l+s+s2}{\PYGZdq{}}\PYG{l+s+s2}{primary must be total\PYGZus{}cost or emissions}\PYG{l+s+s2}{\PYGZdq{}}\PYG{p}{)}
    \PYG{k}{if} \PYG{o+ow}{not} \PYG{n}{np}\PYG{o}{.}\PYG{n}{isfinite}\PYG{p}{(}\PYG{n}{absolute\PYGZus{}tolerance}\PYG{p}{)} \PYG{o+ow}{or} \PYG{o+ow}{not} \PYG{n}{np}\PYG{o}{.}\PYG{n}{isfinite}\PYG{p}{(}\PYG{n}{relative\PYGZus{}tolerance}\PYG{p}{)} \PYG{o+ow}{or} \PYG{n}{absolute\PYGZus{}tolerance} \PYG{o}{\PYGZlt{}} \PYG{l+m+mi}{0} \PYG{o+ow}{or} \PYG{n}{relative\PYGZus{}tolerance} \PYG{o}{\PYGZlt{}} \PYG{l+m+mi}{0}\PYG{p}{:}
        \PYG{k}{raise} \PYG{n+ne}{ValueError}\PYG{p}{(}\PYG{l+s+s2}{\PYGZdq{}}\PYG{l+s+s2}{lexicographic tolerances must be finite and nonnegative}\PYG{l+s+s2}{\PYGZdq{}}\PYG{p}{)}
    \PYG{c+c1}{\PYGZsh{} \CJKCode{第一阶段先取得实际可行候选值，保留其状态、间隙与下界。}}
    \PYG{n}{first} \PYG{o}{=} \PYG{n+nb+bp}{self}\PYG{o}{.}\PYG{n}{solve}\PYG{p}{(}\PYG{n}{mip\PYGZus{}gap}\PYG{o}{=}\PYG{n}{mip\PYGZus{}gap}\PYG{p}{,} \PYG{n}{time\PYGZus{}limit}\PYG{o}{=}\PYG{n}{time\PYGZus{}limit}\PYG{p}{,} \PYG{n}{objective\PYGZus{}mode}\PYG{o}{=}\PYG{n}{primary}\PYG{p}{)}
    \PYG{c+c1}{\PYGZsh{} \CJKCode{排放表达包含登记常量；添加上界时需扣除常量。}}
    \PYG{k}{if} \PYG{n}{primary} \PYG{o}{==} \PYG{l+s+s2}{\PYGZdq{}}\PYG{l+s+s2}{emissions}\PYG{l+s+s2}{\PYGZdq{}}\PYG{p}{:}
        \PYG{n}{coefficients}\PYG{p}{,} \PYG{n}{constant} \PYG{o}{=} \PYG{n+nb+bp}{self}\PYG{o}{.}\PYG{n}{total\PYGZus{}service\PYGZus{}emissions\PYGZus{}expression}\PYG{p}{(}\PYG{p}{)}
        \PYG{n}{value} \PYG{o}{=} \PYG{n+nb}{float}\PYG{p}{(}\PYG{n}{first}\PYG{o}{.}\PYG{n}{summary}\PYG{p}{[}\PYG{l+s+s2}{\PYGZdq{}}\PYG{l+s+s2}{total\PYGZus{}service\PYGZus{}emissions\PYGZus{}kg\PYGZus{}per\PYGZus{}year}\PYG{l+s+s2}{\PYGZdq{}}\PYG{p}{]}\PYG{p}{)}
        \PYG{n}{secondary} \PYG{o}{=} \PYG{l+s+s2}{\PYGZdq{}}\PYG{l+s+s2}{total\PYGZus{}cost}\PYG{l+s+s2}{\PYGZdq{}}
    \PYG{k}{else}\PYG{p}{:}
        \PYG{c+c1}{\PYGZsh{} \CJKCode{费用主目标使用当前原目标系数，不以事后费用统计替代。}}
        \PYG{n}{coefficients} \PYG{o}{=} \PYG{p}{\PYGZob{}}\PYG{n}{i}\PYG{p}{:}\PYG{n}{c} \PYG{k}{for} \PYG{n}{i}\PYG{p}{,}\PYG{n}{c} \PYG{o+ow}{in} \PYG{n+nb}{enumerate}\PYG{p}{(}\PYG{n+nb+bp}{self}\PYG{o}{.}\PYG{n}{problem}\PYG{o}{.}\PYG{n}{objective}\PYG{p}{)} \PYG{k}{if} \PYG{n}{c}\PYG{p}{\PYGZcb{}}
        \PYG{n}{constant} \PYG{o}{=} \PYG{l+m+mf}{0.0}
        \PYG{n}{value} \PYG{o}{=} \PYG{n+nb+bp}{self}\PYG{o}{.}\PYG{n}{\PYGZus{}dot}\PYG{p}{(}\PYG{n}{coefficients}\PYG{p}{,} \PYG{n}{first}\PYG{o}{.}\PYG{n}{values}\PYG{p}{)}
        \PYG{n}{secondary} \PYG{o}{=} \PYG{l+s+s2}{\PYGZdq{}}\PYG{l+s+s2}{emissions}\PYG{l+s+s2}{\PYGZdq{}}
    \PYG{c+c1}{\PYGZsh{} \CJKCode{上界为首阶段值加绝对及相对容差；允许主目标进一步改善，不是双侧邻域。}}
    \PYG{n}{allowance} \PYG{o}{=} \PYG{n}{absolute\PYGZus{}tolerance} \PYG{o}{+} \PYG{n}{relative\PYGZus{}tolerance} \PYG{o}{*} \PYG{n+nb}{abs}\PYG{p}{(}\PYG{n}{value}\PYG{p}{)}
    \PYG{n+nb+bp}{self}\PYG{o}{.}\PYG{n}{problem}\PYG{o}{.}\PYG{n}{add\PYGZus{}constraint}\PYG{p}{(}\PYG{n}{coefficients}\PYG{p}{,} \PYG{n}{ub}\PYG{o}{=}\PYG{n}{value}\PYG{o}{+}\PYG{n}{allowance}\PYG{o}{\PYGZhy{}}\PYG{n}{constant}\PYG{p}{)}
    \PYG{n}{failed\PYGZus{}second} \PYG{o}{=} \PYG{k+kc}{None}
    \PYG{c+c1}{\PYGZsh{} \CJKCode{第二阶段用同样的预算求次目标；不据此宣称全局词典序最优。}}
    \PYG{k}{try}\PYG{p}{:}
        \PYG{n}{second} \PYG{o}{=} \PYG{n+nb+bp}{self}\PYG{o}{.}\PYG{n}{solve}\PYG{p}{(}\PYG{n}{mip\PYGZus{}gap}\PYG{o}{=}\PYG{n}{mip\PYGZus{}gap}\PYG{p}{,} \PYG{n}{time\PYGZus{}limit}\PYG{o}{=}\PYG{n}{time\PYGZus{}limit}\PYG{p}{,} \PYG{n}{objective\PYGZus{}mode}\PYG{o}{=}\PYG{n}{secondary}\PYG{p}{)}
        \PYG{n}{endpoint\PYGZus{}summary} \PYG{o}{=} \PYG{n}{second}\PYG{o}{.}\PYG{n}{summary}
    \PYG{c+c1}{\PYGZsh{} \CJKCode{次阶段无可行解时保留失败摘要，后文继续补足首阶段证据。}}
    \PYG{k}{except} \PYG{n}{SolveFailure} \PYG{k}{as} \PYG{n}{error}\PYG{p}{:}
        \PYG{n}{failed\PYGZus{}second} \PYG{o}{=} \PYG{n}{error}
        \PYG{n}{endpoint\PYGZus{}summary} \PYG{o}{=} \PYG{n}{error}\PYG{o}{.}\PYG{n}{summary}
    \PYG{c+c1}{\PYGZsh{} \CJKCode{同时报告首阶段可行值、容差和终止状态，防止只展示次阶段结果。}}
    \PYG{n}{endpoint\PYGZus{}summary}\PYG{o}{.}\PYG{n}{update}\PYG{p}{(}\PYG{p}{\PYGZob{}}
        \PYG{l+s+s2}{\PYGZdq{}}\PYG{l+s+s2}{endpoint\PYGZus{}method}\PYG{l+s+s2}{\PYGZdq{}}\PYG{p}{:} \PYG{l+s+s2}{\PYGZdq{}}\PYG{l+s+s2}{two\PYGZus{}stage\PYGZus{}lexicographic\PYGZus{}with\PYGZus{}reported\PYGZus{}tolerance}\PYG{l+s+s2}{\PYGZdq{}}\PYG{p}{,}
        \PYG{l+s+s2}{\PYGZdq{}}\PYG{l+s+s2}{primary\PYGZus{}objective}\PYG{l+s+s2}{\PYGZdq{}}\PYG{p}{:} \PYG{n}{primary}\PYG{p}{,} \PYG{l+s+s2}{\PYGZdq{}}\PYG{l+s+s2}{primary\PYGZus{}incumbent\PYGZus{}value}\PYG{l+s+s2}{\PYGZdq{}}\PYG{p}{:} \PYG{n}{value}\PYG{p}{,}
        \PYG{l+s+s2}{\PYGZdq{}}\PYG{l+s+s2}{primary\PYGZus{}objective\PYGZus{}allowance}\PYG{l+s+s2}{\PYGZdq{}}\PYG{p}{:} \PYG{n}{allowance}\PYG{p}{,}
        \PYG{l+s+s2}{\PYGZdq{}}\PYG{l+s+s2}{primary\PYGZus{}solver\PYGZus{}status}\PYG{l+s+s2}{\PYGZdq{}}\PYG{p}{:} \PYG{n}{first}\PYG{o}{.}\PYG{n}{summary}\PYG{p}{[}\PYG{l+s+s2}{\PYGZdq{}}\PYG{l+s+s2}{solver\PYGZus{}status}\PYG{l+s+s2}{\PYGZdq{}}\PYG{p}{]}\PYG{p}{,}
        \PYG{l+s+s2}{\PYGZdq{}}\PYG{l+s+s2}{primary\PYGZus{}mip\PYGZus{}gap}\PYG{l+s+s2}{\PYGZdq{}}\PYG{p}{:} \PYG{n}{first}\PYG{o}{.}\PYG{n}{summary}\PYG{o}{.}\PYG{n}{get}\PYG{p}{(}\PYG{l+s+s2}{\PYGZdq{}}\PYG{l+s+s2}{mip\PYGZus{}gap}\PYG{l+s+s2}{\PYGZdq{}}\PYG{p}{)}\PYG{p}{,}
        \PYG{l+s+s2}{\PYGZdq{}}\PYG{l+s+s2}{primary\PYGZus{}numerical\PYGZus{}audit\PYGZus{}passed}\PYG{l+s+s2}{\PYGZdq{}}\PYG{p}{:} \PYG{n}{first}\PYG{o}{.}\PYG{n}{summary}\PYG{p}{[}\PYG{l+s+s2}{\PYGZdq{}}\PYG{l+s+s2}{numerical\PYGZus{}audit\PYGZus{}passed}\PYG{l+s+s2}{\PYGZdq{}}\PYG{p}{]}\PYG{p}{,}
        \PYG{l+s+s2}{\PYGZdq{}}\PYG{l+s+s2}{primary\PYGZus{}solver\PYGZus{}message}\PYG{l+s+s2}{\PYGZdq{}}\PYG{p}{:} \PYG{n}{first}\PYG{o}{.}\PYG{n}{summary}\PYG{p}{[}\PYG{l+s+s2}{\PYGZdq{}}\PYG{l+s+s2}{solver\PYGZus{}message}\PYG{l+s+s2}{\PYGZdq{}}\PYG{p}{]}\PYG{p}{,}
        \PYG{l+s+s2}{\PYGZdq{}}\PYG{l+s+s2}{primary\PYGZus{}solver\PYGZus{}objective\PYGZus{}value}\PYG{l+s+s2}{\PYGZdq{}}\PYG{p}{:} \PYG{n}{first}\PYG{o}{.}\PYG{n}{summary}\PYG{p}{[}\PYG{l+s+s2}{\PYGZdq{}}\PYG{l+s+s2}{objective\PYGZus{}value}\PYG{l+s+s2}{\PYGZdq{}}\PYG{p}{]}\PYG{p}{,}
        \PYG{l+s+s2}{\PYGZdq{}}\PYG{l+s+s2}{primary\PYGZus{}mip\PYGZus{}dual\PYGZus{}bound}\PYG{l+s+s2}{\PYGZdq{}}\PYG{p}{:} \PYG{n}{first}\PYG{o}{.}\PYG{n}{summary}\PYG{o}{.}\PYG{n}{get}\PYG{p}{(}\PYG{l+s+s2}{\PYGZdq{}}\PYG{l+s+s2}{mip\PYGZus{}dual\PYGZus{}bound}\PYG{l+s+s2}{\PYGZdq{}}\PYG{p}{)}\PYG{p}{,}
        \PYG{l+s+s2}{\PYGZdq{}}\PYG{l+s+s2}{primary\PYGZus{}solver\PYGZus{}requested\PYGZus{}mip\PYGZus{}gap}\PYG{l+s+s2}{\PYGZdq{}}\PYG{p}{:} \PYG{n}{first}\PYG{o}{.}\PYG{n}{summary}\PYG{p}{[}\PYG{l+s+s2}{\PYGZdq{}}\PYG{l+s+s2}{solver\PYGZus{}requested\PYGZus{}mip\PYGZus{}gap}\PYG{l+s+s2}{\PYGZdq{}}\PYG{p}{]}\PYG{p}{,}
        \PYG{l+s+s2}{\PYGZdq{}}\PYG{l+s+s2}{primary\PYGZus{}solver\PYGZus{}time\PYGZus{}limit\PYGZus{}seconds}\PYG{l+s+s2}{\PYGZdq{}}\PYG{p}{:} \PYG{n}{first}\PYG{o}{.}\PYG{n}{summary}\PYG{p}{[}\PYG{l+s+s2}{\PYGZdq{}}\PYG{l+s+s2}{solver\PYGZus{}time\PYGZus{}limit\PYGZus{}seconds}\PYG{l+s+s2}{\PYGZdq{}}\PYG{p}{]}\PYG{p}{,}
        \PYG{l+s+s2}{\PYGZdq{}}\PYG{l+s+s2}{primary\PYGZus{}solver\PYGZus{}elapsed\PYGZus{}seconds}\PYG{l+s+s2}{\PYGZdq{}}\PYG{p}{:} \PYG{n}{first}\PYG{o}{.}\PYG{n}{summary}\PYG{p}{[}\PYG{l+s+s2}{\PYGZdq{}}\PYG{l+s+s2}{solver\PYGZus{}elapsed\PYGZus{}seconds}\PYG{l+s+s2}{\PYGZdq{}}\PYG{p}{]}\PYG{p}{,}
        \PYG{l+s+s2}{\PYGZdq{}}\PYG{l+s+s2}{endpoint\PYGZus{}solver\PYGZus{}elapsed\PYGZus{}seconds}\PYG{l+s+s2}{\PYGZdq{}}\PYG{p}{:} \PYG{n}{first}\PYG{o}{.}\PYG{n}{summary}\PYG{p}{[}\PYG{l+s+s2}{\PYGZdq{}}\PYG{l+s+s2}{solver\PYGZus{}elapsed\PYGZus{}seconds}\PYG{l+s+s2}{\PYGZdq{}}\PYG{p}{]}\PYG{o}{+}\PYG{n}{endpoint\PYGZus{}summary}\PYG{p}{[}\PYG{l+s+s2}{\PYGZdq{}}\PYG{l+s+s2}{solver\PYGZus{}elapsed\PYGZus{}seconds}\PYG{l+s+s2}{\PYGZdq{}}\PYG{p}{]}\PYG{p}{,}
    \PYG{p}{\PYGZcb{}}\PYG{p}{)}
    \PYG{c+c1}{\PYGZsh{} \CJKCode{两阶段证据分别标识}primary\CJKCode{和}secondary\CJKCode{，可追踪求解身份。}}
    \PYG{n}{stages} \PYG{o}{=} \PYG{p}{[}\PYG{n}{deepcopy}\PYG{p}{(}\PYG{n}{first}\PYG{o}{.}\PYG{n}{summary}\PYG{p}{[}\PYG{l+s+s2}{\PYGZdq{}}\PYG{l+s+s2}{solver\PYGZus{}stages}\PYG{l+s+s2}{\PYGZdq{}}\PYG{p}{]}\PYG{p}{[}\PYG{l+m+mi}{0}\PYG{p}{]}\PYG{p}{)}\PYG{p}{,} \PYG{n}{deepcopy}\PYG{p}{(}\PYG{n}{endpoint\PYGZus{}summary}\PYG{p}{[}\PYG{l+s+s2}{\PYGZdq{}}\PYG{l+s+s2}{solver\PYGZus{}stages}\PYG{l+s+s2}{\PYGZdq{}}\PYG{p}{]}\PYG{p}{[}\PYG{l+m+mi}{0}\PYG{p}{]}\PYG{p}{)}\PYG{p}{]}
    \PYG{n}{stages}\PYG{p}{[}\PYG{l+m+mi}{0}\PYG{p}{]}\PYG{p}{[}\PYG{l+s+s2}{\PYGZdq{}}\PYG{l+s+s2}{stage\PYGZus{}id}\PYG{l+s+s2}{\PYGZdq{}}\PYG{p}{]}\PYG{p}{,} \PYG{n}{stages}\PYG{p}{[}\PYG{l+m+mi}{1}\PYG{p}{]}\PYG{p}{[}\PYG{l+s+s2}{\PYGZdq{}}\PYG{l+s+s2}{stage\PYGZus{}id}\PYG{l+s+s2}{\PYGZdq{}}\PYG{p}{]} \PYG{o}{=} \PYG{l+s+s2}{\PYGZdq{}}\PYG{l+s+s2}{primary}\PYG{l+s+s2}{\PYGZdq{}}\PYG{p}{,} \PYG{l+s+s2}{\PYGZdq{}}\PYG{l+s+s2}{secondary}\PYG{l+s+s2}{\PYGZdq{}}
    \PYG{n}{endpoint\PYGZus{}summary}\PYG{p}{[}\PYG{l+s+s2}{\PYGZdq{}}\PYG{l+s+s2}{solver\PYGZus{}stages}\PYG{l+s+s2}{\PYGZdq{}}\PYG{p}{]} \PYG{o}{=} \PYG{n}{stages}
    \PYG{n}{endpoint\PYGZus{}summary}\PYG{p}{[}\PYG{l+s+s2}{\PYGZdq{}}\PYG{l+s+s2}{solution\PYGZus{}quality}\PYG{l+s+s2}{\PYGZdq{}}\PYG{p}{]} \PYG{o}{=} \PYG{n}{solution\PYGZus{}quality\PYGZus{}from\PYGZus{}summary}\PYG{p}{(}
        \PYG{n}{endpoint\PYGZus{}summary}\PYG{p}{,} \PYG{n}{scope}\PYG{o}{=}\PYG{n}{first}\PYG{o}{.}\PYG{n}{summary}\PYG{p}{[}\PYG{l+s+s2}{\PYGZdq{}}\PYG{l+s+s2}{solution\PYGZus{}quality}\PYG{l+s+s2}{\PYGZdq{}}\PYG{p}{]}\PYG{p}{[}\PYG{l+s+s2}{\PYGZdq{}}\PYG{l+s+s2}{scope}\PYG{l+s+s2}{\PYGZdq{}}\PYG{p}{]}\PYG{p}{)}
    \PYG{c+c1}{\PYGZsh{} \CJKCode{次阶段失败在补足证据后重新抛出，不静默返回首阶段当作成功端点。}}
    \PYG{k}{if} \PYG{n}{failed\PYGZus{}second} \PYG{o+ow}{is} \PYG{o+ow}{not} \PYG{k+kc}{None}\PYG{p}{:}
        \PYG{n}{failed\PYGZus{}second}\PYG{o}{.}\PYG{n}{solution\PYGZus{}quality} \PYG{o}{=} \PYG{n}{endpoint\PYGZus{}summary}\PYG{p}{[}\PYG{l+s+s2}{\PYGZdq{}}\PYG{l+s+s2}{solution\PYGZus{}quality}\PYG{l+s+s2}{\PYGZdq{}}\PYG{p}{]}
        \PYG{k}{raise} \PYG{n}{failed\PYGZus{}second}
    \PYG{k}{return} \PYG{n}{second}
\end{Verbatim}

\endgroup

\noindent\textbf{结果核对。}联合查看主目标可行值、允许容差、下界和两次求解状态。次阶段失败要按失败处理；首阶段未取得全局证明时，不能只以次阶段目标值命名全局端点。

\Needspace{12\baselineskip}
\subsection*{10．独立碳帽采样与非支配筛选}
完整独立扫描接口。清除原碳帽、碳价置零，再按碳上限逐档重新构模和检查。

\noindent\textbf{输入：}同一配置和时间序列、至少两档采样；投资自由度由输入配置决定。\par
\noindent\textbf{输出：}样本内非支配集合、全部采样证据、边际减排费用及归一化折中建议。\par

\noindent\textbf{实现思路。}扫描采用同一基础配置和时序，每个碳帽重新建立模型，防止上一档约束残留。先清除碳价与原碳帽，再由两个端点确定采样范围；每个候选必须同时通过数值检查和实际碳帽检查。非支配筛选在已取得的可行点中进行，完整采样记录保留在结果属性中，方便核对被去重或支配的点。

{\small 来源：\path{experiments.py}，原第351--411行；61行原代码、16行补充注释。}\par

\begingroup
\setstretch{1.0}
\begin{Verbatim}[commandchars=\\\{\},numbers=left,firstnumber=1,stepnumber=1,fontsize=\fontsize{9.5}{11.5}\selectfont,breaklines=true,breakanywhere=true,frame=single,framesep=3mm,numbersep=5pt,rulecolor=\color{black},xleftmargin=18pt]
\PYG{c+c1}{\PYGZsh{} \CJKCode{独立接口不自动运行于三城单次入口，也不代表历史扫描使用相同规则。}}
\PYG{k}{def}\PYG{+w}{ }\PYG{n+nf}{epsilon\PYGZus{}constraint\PYGZus{}pareto}\PYG{p}{(}
    \PYG{n}{config\PYGZus{}path}\PYG{p}{:} \PYG{n+nb}{str} \PYG{o}{|} \PYG{n}{Path}\PYG{p}{,}
    \PYG{n}{timeseries\PYGZus{}path}\PYG{p}{:} \PYG{n+nb}{str} \PYG{o}{|} \PYG{n}{Path}\PYG{p}{,}
    \PYG{n}{points}\PYG{p}{:} \PYG{n+nb}{int} \PYG{o}{=} \PYG{l+m+mi}{7}\PYG{p}{,}
\PYG{p}{)} \PYG{o}{\PYGZhy{}}\PYG{o}{\PYGZgt{}} \PYG{n}{pd}\PYG{o}{.}\PYG{n}{DataFrame}\PYG{p}{:}
\PYG{+w}{    }\PYG{l+s+sd}{\PYGZdq{}\PYGZdq{}\PYGZdq{}Sample a nondominated system cost/carbon frontier with explicit caps.}

\PYG{l+s+sd}{    The investment and operating model is unchanged across points. A hard reliability}
\PYG{l+s+sd}{    constraint prevents the emissions objective from creating artificial low\PYGZhy{}carbon}
\PYG{l+s+sd}{    solutions through heat curtailment.}
\PYG{l+s+sd}{    \PYGZdq{}\PYGZdq{}\PYGZdq{}}
    \PYG{k}{if} \PYG{n}{points} \PYG{o}{\PYGZlt{}} \PYG{l+m+mi}{2}\PYG{p}{:}
        \PYG{k}{raise} \PYG{n+ne}{ValueError}\PYG{p}{(}\PYG{l+s+s2}{\PYGZdq{}}\PYG{l+s+s2}{points \CJKCode{至少为} 2}\PYG{l+s+s2}{\PYGZdq{}}\PYG{p}{)}
    \PYG{c+c1}{\PYGZsh{} \CJKCode{配置与时序只读取一次；各档复制同一基础配置，保持输入一致。}}
    \PYG{n}{base\PYGZus{}config}\PYG{p}{,} \PYG{n}{frame} \PYG{o}{=} \PYG{n}{load\PYGZus{}inputs}\PYG{p}{(}\PYG{n}{config\PYGZus{}path}\PYG{p}{,} \PYG{n}{timeseries\PYGZus{}path}\PYG{p}{)}
    \PYG{n}{base\PYGZus{}config} \PYG{o}{=} \PYG{n}{copy}\PYG{o}{.}\PYG{n}{deepcopy}\PYG{p}{(}\PYG{n}{base\PYGZus{}config}\PYG{p}{)}
    \PYG{c+c1}{\PYGZsh{} \CJKCode{清除碳价与原碳帽；其余启用约束和偏好仍按配置保留。}}
    \PYG{n}{base\PYGZus{}config}\PYG{p}{[}\PYG{l+s+s2}{\PYGZdq{}}\PYG{l+s+s2}{planning}\PYG{l+s+s2}{\PYGZdq{}}\PYG{p}{]}\PYG{p}{[}\PYG{l+s+s2}{\PYGZdq{}}\PYG{l+s+s2}{carbon\PYGZus{}price\PYGZus{}yuan\PYGZus{}per\PYGZus{}kg}\PYG{l+s+s2}{\PYGZdq{}}\PYG{p}{]} \PYG{o}{=} \PYG{l+m+mf}{0.0}
    \PYG{n}{base\PYGZus{}config}\PYG{p}{[}\PYG{l+s+s2}{\PYGZdq{}}\PYG{l+s+s2}{planning}\PYG{l+s+s2}{\PYGZdq{}}\PYG{p}{]}\PYG{o}{.}\PYG{n}{pop}\PYG{p}{(}\PYG{l+s+s2}{\PYGZdq{}}\PYG{l+s+s2}{emissions\PYGZus{}cap\PYGZus{}kg\PYGZus{}per\PYGZus{}year}\PYG{l+s+s2}{\PYGZdq{}}\PYG{p}{,} \PYG{k+kc}{None}\PYG{p}{)}

    \PYG{c+c1}{\PYGZsh{} \CJKCode{两个端点分别新建模型，避免前一端点的主目标约束污染另一端点。}}
    \PYG{n}{cost\PYGZus{}optimum} \PYG{o}{=} \PYG{n}{HeatPlanningModel}\PYG{p}{(}\PYG{n}{copy}\PYG{o}{.}\PYG{n}{deepcopy}\PYG{p}{(}\PYG{n}{base\PYGZus{}config}\PYG{p}{)}\PYG{p}{,} \PYG{n}{frame}\PYG{p}{)}\PYG{o}{.}\PYG{n}{solve\PYGZus{}lexicographic}\PYG{p}{(}\PYG{n}{primary}\PYG{o}{=}\PYG{l+s+s2}{\PYGZdq{}}\PYG{l+s+s2}{total\PYGZus{}cost}\PYG{l+s+s2}{\PYGZdq{}}\PYG{p}{)}\PYG{o}{.}\PYG{n}{summary}
    \PYG{n}{emission\PYGZus{}optimum} \PYG{o}{=} \PYG{n}{HeatPlanningModel}\PYG{p}{(}\PYG{n}{copy}\PYG{o}{.}\PYG{n}{deepcopy}\PYG{p}{(}\PYG{n}{base\PYGZus{}config}\PYG{p}{)}\PYG{p}{,} \PYG{n}{frame}\PYG{p}{)}\PYG{o}{.}\PYG{n}{solve\PYGZus{}lexicographic}\PYG{p}{(}\PYG{n}{primary}\PYG{o}{=}\PYG{l+s+s2}{\PYGZdq{}}\PYG{l+s+s2}{emissions}\PYG{l+s+s2}{\PYGZdq{}}\PYG{p}{)}\PYG{o}{.}\PYG{n}{summary}
    \PYG{c+c1}{\PYGZsh{} \CJKCode{采用冷热服务及共享电力口径的登记系统总碳，单位}kg CO2\CJKCode{。}}
    \PYG{n}{emission\PYGZus{}column} \PYG{o}{=} \PYG{l+s+s2}{\PYGZdq{}}\PYG{l+s+s2}{total\PYGZus{}service\PYGZus{}emissions\PYGZus{}kg\PYGZus{}per\PYGZus{}year}\PYG{l+s+s2}{\PYGZdq{}}
    \PYG{n}{min\PYGZus{}emissions} \PYG{o}{=} \PYG{n+nb}{float}\PYG{p}{(}\PYG{n}{emission\PYGZus{}optimum}\PYG{p}{[}\PYG{n}{emission\PYGZus{}column}\PYG{p}{]}\PYG{p}{)}
    \PYG{n}{cost\PYGZus{}optimum\PYGZus{}emissions} \PYG{o}{=} \PYG{n+nb}{float}\PYG{p}{(}\PYG{n}{cost\PYGZus{}optimum}\PYG{p}{[}\PYG{n}{emission\PYGZus{}column}\PYG{p}{]}\PYG{p}{)}
    \PYG{c+c1}{\PYGZsh{} \CJKCode{端点次序超出容差则报错，不伪造可用扫描区间。}}
    \PYG{k}{if} \PYG{n}{min\PYGZus{}emissions} \PYG{o}{\PYGZgt{}} \PYG{n}{cost\PYGZus{}optimum\PYGZus{}emissions} \PYG{o}{+} \PYG{l+m+mf}{1e\PYGZhy{}6} \PYG{o}{+} \PYG{l+m+mf}{1e\PYGZhy{}8}\PYG{o}{*}\PYG{n+nb}{abs}\PYG{p}{(}\PYG{n}{cost\PYGZus{}optimum\PYGZus{}emissions}\PYG{p}{)}\PYG{p}{:}
        \PYG{k}{raise} \PYG{n+ne}{RuntimeError}\PYG{p}{(}\PYG{l+s+s2}{\PYGZdq{}}\PYG{l+s+s2}{Carbon endpoint exceeds cost endpoint; improve solve budgets before defining an epsilon range}\PYG{l+s+s2}{\PYGZdq{}}\PYG{p}{)}
    \PYG{n}{caps} \PYG{o}{=} \PYG{n}{np}\PYG{o}{.}\PYG{n}{linspace}\PYG{p}{(}\PYG{n}{min\PYGZus{}emissions}\PYG{p}{,} \PYG{n}{cost\PYGZus{}optimum\PYGZus{}emissions}\PYG{p}{,} \PYG{n}{points}\PYG{p}{)}

    \PYG{c+c1}{\PYGZsh{} \CJKCode{逐档保存全部可行候选，再统一进行非支配筛选。}}
    \PYG{n}{rows} \PYG{o}{=} \PYG{p}{[}\PYG{p}{]}
    \PYG{c+c1}{\PYGZsh{} \CJKCode{每档均由基础配置开始；本函数没有调用投资冻结。}}
    \PYG{k}{for} \PYG{n}{cap} \PYG{o+ow}{in} \PYG{n}{caps}\PYG{p}{:}
        \PYG{n}{config} \PYG{o}{=} \PYG{n}{copy}\PYG{o}{.}\PYG{n}{deepcopy}\PYG{p}{(}\PYG{n}{base\PYGZus{}config}\PYG{p}{)}
        \PYG{c+c1}{\PYGZsh{} A tiny numerical allowance avoids rejecting the emissions\PYGZhy{}optimal endpoint.}
        \PYG{n}{effective\PYGZus{}cap} \PYG{o}{=} \PYG{n+nb}{float}\PYG{p}{(}\PYG{n}{cap}\PYG{p}{)} \PYG{o}{+} \PYG{l+m+mf}{1e\PYGZhy{}7} \PYG{o}{+} \PYG{n+nb}{abs}\PYG{p}{(}\PYG{n+nb}{float}\PYG{p}{(}\PYG{n}{cap}\PYG{p}{)}\PYG{p}{)}\PYG{o}{*}\PYG{l+m+mf}{1e\PYGZhy{}8}
        \PYG{c+c1}{\PYGZsh{} \CJKCode{实际碳帽包含微小数值余量，同时保存请求值供审查。}}
        \PYG{n}{config}\PYG{p}{[}\PYG{l+s+s2}{\PYGZdq{}}\PYG{l+s+s2}{planning}\PYG{l+s+s2}{\PYGZdq{}}\PYG{p}{]}\PYG{p}{[}\PYG{l+s+s2}{\PYGZdq{}}\PYG{l+s+s2}{emissions\PYGZus{}cap\PYGZus{}kg\PYGZus{}per\PYGZus{}year}\PYG{l+s+s2}{\PYGZdq{}}\PYG{p}{]} \PYG{o}{=} \PYG{n}{effective\PYGZus{}cap}
        \PYG{c+c1}{\PYGZsh{} \CJKCode{自由投资输入时每档重求投资和运行；已固定配置才只检验固定资产。}}
        \PYG{n}{summary} \PYG{o}{=} \PYG{n}{HeatPlanningModel}\PYG{p}{(}\PYG{n}{config}\PYG{p}{,} \PYG{n}{frame}\PYG{p}{)}\PYG{o}{.}\PYG{n}{solve}\PYG{p}{(}\PYG{n}{objective\PYGZus{}mode}\PYG{o}{=}\PYG{l+s+s2}{\PYGZdq{}}\PYG{l+s+s2}{total\PYGZus{}cost}\PYG{l+s+s2}{\PYGZdq{}}\PYG{p}{)}\PYG{o}{.}\PYG{n}{summary}
        \PYG{c+c1}{\PYGZsh{} \CJKCode{数值或碳帽复核失败立即终止，不能把失败点加入前沿。}}
        \PYG{k}{if} \PYG{o+ow}{not} \PYG{n}{summary}\PYG{p}{[}\PYG{l+s+s2}{\PYGZdq{}}\PYG{l+s+s2}{numerical\PYGZus{}audit\PYGZus{}passed}\PYG{l+s+s2}{\PYGZdq{}}\PYG{p}{]} \PYG{o+ow}{or} \PYG{o+ow}{not} \PYG{n}{summary}\PYG{p}{[}\PYG{l+s+s2}{\PYGZdq{}}\PYG{l+s+s2}{emissions\PYGZus{}cap\PYGZus{}pass}\PYG{l+s+s2}{\PYGZdq{}}\PYG{p}{]}\PYG{p}{:}
            \PYG{k}{raise} \PYG{n+ne}{RuntimeError}\PYG{p}{(}\PYG{l+s+s2}{\PYGZdq{}}\PYG{l+s+s2}{A sampled frontier incumbent failed independent numerical or carbon\PYGZhy{}cap checks}\PYG{l+s+s2}{\PYGZdq{}}\PYG{p}{)}
        \PYG{n}{rows}\PYG{o}{.}\PYG{n}{append}\PYG{p}{(}\PYG{p}{\PYGZob{}}\PYG{l+s+s2}{\PYGZdq{}}\PYG{l+s+s2}{requested\PYGZus{}emissions\PYGZus{}cap\PYGZus{}kg\PYGZus{}per\PYGZus{}year}\PYG{l+s+s2}{\PYGZdq{}}\PYG{p}{:} \PYG{n+nb}{float}\PYG{p}{(}\PYG{n}{cap}\PYG{p}{)}\PYG{p}{,} \PYG{o}{*}\PYG{o}{*}\PYG{n}{summary}\PYG{p}{\PYGZcb{}}\PYG{p}{)}
    \PYG{n}{sampled} \PYG{o}{=} \PYG{n}{pd}\PYG{o}{.}\PYG{n}{json\PYGZus{}normalize}\PYG{p}{(}\PYG{n}{rows}\PYG{p}{)}
    \PYG{c+c1}{\PYGZsh{} \CJKCode{非支配关系只针对已采样可行解，不是完整连续}Pareto\CJKCode{前沿。}}
    \PYG{n}{result} \PYG{o}{=} \PYG{n}{nondominated\PYGZus{}cost\PYGZus{}emission\PYGZus{}points}\PYG{p}{(}\PYG{n}{sampled}\PYG{p}{)}
    \PYG{c+c1}{\PYGZsh{} \CJKCode{保留筛选前的完整分母，便于追查重复及被支配样本。}}
    \PYG{n}{result}\PYG{o}{.}\PYG{n}{attrs}\PYG{p}{[}\PYG{l+s+s2}{\PYGZdq{}}\PYG{l+s+s2}{all\PYGZus{}sampled\PYGZus{}points}\PYG{l+s+s2}{\PYGZdq{}}\PYG{p}{]} \PYG{o}{=} \PYG{n}{sampled}\PYG{o}{.}\PYG{n}{to\PYGZus{}dict}\PYG{p}{(}\PYG{n}{orient}\PYG{o}{=}\PYG{l+s+s2}{\PYGZdq{}}\PYG{l+s+s2}{records}\PYG{l+s+s2}{\PYGZdq{}}\PYG{p}{)}
    \PYG{n}{result}\PYG{o}{.}\PYG{n}{attrs}\PYG{p}{[}\PYG{l+s+s2}{\PYGZdq{}}\PYG{l+s+s2}{cost\PYGZus{}endpoint}\PYG{l+s+s2}{\PYGZdq{}}\PYG{p}{]} \PYG{o}{=} \PYG{n}{cost\PYGZus{}optimum}
    \PYG{n}{result}\PYG{o}{.}\PYG{n}{attrs}\PYG{p}{[}\PYG{l+s+s2}{\PYGZdq{}}\PYG{l+s+s2}{carbon\PYGZus{}endpoint}\PYG{l+s+s2}{\PYGZdq{}}\PYG{p}{]} \PYG{o}{=} \PYG{n}{emission\PYGZus{}optimum}
    \PYG{n}{result}\PYG{p}{[}\PYG{l+s+s2}{\PYGZdq{}}\PYG{l+s+s2}{frontier\PYGZus{}scope}\PYG{l+s+s2}{\PYGZdq{}}\PYG{p}{]} \PYG{o}{=} \PYG{l+s+s2}{\PYGZdq{}}\PYG{l+s+s2}{nondominated\PYGZus{}subset\PYGZus{}of\PYGZus{}sampled\PYGZus{}feasible\PYGZus{}solutions}\PYG{l+s+s2}{\PYGZdq{}}
    \PYG{n}{result}\PYG{p}{[}\PYG{l+s+s2}{\PYGZdq{}}\PYG{l+s+s2}{sampled\PYGZus{}point\PYGZus{}count}\PYG{l+s+s2}{\PYGZdq{}}\PYG{p}{]} \PYG{o}{=} \PYG{n+nb}{len}\PYG{p}{(}\PYG{n}{rows}\PYG{p}{)}
    \PYG{n}{result}\PYG{p}{[}\PYG{l+s+s2}{\PYGZdq{}}\PYG{l+s+s2}{removed\PYGZus{}duplicate\PYGZus{}or\PYGZus{}dominated\PYGZus{}count}\PYG{l+s+s2}{\PYGZdq{}}\PYG{p}{]} \PYG{o}{=} \PYG{n+nb}{len}\PYG{p}{(}\PYG{n}{rows}\PYG{p}{)}\PYG{o}{\PYGZhy{}}\PYG{n+nb}{len}\PYG{p}{(}\PYG{n}{result}\PYG{p}{)}
    \PYG{n}{result}\PYG{p}{[}\PYG{l+s+s2}{\PYGZdq{}}\PYG{l+s+s2}{marginal\PYGZus{}abatement\PYGZus{}cost\PYGZus{}yuan\PYGZus{}per\PYGZus{}kg}\PYG{l+s+s2}{\PYGZdq{}}\PYG{p}{]} \PYG{o}{=} \PYG{n}{np}\PYG{o}{.}\PYG{n}{nan}
    \PYG{c+c1}{\PYGZsh{} \CJKCode{相邻点费用差除以减排量，得到元}/kg\CJKCode{的边际减排费用。}}
    \PYG{k}{for} \PYG{n}{i} \PYG{o+ow}{in} \PYG{n+nb}{range}\PYG{p}{(}\PYG{n+nb}{len}\PYG{p}{(}\PYG{n}{result}\PYG{p}{)} \PYG{o}{\PYGZhy{}} \PYG{l+m+mi}{1}\PYG{p}{)}\PYG{p}{:}
        \PYG{n}{delta\PYGZus{}cost} \PYG{o}{=} \PYG{n}{result}\PYG{o}{.}\PYG{n}{loc}\PYG{p}{[}\PYG{n}{i}\PYG{p}{,} \PYG{l+s+s2}{\PYGZdq{}}\PYG{l+s+s2}{objective\PYGZus{}yuan\PYGZus{}per\PYGZus{}year}\PYG{l+s+s2}{\PYGZdq{}}\PYG{p}{]} \PYG{o}{\PYGZhy{}} \PYG{n}{result}\PYG{o}{.}\PYG{n}{loc}\PYG{p}{[}\PYG{n}{i} \PYG{o}{+} \PYG{l+m+mi}{1}\PYG{p}{,} \PYG{l+s+s2}{\PYGZdq{}}\PYG{l+s+s2}{objective\PYGZus{}yuan\PYGZus{}per\PYGZus{}year}\PYG{l+s+s2}{\PYGZdq{}}\PYG{p}{]}
        \PYG{n}{delta\PYGZus{}emission} \PYG{o}{=} \PYG{n}{result}\PYG{o}{.}\PYG{n}{loc}\PYG{p}{[}\PYG{n}{i} \PYG{o}{+} \PYG{l+m+mi}{1}\PYG{p}{,} \PYG{n}{emission\PYGZus{}column}\PYG{p}{]} \PYG{o}{\PYGZhy{}} \PYG{n}{result}\PYG{o}{.}\PYG{n}{loc}\PYG{p}{[}\PYG{n}{i}\PYG{p}{,} \PYG{n}{emission\PYGZus{}column}\PYG{p}{]}
        \PYG{c+c1}{\PYGZsh{} \CJKCode{减排差须超过数值容差，避免用近零分母放大比值。}}
        \PYG{k}{if} \PYG{n}{delta\PYGZus{}emission} \PYG{o}{\PYGZgt{}} \PYG{l+m+mf}{1e\PYGZhy{}6} \PYG{o}{+} \PYG{l+m+mf}{1e\PYGZhy{}8}\PYG{o}{*}\PYG{n+nb}{abs}\PYG{p}{(}\PYG{n}{result}\PYG{o}{.}\PYG{n}{loc}\PYG{p}{[}\PYG{n}{i} \PYG{o}{+} \PYG{l+m+mi}{1}\PYG{p}{,} \PYG{n}{emission\PYGZus{}column}\PYG{p}{]}\PYG{p}{)}\PYG{p}{:}
            \PYG{n}{result}\PYG{o}{.}\PYG{n}{loc}\PYG{p}{[}\PYG{n}{i}\PYG{p}{,} \PYG{l+s+s2}{\PYGZdq{}}\PYG{l+s+s2}{marginal\PYGZus{}abatement\PYGZus{}cost\PYGZus{}yuan\PYGZus{}per\PYGZus{}kg}\PYG{l+s+s2}{\PYGZdq{}}\PYG{p}{]} \PYG{o}{=} \PYG{n}{delta\PYGZus{}cost} \PYG{o}{/} \PYG{n}{delta\PYGZus{}emission}

    \PYG{c+c1}{\PYGZsh{} \CJKCode{理想点距离依赖归一化规则，折中点不是额外物理或工程验收。}}
    \PYG{n}{cost\PYGZus{}span} \PYG{o}{=} \PYG{n}{result}\PYG{p}{[}\PYG{l+s+s2}{\PYGZdq{}}\PYG{l+s+s2}{objective\PYGZus{}yuan\PYGZus{}per\PYGZus{}year}\PYG{l+s+s2}{\PYGZdq{}}\PYG{p}{]}\PYG{o}{.}\PYG{n}{max}\PYG{p}{(}\PYG{p}{)} \PYG{o}{\PYGZhy{}} \PYG{n}{result}\PYG{p}{[}\PYG{l+s+s2}{\PYGZdq{}}\PYG{l+s+s2}{objective\PYGZus{}yuan\PYGZus{}per\PYGZus{}year}\PYG{l+s+s2}{\PYGZdq{}}\PYG{p}{]}\PYG{o}{.}\PYG{n}{min}\PYG{p}{(}\PYG{p}{)}
    \PYG{n}{emission\PYGZus{}span} \PYG{o}{=} \PYG{n}{result}\PYG{p}{[}\PYG{n}{emission\PYGZus{}column}\PYG{p}{]}\PYG{o}{.}\PYG{n}{max}\PYG{p}{(}\PYG{p}{)} \PYG{o}{\PYGZhy{}} \PYG{n}{result}\PYG{p}{[}\PYG{n}{emission\PYGZus{}column}\PYG{p}{]}\PYG{o}{.}\PYG{n}{min}\PYG{p}{(}\PYG{p}{)}
    \PYG{n}{normalized\PYGZus{}cost} \PYG{o}{=} \PYG{p}{(}\PYG{n}{result}\PYG{p}{[}\PYG{l+s+s2}{\PYGZdq{}}\PYG{l+s+s2}{objective\PYGZus{}yuan\PYGZus{}per\PYGZus{}year}\PYG{l+s+s2}{\PYGZdq{}}\PYG{p}{]} \PYG{o}{\PYGZhy{}} \PYG{n}{result}\PYG{p}{[}\PYG{l+s+s2}{\PYGZdq{}}\PYG{l+s+s2}{objective\PYGZus{}yuan\PYGZus{}per\PYGZus{}year}\PYG{l+s+s2}{\PYGZdq{}}\PYG{p}{]}\PYG{o}{.}\PYG{n}{min}\PYG{p}{(}\PYG{p}{)}\PYG{p}{)} \PYG{o}{/} \PYG{n+nb}{max}\PYG{p}{(}\PYG{n}{cost\PYGZus{}span}\PYG{p}{,} \PYG{l+m+mf}{1e\PYGZhy{}9}\PYG{p}{)}
    \PYG{n}{normalized\PYGZus{}emission} \PYG{o}{=} \PYG{p}{(}\PYG{n}{result}\PYG{p}{[}\PYG{n}{emission\PYGZus{}column}\PYG{p}{]} \PYG{o}{\PYGZhy{}} \PYG{n}{result}\PYG{p}{[}\PYG{n}{emission\PYGZus{}column}\PYG{p}{]}\PYG{o}{.}\PYG{n}{min}\PYG{p}{(}\PYG{p}{)}\PYG{p}{)} \PYG{o}{/} \PYG{n+nb}{max}\PYG{p}{(}\PYG{n}{emission\PYGZus{}span}\PYG{p}{,} \PYG{l+m+mf}{1e\PYGZhy{}9}\PYG{p}{)}
    \PYG{n}{result}\PYG{p}{[}\PYG{l+s+s2}{\PYGZdq{}}\PYG{l+s+s2}{distance\PYGZus{}to\PYGZus{}ideal}\PYG{l+s+s2}{\PYGZdq{}}\PYG{p}{]} \PYG{o}{=} \PYG{n}{np}\PYG{o}{.}\PYG{n}{sqrt}\PYG{p}{(}\PYG{n}{normalized\PYGZus{}cost}\PYG{o}{*}\PYG{o}{*}\PYG{l+m+mi}{2} \PYG{o}{+} \PYG{n}{normalized\PYGZus{}emission}\PYG{o}{*}\PYG{o}{*}\PYG{l+m+mi}{2}\PYG{p}{)}
    \PYG{n}{result}\PYG{p}{[}\PYG{l+s+s2}{\PYGZdq{}}\PYG{l+s+s2}{recommended\PYGZus{}compromise}\PYG{l+s+s2}{\PYGZdq{}}\PYG{p}{]} \PYG{o}{=} \PYG{k+kc}{False}
    \PYG{n}{result}\PYG{o}{.}\PYG{n}{loc}\PYG{p}{[}\PYG{n}{result}\PYG{p}{[}\PYG{l+s+s2}{\PYGZdq{}}\PYG{l+s+s2}{distance\PYGZus{}to\PYGZus{}ideal}\PYG{l+s+s2}{\PYGZdq{}}\PYG{p}{]}\PYG{o}{.}\PYG{n}{idxmin}\PYG{p}{(}\PYG{p}{)}\PYG{p}{,} \PYG{l+s+s2}{\PYGZdq{}}\PYG{l+s+s2}{recommended\PYGZus{}compromise}\PYG{l+s+s2}{\PYGZdq{}}\PYG{p}{]} \PYG{o}{=} \PYG{k+kc}{True}
    \PYG{n}{result}\PYG{p}{[}\PYG{l+s+s2}{\PYGZdq{}}\PYG{l+s+s2}{recommendation\PYGZus{}method}\PYG{l+s+s2}{\PYGZdq{}}\PYG{p}{]} \PYG{o}{=} \PYG{l+s+s2}{\PYGZdq{}}\PYG{l+s+s2}{closest\PYGZus{}normalized\PYGZus{}ideal\PYGZus{}point}\PYG{l+s+s2}{\PYGZdq{}}
    \PYG{k}{return} \PYG{n}{result}
\end{Verbatim}

\endgroup

\noindent\textbf{结果核对。}相邻费用差除以减排量得到边际减排费用；推荐折中点由归一化理想点距离选取。投资是否重新优化仍取决于传入配置是否已冻结。

\Needspace{12\baselineskip}
\subsection*{11．代表日硬能量约束可行性诊断}
完整硬约束诊断函数，同时要求能量等式、代表天数守恒和极端日下界；与软权重候选生成分别登记。

\noindent\textbf{输入：}代表日能量矩阵M、目标b、原权重和权重下界（天）。\par
\noindent\textbf{输出：}经过独立复算的可行权重见证，或明确的不可行／求解失败状态。\par

\noindent\textbf{实现思路。}这段代码检查能否同时满足多列能量等式、代表天数守恒和各代表日下界。通过$w=\ell+z$把下界移入右端，求解非负变量$z$的零目标线性规划；能量行按目标量级缩放，以减小数值量级差异。取得候选后恢复原始权重，并在原单位下重新核对能量、总天数和下界。

{\small 来源：\path{representative_adoption.py}，原第48--94行；47行原代码、12行补充注释。}\par

\begingroup
\setstretch{1.0}
\begin{Verbatim}[commandchars=\\\{\},numbers=left,firstnumber=1,stepnumber=1,fontsize=\fontsize{9.5}{11.5}\selectfont,breaklines=true,breakanywhere=true,frame=single,framesep=3mm,numbersep=5pt,rulecolor=\color{black},xleftmargin=18pt]
\PYG{c+c1}{\PYGZsh{} \CJKCode{诊断}Mw=b\CJKCode{是否可行，不是}scenarios.py\CJKCode{中的软能量拟合}SLSQP\CJKCode{。}}
\PYG{k}{def}\PYG{+w}{ }\PYG{n+nf}{diagnose\PYGZus{}exact\PYGZus{}weight\PYGZus{}feasibility}\PYG{p}{(}
    \PYG{n}{design}\PYG{p}{,} \PYG{n}{targets}\PYG{p}{,} \PYG{n}{original\PYGZus{}weights}\PYG{p}{,} \PYG{n}{lower\PYGZus{}bounds}\PYG{p}{,} \PYG{o}{*}\PYG{p}{,} \PYG{n}{total\PYGZus{}weight}\PYG{o}{=}\PYG{k+kc}{None}\PYG{p}{,}
    \PYG{n}{numerical\PYGZus{}tolerance}\PYG{p}{:} \PYG{n+nb}{float} \PYG{o}{=} \PYG{l+m+mf}{1e\PYGZhy{}8}\PYG{p}{,}
\PYG{p}{)} \PYG{o}{\PYGZhy{}}\PYG{o}{\PYGZgt{}} \PYG{n+nb}{dict}\PYG{p}{[}\PYG{n+nb}{str}\PYG{p}{,} \PYG{n}{Any}\PYG{p}{]}\PYG{p}{:}
\PYG{+w}{    }\PYG{l+s+sd}{\PYGZdq{}\PYGZdq{}\PYGZdq{}Diagnose M w = b, sum(w) = W, w \PYGZgt{}= lower using a standard LP.}

\PYG{l+s+sd}{    Feasible weights are a witness only. An infeasible result is explicit;}
\PYG{l+s+sd}{    numerical/solver failures are never reported as a proof of infeasibility.}
\PYG{l+s+sd}{    Rows are scaled to protect the LP from energy\PYGZhy{}unit magnitude differences.}
\PYG{l+s+sd}{    The witness is independently recomputed in original units afterwards.}
\PYG{l+s+sd}{    \PYGZdq{}\PYGZdq{}\PYGZdq{}}
    \PYG{c+c1}{\PYGZsh{} \CJKCode{矩阵列对应代表日，目标为声明能量口径；权重单位为天。}}
    \PYG{n}{matrix}\PYG{p}{,} \PYG{n}{target}\PYG{p}{,} \PYG{n}{initial}\PYG{p}{,} \PYG{n}{lower}\PYG{p}{,} \PYG{n}{total} \PYG{o}{=} \PYG{n}{\PYGZus{}inputs}\PYG{p}{(}\PYG{n}{design}\PYG{p}{,} \PYG{n}{targets}\PYG{p}{,} \PYG{n}{original\PYGZus{}weights}\PYG{p}{,} \PYG{n}{lower\PYGZus{}bounds}\PYG{p}{,} \PYG{n}{total\PYGZus{}weight}\PYG{p}{)}
    \PYG{n}{tolerance} \PYG{o}{=} \PYG{n+nb}{float}\PYG{p}{(}\PYG{n}{numerical\PYGZus{}tolerance}\PYG{p}{)}
    \PYG{k}{if} \PYG{o+ow}{not} \PYG{n}{np}\PYG{o}{.}\PYG{n}{isfinite}\PYG{p}{(}\PYG{n}{tolerance}\PYG{p}{)} \PYG{o+ow}{or} \PYG{n}{tolerance} \PYG{o}{\PYGZlt{}}\PYG{o}{=} \PYG{l+m+mi}{0}\PYG{p}{:}
        \PYG{k}{raise} \PYG{n+ne}{ValueError}\PYG{p}{(}\PYG{l+s+s2}{\PYGZdq{}}\PYG{l+s+s2}{numerical\PYGZus{}tolerance must be positive and finite}\PYG{l+s+s2}{\PYGZdq{}}\PYG{p}{)}
    \PYG{c+c1}{\PYGZsh{} \CJKCode{这项诊断不修改来源数据，也不证明工程运行可行。}}
    \PYG{n}{report} \PYG{o}{=} \PYG{p}{\PYGZob{}}\PYG{l+s+s2}{\PYGZdq{}}\PYG{l+s+s2}{method}\PYG{l+s+s2}{\PYGZdq{}}\PYG{p}{:} \PYG{l+s+s2}{\PYGZdq{}}\PYG{l+s+s2}{standard\PYGZus{}linear\PYGZus{}feasibility\PYGZus{}HiGHS}\PYG{l+s+s2}{\PYGZdq{}}\PYG{p}{,}
              \PYG{l+s+s2}{\PYGZdq{}}\PYG{l+s+s2}{equations}\PYG{l+s+s2}{\PYGZdq{}}\PYG{p}{:} \PYG{l+s+s2}{\PYGZdq{}}\PYG{l+s+s2}{M w = b; sum(w) = W; w \PYGZgt{}= lower}\PYG{l+s+s2}{\PYGZdq{}}\PYG{p}{,}
              \PYG{l+s+s2}{\PYGZdq{}}\PYG{l+s+s2}{total\PYGZus{}weight\PYGZus{}days}\PYG{l+s+s2}{\PYGZdq{}}\PYG{p}{:} \PYG{n}{total}\PYG{p}{,} \PYG{l+s+s2}{\PYGZdq{}}\PYG{l+s+s2}{lower\PYGZus{}bound\PYGZus{}sum\PYGZus{}days}\PYG{l+s+s2}{\PYGZdq{}}\PYG{p}{:} \PYG{n+nb}{float}\PYG{p}{(}\PYG{n}{lower}\PYG{o}{.}\PYG{n}{sum}\PYG{p}{(}\PYG{p}{)}\PYG{p}{)}\PYG{p}{,}
              \PYG{l+s+s2}{\PYGZdq{}}\PYG{l+s+s2}{zero\PYGZus{}target\PYGZus{}rows}\PYG{l+s+s2}{\PYGZdq{}}\PYG{p}{:} \PYG{n}{np}\PYG{o}{.}\PYG{n}{flatnonzero}\PYG{p}{(}\PYG{n}{target} \PYG{o}{==} \PYG{l+m+mi}{0}\PYG{p}{)}\PYG{o}{.}\PYG{n}{tolist}\PYG{p}{(}\PYG{p}{)}\PYG{p}{,}
              \PYG{l+s+s2}{\PYGZdq{}}\PYG{l+s+s2}{source\PYGZus{}data\PYGZus{}modified}\PYG{l+s+s2}{\PYGZdq{}}\PYG{p}{:} \PYG{k+kc}{False}\PYG{p}{,} \PYG{l+s+s2}{\PYGZdq{}}\PYG{l+s+s2}{engineering\PYGZus{}operation\PYGZus{}validated}\PYG{l+s+s2}{\PYGZdq{}}\PYG{p}{:} \PYG{k+kc}{False}\PYG{p}{\PYGZcb{}}
    \PYG{c+c1}{\PYGZsh{} \CJKCode{所有下界之和超过总天数即可判不可行，不需优化求解。}}
    \PYG{k}{if} \PYG{n}{lower}\PYG{o}{.}\PYG{n}{sum}\PYG{p}{(}\PYG{p}{)} \PYG{o}{\PYGZgt{}} \PYG{n}{total} \PYG{o}{+} \PYG{n}{tolerance}\PYG{p}{:}
        \PYG{k}{return} \PYG{p}{\PYGZob{}}\PYG{o}{*}\PYG{o}{*}\PYG{n}{report}\PYG{p}{,} \PYG{l+s+s2}{\PYGZdq{}}\PYG{l+s+s2}{status}\PYG{l+s+s2}{\PYGZdq{}}\PYG{p}{:} \PYG{l+s+s2}{\PYGZdq{}}\PYG{l+s+s2}{infeasible}\PYG{l+s+s2}{\PYGZdq{}}\PYG{p}{,} \PYG{l+s+s2}{\PYGZdq{}}\PYG{l+s+s2}{reason}\PYG{l+s+s2}{\PYGZdq{}}\PYG{p}{:} \PYG{l+s+s2}{\PYGZdq{}}\PYG{l+s+s2}{lower\PYGZus{}bounds\PYGZus{}exceed\PYGZus{}total\PYGZus{}weight}\PYG{l+s+s2}{\PYGZdq{}}\PYG{p}{,} \PYG{l+s+s2}{\PYGZdq{}}\PYG{l+s+s2}{witness\PYGZus{}weights}\PYG{l+s+s2}{\PYGZdq{}}\PYG{p}{:} \PYG{k+kc}{None}\PYG{p}{\PYGZcb{}}
    \PYG{c+c1}{\PYGZsh{} \CJKCode{先按能量量级缩放，并用}w=lower+z\CJKCode{转成非负变量。}}
    \PYG{n}{energy\PYGZus{}scale} \PYG{o}{=} \PYG{n}{np}\PYG{o}{.}\PYG{n}{maximum}\PYG{p}{(}\PYG{n}{np}\PYG{o}{.}\PYG{n}{abs}\PYG{p}{(}\PYG{n}{target}\PYG{p}{)}\PYG{p}{,} \PYG{l+m+mf}{1.0}\PYG{p}{)}
    \PYG{c+c1}{\PYGZsh{} Shift w = lower + z, z \PYGZgt{}= 0, and scale the conservation equation too.}
    \PYG{n}{scale} \PYG{o}{=} \PYG{n}{np}\PYG{o}{.}\PYG{n}{concatenate}\PYG{p}{(}\PYG{p}{(}\PYG{n}{energy\PYGZus{}scale}\PYG{p}{,} \PYG{p}{[}\PYG{n+nb}{max}\PYG{p}{(}\PYG{n}{total}\PYG{p}{,} \PYG{l+m+mf}{1.0}\PYG{p}{)}\PYG{p}{]}\PYG{p}{)}\PYG{p}{)}
    \PYG{c+c1}{\PYGZsh{} \CJKCode{最后一行是总天数守恒，与每列负荷能量等式共同约束。}}
    \PYG{n}{equalities} \PYG{o}{=} \PYG{n}{np}\PYG{o}{.}\PYG{n}{vstack}\PYG{p}{(}\PYG{p}{(}\PYG{n}{matrix}\PYG{p}{,} \PYG{n}{np}\PYG{o}{.}\PYG{n}{ones}\PYG{p}{(}\PYG{n+nb}{len}\PYG{p}{(}\PYG{n}{initial}\PYG{p}{)}\PYG{p}{)}\PYG{p}{)}\PYG{p}{)} \PYG{o}{/} \PYG{n}{scale}\PYG{p}{[}\PYG{p}{:}\PYG{p}{,} \PYG{k+kc}{None}\PYG{p}{]}
    \PYG{n}{rhs} \PYG{o}{=} \PYG{n}{np}\PYG{o}{.}\PYG{n}{concatenate}\PYG{p}{(}\PYG{p}{(}\PYG{n}{target} \PYG{o}{\PYGZhy{}} \PYG{n}{matrix} \PYG{o}{@} \PYG{n}{lower}\PYG{p}{,} \PYG{p}{[}\PYG{n}{total} \PYG{o}{\PYGZhy{}} \PYG{n}{lower}\PYG{o}{.}\PYG{n}{sum}\PYG{p}{(}\PYG{p}{)}\PYG{p}{]}\PYG{p}{)}\PYG{p}{)} \PYG{o}{/} \PYG{n}{scale}
    \PYG{c+c1}{\PYGZsh{} \CJKCode{零目标线性规划只寻找可行见证，能量与总天数均为等式。}}
    \PYG{n}{result} \PYG{o}{=} \PYG{n}{linprog}\PYG{p}{(}\PYG{n}{np}\PYG{o}{.}\PYG{n}{zeros}\PYG{p}{(}\PYG{n+nb}{len}\PYG{p}{(}\PYG{n}{initial}\PYG{p}{)}\PYG{p}{)}\PYG{p}{,} \PYG{n}{A\PYGZus{}eq}\PYG{o}{=}\PYG{n}{equalities}\PYG{p}{,} \PYG{n}{b\PYGZus{}eq}\PYG{o}{=}\PYG{n}{rhs}\PYG{p}{,}
                     \PYG{n}{bounds}\PYG{o}{=}\PYG{p}{[}\PYG{p}{(}\PYG{l+m+mf}{0.0}\PYG{p}{,} \PYG{k+kc}{None}\PYG{p}{)}\PYG{p}{]} \PYG{o}{*} \PYG{n+nb}{len}\PYG{p}{(}\PYG{n}{initial}\PYG{p}{)}\PYG{p}{,} \PYG{n}{method}\PYG{o}{=}\PYG{l+s+s2}{\PYGZdq{}}\PYG{l+s+s2}{highs}\PYG{l+s+s2}{\PYGZdq{}}\PYG{p}{)}
    \PYG{n}{report}\PYG{o}{.}\PYG{n}{update}\PYG{p}{(}\PYG{n}{solver\PYGZus{}status}\PYG{o}{=}\PYG{n+nb}{int}\PYG{p}{(}\PYG{n}{result}\PYG{o}{.}\PYG{n}{status}\PYG{p}{)}\PYG{p}{,} \PYG{n}{solver\PYGZus{}message}\PYG{o}{=}\PYG{n+nb}{str}\PYG{p}{(}\PYG{n}{result}\PYG{o}{.}\PYG{n}{message}\PYG{p}{)}\PYG{p}{)}
    \PYG{c+c1}{\PYGZsh{} \CJKCode{仅明确的}LP\CJKCode{不可行状态作为不可行结论。}}
    \PYG{k}{if} \PYG{n}{result}\PYG{o}{.}\PYG{n}{status} \PYG{o}{==} \PYG{l+m+mi}{2}\PYG{p}{:}
        \PYG{k}{return} \PYG{p}{\PYGZob{}}\PYG{o}{*}\PYG{o}{*}\PYG{n}{report}\PYG{p}{,} \PYG{l+s+s2}{\PYGZdq{}}\PYG{l+s+s2}{status}\PYG{l+s+s2}{\PYGZdq{}}\PYG{p}{:} \PYG{l+s+s2}{\PYGZdq{}}\PYG{l+s+s2}{infeasible}\PYG{l+s+s2}{\PYGZdq{}}\PYG{p}{,} \PYG{l+s+s2}{\PYGZdq{}}\PYG{l+s+s2}{reason}\PYG{l+s+s2}{\PYGZdq{}}\PYG{p}{:} \PYG{l+s+s2}{\PYGZdq{}}\PYG{l+s+s2}{LP\PYGZus{}infeasible}\PYG{l+s+s2}{\PYGZdq{}}\PYG{p}{,} \PYG{l+s+s2}{\PYGZdq{}}\PYG{l+s+s2}{witness\PYGZus{}weights}\PYG{l+s+s2}{\PYGZdq{}}\PYG{p}{:} \PYG{k+kc}{None}\PYG{p}{\PYGZcb{}}
    \PYG{c+c1}{\PYGZsh{} \CJKCode{求解或数值失败单独登记，不冒称数学上已经证伪。}}
    \PYG{k}{if} \PYG{o+ow}{not} \PYG{n}{result}\PYG{o}{.}\PYG{n}{success} \PYG{o+ow}{or} \PYG{n}{result}\PYG{o}{.}\PYG{n}{x} \PYG{o+ow}{is} \PYG{k+kc}{None}\PYG{p}{:}
        \PYG{k}{return} \PYG{p}{\PYGZob{}}\PYG{o}{*}\PYG{o}{*}\PYG{n}{report}\PYG{p}{,} \PYG{l+s+s2}{\PYGZdq{}}\PYG{l+s+s2}{status}\PYG{l+s+s2}{\PYGZdq{}}\PYG{p}{:} \PYG{l+s+s2}{\PYGZdq{}}\PYG{l+s+s2}{solver\PYGZus{}failure}\PYG{l+s+s2}{\PYGZdq{}}\PYG{p}{,} \PYG{l+s+s2}{\PYGZdq{}}\PYG{l+s+s2}{reason}\PYG{l+s+s2}{\PYGZdq{}}\PYG{p}{:} \PYG{l+s+s2}{\PYGZdq{}}\PYG{l+s+s2}{LP\PYGZus{}did\PYGZus{}not\PYGZus{}return\PYGZus{}a\PYGZus{}feasible\PYGZus{}witness}\PYG{l+s+s2}{\PYGZdq{}}\PYG{p}{,} \PYG{l+s+s2}{\PYGZdq{}}\PYG{l+s+s2}{witness\PYGZus{}weights}\PYG{l+s+s2}{\PYGZdq{}}\PYG{p}{:} \PYG{k+kc}{None}\PYG{p}{\PYGZcb{}}
    \PYG{c+c1}{\PYGZsh{} \CJKCode{恢复原单位权重并重新计算能量、天数与下界残差。}}
    \PYG{n}{witness} \PYG{o}{=} \PYG{n}{lower} \PYG{o}{+} \PYG{n}{np}\PYG{o}{.}\PYG{n}{asarray}\PYG{p}{(}\PYG{n}{result}\PYG{o}{.}\PYG{n}{x}\PYG{p}{,} \PYG{n}{dtype}\PYG{o}{=}\PYG{n+nb}{float}\PYG{p}{)}
    \PYG{c+c1}{\PYGZsh{} \CJKCode{见证恢复后使用原矩阵复算，不只信任缩放后的求解值。}}
    \PYG{n}{energy\PYGZus{}residual} \PYG{o}{=} \PYG{n}{matrix} \PYG{o}{@} \PYG{n}{witness} \PYG{o}{\PYGZhy{}} \PYG{n}{target}
    \PYG{n}{conservation} \PYG{o}{=} \PYG{n+nb}{abs}\PYG{p}{(}\PYG{n+nb}{float}\PYG{p}{(}\PYG{n}{witness}\PYG{o}{.}\PYG{n}{sum}\PYG{p}{(}\PYG{p}{)} \PYG{o}{\PYGZhy{}} \PYG{n}{total}\PYG{p}{)}\PYG{p}{)}
    \PYG{n}{bound\PYGZus{}violation} \PYG{o}{=} \PYG{n+nb}{float}\PYG{p}{(}\PYG{n}{np}\PYG{o}{.}\PYG{n}{maximum}\PYG{p}{(}\PYG{n}{lower} \PYG{o}{\PYGZhy{}} \PYG{n}{witness}\PYG{p}{,} \PYG{l+m+mi}{0}\PYG{p}{)}\PYG{o}{.}\PYG{n}{max}\PYG{p}{(}\PYG{p}{)}\PYG{p}{)}
    \PYG{n}{normalized\PYGZus{}residual} \PYG{o}{=} \PYG{n+nb}{float}\PYG{p}{(}\PYG{n}{np}\PYG{o}{.}\PYG{n}{max}\PYG{p}{(}\PYG{n}{np}\PYG{o}{.}\PYG{n}{abs}\PYG{p}{(}\PYG{n}{energy\PYGZus{}residual}\PYG{p}{)} \PYG{o}{/} \PYG{n}{energy\PYGZus{}scale}\PYG{p}{)}\PYG{p}{)}
    \PYG{c+c1}{\PYGZsh{} \CJKCode{见证必须通过独立有限性和残差检查；通过也不自动代表已采用。}}
    \PYG{n}{verified} \PYG{o}{=} \PYG{p}{(}\PYG{n}{np}\PYG{o}{.}\PYG{n}{isfinite}\PYG{p}{(}\PYG{n}{witness}\PYG{p}{)}\PYG{o}{.}\PYG{n}{all}\PYG{p}{(}\PYG{p}{)} \PYG{o+ow}{and} \PYG{n}{normalized\PYGZus{}residual} \PYG{o}{\PYGZlt{}}\PYG{o}{=} \PYG{n}{tolerance}
                \PYG{o+ow}{and} \PYG{n}{conservation} \PYG{o}{\PYGZlt{}}\PYG{o}{=} \PYG{n}{tolerance} \PYG{o}{*} \PYG{n+nb}{max}\PYG{p}{(}\PYG{n}{total}\PYG{p}{,} \PYG{l+m+mf}{1.0}\PYG{p}{)} \PYG{o+ow}{and} \PYG{n}{bound\PYGZus{}violation} \PYG{o}{\PYGZlt{}}\PYG{o}{=} \PYG{n}{tolerance}\PYG{p}{)}
    \PYG{n}{report}\PYG{o}{.}\PYG{n}{update}\PYG{p}{(}\PYG{n}{status}\PYG{o}{=}\PYG{l+s+s2}{\PYGZdq{}}\PYG{l+s+s2}{feasible}\PYG{l+s+s2}{\PYGZdq{}} \PYG{k}{if} \PYG{n}{verified} \PYG{k}{else} \PYG{l+s+s2}{\PYGZdq{}}\PYG{l+s+s2}{solver\PYGZus{}failure}\PYG{l+s+s2}{\PYGZdq{}}\PYG{p}{,}
                  \PYG{n}{reason}\PYG{o}{=}\PYG{l+s+s2}{\PYGZdq{}}\PYG{l+s+s2}{verified\PYGZus{}LP\PYGZus{}witness}\PYG{l+s+s2}{\PYGZdq{}} \PYG{k}{if} \PYG{n}{verified} \PYG{k}{else} \PYG{l+s+s2}{\PYGZdq{}}\PYG{l+s+s2}{LP\PYGZus{}witness\PYGZus{}failed\PYGZus{}independent\PYGZus{}residual\PYGZus{}check}\PYG{l+s+s2}{\PYGZdq{}}\PYG{p}{,}
                  \PYG{n}{witness\PYGZus{}weights}\PYG{o}{=}\PYG{n}{witness}\PYG{o}{.}\PYG{n}{tolist}\PYG{p}{(}\PYG{p}{)}\PYG{p}{,} \PYG{n}{energy\PYGZus{}residual\PYGZus{}original\PYGZus{}units}\PYG{o}{=}\PYG{n}{energy\PYGZus{}residual}\PYG{o}{.}\PYG{n}{tolist}\PYG{p}{(}\PYG{p}{)}\PYG{p}{,}
                  \PYG{n}{maximum\PYGZus{}normalized\PYGZus{}energy\PYGZus{}residual}\PYG{o}{=}\PYG{n}{normalized\PYGZus{}residual}\PYG{p}{,}
                  \PYG{n}{conservation\PYGZus{}residual\PYGZus{}days}\PYG{o}{=}\PYG{n}{conservation}\PYG{p}{,} \PYG{n}{maximum\PYGZus{}lower\PYGZus{}bound\PYGZus{}violation}\PYG{o}{=}\PYG{n}{bound\PYGZus{}violation}\PYG{p}{)}
    \PYG{k}{return} \PYG{n}{report}
\end{Verbatim}

\endgroup

\noindent\textbf{结果核对。}该函数给出可行权重见证，不负责最小化分布距离。其\path{status}区分可行、不可行及求解失败；最小$W_1$的具体优化实现见后面的新增代码。

\Needspace{12\baselineskip}
\subsection*{12．有限模式管网反馈与失败处理}
完整有限目录反馈函数。适用域为登记的独立泵送星形支路；发现不合格模式后加禁用约束并重求。

\noindent\textbf{输入：}支持范围内的配置、连续时序、求解预算、轮数上限和可选逐轮记录回调。\par
\noindent\textbf{输出：}通过硬校核的候选及逐轮历史；失败则随异常保留原因和已有割约束。\par

\noindent\textbf{实现思路。}反馈循环在登记的有限模式目录内工作：求得候选后，独立检查其水力、温度与服务条件；将新增不合格模式在全部时段的选择变量固定为零，再求下一轮。已排除模式集合防止重复加割；通过、异常、无法产生新割和轮数耗尽都有独立出口，逐轮回调可保存问题定位所需的中间结果。

{\small 来源：\path{network_feedback.py}，原第440--481行；42行原代码、15行补充注释。}\par

\begingroup
\setstretch{1.0}
\begin{Verbatim}[commandchars=\\\{\},numbers=left,firstnumber=1,stepnumber=1,fontsize=\fontsize{9.5}{11.5}\selectfont,breaklines=true,breakanywhere=true,frame=single,framesep=3mm,numbersep=5pt,rulecolor=\color{black},xleftmargin=18pt]
\PYG{c+c1}{\PYGZsh{} \CJKCode{此接口不是三城聚合网络的通用连续非线性水力优化器。}}
\PYG{k}{def}\PYG{+w}{ }\PYG{n+nf}{solve\PYGZus{}with\PYGZus{}network\PYGZus{}feedback}\PYG{p}{(}\PYG{n}{config}\PYG{p}{,} \PYG{n}{timeseries}\PYG{p}{,} \PYG{o}{*}\PYG{p}{,} \PYG{n}{mip\PYGZus{}gap}\PYG{o}{=}\PYG{l+m+mf}{0.}\PYG{p}{,} \PYG{n}{time\PYGZus{}limit}\PYG{o}{=}\PYG{l+m+mf}{60.}\PYG{p}{,}
                                \PYG{n}{max\PYGZus{}rounds}\PYG{o}{=}\PYG{k+kc}{None}\PYG{p}{,} \PYG{n}{on\PYGZus{}iteration}\PYG{p}{:} \PYG{n}{Callable} \PYG{o}{|} \PYG{k+kc}{None} \PYG{o}{=} \PYG{k+kc}{None}\PYG{p}{)}\PYG{p}{:}
    \PYG{c+c1}{\PYGZsh{} \CJKCode{模式目录和适用配置在构模入口先验证，禁用模式必须来自该有限集合。}}
    \PYG{n}{model}\PYG{p}{,} \PYG{n}{catalog}\PYG{p}{,} \PYG{n}{flags} \PYG{o}{=} \PYG{n}{build\PYGZus{}feedback\PYGZus{}model}\PYG{p}{(}\PYG{n}{config}\PYG{p}{,} \PYG{n}{timeseries}\PYG{p}{)}
    \PYG{c+c1}{\PYGZsh{} \CJKCode{默认轮数有限；不会承诺任意网络的全局收敛。}}
    \PYG{k}{if} \PYG{n}{max\PYGZus{}rounds} \PYG{o+ow}{is} \PYG{k+kc}{None}\PYG{p}{:}
        \PYG{n}{max\PYGZus{}rounds} \PYG{o}{=} \PYG{n+nb}{len}\PYG{p}{(}\PYG{n}{catalog}\PYG{p}{)} \PYG{o}{+} \PYG{l+m+mi}{1}
    \PYG{c+c1}{\PYGZsh{} \CJKCode{拒绝无效轮数、时间预算和间隙。}}
    \PYG{k}{if} \PYG{n+nb}{isinstance}\PYG{p}{(}\PYG{n}{max\PYGZus{}rounds}\PYG{p}{,} \PYG{n+nb}{bool}\PYG{p}{)} \PYG{o+ow}{or} \PYG{n+nb}{int}\PYG{p}{(}\PYG{n}{max\PYGZus{}rounds}\PYG{p}{)} \PYG{o}{!=} \PYG{n}{max\PYGZus{}rounds} \PYG{o+ow}{or} \PYG{n}{max\PYGZus{}rounds} \PYG{o}{\PYGZlt{}} \PYG{l+m+mi}{1}\PYG{p}{:}
        \PYG{k}{raise} \PYG{n+ne}{ValueError}\PYG{p}{(}\PYG{l+s+s2}{\PYGZdq{}}\PYG{l+s+s2}{max\PYGZus{}rounds must be a positive integer}\PYG{l+s+s2}{\PYGZdq{}}\PYG{p}{)}
    \PYG{k}{if} \PYG{o+ow}{not} \PYG{n}{math}\PYG{o}{.}\PYG{n}{isfinite}\PYG{p}{(}\PYG{n}{time\PYGZus{}limit}\PYG{p}{)} \PYG{o+ow}{or} \PYG{n}{time\PYGZus{}limit} \PYG{o}{\PYGZlt{}}\PYG{o}{=} \PYG{l+m+mi}{0} \PYG{o+ow}{or} \PYG{o+ow}{not} \PYG{n}{math}\PYG{o}{.}\PYG{n}{isfinite}\PYG{p}{(}\PYG{n}{mip\PYGZus{}gap}\PYG{p}{)} \PYG{o+ow}{or} \PYG{o+ow}{not} \PYG{l+m+mi}{0} \PYG{o}{\PYGZlt{}}\PYG{o}{=} \PYG{n}{mip\PYGZus{}gap} \PYG{o}{\PYGZlt{}} \PYG{l+m+mi}{1}\PYG{p}{:}
        \PYG{k}{raise} \PYG{n+ne}{ValueError}\PYG{p}{(}\PYG{l+s+s2}{\PYGZdq{}}\PYG{l+s+s2}{invalid solve budget/gap}\PYG{l+s+s2}{\PYGZdq{}}\PYG{p}{)}
    \PYG{c+c1}{\PYGZsh{} \CJKCode{逐轮历史及已排除模式单独记录，便于辨识失败发生在哪一步。}}
    \PYG{n}{history}\PYG{p}{,} \PYG{n}{eliminated} \PYG{o}{=} \PYG{p}{[}\PYG{p}{]}\PYG{p}{,} \PYG{n+nb}{set}\PYG{p}{(}\PYG{p}{)}
    \PYG{c+c1}{\PYGZsh{} \CJKCode{每轮在已有割约束下求解，然后接受独立模式校核。}}
    \PYG{k}{for} \PYG{n}{iteration} \PYG{o+ow}{in} \PYG{n+nb}{range}\PYG{p}{(}\PYG{n+nb}{int}\PYG{p}{(}\PYG{n}{max\PYGZus{}rounds}\PYG{p}{)}\PYG{p}{)}\PYG{p}{:}
        \PYG{k}{try}\PYG{p}{:}
            \PYG{c+c1}{\PYGZsh{} \CJKCode{费用最小可行候选仍需通过温度、水力及服务范围检查。}}
            \PYG{n}{solution} \PYG{o}{=} \PYG{n}{model}\PYG{o}{.}\PYG{n}{solve}\PYG{p}{(}\PYG{n}{time\PYGZus{}limit}\PYG{o}{=}\PYG{n}{time\PYGZus{}limit}\PYG{p}{,} \PYG{n}{mip\PYGZus{}gap}\PYG{o}{=}\PYG{n}{mip\PYGZus{}gap}\PYG{p}{)}
            \PYG{n}{dispatch}\PYG{p}{,} \PYG{n}{audit} \PYG{o}{=} \PYG{n}{audit\PYGZus{}feedback\PYGZus{}solution}\PYG{p}{(}\PYG{n}{config}\PYG{p}{,} \PYG{n}{model}\PYG{p}{,} \PYG{n}{solution}\PYG{p}{,} \PYG{n}{catalog}\PYG{p}{)}
        \PYG{c+c1}{\PYGZsh{} \CJKCode{求解或审计异常保存至}history\CJKCode{，再抛出}FeedbackFailure\CJKCode{。}}
        \PYG{k}{except} \PYG{n+ne}{Exception} \PYG{k}{as} \PYG{n}{error}\PYG{p}{:}
            \PYG{n}{record} \PYG{o}{=} \PYG{p}{\PYGZob{}}\PYG{l+s+s2}{\PYGZdq{}}\PYG{l+s+s2}{iteration}\PYG{l+s+s2}{\PYGZdq{}}\PYG{p}{:} \PYG{n}{iteration}\PYG{p}{,} \PYG{l+s+s2}{\PYGZdq{}}\PYG{l+s+s2}{status}\PYG{l+s+s2}{\PYGZdq{}}\PYG{p}{:} \PYG{l+s+s2}{\PYGZdq{}}\PYG{l+s+s2}{failed}\PYG{l+s+s2}{\PYGZdq{}}\PYG{p}{,} \PYG{l+s+s2}{\PYGZdq{}}\PYG{l+s+s2}{type}\PYG{l+s+s2}{\PYGZdq{}}\PYG{p}{:} \PYG{n+nb}{type}\PYG{p}{(}\PYG{n}{error}\PYG{p}{)}\PYG{o}{.}\PYG{n+nv+vm}{\PYGZus{}\PYGZus{}name\PYGZus{}\PYGZus{}}\PYG{p}{,} \PYG{l+s+s2}{\PYGZdq{}}\PYG{l+s+s2}{reason}\PYG{l+s+s2}{\PYGZdq{}}\PYG{p}{:} \PYG{n+nb}{str}\PYG{p}{(}\PYG{n}{error}\PYG{p}{)}\PYG{p}{,}
                      \PYG{l+s+s2}{\PYGZdq{}}\PYG{l+s+s2}{solver\PYGZus{}evidence}\PYG{l+s+s2}{\PYGZdq{}}\PYG{p}{:} \PYG{n+nb}{getattr}\PYG{p}{(}\PYG{n}{error}\PYG{p}{,} \PYG{l+s+s2}{\PYGZdq{}}\PYG{l+s+s2}{summary}\PYG{l+s+s2}{\PYGZdq{}}\PYG{p}{,} \PYG{k+kc}{None}\PYG{p}{)}\PYG{p}{,} \PYG{l+s+s2}{\PYGZdq{}}\PYG{l+s+s2}{previous\PYGZus{}cuts}\PYG{l+s+s2}{\PYGZdq{}}\PYG{p}{:} \PYG{n+nb}{sorted}\PYG{p}{(}\PYG{n}{eliminated}\PYG{p}{)}\PYG{p}{\PYGZcb{}}
            \PYG{n}{history}\PYG{o}{.}\PYG{n}{append}\PYG{p}{(}\PYG{n}{record}\PYG{p}{)}
            \PYG{c+c1}{\PYGZsh{} \CJKCode{回调可逐轮落盘，让拒绝方案与失败原因可追溯。}}
            \PYG{k}{if} \PYG{n}{on\PYGZus{}iteration}\PYG{p}{:}
                \PYG{n}{on\PYGZus{}iteration}\PYG{p}{(}\PYG{n}{record}\PYG{p}{,} \PYG{k+kc}{None}\PYG{p}{,} \PYG{k+kc}{None}\PYG{p}{)}
            \PYG{k}{raise} \PYG{n}{FeedbackFailure}\PYG{p}{(}\PYG{n+nb}{str}\PYG{p}{(}\PYG{n}{error}\PYG{p}{)}\PYG{p}{,} \PYG{n}{history}\PYG{p}{)} \PYG{k+kn}{from}\PYG{+w}{ }\PYG{n+nn}{error}
        \PYG{c+c1}{\PYGZsh{} invalid\PYGZus{}modes\CJKCode{来自独立硬校核，不由软惩罚值直接替代。}}
        \PYG{n}{invalid} \PYG{o}{=} \PYG{n+nb}{set}\PYG{p}{(}\PYG{n}{audit}\PYG{p}{[}\PYG{l+s+s2}{\PYGZdq{}}\PYG{l+s+s2}{invalid\PYGZus{}modes}\PYG{l+s+s2}{\PYGZdq{}}\PYG{p}{]}\PYG{p}{)}
        \PYG{n}{record} \PYG{o}{=} \PYG{p}{\PYGZob{}}\PYG{l+s+s2}{\PYGZdq{}}\PYG{l+s+s2}{iteration}\PYG{l+s+s2}{\PYGZdq{}}\PYG{p}{:} \PYG{n}{iteration}\PYG{p}{,} \PYG{l+s+s2}{\PYGZdq{}}\PYG{l+s+s2}{status}\PYG{l+s+s2}{\PYGZdq{}}\PYG{p}{:} \PYG{l+s+s2}{\PYGZdq{}}\PYG{l+s+s2}{accepted}\PYG{l+s+s2}{\PYGZdq{}} \PYG{k}{if} \PYG{n}{audit}\PYG{p}{[}\PYG{l+s+s2}{\PYGZdq{}}\PYG{l+s+s2}{passed}\PYG{l+s+s2}{\PYGZdq{}}\PYG{p}{]} \PYG{k}{else} \PYG{l+s+s2}{\PYGZdq{}}\PYG{l+s+s2}{rejected}\PYG{l+s+s2}{\PYGZdq{}}\PYG{p}{,}
            \PYG{l+s+s2}{\PYGZdq{}}\PYG{l+s+s2}{objective\PYGZus{}yuan}\PYG{l+s+s2}{\PYGZdq{}}\PYG{p}{:} \PYG{n}{solution}\PYG{o}{.}\PYG{n}{summary}\PYG{p}{[}\PYG{l+s+s2}{\PYGZdq{}}\PYG{l+s+s2}{objective\PYGZus{}yuan\PYGZus{}per\PYGZus{}year}\PYG{l+s+s2}{\PYGZdq{}}\PYG{p}{]}\PYG{p}{,}
            \PYG{l+s+s2}{\PYGZdq{}}\PYG{l+s+s2}{solver\PYGZus{}status}\PYG{l+s+s2}{\PYGZdq{}}\PYG{p}{:} \PYG{n}{solution}\PYG{o}{.}\PYG{n}{summary}\PYG{p}{[}\PYG{l+s+s2}{\PYGZdq{}}\PYG{l+s+s2}{solver\PYGZus{}status}\PYG{l+s+s2}{\PYGZdq{}}\PYG{p}{]}\PYG{p}{,} \PYG{l+s+s2}{\PYGZdq{}}\PYG{l+s+s2}{mip\PYGZus{}gap}\PYG{l+s+s2}{\PYGZdq{}}\PYG{p}{:} \PYG{n}{solution}\PYG{o}{.}\PYG{n}{summary}\PYG{o}{.}\PYG{n}{get}\PYG{p}{(}\PYG{l+s+s2}{\PYGZdq{}}\PYG{l+s+s2}{mip\PYGZus{}gap}\PYG{l+s+s2}{\PYGZdq{}}\PYG{p}{)}\PYG{p}{,}
            \PYG{l+s+s2}{\PYGZdq{}}\PYG{l+s+s2}{solution\PYGZus{}quality}\PYG{l+s+s2}{\PYGZdq{}}\PYG{p}{:} \PYG{n}{solution}\PYG{o}{.}\PYG{n}{summary}\PYG{p}{[}\PYG{l+s+s2}{\PYGZdq{}}\PYG{l+s+s2}{solution\PYGZus{}quality}\PYG{l+s+s2}{\PYGZdq{}}\PYG{p}{]}\PYG{p}{,} \PYG{l+s+s2}{\PYGZdq{}}\PYG{l+s+s2}{audit}\PYG{l+s+s2}{\PYGZdq{}}\PYG{p}{:} \PYG{n}{audit}\PYG{p}{,}
            \PYG{l+s+s2}{\PYGZdq{}}\PYG{l+s+s2}{added\PYGZus{}mode\PYGZus{}cuts}\PYG{l+s+s2}{\PYGZdq{}}\PYG{p}{:} \PYG{n+nb}{sorted}\PYG{p}{(}\PYG{n}{invalid} \PYG{o}{\PYGZhy{}} \PYG{n}{eliminated}\PYG{p}{)}\PYG{p}{,} \PYG{l+s+s2}{\PYGZdq{}}\PYG{l+s+s2}{previous\PYGZus{}cuts}\PYG{l+s+s2}{\PYGZdq{}}\PYG{p}{:} \PYG{n+nb}{sorted}\PYG{p}{(}\PYG{n}{eliminated}\PYG{p}{)}\PYG{p}{\PYGZcb{}}
        \PYG{n}{history}\PYG{o}{.}\PYG{n}{append}\PYG{p}{(}\PYG{n}{record}\PYG{p}{)}
        \PYG{k}{if} \PYG{n}{on\PYGZus{}iteration}\PYG{p}{:}
            \PYG{n}{on\PYGZus{}iteration}\PYG{p}{(}\PYG{n}{record}\PYG{p}{,} \PYG{n}{solution}\PYG{p}{,} \PYG{n}{dispatch}\PYG{p}{)}
        \PYG{c+c1}{\PYGZsh{} \CJKCode{全部校核通过才返回；最优性范围限于有限目录和登记容差。}}
        \PYG{k}{if} \PYG{n}{audit}\PYG{p}{[}\PYG{l+s+s2}{\PYGZdq{}}\PYG{l+s+s2}{passed}\PYG{l+s+s2}{\PYGZdq{}}\PYG{p}{]}\PYG{p}{:}
            \PYG{n}{solution}\PYG{o}{.}\PYG{n}{summary}\PYG{p}{[}\PYG{l+s+s2}{\PYGZdq{}}\PYG{l+s+s2}{network\PYGZus{}feedback}\PYG{l+s+s2}{\PYGZdq{}}\PYG{p}{]} \PYG{o}{=} \PYG{p}{\PYGZob{}}\PYG{l+s+s2}{\PYGZdq{}}\PYG{l+s+s2}{iterations}\PYG{l+s+s2}{\PYGZdq{}}\PYG{p}{:} \PYG{n}{iteration} \PYG{o}{+} \PYG{l+m+mi}{1}\PYG{p}{,} \PYG{l+s+s2}{\PYGZdq{}}\PYG{l+s+s2}{eliminated\PYGZus{}modes}\PYG{l+s+s2}{\PYGZdq{}}\PYG{p}{:} \PYG{n+nb}{sorted}\PYG{p}{(}\PYG{n}{eliminated}\PYG{p}{)}\PYG{p}{,}
                \PYG{l+s+s2}{\PYGZdq{}}\PYG{l+s+s2}{audit}\PYG{l+s+s2}{\PYGZdq{}}\PYG{p}{:} \PYG{n}{audit}\PYG{p}{,} \PYG{l+s+s2}{\PYGZdq{}}\PYG{l+s+s2}{optimality\PYGZus{}domain}\PYG{l+s+s2}{\PYGZdq{}}\PYG{p}{:} \PYG{l+s+s2}{\PYGZdq{}}\PYG{l+s+s2}{finite catalog MILP within recorded solver tolerances; not continuous nonlinear network optimum}\PYG{l+s+s2}{\PYGZdq{}}\PYG{p}{,}
                \PYG{l+s+s2}{\PYGZdq{}}\PYG{l+s+s2}{soft\PYGZus{}penalty\PYGZus{}is\PYGZus{}not\PYGZus{}cash\PYGZus{}cost}\PYG{l+s+s2}{\PYGZdq{}}\PYG{p}{:} \PYG{k+kc}{True}\PYG{p}{\PYGZcb{}}
            \PYG{k}{return} \PYG{n}{FeedbackArtifacts}\PYG{p}{(}\PYG{n}{model}\PYG{p}{,} \PYG{n}{solution}\PYG{p}{,} \PYG{n}{dispatch}\PYG{p}{,} \PYG{n}{history}\PYG{p}{,} \PYG{n}{audit}\PYG{p}{,} \PYG{n}{catalog}\PYG{p}{)}
        \PYG{c+c1}{\PYGZsh{} \CJKCode{没有可新增禁用模式时明确失败，避免重复无效求解。}}
        \PYG{k}{if} \PYG{o+ow}{not} \PYG{n}{invalid} \PYG{o}{\PYGZhy{}} \PYG{n}{eliminated}\PYG{p}{:}
            \PYG{k}{raise} \PYG{n}{FeedbackFailure}\PYG{p}{(}\PYG{l+s+s2}{\PYGZdq{}}\PYG{l+s+s2}{non\PYGZhy{}engineering audit failure; no valid mode cut available}\PYG{l+s+s2}{\PYGZdq{}}\PYG{p}{,} \PYG{n}{history}\PYG{p}{)}
        \PYG{c+c1}{\PYGZsh{} \CJKCode{对每个时段将不合格模式的选择变量置零，再重新求解。}}
        \PYG{k}{for} \PYG{n}{virtual} \PYG{o+ow}{in} \PYG{n+nb}{sorted}\PYG{p}{(}\PYG{n}{invalid} \PYG{o}{\PYGZhy{}} \PYG{n}{eliminated}\PYG{p}{)}\PYG{p}{:}
            \PYG{k}{for} \PYG{n}{t} \PYG{o+ow}{in} \PYG{n}{model}\PYG{o}{.}\PYG{n}{times}\PYG{p}{:}
                \PYG{n}{model}\PYG{o}{.}\PYG{n}{problem}\PYG{o}{.}\PYG{n}{add\PYGZus{}constraint}\PYG{p}{(}\PYG{p}{\PYGZob{}}\PYG{n}{flags}\PYG{p}{[}\PYG{n}{virtual}\PYG{p}{,} \PYG{n}{t}\PYG{p}{]}\PYG{p}{:} \PYG{l+m+mf}{1.}\PYG{p}{\PYGZcb{}}\PYG{p}{,} \PYG{n}{ub}\PYG{o}{=}\PYG{l+m+mf}{0.}\PYG{p}{)}
        \PYG{c+c1}{\PYGZsh{} \CJKCode{登记已排除集合，下轮只对首次发现的不合格模式加割。}}
        \PYG{n}{eliminated}\PYG{o}{.}\PYG{n}{update}\PYG{p}{(}\PYG{n}{invalid}\PYG{p}{)}
    \PYG{c+c1}{\PYGZsh{} \CJKCode{轮数耗尽仍未验收则失败，不把最后一轮写成合格方案。}}
    \PYG{k}{raise} \PYG{n}{FeedbackFailure}\PYG{p}{(}\PYG{l+s+s2}{\PYGZdq{}}\PYG{l+s+s2}{feedback iteration limit reached without engineering acceptance}\PYG{l+s+s2}{\PYGZdq{}}\PYG{p}{,} \PYG{n}{history}\PYG{p}{)}
\end{Verbatim}

\endgroup

\noindent\textbf{结果核对。}比较每轮\path{audit}、\path{added_mode_cuts}与\path{previous_cuts}，确认新割有物理依据。最终接收条件是硬校核通过，软惩罚只参与目标。

\Needspace{12\baselineskip}
\subsection*{13．热源出力的上游余热共用约束}
完整的共享上游热资源装配函数。把直供换热器与热泵已有出力变量共同映射到同一资源池，约束逐时实际取热，而不是让每台设备分别重复使用全部余热。

\noindent\textbf{输入：}明确登记的资源池名称、成员热源、可用热功率常量或时序列（kW）、传热效率（无量纲），以及热源逐温位出力变量与逐时COP。\par
\noindent\textbf{输出：}共享资源取热变量\path{shared_heat_resource_use}（kW）、等式连接及可用功率上界；资源行同时登记以供结果核验。\par

\noindent\textbf{实现思路。}本函数把正文热源出力$Q_{s,g,t}$（kW）折算为资源池上游取热。直供项为$Q/\eta$，热泵项为$Q(1-1/\mathrm{COP})/\eta$；$\eta$为无量纲传热效率，压缩机电功不算余热。程序用取热变量与全部成员、全部温位的折算和建立等式，再以逐时可用热功率限制其上界，因此换热器与热泵不能重复占用同一资源。该约束衔接已有热源出力和节点热平衡，仅在配置登记资源池时启用。

{\small 来源：\path{physical_extensions.py}，原第25--79行；55行原代码、14行补充注释。}\par

\begingroup
\setstretch{1.0}
\begin{Verbatim}[commandchars=\\\{\},numbers=left,firstnumber=1,stepnumber=1,fontsize=\fontsize{9.5}{11.5}\selectfont,breaklines=true,breakanywhere=true,frame=single,framesep=3mm,numbersep=5pt,rulecolor=\color{black},xleftmargin=18pt]
\PYG{c+c1}{\PYGZsh{} \CJKCode{本函数使用已注册热源出力，不新增热源设备。}}
\PYG{k}{def}\PYG{+w}{ }\PYG{n+nf}{add\PYGZus{}shared\PYGZus{}heat\PYGZus{}resources}\PYG{p}{(}\PYG{n}{model}\PYG{p}{)}\PYG{p}{:}
\PYG{+w}{    }\PYG{l+s+sd}{\PYGZdq{}\PYGZdq{}\PYGZdq{}Bound upstream heat drawn jointly by direct exchangers and heat pumps.}

\PYG{l+s+sd}{    Q\PYGZus{}upstream = Q\PYGZus{}direct / eta + Q\PYGZus{}HP * (1 \PYGZhy{} 1/COP) / eta.}
\PYG{l+s+sd}{    eta is the explicitly declared resource\PYGZhy{}to\PYGZhy{}device transfer efficiency.}
\PYG{l+s+sd}{    Compressor electricity is excluded. Cooling/rejection is not allocated to}
\PYG{l+s+sd}{    this pool; such sources require a separate, explicitly defined resource model.}
\PYG{l+s+sd}{    \PYGZdq{}\PYGZdq{}\PYGZdq{}}
    \PYG{c+c1}{\PYGZsh{} \CJKCode{资源行保留可用上界、出力系数与取热变量。}}
    \PYG{n}{model}\PYG{o}{.}\PYG{n}{shared\PYGZus{}heat\PYGZus{}resource\PYGZus{}rows} \PYG{o}{=} \PYG{p}{[}\PYG{p}{]}
    \PYG{n}{names}\PYG{p}{,} \PYG{n}{assigned} \PYG{o}{=} \PYG{n+nb}{set}\PYG{p}{(}\PYG{p}{)}\PYG{p}{,} \PYG{n+nb}{set}\PYG{p}{(}\PYG{p}{)}
    \PYG{c+c1}{\PYGZsh{} \CJKCode{只有配置登记的资源池才会增加这组约束。}}
    \PYG{k}{for} \PYG{n}{pool} \PYG{o+ow}{in} \PYG{n}{model}\PYG{o}{.}\PYG{n}{config}\PYG{o}{.}\PYG{n}{get}\PYG{p}{(}\PYG{l+s+s2}{\PYGZdq{}}\PYG{l+s+s2}{shared\PYGZus{}heat\PYGZus{}resources}\PYG{l+s+s2}{\PYGZdq{}}\PYG{p}{,} \PYG{p}{[}\PYG{p}{]}\PYG{p}{)}\PYG{p}{:}
        \PYG{n}{name} \PYG{o}{=} \PYG{n+nb}{str}\PYG{p}{(}\PYG{n}{pool}\PYG{p}{[}\PYG{l+s+s2}{\PYGZdq{}}\PYG{l+s+s2}{name}\PYG{l+s+s2}{\PYGZdq{}}\PYG{p}{]}\PYG{p}{)}
        \PYG{n}{members} \PYG{o}{=} \PYG{n+nb}{list}\PYG{p}{(}\PYG{n}{pool}\PYG{p}{[}\PYG{l+s+s2}{\PYGZdq{}}\PYG{l+s+s2}{source\PYGZus{}names}\PYG{l+s+s2}{\PYGZdq{}}\PYG{p}{]}\PYG{p}{)}
        \PYG{c+c1}{\PYGZsh{} \CJKCode{资源名称和成员须唯一，避免资源账本重复。}}
        \PYG{k}{if} \PYG{o+ow}{not} \PYG{n}{name}\PYG{o}{.}\PYG{n}{strip}\PYG{p}{(}\PYG{p}{)} \PYG{o+ow}{or} \PYG{n}{name} \PYG{o+ow}{in} \PYG{n}{names} \PYG{o+ow}{or} \PYG{o+ow}{not} \PYG{n}{members} \PYG{o+ow}{or} \PYG{n+nb}{len}\PYG{p}{(}\PYG{n}{members}\PYG{p}{)} \PYG{o}{!=} \PYG{n+nb}{len}\PYG{p}{(}\PYG{n+nb}{set}\PYG{p}{(}\PYG{n}{members}\PYG{p}{)}\PYG{p}{)}\PYG{p}{:}
            \PYG{k}{raise} \PYG{n+ne}{ValueError}\PYG{p}{(}\PYG{l+s+s2}{\PYGZdq{}}\PYG{l+s+s2}{shared heat resource names/members must be unique and nonempty}\PYG{l+s+s2}{\PYGZdq{}}\PYG{p}{)}
        \PYG{n}{names}\PYG{o}{.}\PYG{n}{add}\PYG{p}{(}\PYG{n}{name}\PYG{p}{)}
        \PYG{c+c1}{\PYGZsh{} \CJKCode{同一热源不可分属多个资源池而重复计量。}}
        \PYG{k}{if} \PYG{n+nb}{set}\PYG{p}{(}\PYG{n}{members}\PYG{p}{)} \PYG{o}{\PYGZhy{}} \PYG{n+nb}{set}\PYG{p}{(}\PYG{n}{model}\PYG{o}{.}\PYG{n}{sources}\PYG{p}{)} \PYG{o+ow}{or} \PYG{n}{assigned}\PYG{o}{.}\PYG{n}{intersection}\PYG{p}{(}\PYG{n}{members}\PYG{p}{)}\PYG{p}{:}
            \PYG{k}{raise} \PYG{n+ne}{ValueError}\PYG{p}{(}\PYG{l+s+sa}{f}\PYG{l+s+s2}{\PYGZdq{}}\PYG{l+s+s2}{resource }\PYG{l+s+si}{\PYGZob{}}\PYG{n}{name}\PYG{l+s+si}{\PYGZcb{}}\PYG{l+s+s2}{: unknown source or source assigned to multiple pools}\PYG{l+s+s2}{\PYGZdq{}}\PYG{p}{)}
        \PYG{n}{assigned}\PYG{o}{.}\PYG{n}{update}\PYG{p}{(}\PYG{n}{members}\PYG{p}{)}
        \PYG{c+c1}{\PYGZsh{} \CJKCode{可用热功率以}kW\CJKCode{登记，常量与时序二选一。}}
        \PYG{n}{scalar} \PYG{o}{=} \PYG{l+s+s2}{\PYGZdq{}}\PYG{l+s+s2}{available\PYGZus{}heat\PYGZus{}kw}\PYG{l+s+s2}{\PYGZdq{}} \PYG{o+ow}{in} \PYG{n}{pool}
        \PYG{n}{column} \PYG{o}{=} \PYG{l+s+s2}{\PYGZdq{}}\PYG{l+s+s2}{available\PYGZus{}heat\PYGZus{}column}\PYG{l+s+s2}{\PYGZdq{}} \PYG{o+ow}{in} \PYG{n}{pool}
        \PYG{k}{if} \PYG{n}{scalar} \PYG{o}{==} \PYG{n}{column}\PYG{p}{:}
            \PYG{k}{raise} \PYG{n+ne}{ValueError}\PYG{p}{(}\PYG{l+s+sa}{f}\PYG{l+s+s2}{\PYGZdq{}}\PYG{l+s+s2}{resource }\PYG{l+s+si}{\PYGZob{}}\PYG{n}{name}\PYG{l+s+si}{\PYGZcb{}}\PYG{l+s+s2}{: specify exactly one available\PYGZus{}heat\PYGZus{}kw/column}\PYG{l+s+s2}{\PYGZdq{}}\PYG{p}{)}
        \PYG{c+c1}{\PYGZsh{} \CJKCode{效率无量纲，输出热量折回上游时需除以它。}}
        \PYG{n}{efficiency} \PYG{o}{=} \PYG{n}{\PYGZus{}nonnegative}\PYG{p}{(}\PYG{n}{pool}\PYG{p}{[}\PYG{l+s+s2}{\PYGZdq{}}\PYG{l+s+s2}{transfer\PYGZus{}efficiency}\PYG{l+s+s2}{\PYGZdq{}}\PYG{p}{]}\PYG{p}{,} \PYG{l+s+sa}{f}\PYG{l+s+s2}{\PYGZdq{}}\PYG{l+s+s2}{resource }\PYG{l+s+si}{\PYGZob{}}\PYG{n}{name}\PYG{l+s+si}{\PYGZcb{}}\PYG{l+s+s2}{ efficiency}\PYG{l+s+s2}{\PYGZdq{}}\PYG{p}{)}
        \PYG{k}{if} \PYG{o+ow}{not} \PYG{l+m+mi}{0} \PYG{o}{\PYGZlt{}} \PYG{n}{efficiency} \PYG{o}{\PYGZlt{}}\PYG{o}{=} \PYG{l+m+mi}{1}\PYG{p}{:}
            \PYG{k}{raise} \PYG{n+ne}{ValueError}\PYG{p}{(}\PYG{l+s+s2}{\PYGZdq{}}\PYG{l+s+s2}{resource transfer\PYGZus{}efficiency must be in (0, 1]}\PYG{l+s+s2}{\PYGZdq{}}\PYG{p}{)}
        \PYG{k}{for} \PYG{n}{source\PYGZus{}name} \PYG{o+ow}{in} \PYG{n}{members}\PYG{p}{:}
            \PYG{n}{source} \PYG{o}{=} \PYG{n}{model}\PYG{o}{.}\PYG{n}{sources}\PYG{p}{[}\PYG{n}{source\PYGZus{}name}\PYG{p}{]}
            \PYG{c+c1}{\PYGZsh{} \CJKCode{此资源口径仅接受直供换热器与供热热泵。}}
            \PYG{k}{if} \PYG{n}{source}\PYG{o}{.}\PYG{n}{get}\PYG{p}{(}\PYG{l+s+s2}{\PYGZdq{}}\PYG{l+s+s2}{type}\PYG{l+s+s2}{\PYGZdq{}}\PYG{p}{)} \PYG{o+ow}{not} \PYG{o+ow}{in} \PYG{p}{\PYGZob{}}\PYG{l+s+s2}{\PYGZdq{}}\PYG{l+s+s2}{direct\PYGZus{}heat}\PYG{l+s+s2}{\PYGZdq{}}\PYG{p}{,} \PYG{l+s+s2}{\PYGZdq{}}\PYG{l+s+s2}{heat\PYGZus{}pump}\PYG{l+s+s2}{\PYGZdq{}}\PYG{p}{\PYGZcb{}}\PYG{p}{:}
                \PYG{k}{raise} \PYG{n+ne}{ValueError}\PYG{p}{(}\PYG{l+s+sa}{f}\PYG{l+s+s2}{\PYGZdq{}}\PYG{l+s+s2}{resource }\PYG{l+s+si}{\PYGZob{}}\PYG{n}{name}\PYG{l+s+si}{\PYGZcb{}}\PYG{l+s+s2}{: only direct\PYGZus{}heat and heat\PYGZus{}pump are supported}\PYG{l+s+s2}{\PYGZdq{}}\PYG{p}{)}
            \PYG{c+c1}{\PYGZsh{} \CJKCode{制冷或排热设备须另建资源守恒，不能套用。}}
            \PYG{k}{if} \PYG{n+nb}{any}\PYG{p}{(}\PYG{n}{key}\PYG{p}{[}\PYG{l+m+mi}{0}\PYG{p}{]} \PYG{o}{==} \PYG{n}{source\PYGZus{}name} \PYG{k}{for} \PYG{n}{key} \PYG{o+ow}{in} \PYG{n}{model}\PYG{o}{.}\PYG{n}{variables}\PYG{p}{[}\PYG{l+s+s2}{\PYGZdq{}}\PYG{l+s+s2}{cooling\PYGZus{}output}\PYG{l+s+s2}{\PYGZdq{}}\PYG{p}{]}\PYG{p}{)} \PYG{o+ow}{or} \PYG{n}{source}\PYG{o}{.}\PYG{n}{get}\PYG{p}{(}\PYG{l+s+s2}{\PYGZdq{}}\PYG{l+s+s2}{heat\PYGZus{}rejection\PYGZus{}enabled}\PYG{l+s+s2}{\PYGZdq{}}\PYG{p}{)}\PYG{p}{:}
                \PYG{k}{raise} \PYG{n+ne}{ValueError}\PYG{p}{(}\PYG{l+s+sa}{f}\PYG{l+s+s2}{\PYGZdq{}}\PYG{l+s+s2}{resource }\PYG{l+s+si}{\PYGZob{}}\PYG{n}{name}\PYG{l+s+si}{\PYGZcb{}}\PYG{l+s+s2}{: cooling/rejection needs a separate resource balance}\PYG{l+s+s2}{\PYGZdq{}}\PYG{p}{)}
        \PYG{k}{for} \PYG{n}{t} \PYG{o+ow}{in} \PYG{n}{model}\PYG{o}{.}\PYG{n}{times}\PYG{p}{:}
            \PYG{n}{available} \PYG{o}{=} \PYG{n}{\PYGZus{}nonnegative}\PYG{p}{(}
                \PYG{n}{model}\PYG{o}{.}\PYG{n}{\PYGZus{}ts\PYGZus{}value}\PYG{p}{(}\PYG{n}{t}\PYG{p}{,} \PYG{n}{pool}\PYG{p}{[}\PYG{l+s+s2}{\PYGZdq{}}\PYG{l+s+s2}{available\PYGZus{}heat\PYGZus{}column}\PYG{l+s+s2}{\PYGZdq{}}\PYG{p}{]}\PYG{p}{)} \PYG{k}{if} \PYG{n}{column} \PYG{k}{else} \PYG{n}{pool}\PYG{p}{[}\PYG{l+s+s2}{\PYGZdq{}}\PYG{l+s+s2}{available\PYGZus{}heat\PYGZus{}kw}\PYG{l+s+s2}{\PYGZdq{}}\PYG{p}{]}\PYG{p}{,}
                \PYG{l+s+sa}{f}\PYG{l+s+s2}{\PYGZdq{}}\PYG{l+s+s2}{resource }\PYG{l+s+si}{\PYGZob{}}\PYG{n}{name}\PYG{l+s+si}{\PYGZcb{}}\PYG{l+s+s2}{ available heat at }\PYG{l+s+si}{\PYGZob{}}\PYG{n}{t}\PYG{l+s+si}{\PYGZcb{}}\PYG{l+s+s2}{\PYGZdq{}}\PYG{p}{,}
            \PYG{p}{)}
            \PYG{n}{terms} \PYG{o}{=} \PYG{p}{\PYGZob{}}\PYG{p}{\PYGZcb{}}
            \PYG{k}{for} \PYG{n}{source\PYGZus{}name} \PYG{o+ow}{in} \PYG{n}{members}\PYG{p}{:}
                \PYG{n}{source} \PYG{o}{=} \PYG{n}{model}\PYG{o}{.}\PYG{n}{sources}\PYG{p}{[}\PYG{n}{source\PYGZus{}name}\PYG{p}{]}
                \PYG{k}{for} \PYG{n}{grade\PYGZus{}name}\PYG{p}{,} \PYG{n}{grade} \PYG{o+ow}{in} \PYG{n}{model}\PYG{o}{.}\PYG{n}{grades}\PYG{o}{.}\PYG{n}{items}\PYG{p}{(}\PYG{p}{)}\PYG{p}{:}
                    \PYG{c+c1}{\PYGZsh{} \CJKCode{按温位汇总同一成员已有的供热功率变量。}}
                    \PYG{n}{idx} \PYG{o}{=} \PYG{n}{model}\PYG{o}{.}\PYG{n}{variables}\PYG{p}{[}\PYG{l+s+s2}{\PYGZdq{}}\PYG{l+s+s2}{source\PYGZus{}heat}\PYG{l+s+s2}{\PYGZdq{}}\PYG{p}{]}\PYG{o}{.}\PYG{n}{get}\PYG{p}{(}\PYG{p}{(}\PYG{n}{source\PYGZus{}name}\PYG{p}{,} \PYG{n}{grade\PYGZus{}name}\PYG{p}{,} \PYG{n}{t}\PYG{p}{)}\PYG{p}{)}
                    \PYG{k}{if} \PYG{n}{idx} \PYG{o+ow}{is} \PYG{k+kc}{None}\PYG{p}{:}
                        \PYG{k}{continue}
                    \PYG{c+c1}{\PYGZsh{} \CJKCode{直供取热比例为}1\CJKCode{；热泵需扣除压缩机电功。}}
                    \PYG{n}{fraction} \PYG{o}{=} \PYG{l+m+mf}{1.0}
                    \PYG{k}{if} \PYG{n}{source}\PYG{o}{.}\PYG{n}{get}\PYG{p}{(}\PYG{l+s+s2}{\PYGZdq{}}\PYG{l+s+s2}{type}\PYG{l+s+s2}{\PYGZdq{}}\PYG{p}{)} \PYG{o}{==} \PYG{l+s+s2}{\PYGZdq{}}\PYG{l+s+s2}{heat\PYGZus{}pump}\PYG{l+s+s2}{\PYGZdq{}}\PYG{p}{:}
                        \PYG{n}{cop} \PYG{o}{=} \PYG{n+nb}{float}\PYG{p}{(}\PYG{n}{model}\PYG{o}{.}\PYG{n}{\PYGZus{}cop}\PYG{p}{(}\PYG{n}{source}\PYG{p}{,} \PYG{n}{grade}\PYG{p}{,} \PYG{n}{t}\PYG{p}{)}\PYG{p}{)}
                        \PYG{k}{if} \PYG{o+ow}{not} \PYG{n}{isfinite}\PYG{p}{(}\PYG{n}{cop}\PYG{p}{)} \PYG{o+ow}{or} \PYG{n}{cop} \PYG{o}{\PYGZlt{}}\PYG{o}{=} \PYG{l+m+mi}{1}\PYG{p}{:}
                            \PYG{k}{raise} \PYG{n+ne}{ValueError}\PYG{p}{(}\PYG{l+s+s2}{\PYGZdq{}}\PYG{l+s+s2}{a resource heat pump requires finite COP \PYGZgt{} 1}\PYG{l+s+s2}{\PYGZdq{}}\PYG{p}{)}
                        \PYG{c+c1}{\PYGZsh{} \CJKCode{热泵上游取热为供热功率乘以}1\CJKCode{减}COP\CJKCode{倒数。}}
                        \PYG{n}{fraction} \PYG{o}{\PYGZhy{}}\PYG{o}{=} \PYG{l+m+mf}{1.0} \PYG{o}{/} \PYG{n}{cop}
                    \PYG{c+c1}{\PYGZsh{} \CJKCode{再除传热效率，得到每}kW\CJKCode{供热对应的上游取热。}}
                    \PYG{n}{model}\PYG{o}{.}\PYG{n}{\PYGZus{}add}\PYG{p}{(}\PYG{n}{terms}\PYG{p}{,} \PYG{n}{idx}\PYG{p}{,} \PYG{n}{fraction} \PYG{o}{/} \PYG{n}{efficiency}\PYG{p}{)}
            \PYG{c+c1}{\PYGZsh{} \CJKCode{取热变量上界为逐时可用资源，等式绑定总取热。}}
            \PYG{n}{use} \PYG{o}{=} \PYG{n}{model}\PYG{o}{.}\PYG{n}{problem}\PYG{o}{.}\PYG{n}{add\PYGZus{}var}\PYG{p}{(}\PYG{l+s+sa}{f}\PYG{l+s+s2}{\PYGZdq{}}\PYG{l+s+s2}{shared\PYGZus{}heat\PYGZus{}resource\PYGZus{}use[}\PYG{l+s+si}{\PYGZob{}}\PYG{n}{name}\PYG{l+s+si}{\PYGZcb{}}\PYG{l+s+s2}{,}\PYG{l+s+si}{\PYGZob{}}\PYG{n}{t}\PYG{l+s+si}{\PYGZcb{}}\PYG{l+s+s2}{]}\PYG{l+s+s2}{\PYGZdq{}}\PYG{p}{,} \PYG{n}{ub}\PYG{o}{=}\PYG{n}{available}\PYG{p}{)}
            \PYG{n}{model}\PYG{o}{.}\PYG{n}{variables}\PYG{p}{[}\PYG{l+s+s2}{\PYGZdq{}}\PYG{l+s+s2}{shared\PYGZus{}heat\PYGZus{}resource\PYGZus{}use}\PYG{l+s+s2}{\PYGZdq{}}\PYG{p}{]}\PYG{p}{[}\PYG{p}{(}\PYG{n}{name}\PYG{p}{,} \PYG{n}{t}\PYG{p}{)}\PYG{p}{]} \PYG{o}{=} \PYG{n}{use}
            \PYG{n}{model}\PYG{o}{.}\PYG{n}{problem}\PYG{o}{.}\PYG{n}{add\PYGZus{}constraint}\PYG{p}{(}\PYG{p}{\PYGZob{}}\PYG{n}{use}\PYG{p}{:} \PYG{l+m+mf}{1.0}\PYG{p}{,} \PYG{o}{*}\PYG{o}{*}\PYG{p}{\PYGZob{}}\PYG{n}{i}\PYG{p}{:} \PYG{o}{\PYGZhy{}}\PYG{n}{c} \PYG{k}{for} \PYG{n}{i}\PYG{p}{,} \PYG{n}{c} \PYG{o+ow}{in} \PYG{n}{terms}\PYG{o}{.}\PYG{n}{items}\PYG{p}{(}\PYG{p}{)}\PYG{p}{\PYGZcb{}}\PYG{p}{\PYGZcb{}}\PYG{p}{,} \PYG{n}{lb}\PYG{o}{=}\PYG{l+m+mf}{0.0}\PYG{p}{,} \PYG{n}{ub}\PYG{o}{=}\PYG{l+m+mf}{0.0}\PYG{p}{)}
            \PYG{n}{model}\PYG{o}{.}\PYG{n}{shared\PYGZus{}heat\PYGZus{}resource\PYGZus{}rows}\PYG{o}{.}\PYG{n}{append}\PYG{p}{(}\PYG{p}{(}\PYG{n}{name}\PYG{p}{,} \PYG{n}{t}\PYG{p}{,} \PYG{n}{available}\PYG{p}{,} \PYG{n}{terms}\PYG{p}{,} \PYG{n}{use}\PYG{p}{)}\PYG{p}{)}
\end{Verbatim}

\endgroup

\noindent\textbf{结果核对。}核验\path{shared_heat_resource_use}的逐时kW值及资源行登记的可用上界、COP和折算系数；适用于明确登记的直供与供热热泵共享资源池。制冷、排热支路不属于本口径。

\Needspace{12\baselineskip}
\subsection*{14．逐楼RC状态、舒适边界与节点交付量}
完整的逐楼1R1C约束装配函数，与前述精确离散系数函数互补。室温变量直接承载逐时舒适上下界，实际冷热交付量进入热状态递推，并按节点温位聚合后交给主模型平衡。

\noindent\textbf{输入：}通过校验的逐楼热损失系数、有效热容、离散系数a和b、室外温度与得热时序、舒适温度及冷热功率上界、分组边界和互斥模式。\par
\noindent\textbf{输出：}逐楼末态室温（摄氏度）、实际供热和制冷功率（kW）、可选冷热模式二元变量、节点交付量映射以及RC递推等式与末态边界。\par

\noindent\textbf{实现思路。}逐楼末态室温$T_t$（摄氏度）受逐时舒适上下界约束，供热$Q_h$与制冷$Q_c$均为kW。以此前计算的$a$（无量纲）、$b$（K/kW）装配$T_t-aT_{t-1}-bQ_h+bQ_c=b(HT_{\mathrm{out}}+G)$；$H$为kW/K，$G$为内得热与太阳得热之和。固定首态移入右端，循环组连接组内前驱，固定末态另受温度界约束。聚合交付量接入正文节点平衡；冷热互斥二元变量仅按配置启用。

{\small 来源：\path{building_thermal.py}，原第463--525行；63行原代码、15行补充注释。}\par

\begingroup
\setstretch{1.0}
\begin{Verbatim}[commandchars=\\\{\},numbers=left,firstnumber=1,stepnumber=1,fontsize=\fontsize{9.5}{11.5}\selectfont,breaklines=true,breakanywhere=true,frame=single,framesep=3mm,numbersep=5pt,rulecolor=\color{black},xleftmargin=18pt]
\PYG{c+c1}{\PYGZsh{} \CJKCode{逐楼热状态是独立}RC\CJKCode{分支，节点平衡须消费其交付量。}}
\PYG{k}{def}\PYG{+w}{ }\PYG{n+nf}{add\PYGZus{}building\PYGZus{}thermal\PYGZus{}constraints}\PYG{p}{(}\PYG{n}{model}\PYG{p}{:} \PYG{n}{Any}\PYG{p}{)} \PYG{o}{\PYGZhy{}}\PYG{o}{\PYGZgt{}} \PYG{n}{BuildingThermalBlock} \PYG{o}{|} \PYG{k+kc}{None}\PYG{p}{:}
\PYG{+w}{    }\PYG{l+s+sd}{\PYGZdq{}\PYGZdq{}\PYGZdq{}Register on model.problem / model.variables and return load variable maps.}

\PYG{l+s+sd}{    Needs config, ts, dt, problem and variables, optionally boundary\PYGZus{}groups.}
\PYG{l+s+sd}{    Adds no source equipment, costs, unmet\PYGZhy{}load variables or reliability credits.}
\PYG{l+s+sd}{    The caller MUST consume node\PYGZus{}heat/node\PYGZus{}cooling in physical supply balances.}
\PYG{l+s+sd}{    \PYGZdq{}\PYGZdq{}\PYGZdq{}}
    \PYG{c+c1}{\PYGZsh{} \CJKCode{先校验逐楼输入与分组边界，未启用则返回空。}}
    \PYG{n}{inputs} \PYG{o}{=} \PYG{n}{validate\PYGZus{}building\PYGZus{}thermal\PYGZus{}inputs}\PYG{p}{(}\PYG{n}{model}\PYG{o}{.}\PYG{n}{config}\PYG{p}{,} \PYG{n}{model}\PYG{o}{.}\PYG{n}{ts}\PYG{p}{,} \PYG{n}{dt\PYGZus{}hours}\PYG{o}{=}\PYG{n}{model}\PYG{o}{.}\PYG{n}{dt}\PYG{p}{,}
                                             \PYG{n}{boundary\PYGZus{}groups}\PYG{o}{=}\PYG{n+nb}{getattr}\PYG{p}{(}\PYG{n}{model}\PYG{p}{,} \PYG{l+s+s2}{\PYGZdq{}}\PYG{l+s+s2}{boundary\PYGZus{}groups}\PYG{l+s+s2}{\PYGZdq{}}\PYG{p}{,} \PYG{k+kc}{None}\PYG{p}{)}\PYG{p}{)}
    \PYG{k}{if} \PYG{n}{inputs} \PYG{o+ow}{is} \PYG{k+kc}{None}\PYG{p}{:}
        \PYG{k}{return} \PYG{k+kc}{None}
    \PYG{n}{categories} \PYG{o}{=} \PYG{p}{(}\PYG{l+s+s2}{\PYGZdq{}}\PYG{l+s+s2}{building\PYGZus{}temperature}\PYG{l+s+s2}{\PYGZdq{}}\PYG{p}{,} \PYG{l+s+s2}{\PYGZdq{}}\PYG{l+s+s2}{building\PYGZus{}heat}\PYG{l+s+s2}{\PYGZdq{}}\PYG{p}{,} \PYG{l+s+s2}{\PYGZdq{}}\PYG{l+s+s2}{building\PYGZus{}cooling}\PYG{l+s+s2}{\PYGZdq{}}\PYG{p}{,} \PYG{l+s+s2}{\PYGZdq{}}\PYG{l+s+s2}{building\PYGZus{}heat\PYGZus{}mode}\PYG{l+s+s2}{\PYGZdq{}}\PYG{p}{,}
                  \PYG{l+s+s2}{\PYGZdq{}}\PYG{l+s+s2}{building\PYGZus{}heat\PYGZus{}by\PYGZus{}node}\PYG{l+s+s2}{\PYGZdq{}}\PYG{p}{,} \PYG{l+s+s2}{\PYGZdq{}}\PYG{l+s+s2}{building\PYGZus{}cooling\PYGZus{}by\PYGZus{}node}\PYG{l+s+s2}{\PYGZdq{}}\PYG{p}{)}
    \PYG{c+c1}{\PYGZsh{} \CJKCode{拒绝重复注册，避免同一建筑出现两套状态变量。}}
    \PYG{k}{if} \PYG{n+nb}{any}\PYG{p}{(}\PYG{n}{model}\PYG{o}{.}\PYG{n}{variables}\PYG{o}{.}\PYG{n}{get}\PYG{p}{(}\PYG{n}{category}\PYG{p}{)} \PYG{k}{for} \PYG{n}{category} \PYG{o+ow}{in} \PYG{n}{categories}\PYG{p}{)}\PYG{p}{:}
        \PYG{k}{raise} \PYG{n+ne}{ValueError}\PYG{p}{(}\PYG{l+s+s2}{\PYGZdq{}}\PYG{l+s+s2}{building thermal variables already registered}\PYG{l+s+s2}{\PYGZdq{}}\PYG{p}{)}
    \PYG{n}{maps} \PYG{o}{=} \PYG{p}{\PYGZob{}}\PYG{n}{category}\PYG{p}{:} \PYG{p}{\PYGZob{}}\PYG{p}{\PYGZcb{}} \PYG{k}{for} \PYG{n}{category} \PYG{o+ow}{in} \PYG{n}{categories}\PYG{p}{\PYGZcb{}}
    \PYG{n}{temperature}\PYG{p}{,} \PYG{n}{heat}\PYG{p}{,} \PYG{n}{cooling}\PYG{p}{,} \PYG{n}{mode}\PYG{p}{,} \PYG{n}{node\PYGZus{}heat}\PYG{p}{,} \PYG{n}{node\PYGZus{}cooling} \PYG{o}{=} \PYG{p}{(}\PYG{n}{maps}\PYG{p}{[}\PYG{n}{category}\PYG{p}{]} \PYG{k}{for} \PYG{n}{category} \PYG{o+ow}{in} \PYG{n}{categories}\PYG{p}{)}
    \PYG{n}{p} \PYG{o}{=} \PYG{n}{model}\PYG{o}{.}\PYG{n}{problem}
    \PYG{n}{heat\PYGZus{}members}\PYG{p}{:} \PYG{n+nb}{dict}\PYG{p}{[}\PYG{n+nb}{tuple}\PYG{p}{[}\PYG{n+nb}{str}\PYG{p}{,} \PYG{n+nb}{str}\PYG{p}{,} \PYG{n+nb}{int}\PYG{p}{]}\PYG{p}{,} \PYG{n+nb}{list}\PYG{p}{[}\PYG{n+nb}{int}\PYG{p}{]}\PYG{p}{]} \PYG{o}{=} \PYG{n}{defaultdict}\PYG{p}{(}\PYG{n+nb}{list}\PYG{p}{)}
    \PYG{n}{cooling\PYGZus{}members}\PYG{p}{:} \PYG{n+nb}{dict}\PYG{p}{[}\PYG{n+nb}{tuple}\PYG{p}{[}\PYG{n+nb}{str}\PYG{p}{,} \PYG{n+nb}{int}\PYG{p}{]}\PYG{p}{,} \PYG{n+nb}{list}\PYG{p}{[}\PYG{n+nb}{int}\PYG{p}{]}\PYG{p}{]} \PYG{o}{=} \PYG{n}{defaultdict}\PYG{p}{(}\PYG{n+nb}{list}\PYG{p}{)}
    \PYG{k}{for} \PYG{n}{building} \PYG{o+ow}{in} \PYG{n}{inputs}\PYG{o}{.}\PYG{n}{buildings}\PYG{p}{:}
        \PYG{k}{for} \PYG{n}{t} \PYG{o+ow}{in} \PYG{n+nb}{range}\PYG{p}{(}\PYG{n+nb}{len}\PYG{p}{(}\PYG{n}{inputs}\PYG{o}{.}\PYG{n}{frame}\PYG{p}{)}\PYG{p}{)}\PYG{p}{:}
            \PYG{n}{key} \PYG{o}{=} \PYG{n}{building}\PYG{o}{.}\PYG{n}{name}\PYG{p}{,} \PYG{n}{t}
            \PYG{c+c1}{\PYGZsh{} \CJKCode{室温单位为摄氏度，逐时舒适区间直接作为变量界。}}
            \PYG{n}{temperature}\PYG{p}{[}\PYG{n}{key}\PYG{p}{]} \PYG{o}{=} \PYG{n}{p}\PYG{o}{.}\PYG{n}{add\PYGZus{}var}\PYG{p}{(}\PYG{l+s+sa}{f}\PYG{l+s+s2}{\PYGZdq{}}\PYG{l+s+s2}{building\PYGZus{}temperature[}\PYG{l+s+si}{\PYGZob{}}\PYG{n}{building}\PYG{o}{.}\PYG{n}{name}\PYG{l+s+si}{\PYGZcb{}}\PYG{l+s+s2}{,}\PYG{l+s+si}{\PYGZob{}}\PYG{n}{t}\PYG{l+s+si}{\PYGZcb{}}\PYG{l+s+s2}{]}\PYG{l+s+s2}{\PYGZdq{}}\PYG{p}{,} \PYG{n}{lb}\PYG{o}{=}\PYG{n}{building}\PYG{o}{.}\PYG{n}{minimum}\PYG{p}{[}\PYG{n}{t}\PYG{p}{]}\PYG{p}{,} \PYG{n}{ub}\PYG{o}{=}\PYG{n}{building}\PYG{o}{.}\PYG{n}{maximum}\PYG{p}{[}\PYG{n}{t}\PYG{p}{]}\PYG{p}{)}
            \PYG{c+c1}{\PYGZsh{} \CJKCode{冷热交付均为非负}kW\CJKCode{，各自受末端功率上界限制。}}
            \PYG{n}{heat}\PYG{p}{[}\PYG{n}{key}\PYG{p}{]} \PYG{o}{=} \PYG{n}{p}\PYG{o}{.}\PYG{n}{add\PYGZus{}var}\PYG{p}{(}\PYG{l+s+sa}{f}\PYG{l+s+s2}{\PYGZdq{}}\PYG{l+s+s2}{building\PYGZus{}heat[}\PYG{l+s+si}{\PYGZob{}}\PYG{n}{building}\PYG{o}{.}\PYG{n}{name}\PYG{l+s+si}{\PYGZcb{}}\PYG{l+s+s2}{,}\PYG{l+s+si}{\PYGZob{}}\PYG{n}{t}\PYG{l+s+si}{\PYGZcb{}}\PYG{l+s+s2}{]}\PYG{l+s+s2}{\PYGZdq{}}\PYG{p}{,} \PYG{n}{ub}\PYG{o}{=}\PYG{n}{building}\PYG{o}{.}\PYG{n}{heat\PYGZus{}upper}\PYG{p}{[}\PYG{n}{t}\PYG{p}{]}\PYG{p}{)}
            \PYG{n}{cooling}\PYG{p}{[}\PYG{n}{key}\PYG{p}{]} \PYG{o}{=} \PYG{n}{p}\PYG{o}{.}\PYG{n}{add\PYGZus{}var}\PYG{p}{(}\PYG{l+s+sa}{f}\PYG{l+s+s2}{\PYGZdq{}}\PYG{l+s+s2}{building\PYGZus{}cooling[}\PYG{l+s+si}{\PYGZob{}}\PYG{n}{building}\PYG{o}{.}\PYG{n}{name}\PYG{l+s+si}{\PYGZcb{}}\PYG{l+s+s2}{,}\PYG{l+s+si}{\PYGZob{}}\PYG{n}{t}\PYG{l+s+si}{\PYGZcb{}}\PYG{l+s+s2}{]}\PYG{l+s+s2}{\PYGZdq{}}\PYG{p}{,} \PYG{n}{ub}\PYG{o}{=}\PYG{n}{building}\PYG{o}{.}\PYG{n}{cooling\PYGZus{}upper}\PYG{p}{[}\PYG{n}{t}\PYG{p}{]}\PYG{p}{)}
            \PYG{c+c1}{\PYGZsh{} \CJKCode{按建筑所属节点和供热温位收集实际交付功率。}}
            \PYG{n}{heat\PYGZus{}members}\PYG{p}{[}\PYG{n}{building}\PYG{o}{.}\PYG{n}{node}\PYG{p}{,} \PYG{n}{building}\PYG{o}{.}\PYG{n}{heat\PYGZus{}grade}\PYG{p}{,} \PYG{n}{t}\PYG{p}{]}\PYG{o}{.}\PYG{n}{append}\PYG{p}{(}\PYG{n}{heat}\PYG{p}{[}\PYG{n}{key}\PYG{p}{]}\PYG{p}{)}
            \PYG{n}{cooling\PYGZus{}members}\PYG{p}{[}\PYG{n}{building}\PYG{o}{.}\PYG{n}{node}\PYG{p}{,} \PYG{n}{t}\PYG{p}{]}\PYG{o}{.}\PYG{n}{append}\PYG{p}{(}\PYG{n}{cooling}\PYG{p}{[}\PYG{n}{key}\PYG{p}{]}\PYG{p}{)}
            \PYG{c+c1}{\PYGZsh{} \CJKCode{仅互斥模式且冷热均可用时引入二元开关。}}
            \PYG{k}{if} \PYG{n}{inputs}\PYG{o}{.}\PYG{n}{simultaneous\PYGZus{}mode} \PYG{o}{==} \PYG{l+s+s2}{\PYGZdq{}}\PYG{l+s+s2}{exclusive}\PYG{l+s+s2}{\PYGZdq{}} \PYG{o+ow}{and} \PYG{n}{building}\PYG{o}{.}\PYG{n}{heat\PYGZus{}upper}\PYG{p}{[}\PYG{n}{t}\PYG{p}{]} \PYG{o}{\PYGZgt{}} \PYG{l+m+mi}{0} \PYG{o+ow}{and} \PYG{n}{building}\PYG{o}{.}\PYG{n}{cooling\PYGZus{}upper}\PYG{p}{[}\PYG{n}{t}\PYG{p}{]} \PYG{o}{\PYGZgt{}} \PYG{l+m+mi}{0}\PYG{p}{:}
                \PYG{n}{mode}\PYG{p}{[}\PYG{n}{key}\PYG{p}{]} \PYG{o}{=} \PYG{n}{p}\PYG{o}{.}\PYG{n}{add\PYGZus{}var}\PYG{p}{(}\PYG{l+s+sa}{f}\PYG{l+s+s2}{\PYGZdq{}}\PYG{l+s+s2}{building\PYGZus{}heat\PYGZus{}mode[}\PYG{l+s+si}{\PYGZob{}}\PYG{n}{building}\PYG{o}{.}\PYG{n}{name}\PYG{l+s+si}{\PYGZcb{}}\PYG{l+s+s2}{,}\PYG{l+s+si}{\PYGZob{}}\PYG{n}{t}\PYG{l+s+si}{\PYGZcb{}}\PYG{l+s+s2}{]}\PYG{l+s+s2}{\PYGZdq{}}\PYG{p}{,} \PYG{n}{ub}\PYG{o}{=}\PYG{l+m+mf}{1.0}\PYG{p}{,} \PYG{n}{integer}\PYG{o}{=}\PYG{k+kc}{True}\PYG{p}{)}
                \PYG{c+c1}{\PYGZsh{} \CJKCode{模式为}1\CJKCode{允许供热；模式为}0\CJKCode{允许制冷。}}
                \PYG{n}{p}\PYG{o}{.}\PYG{n}{add\PYGZus{}constraint}\PYG{p}{(}\PYG{p}{\PYGZob{}}\PYG{n}{heat}\PYG{p}{[}\PYG{n}{key}\PYG{p}{]}\PYG{p}{:} \PYG{l+m+mf}{1.0}\PYG{p}{,} \PYG{n}{mode}\PYG{p}{[}\PYG{n}{key}\PYG{p}{]}\PYG{p}{:} \PYG{o}{\PYGZhy{}}\PYG{n}{building}\PYG{o}{.}\PYG{n}{heat\PYGZus{}upper}\PYG{p}{[}\PYG{n}{t}\PYG{p}{]}\PYG{p}{\PYGZcb{}}\PYG{p}{,} \PYG{n}{ub}\PYG{o}{=}\PYG{l+m+mf}{0.0}\PYG{p}{)}
                \PYG{n}{p}\PYG{o}{.}\PYG{n}{add\PYGZus{}constraint}\PYG{p}{(}\PYG{p}{\PYGZob{}}\PYG{n}{cooling}\PYG{p}{[}\PYG{n}{key}\PYG{p}{]}\PYG{p}{:} \PYG{l+m+mf}{1.0}\PYG{p}{,} \PYG{n}{mode}\PYG{p}{[}\PYG{n}{key}\PYG{p}{]}\PYG{p}{:} \PYG{n}{building}\PYG{o}{.}\PYG{n}{cooling\PYGZus{}upper}\PYG{p}{[}\PYG{n}{t}\PYG{p}{]}\PYG{p}{\PYGZcb{}}\PYG{p}{,} \PYG{n}{ub}\PYG{o}{=}\PYG{n}{building}\PYG{o}{.}\PYG{n}{cooling\PYGZus{}upper}\PYG{p}{[}\PYG{n}{t}\PYG{p}{]}\PYG{p}{)}
        \PYG{c+c1}{\PYGZsh{} \CJKCode{状态递推只连接所属边界组，避免跨场景串联。}}
        \PYG{k}{for} \PYG{n}{group}\PYG{p}{,} \PYG{n}{times} \PYG{o+ow}{in} \PYG{n}{inputs}\PYG{o}{.}\PYG{n}{boundary\PYGZus{}groups}\PYG{o}{.}\PYG{n}{items}\PYG{p}{(}\PYG{p}{)}\PYG{p}{:}
            \PYG{n}{condition} \PYG{o}{=} \PYG{n}{building}\PYG{o}{.}\PYG{n}{boundary}\PYG{p}{[}\PYG{n}{group}\PYG{p}{]}
            \PYG{k}{for} \PYG{n}{position}\PYG{p}{,} \PYG{n}{t} \PYG{o+ow}{in} \PYG{n+nb}{enumerate}\PYG{p}{(}\PYG{n}{times}\PYG{p}{)}\PYG{p}{:}
                \PYG{n}{key} \PYG{o}{=} \PYG{n}{building}\PYG{o}{.}\PYG{n}{name}\PYG{p}{,} \PYG{n}{t}
                \PYG{c+c1}{\PYGZsh{} a\CJKCode{无量纲，}b\CJKCode{为}K\CJKCode{每}kW\CJKCode{；供热升温而制冷降温。}}
                \PYG{n}{row} \PYG{o}{=} \PYG{p}{\PYGZob{}}\PYG{n}{temperature}\PYG{p}{[}\PYG{n}{key}\PYG{p}{]}\PYG{p}{:} \PYG{l+m+mf}{1.0}\PYG{p}{,} \PYG{n}{heat}\PYG{p}{[}\PYG{n}{key}\PYG{p}{]}\PYG{p}{:} \PYG{o}{\PYGZhy{}}\PYG{n}{building}\PYG{o}{.}\PYG{n}{b}\PYG{p}{,} \PYG{n}{cooling}\PYG{p}{[}\PYG{n}{key}\PYG{p}{]}\PYG{p}{:} \PYG{n}{building}\PYG{o}{.}\PYG{n}{b}\PYG{p}{\PYGZcb{}}
                \PYG{c+c1}{\PYGZsh{} H\CJKCode{乘室外温度与内、太阳得热合成递推常量项。}}
                \PYG{n}{rhs} \PYG{o}{=} \PYG{n}{building}\PYG{o}{.}\PYG{n}{b} \PYG{o}{*} \PYG{p}{(}\PYG{n}{building}\PYG{o}{.}\PYG{n}{heat\PYGZus{}loss\PYGZus{}kw\PYGZus{}per\PYGZus{}c} \PYG{o}{*} \PYG{n}{inputs}\PYG{o}{.}\PYG{n}{outdoor\PYGZus{}temperature}\PYG{p}{[}\PYG{n}{t}\PYG{p}{]}
                                     \PYG{o}{+} \PYG{n}{building}\PYG{o}{.}\PYG{n}{internal\PYGZus{}gain}\PYG{p}{[}\PYG{n}{t}\PYG{p}{]} \PYG{o}{+} \PYG{n}{building}\PYG{o}{.}\PYG{n}{solar\PYGZus{}gain}\PYG{p}{[}\PYG{n}{t}\PYG{p}{]}\PYG{p}{)}
                \PYG{c+c1}{\PYGZsh{} \CJKCode{固定组首时段把初始室温乘}a\CJKCode{移到等式右端。}}
                \PYG{k}{if} \PYG{n}{position} \PYG{o}{==} \PYG{l+m+mi}{0} \PYG{o+ow}{and} \PYG{n}{condition}\PYG{o}{.}\PYG{n}{mode} \PYG{o}{==} \PYG{l+s+s2}{\PYGZdq{}}\PYG{l+s+s2}{fixed}\PYG{l+s+s2}{\PYGZdq{}}\PYG{p}{:}
                    \PYG{n}{rhs} \PYG{o}{+}\PYG{o}{=} \PYG{n}{building}\PYG{o}{.}\PYG{n}{a} \PYG{o}{*} \PYG{n}{condition}\PYG{o}{.}\PYG{n}{initial\PYGZus{}temperature\PYGZus{}c}
                \PYG{k}{else}\PYG{p}{:}
                    \PYG{n}{previous} \PYG{o}{=} \PYG{n}{temperature}\PYG{p}{[}\PYG{n}{building}\PYG{o}{.}\PYG{n}{name}\PYG{p}{,} \PYG{n}{times}\PYG{p}{[}\PYG{n}{position} \PYG{o}{\PYGZhy{}} \PYG{l+m+mi}{1}\PYG{p}{]}\PYG{p}{]}
                    \PYG{c+c1}{\PYGZsh{} Addition is essential for single\PYGZhy{}step cyclic groups (same index).}
                    \PYG{c+c1}{\PYGZsh{} \CJKCode{循环组接组内前驱，单时段同索引须累加系数。}}
                    \PYG{n}{row}\PYG{p}{[}\PYG{n}{previous}\PYG{p}{]} \PYG{o}{=} \PYG{n}{row}\PYG{o}{.}\PYG{n}{get}\PYG{p}{(}\PYG{n}{previous}\PYG{p}{,} \PYG{l+m+mf}{0.0}\PYG{p}{)} \PYG{o}{\PYGZhy{}} \PYG{n}{building}\PYG{o}{.}\PYG{n}{a}
                \PYG{n}{p}\PYG{o}{.}\PYG{n}{add\PYGZus{}constraint}\PYG{p}{(}\PYG{n}{row}\PYG{p}{,} \PYG{n}{lb}\PYG{o}{=}\PYG{n}{rhs}\PYG{p}{,} \PYG{n}{ub}\PYG{o}{=}\PYG{n}{rhs}\PYG{p}{)}
            \PYG{c+c1}{\PYGZsh{} \CJKCode{固定边界另约束末态室温，不强制循环相等。}}
            \PYG{k}{if} \PYG{n}{condition}\PYG{o}{.}\PYG{n}{mode} \PYG{o}{==} \PYG{l+s+s2}{\PYGZdq{}}\PYG{l+s+s2}{fixed}\PYG{l+s+s2}{\PYGZdq{}}\PYG{p}{:}
                \PYG{n}{p}\PYG{o}{.}\PYG{n}{add\PYGZus{}constraint}\PYG{p}{(}\PYG{p}{\PYGZob{}}\PYG{n}{temperature}\PYG{p}{[}\PYG{n}{building}\PYG{o}{.}\PYG{n}{name}\PYG{p}{,} \PYG{n}{times}\PYG{p}{[}\PYG{o}{\PYGZhy{}}\PYG{l+m+mi}{1}\PYG{p}{]}\PYG{p}{]}\PYG{p}{:} \PYG{l+m+mf}{1.0}\PYG{p}{\PYGZcb{}}\PYG{p}{,}
                                 \PYG{n}{lb}\PYG{o}{=}\PYG{n}{condition}\PYG{o}{.}\PYG{n}{final\PYGZus{}min\PYGZus{}temperature\PYGZus{}c}\PYG{p}{,}
                                 \PYG{n}{ub}\PYG{o}{=}\PYG{n}{condition}\PYG{o}{.}\PYG{n}{final\PYGZus{}max\PYGZus{}temperature\PYGZus{}c} \PYG{k}{if} \PYG{n}{condition}\PYG{o}{.}\PYG{n}{final\PYGZus{}max\PYGZus{}temperature\PYGZus{}c} \PYG{o+ow}{is} \PYG{o+ow}{not} \PYG{k+kc}{None} \PYG{k}{else} \PYG{n}{np}\PYG{o}{.}\PYG{n}{inf}\PYG{p}{)}
    \PYG{c+c1}{\PYGZsh{} \CJKCode{多楼以等式汇总交付量，单楼直接复用变量索引。}}
    \PYG{k}{for} \PYG{n}{category}\PYG{p}{,} \PYG{n}{members}\PYG{p}{,} \PYG{n}{target} \PYG{o+ow}{in} \PYG{p}{(}\PYG{p}{(}\PYG{l+s+s2}{\PYGZdq{}}\PYG{l+s+s2}{building\PYGZus{}heat\PYGZus{}by\PYGZus{}node}\PYG{l+s+s2}{\PYGZdq{}}\PYG{p}{,} \PYG{n}{heat\PYGZus{}members}\PYG{p}{,} \PYG{n}{node\PYGZus{}heat}\PYG{p}{)}\PYG{p}{,}
                                      \PYG{p}{(}\PYG{l+s+s2}{\PYGZdq{}}\PYG{l+s+s2}{building\PYGZus{}cooling\PYGZus{}by\PYGZus{}node}\PYG{l+s+s2}{\PYGZdq{}}\PYG{p}{,} \PYG{n}{cooling\PYGZus{}members}\PYG{p}{,} \PYG{n}{node\PYGZus{}cooling}\PYG{p}{)}\PYG{p}{)}\PYG{p}{:}
        \PYG{k}{for} \PYG{n}{key}\PYG{p}{,} \PYG{n}{indices} \PYG{o+ow}{in} \PYG{n}{members}\PYG{o}{.}\PYG{n}{items}\PYG{p}{(}\PYG{p}{)}\PYG{p}{:}
            \PYG{k}{if} \PYG{n+nb}{len}\PYG{p}{(}\PYG{n}{indices}\PYG{p}{)} \PYG{o}{==} \PYG{l+m+mi}{1}\PYG{p}{:}
                \PYG{n}{target}\PYG{p}{[}\PYG{n}{key}\PYG{p}{]} \PYG{o}{=} \PYG{n}{indices}\PYG{p}{[}\PYG{l+m+mi}{0}\PYG{p}{]}  \PYG{c+c1}{\PYGZsh{} Alias avoids a redundant row/variable.}
            \PYG{k}{else}\PYG{p}{:}
                \PYG{n}{aggregate} \PYG{o}{=} \PYG{n}{p}\PYG{o}{.}\PYG{n}{add\PYGZus{}var}\PYG{p}{(}\PYG{l+s+sa}{f}\PYG{l+s+s2}{\PYGZdq{}}\PYG{l+s+si}{\PYGZob{}}\PYG{n}{category}\PYG{l+s+si}{\PYGZcb{}}\PYG{l+s+s2}{[}\PYG{l+s+si}{\PYGZob{}}\PYG{l+s+s1}{\PYGZsq{}}\PYG{l+s+s1}{,}\PYG{l+s+s1}{\PYGZsq{}}\PYG{o}{.}\PYG{n}{join}\PYG{p}{(}\PYG{n+nb}{map}\PYG{p}{(}\PYG{n+nb}{str}\PYG{p}{,}\PYG{+w}{ }\PYG{n}{key}\PYG{p}{)}\PYG{p}{)}\PYG{l+s+si}{\PYGZcb{}}\PYG{l+s+s2}{]}\PYG{l+s+s2}{\PYGZdq{}}\PYG{p}{,} \PYG{n}{ub}\PYG{o}{=}\PYG{n+nb}{sum}\PYG{p}{(}\PYG{n}{p}\PYG{o}{.}\PYG{n}{upper}\PYG{p}{[}\PYG{n}{i}\PYG{p}{]} \PYG{k}{for} \PYG{n}{i} \PYG{o+ow}{in} \PYG{n}{indices}\PYG{p}{)}\PYG{p}{)}
                \PYG{n}{target}\PYG{p}{[}\PYG{n}{key}\PYG{p}{]} \PYG{o}{=} \PYG{n}{aggregate}
                \PYG{n}{row} \PYG{o}{=} \PYG{p}{\PYGZob{}}\PYG{n}{index}\PYG{p}{:} \PYG{o}{\PYGZhy{}}\PYG{l+m+mf}{1.0} \PYG{k}{for} \PYG{n}{index} \PYG{o+ow}{in} \PYG{n}{indices}\PYG{p}{\PYGZcb{}}
                \PYG{n}{row}\PYG{p}{[}\PYG{n}{aggregate}\PYG{p}{]} \PYG{o}{=} \PYG{l+m+mf}{1.0}
                \PYG{n}{p}\PYG{o}{.}\PYG{n}{add\PYGZus{}constraint}\PYG{p}{(}\PYG{n}{row}\PYG{p}{,} \PYG{n}{lb}\PYG{o}{=}\PYG{l+m+mf}{0.0}\PYG{p}{,} \PYG{n}{ub}\PYG{o}{=}\PYG{l+m+mf}{0.0}\PYG{p}{)}
    \PYG{n}{model}\PYG{o}{.}\PYG{n}{variables}\PYG{o}{.}\PYG{n}{update}\PYG{p}{(}\PYG{n}{maps}\PYG{p}{)}
    \PYG{k}{return} \PYG{n}{BuildingThermalBlock}\PYG{p}{(}\PYG{n}{inputs}\PYG{p}{,} \PYG{n}{temperature}\PYG{p}{,} \PYG{n}{heat}\PYG{p}{,} \PYG{n}{cooling}\PYG{p}{,} \PYG{n}{mode}\PYG{p}{,} \PYG{n}{node\PYGZus{}heat}\PYG{p}{,} \PYG{n}{node\PYGZus{}cooling}\PYG{p}{)}
\end{Verbatim}

\endgroup

\noindent\textbf{结果核对。}读取逐楼轨迹的\path{temperature_end_c}、实际冷热交付与RC残差，并结合舒适摘要的低温、超温及末态偏差核验；对应登记的RC建筑与边界组，不能用聚合平均温度替代逐楼边界。

\Needspace{12\baselineskip}
\subsection*{15．逐时配电接入容量及可靠容量方法入口}
build方法内部的连续上下文片段。完整展示逐时电网接入容量、可用率及微电网净负荷约束，并保留后续可靠容量方法的参数准备和合法性检查；两种可靠容量方法的实际装配位于后续源码。

\noindent\textbf{输入：}已注册的热泵、电锅炉、循环泵等电负荷系数，既有和新增接入容量上限（kW）、增容造价（元/kW）、备用裕度、可用率时序，以及可选微电网变量。\par
\noindent\textbf{输出：}可选新增接入容量及投资账项、逐时净电负荷容量上界；同时准备可靠本地电源容量行并校验\path{firm_capacity_method}的方法身份。\par

\noindent\textbf{实现思路。}程序读取统一电负荷登记表，将热泵、电锅炉和泵耗等形成$P_t$（kW）。令备用倍率$m=1+r$、接入可用率$\alpha_t\in[0,1]$，逐时约束为$m(P_t+P_{\mathrm{ch}}-P_{\mathrm{local}}-P_{\mathrm{dis}})\le\alpha_t(K_0+K_{\mathrm{new}})$，微电网项仅在其启用时加入。新增容量$K_{\mathrm{new}}$（kW）同时进入投资账本，使正文供热方案受配电接入限制。本节末部只准备可靠容量方法入口，后续静态分支另行装配。

{\small 来源：\path{model.py}，原第2030--2089行；60行原代码、14行补充注释。}\par

\begingroup
\setstretch{1.0}
\begin{Verbatim}[commandchars=\\\{\},numbers=left,firstnumber=1,stepnumber=1,fontsize=\fontsize{9.5}{11.5}\selectfont,breaklines=true,breakanywhere=true,frame=single,framesep=3mm,numbersep=5pt,rulecolor=\color{black},xleftmargin=18pt]
\PYG{c+c1}{\PYGZsh{} Local electricity connection adequacy for heat pumps and circulation}
\PYG{c+c1}{\PYGZsh{} pumps. This prevents a heat\PYGZhy{}only optimum that cannot be connected to}
\PYG{c+c1}{\PYGZsh{} the distribution grid without reinforcement.}
\PYG{c+c1}{\PYGZsh{} \CJKCode{以下依赖}build\CJKCode{内既有变量，配置接入模块才装配。}}
\PYG{n}{grid} \PYG{o}{=} \PYG{n+nb+bp}{self}\PYG{o}{.}\PYG{n}{config}\PYG{o}{.}\PYG{n}{get}\PYG{p}{(}\PYG{l+s+s2}{\PYGZdq{}}\PYG{l+s+s2}{grid\PYGZus{}connection}\PYG{l+s+s2}{\PYGZdq{}}\PYG{p}{)}
\PYG{k}{if} \PYG{n}{grid}\PYG{p}{:}
    \PYG{c+c1}{\PYGZsh{} \CJKCode{既有和新增接入容量均为}kW\CJKCode{，新增量有投资上界。}}
    \PYG{n}{existing\PYGZus{}grid} \PYG{o}{=} \PYG{n+nb}{float}\PYG{p}{(}\PYG{n}{grid}\PYG{o}{.}\PYG{n}{get}\PYG{p}{(}\PYG{l+s+s2}{\PYGZdq{}}\PYG{l+s+s2}{existing\PYGZus{}capacity\PYGZus{}kw}\PYG{l+s+s2}{\PYGZdq{}}\PYG{p}{,} \PYG{l+m+mf}{0.0}\PYG{p}{)}\PYG{p}{)}
    \PYG{n}{max\PYGZus{}new\PYGZus{}grid} \PYG{o}{=} \PYG{n+nb}{float}\PYG{p}{(}\PYG{n}{grid}\PYG{o}{.}\PYG{n}{get}\PYG{p}{(}\PYG{l+s+s2}{\PYGZdq{}}\PYG{l+s+s2}{max\PYGZus{}new\PYGZus{}capacity\PYGZus{}kw}\PYG{l+s+s2}{\PYGZdq{}}\PYG{p}{,} \PYG{l+m+mf}{0.0}\PYG{p}{)}\PYG{p}{)}
    \PYG{n}{grid\PYGZus{}idx} \PYG{o}{=} \PYG{k+kc}{None}
    \PYG{c+c1}{\PYGZsh{} \CJKCode{不允许增容时仅使用既有容量，不建立新增变量。}}
    \PYG{k}{if} \PYG{n}{max\PYGZus{}new\PYGZus{}grid} \PYG{o}{\PYGZgt{}} \PYG{l+m+mi}{0}\PYG{p}{:}
        \PYG{n}{grid\PYGZus{}idx} \PYG{o}{=} \PYG{n}{p}\PYG{o}{.}\PYG{n}{add\PYGZus{}var}\PYG{p}{(}\PYG{l+s+s2}{\PYGZdq{}}\PYG{l+s+s2}{grid\PYGZus{}connection\PYGZus{}capacity}\PYG{l+s+s2}{\PYGZdq{}}\PYG{p}{,} \PYG{n}{ub}\PYG{o}{=}\PYG{n}{max\PYGZus{}new\PYGZus{}grid}\PYG{p}{)}
        \PYG{n+nb+bp}{self}\PYG{o}{.}\PYG{n}{variables}\PYG{p}{[}\PYG{l+s+s2}{\PYGZdq{}}\PYG{l+s+s2}{grid\PYGZus{}connection\PYGZus{}capacity}\PYG{l+s+s2}{\PYGZdq{}}\PYG{p}{]}\PYG{p}{[}\PYG{l+s+s2}{\PYGZdq{}}\PYG{l+s+s2}{system}\PYG{l+s+s2}{\PYGZdq{}}\PYG{p}{]} \PYG{o}{=} \PYG{n}{grid\PYGZus{}idx}
        \PYG{c+c1}{\PYGZsh{} \CJKCode{增容造价按元每}kW\CJKCode{进入既有投资年化账本。}}
        \PYG{n+nb+bp}{self}\PYG{o}{.}\PYG{n}{\PYGZus{}record\PYGZus{}investment\PYGZus{}cost}\PYG{p}{(}
            \PYG{n}{grid\PYGZus{}idx}\PYG{p}{,} \PYG{l+s+s2}{\PYGZdq{}}\PYG{l+s+s2}{grid\PYGZus{}connection\PYGZus{}capacity}\PYG{l+s+s2}{\PYGZdq{}}\PYG{p}{,} \PYG{l+s+s2}{\PYGZdq{}}\PYG{l+s+s2}{system}\PYG{l+s+s2}{\PYGZdq{}}\PYG{p}{,} \PYG{n+nb}{float}\PYG{p}{(}\PYG{n}{grid}\PYG{o}{.}\PYG{n}{get}\PYG{p}{(}\PYG{l+s+s2}{\PYGZdq{}}\PYG{l+s+s2}{capex\PYGZus{}yuan\PYGZus{}per\PYGZus{}kw}\PYG{l+s+s2}{\PYGZdq{}}\PYG{p}{,} \PYG{l+m+mf}{0.0}\PYG{p}{)}\PYG{p}{)}
        \PYG{p}{)}
    \PYG{c+c1}{\PYGZsh{} \CJKCode{裕度无量纲，以}1\CJKCode{加备用比例放大净电负荷。}}
    \PYG{n}{margin} \PYG{o}{=} \PYG{l+m+mf}{1.0} \PYG{o}{+} \PYG{n+nb}{float}\PYG{p}{(}\PYG{n}{grid}\PYG{o}{.}\PYG{n}{get}\PYG{p}{(}\PYG{l+s+s2}{\PYGZdq{}}\PYG{l+s+s2}{reserve\PYGZus{}margin}\PYG{l+s+s2}{\PYGZdq{}}\PYG{p}{,} \PYG{l+m+mf}{0.0}\PYG{p}{)}\PYG{p}{)}
    \PYG{k}{for} \PYG{n}{t} \PYG{o+ow}{in} \PYG{n+nb+bp}{self}\PYG{o}{.}\PYG{n}{times}\PYG{p}{:}
        \PYG{n}{grid\PYGZus{}availability} \PYG{o}{=} \PYG{l+m+mf}{1.0}
        \PYG{c+c1}{\PYGZsh{} \CJKCode{若登记可用率时序，则逐时截取到}0\CJKCode{至}1\CJKCode{。}}
        \PYG{k}{if} \PYG{n}{grid}\PYG{o}{.}\PYG{n}{get}\PYG{p}{(}\PYG{l+s+s2}{\PYGZdq{}}\PYG{l+s+s2}{availability\PYGZus{}column}\PYG{l+s+s2}{\PYGZdq{}}\PYG{p}{)}\PYG{p}{:}
            \PYG{n}{grid\PYGZus{}availability} \PYG{o}{=} \PYG{n+nb+bp}{self}\PYG{o}{.}\PYG{n}{\PYGZus{}clip\PYGZus{}scalar}\PYG{p}{(}
                \PYG{n+nb+bp}{self}\PYG{o}{.}\PYG{n}{\PYGZus{}ts\PYGZus{}value}\PYG{p}{(}\PYG{n}{t}\PYG{p}{,} \PYG{n}{grid}\PYG{p}{[}\PYG{l+s+s2}{\PYGZdq{}}\PYG{l+s+s2}{availability\PYGZus{}column}\PYG{l+s+s2}{\PYGZdq{}}\PYG{p}{]}\PYG{p}{)}\PYG{p}{,}
                \PYG{l+m+mf}{0.0}\PYG{p}{,}
                \PYG{l+m+mf}{1.0}\PYG{p}{,}
            \PYG{p}{)}
        \PYG{c+c1}{\PYGZsh{} \CJKCode{读取统一电负荷行，热泵和泵电耗不另算一遍。}}
        \PYG{n}{row} \PYG{o}{=} \PYG{n+nb+bp}{self}\PYG{o}{.}\PYG{n}{\PYGZus{}electric\PYGZus{}load\PYGZus{}row}\PYG{p}{(}\PYG{n}{t}\PYG{p}{,} \PYG{n}{margin}\PYG{o}{=}\PYG{n}{margin}\PYG{p}{)}
        \PYG{c+c1}{\PYGZsh{} \CJKCode{微电网分支仅在配置启用时修正接入净负荷。}}
        \PYG{k}{if} \PYG{n}{microgrid}\PYG{p}{:}
            \PYG{k}{for} \PYG{n}{category}\PYG{p}{,} \PYG{n}{sign} \PYG{o+ow}{in} \PYG{p}{(}
                \PYG{c+c1}{\PYGZsh{} \CJKCode{储电充电增加负荷，本地发电与放电减少负荷。}}
                \PYG{p}{(}\PYG{l+s+s2}{\PYGZdq{}}\PYG{l+s+s2}{microgrid\PYGZus{}battery\PYGZus{}charge}\PYG{l+s+s2}{\PYGZdq{}}\PYG{p}{,} \PYG{l+m+mf}{1.0}\PYG{p}{)}\PYG{p}{,}
                \PYG{p}{(}\PYG{l+s+s2}{\PYGZdq{}}\PYG{l+s+s2}{microgrid\PYGZus{}renewable\PYGZus{}use}\PYG{l+s+s2}{\PYGZdq{}}\PYG{p}{,} \PYG{o}{\PYGZhy{}}\PYG{l+m+mf}{1.0}\PYG{p}{)}\PYG{p}{,}
                \PYG{p}{(}\PYG{l+s+s2}{\PYGZdq{}}\PYG{l+s+s2}{microgrid\PYGZus{}backup\PYGZus{}generation}\PYG{l+s+s2}{\PYGZdq{}}\PYG{p}{,} \PYG{o}{\PYGZhy{}}\PYG{l+m+mf}{1.0}\PYG{p}{)}\PYG{p}{,}
                \PYG{p}{(}\PYG{l+s+s2}{\PYGZdq{}}\PYG{l+s+s2}{microgrid\PYGZus{}battery\PYGZus{}discharge}\PYG{l+s+s2}{\PYGZdq{}}\PYG{p}{,} \PYG{o}{\PYGZhy{}}\PYG{l+m+mf}{1.0}\PYG{p}{)}\PYG{p}{,}
            \PYG{p}{)}\PYG{p}{:}
                \PYG{n+nb+bp}{self}\PYG{o}{.}\PYG{n}{\PYGZus{}add}\PYG{p}{(}\PYG{n}{row}\PYG{p}{,} \PYG{n+nb+bp}{self}\PYG{o}{.}\PYG{n}{variables}\PYG{p}{[}\PYG{n}{category}\PYG{p}{]}\PYG{p}{[}\PYG{n}{t}\PYG{p}{]}\PYG{p}{,} \PYG{n}{margin} \PYG{o}{*} \PYG{n}{sign}\PYG{p}{)}
        \PYG{c+c1}{\PYGZsh{} \CJKCode{新增容量移到左端且乘可用率，保持线性上界。}}
        \PYG{k}{if} \PYG{n}{grid\PYGZus{}idx} \PYG{o+ow}{is} \PYG{o+ow}{not} \PYG{k+kc}{None}\PYG{p}{:}
            \PYG{n+nb+bp}{self}\PYG{o}{.}\PYG{n}{\PYGZus{}add}\PYG{p}{(}\PYG{n}{row}\PYG{p}{,} \PYG{n}{grid\PYGZus{}idx}\PYG{p}{,} \PYG{o}{\PYGZhy{}}\PYG{n}{grid\PYGZus{}availability}\PYG{p}{)}
        \PYG{c+c1}{\PYGZsh{} \CJKCode{净负荷须小于可用率乘既有加新增接入容量。}}
        \PYG{n}{p}\PYG{o}{.}\PYG{n}{add\PYGZus{}constraint}\PYG{p}{(}\PYG{n}{row}\PYG{p}{,} \PYG{n}{ub}\PYG{o}{=}\PYG{n}{existing\PYGZus{}grid} \PYG{o}{*} \PYG{n}{grid\PYGZus{}availability}\PYG{p}{)}

    \PYG{c+c1}{\PYGZsh{} Ensure that electrically supplied firm heat capacity is actually}
    \PYG{c+c1}{\PYGZsh{} deliverable through the grid connection in a contingency.  The}
    \PYG{c+c1}{\PYGZsh{} legacy method assumes every installed electric heat source runs}
    \PYG{c+c1}{\PYGZsh{} simultaneously.  ``portfolio\PYGZus{}dispatch`` instead selects a}
    \PYG{c+c1}{\PYGZsh{} feasible firm\PYGZhy{}heat portfolio for the peak (and each configured}
    \PYG{c+c1}{\PYGZsh{} source outage), avoiding double\PYGZhy{}counting mutually substitutable}
    \PYG{c+c1}{\PYGZsh{} heat pumps and backup boilers.}
    \PYG{c+c1}{\PYGZsh{} \CJKCode{以下准备静态可靠容量检查；同时率与裕度不同。}}
    \PYG{n}{simultaneity} \PYG{o}{=} \PYG{n+nb}{float}\PYG{p}{(}\PYG{n}{grid}\PYG{o}{.}\PYG{n}{get}\PYG{p}{(}\PYG{l+s+s2}{\PYGZdq{}}\PYG{l+s+s2}{firm\PYGZus{}simultaneity}\PYG{l+s+s2}{\PYGZdq{}}\PYG{p}{,} \PYG{l+m+mf}{1.0}\PYG{p}{)}\PYG{p}{)}
    \PYG{n}{existing\PYGZus{}firm\PYGZus{}local} \PYG{o}{=} \PYG{l+m+mf}{0.0}
    \PYG{n}{firm\PYGZus{}local\PYGZus{}row}\PYG{p}{:} \PYG{n+nb}{dict}\PYG{p}{[}\PYG{n+nb}{int}\PYG{p}{,} \PYG{n+nb}{float}\PYG{p}{]} \PYG{o}{=} \PYG{p}{\PYGZob{}}\PYG{p}{\PYGZcb{}}
    \PYG{k}{if} \PYG{n}{microgrid}\PYG{p}{:}
        \PYG{c+c1}{\PYGZsh{} \CJKCode{本地可再生可靠容量按登记出力因子折算。}}
        \PYG{n}{renewable\PYGZus{}firm} \PYG{o}{=} \PYG{n+nb}{float}\PYG{p}{(}\PYG{n}{microgrid}\PYG{o}{.}\PYG{n}{get}\PYG{p}{(}\PYG{l+s+s2}{\PYGZdq{}}\PYG{l+s+s2}{renewable\PYGZus{}firm\PYGZus{}output\PYGZus{}factor}\PYG{l+s+s2}{\PYGZdq{}}\PYG{p}{,} \PYG{l+m+mf}{0.15}\PYG{p}{)}\PYG{p}{)}
        \PYG{k}{for} \PYG{n}{asset}\PYG{p}{,} \PYG{n}{factor} \PYG{o+ow}{in} \PYG{p}{(}\PYG{p}{(}\PYG{l+s+s2}{\PYGZdq{}}\PYG{l+s+s2}{renewable}\PYG{l+s+s2}{\PYGZdq{}}\PYG{p}{,} \PYG{n}{renewable\PYGZus{}firm}\PYG{p}{)}\PYG{p}{,} \PYG{p}{(}\PYG{l+s+s2}{\PYGZdq{}}\PYG{l+s+s2}{battery\PYGZus{}power}\PYG{l+s+s2}{\PYGZdq{}}\PYG{p}{,} \PYG{l+m+mf}{1.0}\PYG{p}{)}\PYG{p}{,} \PYG{p}{(}\PYG{l+s+s2}{\PYGZdq{}}\PYG{l+s+s2}{backup}\PYG{l+s+s2}{\PYGZdq{}}\PYG{p}{,} \PYG{l+m+mf}{1.0}\PYG{p}{)}\PYG{p}{)}\PYG{p}{:}
            \PYG{n}{idx}\PYG{p}{,} \PYG{n}{existing}\PYG{p}{,} \PYG{n}{\PYGZus{}} \PYG{o}{=} \PYG{n}{microgrid\PYGZus{}specs}\PYG{p}{[}\PYG{n}{asset}\PYG{p}{]}
            \PYG{n}{existing\PYGZus{}firm\PYGZus{}local} \PYG{o}{+}\PYG{o}{=} \PYG{n}{factor} \PYG{o}{*} \PYG{n}{existing}
            \PYG{k}{if} \PYG{n}{idx} \PYG{o+ow}{is} \PYG{o+ow}{not} \PYG{k+kc}{None}\PYG{p}{:}
                \PYG{n}{firm\PYGZus{}local\PYGZus{}row}\PYG{p}{[}\PYG{n}{idx}\PYG{p}{]} \PYG{o}{=} \PYG{o}{\PYGZhy{}}\PYG{n}{factor}

    \PYG{c+c1}{\PYGZsh{} \CJKCode{选择可靠容量方法；本片段尚未装配该方法约束。}}
    \PYG{n}{firm\PYGZus{}method} \PYG{o}{=} \PYG{n+nb}{str}\PYG{p}{(}\PYG{n}{grid}\PYG{o}{.}\PYG{n}{get}\PYG{p}{(}\PYG{l+s+s2}{\PYGZdq{}}\PYG{l+s+s2}{firm\PYGZus{}capacity\PYGZus{}method}\PYG{l+s+s2}{\PYGZdq{}}\PYG{p}{,} \PYG{l+s+s2}{\PYGZdq{}}\PYG{l+s+s2}{all\PYGZus{}installed\PYGZus{}electric}\PYG{l+s+s2}{\PYGZdq{}}\PYG{p}{)}\PYG{p}{)}
    \PYG{k}{if} \PYG{n}{firm\PYGZus{}method} \PYG{o+ow}{not} \PYG{o+ow}{in} \PYG{p}{\PYGZob{}}\PYG{l+s+s2}{\PYGZdq{}}\PYG{l+s+s2}{all\PYGZus{}installed\PYGZus{}electric}\PYG{l+s+s2}{\PYGZdq{}}\PYG{p}{,} \PYG{l+s+s2}{\PYGZdq{}}\PYG{l+s+s2}{portfolio\PYGZus{}dispatch}\PYG{l+s+s2}{\PYGZdq{}}\PYG{p}{\PYGZcb{}}\PYG{p}{:}
        \PYG{k}{raise} \PYG{n+ne}{ValueError}\PYG{p}{(}
            \PYG{l+s+s2}{\PYGZdq{}}\PYG{l+s+s2}{grid\PYGZus{}connection.firm\PYGZus{}capacity\PYGZus{}method \CJKCode{仅支持} }\PYG{l+s+s2}{\PYGZdq{}}
            \PYG{l+s+s2}{\PYGZdq{}}\PYG{l+s+s2}{all\PYGZus{}installed\PYGZus{}electric \CJKCode{或} portfolio\PYGZus{}dispatch}\PYG{l+s+s2}{\PYGZdq{}}
        \PYG{p}{)}
\end{Verbatim}

\endgroup

\noindent\textbf{结果核对。}核验逐时电负荷行、接入可用率和\path{grid_connection_capacity_kw}，并记录\path{firm_grid_capacity_method}；本节证明接入约束的装配方式，可靠容量核验应追踪后续实际选择的分支。

\Needspace{12\baselineskip}
\subsection*{16．硬能量守恒下的最小一维W1权重线性规划}
完整展示实际代表日保持不变、仅优化权重的主算法；以能量和总天数为硬等式，在权重下界内最小化归一化的一维经验分布W1，并独立重算残差。

\noindent\textbf{输入：}完整等间隔实测日序列、已选实际代表日矩阵、原天数权重、逐日权重下界与正数数值容差；代表日值与观测值须同量纲。\par
\noindent\textbf{输出：}通过独立复核的权重、归一化W1及原量纲W1、能量/天数/下界/目标残差；另明确返回不可行、零能量无定义或求解失败身份。\par

\noindent\textbf{实现思路。}设实测经验分布为$F$，第$j$个固定代表日的CDF贡献为$C_{ij}$，原总天数为$W$。代码用$w=\ell+z$把下界转为$z\ge0$，同时硬性保持加权总能量与总天数；两组不等式使辅助变量包住各支撑区间的CDF差，目标按区间宽度和实测均值加权，得到最小归一化一维W1。其优势是守恒条件不会被拟合目标让步，数据与代表日也不被改写；所得最优权重仍须独立重算残差。

{\small 来源：\path{hard_weight_calibration.py}，原第13--97行；85行原代码、15行补充注释。}\par

\begingroup
\setstretch{1.0}
\begin{Verbatim}[commandchars=\\\{\},numbers=left,firstnumber=1,stepnumber=1,fontsize=\fontsize{9.5}{11.5}\selectfont,breaklines=true,breakanywhere=true,frame=single,framesep=3mm,numbersep=5pt,rulecolor=\color{black},xleftmargin=18pt]
\PYG{k}{def}\PYG{+w}{ }\PYG{n+nf}{minimum\PYGZus{}w1\PYGZus{}energy\PYGZus{}weights}\PYG{p}{(}\PYG{n}{observed}\PYG{p}{,} \PYG{n}{representative\PYGZus{}profiles}\PYG{p}{,} \PYG{n}{original\PYGZus{}weights}\PYG{p}{,}
                              \PYG{n}{lower\PYGZus{}bounds}\PYG{p}{,} \PYG{o}{*}\PYG{p}{,} \PYG{n}{numerical\PYGZus{}tolerance}\PYG{o}{=}\PYG{l+m+mf}{1e\PYGZhy{}8}\PYG{p}{)}\PYG{p}{:}
\PYG{+w}{    }\PYG{l+s+sd}{\PYGZdq{}\PYGZdq{}\PYGZdq{}Minimize normalized 1\PYGZhy{}D W1 with exact energy and day conservation.}

\PYG{l+s+sd}{    Each representative row is one equally sampled actual day; observations}
\PYG{l+s+sd}{    comprise complete equally sampled days. All values must be finite and}
\PYG{l+s+sd}{    nonnegative. With zero observed mean the normalized metric is undefined.}
\PYG{l+s+sd}{    Returns explicit success, infeasibility, or failure with no data mutation.}
\PYG{l+s+sd}{    \PYGZdq{}\PYGZdq{}\PYGZdq{}}
    \PYG{c+c1}{\PYGZsh{} \CJKCode{只读入观测与固定代表日；优化对象是天数权重。}}
    \PYG{n}{values} \PYG{o}{=} \PYG{n}{np}\PYG{o}{.}\PYG{n}{asarray}\PYG{p}{(}\PYG{n}{observed}\PYG{p}{,} \PYG{n}{dtype}\PYG{o}{=}\PYG{n+nb}{float}\PYG{p}{)}
    \PYG{n}{profiles} \PYG{o}{=} \PYG{n}{np}\PYG{o}{.}\PYG{n}{asarray}\PYG{p}{(}\PYG{n}{representative\PYGZus{}profiles}\PYG{p}{,} \PYG{n}{dtype}\PYG{o}{=}\PYG{n+nb}{float}\PYG{p}{)}
    \PYG{n}{original} \PYG{o}{=} \PYG{n}{np}\PYG{o}{.}\PYG{n}{asarray}\PYG{p}{(}\PYG{n}{original\PYGZus{}weights}\PYG{p}{,} \PYG{n}{dtype}\PYG{o}{=}\PYG{n+nb}{float}\PYG{p}{)}
    \PYG{n}{lower} \PYG{o}{=} \PYG{n}{np}\PYG{o}{.}\PYG{n}{asarray}\PYG{p}{(}\PYG{n}{lower\PYGZus{}bounds}\PYG{p}{,} \PYG{n}{dtype}\PYG{o}{=}\PYG{n+nb}{float}\PYG{p}{)}
    \PYG{k}{if} \PYG{n}{values}\PYG{o}{.}\PYG{n}{ndim} \PYG{o}{!=} \PYG{l+m+mi}{1} \PYG{o+ow}{or} \PYG{o+ow}{not} \PYG{n}{values}\PYG{o}{.}\PYG{n}{size} \PYG{o+ow}{or} \PYG{n}{profiles}\PYG{o}{.}\PYG{n}{ndim} \PYG{o}{!=} \PYG{l+m+mi}{2} \PYG{o+ow}{or} \PYG{o+ow}{not} \PYG{n+nb}{all}\PYG{p}{(}\PYG{n}{profiles}\PYG{o}{.}\PYG{n}{shape}\PYG{p}{)}\PYG{p}{:}
        \PYG{k}{raise} \PYG{n+ne}{ValueError}\PYG{p}{(}\PYG{l+s+s2}{\PYGZdq{}}\PYG{l+s+s2}{observed must be nonempty 1D and representative profiles nonempty 2D}\PYG{l+s+s2}{\PYGZdq{}}\PYG{p}{)}
    \PYG{c+c1}{\PYGZsh{} \CJKCode{每行是一个实际代表日；观测须由完整同长日组成。}}
    \PYG{n}{m}\PYG{p}{,} \PYG{n}{hours} \PYG{o}{=} \PYG{n}{profiles}\PYG{o}{.}\PYG{n}{shape}
    \PYG{k}{if} \PYG{n}{original}\PYG{o}{.}\PYG{n}{shape} \PYG{o}{!=} \PYG{p}{(}\PYG{n}{m}\PYG{p}{,}\PYG{p}{)} \PYG{o+ow}{or} \PYG{n}{lower}\PYG{o}{.}\PYG{n}{shape} \PYG{o}{!=} \PYG{p}{(}\PYG{n}{m}\PYG{p}{,}\PYG{p}{)} \PYG{o+ow}{or} \PYG{n+nb}{len}\PYG{p}{(}\PYG{n}{values}\PYG{p}{)} \PYG{o}{\PYGZpc{}} \PYG{n}{hours}\PYG{p}{:}
        \PYG{k}{raise} \PYG{n+ne}{ValueError}\PYG{p}{(}\PYG{l+s+s2}{\PYGZdq{}}\PYG{l+s+s2}{weight/profile dimensions or complete\PYGZhy{}day observation length are invalid}\PYG{l+s+s2}{\PYGZdq{}}\PYG{p}{)}
    \PYG{k}{if} \PYG{n+nb}{any}\PYG{p}{(}\PYG{o+ow}{not} \PYG{n}{np}\PYG{o}{.}\PYG{n}{isfinite}\PYG{p}{(}\PYG{n}{x}\PYG{p}{)}\PYG{o}{.}\PYG{n}{all}\PYG{p}{(}\PYG{p}{)} \PYG{o+ow}{or} \PYG{p}{(}\PYG{n}{x} \PYG{o}{\PYGZlt{}} \PYG{l+m+mi}{0}\PYG{p}{)}\PYG{o}{.}\PYG{n}{any}\PYG{p}{(}\PYG{p}{)} \PYG{k}{for} \PYG{n}{x} \PYG{o+ow}{in} \PYG{p}{(}\PYG{n}{values}\PYG{p}{,} \PYG{n}{profiles}\PYG{p}{,} \PYG{n}{original}\PYG{p}{,} \PYG{n}{lower}\PYG{p}{)}\PYG{p}{)}\PYG{p}{:}
        \PYG{k}{raise} \PYG{n+ne}{ValueError}\PYG{p}{(}\PYG{l+s+s2}{\PYGZdq{}}\PYG{l+s+s2}{observations, profiles, weights and lower bounds must be finite and nonnegative}\PYG{l+s+s2}{\PYGZdq{}}\PYG{p}{)}
    \PYG{n}{tolerance} \PYG{o}{=} \PYG{n+nb}{float}\PYG{p}{(}\PYG{n}{numerical\PYGZus{}tolerance}\PYG{p}{)}
    \PYG{k}{if} \PYG{o+ow}{not} \PYG{n}{np}\PYG{o}{.}\PYG{n}{isfinite}\PYG{p}{(}\PYG{n}{tolerance}\PYG{p}{)} \PYG{o+ow}{or} \PYG{n}{tolerance} \PYG{o}{\PYGZlt{}}\PYG{o}{=} \PYG{l+m+mi}{0}\PYG{p}{:}
        \PYG{k}{raise} \PYG{n+ne}{ValueError}\PYG{p}{(}\PYG{l+s+s2}{\PYGZdq{}}\PYG{l+s+s2}{numerical\PYGZus{}tolerance must be positive and finite}\PYG{l+s+s2}{\PYGZdq{}}\PYG{p}{)}
    \PYG{c+c1}{\PYGZsh{} \CJKCode{原权重总天数乘每日采样数须等于观测长度。}}
    \PYG{n}{total} \PYG{o}{=} \PYG{n+nb}{float}\PYG{p}{(}\PYG{n}{original}\PYG{o}{.}\PYG{n}{sum}\PYG{p}{(}\PYG{p}{)}\PYG{p}{)}
    \PYG{k}{if} \PYG{o+ow}{not} \PYG{n}{np}\PYG{o}{.}\PYG{n}{isclose}\PYG{p}{(}\PYG{n}{total} \PYG{o}{*} \PYG{n}{hours}\PYG{p}{,} \PYG{n+nb}{len}\PYG{p}{(}\PYG{n}{values}\PYG{p}{)}\PYG{p}{,} \PYG{n}{rtol}\PYG{o}{=}\PYG{l+m+mi}{0}\PYG{p}{,} \PYG{n}{atol}\PYG{o}{=}\PYG{n}{tolerance}\PYG{p}{)}\PYG{p}{:}
        \PYG{k}{raise} \PYG{n+ne}{ValueError}\PYG{p}{(}\PYG{l+s+s2}{\PYGZdq{}}\PYG{l+s+s2}{original day weight does not represent the observed complete\PYGZhy{}day count}\PYG{l+s+s2}{\PYGZdq{}}\PYG{p}{)}
    \PYG{c+c1}{\PYGZsh{} \CJKCode{报告明确保留未验证工程运行的身份。}}
    \PYG{n}{report} \PYG{o}{=} \PYG{p}{\PYGZob{}}\PYG{l+s+s2}{\PYGZdq{}}\PYG{l+s+s2}{method}\PYG{l+s+s2}{\PYGZdq{}}\PYG{p}{:} \PYG{l+s+s2}{\PYGZdq{}}\PYG{l+s+s2}{standard\PYGZus{}LP\PYGZus{}hard\PYGZus{}energy\PYGZus{}minimum\PYGZus{}1D\PYGZus{}CDF\PYGZus{}W1}\PYG{l+s+s2}{\PYGZdq{}}\PYG{p}{,}
              \PYG{l+s+s2}{\PYGZdq{}}\PYG{l+s+s2}{energy\PYGZus{}matching}\PYG{l+s+s2}{\PYGZdq{}}\PYG{p}{:} \PYG{l+s+s2}{\PYGZdq{}}\PYG{l+s+s2}{hard\PYGZus{}equality\PYGZus{}with\PYGZus{}independent\PYGZus{}residual\PYGZus{}check}\PYG{l+s+s2}{\PYGZdq{}}\PYG{p}{,}
              \PYG{l+s+s2}{\PYGZdq{}}\PYG{l+s+s2}{source\PYGZus{}data\PYGZus{}modified}\PYG{l+s+s2}{\PYGZdq{}}\PYG{p}{:} \PYG{k+kc}{False}\PYG{p}{,} \PYG{l+s+s2}{\PYGZdq{}}\PYG{l+s+s2}{selected\PYGZus{}profiles\PYGZus{}modified}\PYG{l+s+s2}{\PYGZdq{}}\PYG{p}{:} \PYG{k+kc}{False}\PYG{p}{,}
              \PYG{l+s+s2}{\PYGZdq{}}\PYG{l+s+s2}{engineering\PYGZus{}operation\PYGZus{}validated}\PYG{l+s+s2}{\PYGZdq{}}\PYG{p}{:} \PYG{k+kc}{False}\PYG{p}{,} \PYG{l+s+s2}{\PYGZdq{}}\PYG{l+s+s2}{total\PYGZus{}weight\PYGZus{}days}\PYG{l+s+s2}{\PYGZdq{}}\PYG{p}{:} \PYG{n}{total}\PYG{p}{,}
              \PYG{l+s+s2}{\PYGZdq{}}\PYG{l+s+s2}{original\PYGZus{}weights}\PYG{l+s+s2}{\PYGZdq{}}\PYG{p}{:} \PYG{n}{original}\PYG{o}{.}\PYG{n}{tolist}\PYG{p}{(}\PYG{p}{)}\PYG{p}{,} \PYG{l+s+s2}{\PYGZdq{}}\PYG{l+s+s2}{lower\PYGZus{}bounds}\PYG{l+s+s2}{\PYGZdq{}}\PYG{p}{:} \PYG{n}{lower}\PYG{o}{.}\PYG{n}{tolist}\PYG{p}{(}\PYG{p}{)}\PYG{p}{\PYGZcb{}}
    \PYG{c+c1}{\PYGZsh{} \CJKCode{下界总和超过总天数时直接返回不可行。}}
    \PYG{k}{if} \PYG{n}{lower}\PYG{o}{.}\PYG{n}{sum}\PYG{p}{(}\PYG{p}{)} \PYG{o}{\PYGZgt{}} \PYG{n}{total} \PYG{o}{+} \PYG{n}{tolerance}\PYG{p}{:}
        \PYG{k}{return} \PYG{p}{\PYGZob{}}\PYG{o}{*}\PYG{o}{*}\PYG{n}{report}\PYG{p}{,} \PYG{l+s+s2}{\PYGZdq{}}\PYG{l+s+s2}{status}\PYG{l+s+s2}{\PYGZdq{}}\PYG{p}{:} \PYG{l+s+s2}{\PYGZdq{}}\PYG{l+s+s2}{infeasible}\PYG{l+s+s2}{\PYGZdq{}}\PYG{p}{,} \PYG{l+s+s2}{\PYGZdq{}}\PYG{l+s+s2}{reason}\PYG{l+s+s2}{\PYGZdq{}}\PYG{p}{:} \PYG{l+s+s2}{\PYGZdq{}}\PYG{l+s+s2}{lower\PYGZus{}bounds\PYGZus{}exceed\PYGZus{}total\PYGZus{}weight}\PYG{l+s+s2}{\PYGZdq{}}\PYG{p}{,} \PYG{l+s+s2}{\PYGZdq{}}\PYG{l+s+s2}{weights}\PYG{l+s+s2}{\PYGZdq{}}\PYG{p}{:} \PYG{k+kc}{None}\PYG{p}{\PYGZcb{}}
    \PYG{c+c1}{\PYGZsh{} \CJKCode{零均值时归一化}W1\CJKCode{无定义，不能伪造零误差。}}
    \PYG{n}{mean} \PYG{o}{=} \PYG{n+nb}{float}\PYG{p}{(}\PYG{n}{values}\PYG{o}{.}\PYG{n}{mean}\PYG{p}{(}\PYG{p}{)}\PYG{p}{)}
    \PYG{k}{if} \PYG{n}{mean} \PYG{o}{==} \PYG{l+m+mi}{0}\PYG{p}{:}
        \PYG{k}{return} \PYG{p}{\PYGZob{}}\PYG{o}{*}\PYG{o}{*}\PYG{n}{report}\PYG{p}{,} \PYG{l+s+s2}{\PYGZdq{}}\PYG{l+s+s2}{status}\PYG{l+s+s2}{\PYGZdq{}}\PYG{p}{:} \PYG{l+s+s2}{\PYGZdq{}}\PYG{l+s+s2}{zero\PYGZus{}energy\PYGZus{}undefined}\PYG{l+s+s2}{\PYGZdq{}}\PYG{p}{,} \PYG{l+s+s2}{\PYGZdq{}}\PYG{l+s+s2}{reason}\PYG{l+s+s2}{\PYGZdq{}}\PYG{p}{:} \PYG{l+s+s2}{\PYGZdq{}}\PYG{l+s+s2}{normalized\PYGZus{}W1\PYGZus{}and\PYGZus{}relative\PYGZus{}energy\PYGZus{}undefined}\PYG{l+s+s2}{\PYGZdq{}}\PYG{p}{,} \PYG{l+s+s2}{\PYGZdq{}}\PYG{l+s+s2}{weights}\PYG{l+s+s2}{\PYGZdq{}}\PYG{p}{:} \PYG{k+kc}{None}\PYG{p}{\PYGZcb{}}
    \PYG{c+c1}{\PYGZsh{} \CJKCode{合并取值支撑点，将}CDF\CJKCode{差积分化为区间加权和。}}
    \PYG{n}{support} \PYG{o}{=} \PYG{n}{np}\PYG{o}{.}\PYG{n}{unique}\PYG{p}{(}\PYG{n}{np}\PYG{o}{.}\PYG{n}{concatenate}\PYG{p}{(}\PYG{p}{(}\PYG{n}{values}\PYG{p}{,} \PYG{n}{profiles}\PYG{o}{.}\PYG{n}{ravel}\PYG{p}{(}\PYG{p}{)}\PYG{p}{)}\PYG{p}{)}\PYG{p}{)}
    \PYG{n}{left}\PYG{p}{,} \PYG{n}{widths} \PYG{o}{=} \PYG{n}{support}\PYG{p}{[}\PYG{p}{:}\PYG{o}{\PYGZhy{}}\PYG{l+m+mi}{1}\PYG{p}{]}\PYG{p}{,} \PYG{n}{np}\PYG{o}{.}\PYG{n}{diff}\PYG{p}{(}\PYG{n}{support}\PYG{p}{)}
    \PYG{n}{observed\PYGZus{}cdf} \PYG{o}{=} \PYG{n}{np}\PYG{o}{.}\PYG{n}{searchsorted}\PYG{p}{(}\PYG{n}{np}\PYG{o}{.}\PYG{n}{sort}\PYG{p}{(}\PYG{n}{values}\PYG{p}{)}\PYG{p}{,} \PYG{n}{left}\PYG{p}{,} \PYG{n}{side}\PYG{o}{=}\PYG{l+s+s2}{\PYGZdq{}}\PYG{l+s+s2}{right}\PYG{l+s+s2}{\PYGZdq{}}\PYG{p}{)} \PYG{o}{/} \PYG{n+nb}{len}\PYG{p}{(}\PYG{n}{values}\PYG{p}{)}
    \PYG{c+c1}{\PYGZsh{} A day weight contributes count(profile \PYGZlt{}= support) / (H*W) to the CDF.}
    \PYG{c+c1}{\PYGZsh{} \CJKCode{每列给出一个代表日单位天数的}CDF\CJKCode{贡献。}}
    \PYG{n}{cdf} \PYG{o}{=} \PYG{n}{np}\PYG{o}{.}\PYG{n}{column\PYGZus{}stack}\PYG{p}{(}\PYG{p}{[}\PYG{n}{np}\PYG{o}{.}\PYG{n}{searchsorted}\PYG{p}{(}\PYG{n}{np}\PYG{o}{.}\PYG{n}{sort}\PYG{p}{(}\PYG{n}{p}\PYG{p}{)}\PYG{p}{,} \PYG{n}{left}\PYG{p}{,} \PYG{n}{side}\PYG{o}{=}\PYG{l+s+s2}{\PYGZdq{}}\PYG{l+s+s2}{right}\PYG{l+s+s2}{\PYGZdq{}}\PYG{p}{)} \PYG{k}{for} \PYG{n}{p} \PYG{o+ow}{in} \PYG{n}{profiles}\PYG{p}{]}\PYG{p}{)} \PYG{o}{/} \PYG{p}{(}\PYG{n}{hours} \PYG{o}{*} \PYG{n}{total}\PYG{p}{)}
    \PYG{n}{day\PYGZus{}energy}\PYG{p}{,} \PYG{n}{target} \PYG{o}{=} \PYG{n}{profiles}\PYG{o}{.}\PYG{n}{sum}\PYG{p}{(}\PYG{n}{axis}\PYG{o}{=}\PYG{l+m+mi}{1}\PYG{p}{)}\PYG{p}{,} \PYG{n+nb}{float}\PYG{p}{(}\PYG{n}{values}\PYG{o}{.}\PYG{n}{sum}\PYG{p}{(}\PYG{p}{)}\PYG{p}{)}
    \PYG{n}{n} \PYG{o}{=} \PYG{n+nb}{len}\PYG{p}{(}\PYG{n}{left}\PYG{p}{)}
    \PYG{c+c1}{\PYGZsh{} \CJKCode{令权重等于下界加非负余量，显式保证逐日下界。}}
    \PYG{n}{shifted\PYGZus{}target} \PYG{o}{=} \PYG{n}{target} \PYG{o}{\PYGZhy{}} \PYG{n+nb}{float}\PYG{p}{(}\PYG{n}{day\PYGZus{}energy} \PYG{o}{@} \PYG{n}{lower}\PYG{p}{)}
    \PYG{c+c1}{\PYGZsh{} \CJKCode{两行硬等式分别锁定总能量与总天数，并按量级缩放。}}
    \PYG{n}{equality} \PYG{o}{=} \PYG{n}{csr\PYGZus{}matrix}\PYG{p}{(}\PYG{n}{np}\PYG{o}{.}\PYG{n}{vstack}\PYG{p}{(}\PYG{p}{(}\PYG{n}{day\PYGZus{}energy} \PYG{o}{/} \PYG{n+nb}{max}\PYG{p}{(}\PYG{n}{target}\PYG{p}{,} \PYG{l+m+mi}{1}\PYG{p}{)}\PYG{p}{,} \PYG{n}{np}\PYG{o}{.}\PYG{n}{ones}\PYG{p}{(}\PYG{n}{m}\PYG{p}{)} \PYG{o}{/} \PYG{n+nb}{max}\PYG{p}{(}\PYG{n}{total}\PYG{p}{,} \PYG{l+m+mi}{1}\PYG{p}{)}\PYG{p}{)}\PYG{p}{)}\PYG{p}{)}
    \PYG{n}{equality} \PYG{o}{=} \PYG{n}{hstack}\PYG{p}{(}\PYG{p}{(}\PYG{n}{equality}\PYG{p}{,} \PYG{n}{csr\PYGZus{}matrix}\PYG{p}{(}\PYG{p}{(}\PYG{l+m+mi}{2}\PYG{p}{,} \PYG{n}{n}\PYG{p}{)}\PYG{p}{)}\PYG{p}{)}\PYG{p}{,} \PYG{n+nb}{format}\PYG{o}{=}\PYG{l+s+s2}{\PYGZdq{}}\PYG{l+s+s2}{csr}\PYG{l+s+s2}{\PYGZdq{}}\PYG{p}{)}
    \PYG{n}{rhs} \PYG{o}{=} \PYG{n}{np}\PYG{o}{.}\PYG{n}{array}\PYG{p}{(}\PYG{p}{[}\PYG{n}{shifted\PYGZus{}target}\PYG{o}{/}\PYG{n+nb}{max}\PYG{p}{(}\PYG{n}{target}\PYG{p}{,} \PYG{l+m+mi}{1}\PYG{p}{)}\PYG{p}{,} \PYG{p}{(}\PYG{n}{total}\PYG{o}{\PYGZhy{}}\PYG{n}{lower}\PYG{o}{.}\PYG{n}{sum}\PYG{p}{(}\PYG{p}{)}\PYG{p}{)}\PYG{o}{/}\PYG{n+nb}{max}\PYG{p}{(}\PYG{n}{total}\PYG{p}{,} \PYG{l+m+mi}{1}\PYG{p}{)}\PYG{p}{]}\PYG{p}{)}
    \PYG{k}{if} \PYG{n}{n}\PYG{p}{:}
        \PYG{c+c1}{\PYGZsh{} \CJKCode{两组不等式让辅助量包住}CDF\CJKCode{误差的正负两侧。}}
        \PYG{n}{residual} \PYG{o}{=} \PYG{n}{observed\PYGZus{}cdf} \PYG{o}{\PYGZhy{}} \PYG{n}{cdf} \PYG{o}{@} \PYG{n}{lower}
        \PYG{n}{cdf\PYGZus{}sparse} \PYG{o}{=} \PYG{n}{csr\PYGZus{}matrix}\PYG{p}{(}\PYG{n}{cdf}\PYG{p}{)}
        \PYG{n}{absolute} \PYG{o}{=} \PYG{o}{\PYGZhy{}}\PYG{n}{eye}\PYG{p}{(}\PYG{n}{n}\PYG{p}{,} \PYG{n+nb}{format}\PYG{o}{=}\PYG{l+s+s2}{\PYGZdq{}}\PYG{l+s+s2}{csr}\PYG{l+s+s2}{\PYGZdq{}}\PYG{p}{)}
        \PYG{n}{inequalities} \PYG{o}{=} \PYG{n}{vstack}\PYG{p}{(}\PYG{p}{(}\PYG{n}{hstack}\PYG{p}{(}\PYG{p}{(}\PYG{n}{cdf\PYGZus{}sparse}\PYG{p}{,} \PYG{n}{absolute}\PYG{p}{)}\PYG{p}{)}\PYG{p}{,} \PYG{n}{hstack}\PYG{p}{(}\PYG{p}{(}\PYG{o}{\PYGZhy{}}\PYG{n}{cdf\PYGZus{}sparse}\PYG{p}{,} \PYG{n}{absolute}\PYG{p}{)}\PYG{p}{)}\PYG{p}{)}\PYG{p}{,} \PYG{n+nb}{format}\PYG{o}{=}\PYG{l+s+s2}{\PYGZdq{}}\PYG{l+s+s2}{csr}\PYG{l+s+s2}{\PYGZdq{}}\PYG{p}{)}
        \PYG{n}{upper} \PYG{o}{=} \PYG{n}{np}\PYG{o}{.}\PYG{n}{concatenate}\PYG{p}{(}\PYG{p}{(}\PYG{n}{residual}\PYG{p}{,} \PYG{o}{\PYGZhy{}}\PYG{n}{residual}\PYG{p}{)}\PYG{p}{)}
    \PYG{k}{else}\PYG{p}{:}
        \PYG{n}{inequalities}\PYG{p}{,} \PYG{n}{upper} \PYG{o}{=} \PYG{k+kc}{None}\PYG{p}{,} \PYG{k+kc}{None}
    \PYG{c+c1}{\PYGZsh{} \CJKCode{区间宽度除以实测均值是绝对误差辅助量的目标系数。}}
    \PYG{n}{objective} \PYG{o}{=} \PYG{n}{np}\PYG{o}{.}\PYG{n}{concatenate}\PYG{p}{(}\PYG{p}{(}\PYG{n}{np}\PYG{o}{.}\PYG{n}{zeros}\PYG{p}{(}\PYG{n}{m}\PYG{p}{)}\PYG{p}{,} \PYG{n}{widths} \PYG{o}{/} \PYG{n}{mean}\PYG{p}{)}\PYG{p}{)}
    \PYG{n}{result} \PYG{o}{=} \PYG{n}{linprog}\PYG{p}{(}\PYG{n}{objective}\PYG{p}{,} \PYG{n}{A\PYGZus{}ub}\PYG{o}{=}\PYG{n}{inequalities}\PYG{p}{,} \PYG{n}{b\PYGZus{}ub}\PYG{o}{=}\PYG{n}{upper}\PYG{p}{,}
                     \PYG{n}{A\PYGZus{}eq}\PYG{o}{=}\PYG{n}{equality}\PYG{p}{,} \PYG{n}{b\PYGZus{}eq}\PYG{o}{=}\PYG{n}{rhs}\PYG{p}{,}
                     \PYG{n}{bounds}\PYG{o}{=}\PYG{p}{[}\PYG{p}{(}\PYG{l+m+mi}{0}\PYG{p}{,} \PYG{k+kc}{None}\PYG{p}{)}\PYG{p}{]} \PYG{o}{*} \PYG{p}{(}\PYG{n}{m}\PYG{o}{+}\PYG{n}{n}\PYG{p}{)}\PYG{p}{,} \PYG{n}{method}\PYG{o}{=}\PYG{l+s+s2}{\PYGZdq{}}\PYG{l+s+s2}{highs}\PYG{l+s+s2}{\PYGZdq{}}\PYG{p}{)}
    \PYG{n}{report}\PYG{o}{.}\PYG{n}{update}\PYG{p}{(}\PYG{n}{solver\PYGZus{}status}\PYG{o}{=}\PYG{n+nb}{int}\PYG{p}{(}\PYG{n}{result}\PYG{o}{.}\PYG{n}{status}\PYG{p}{)}\PYG{p}{,} \PYG{n}{solver\PYGZus{}message}\PYG{o}{=}\PYG{n+nb}{str}\PYG{p}{(}\PYG{n}{result}\PYG{o}{.}\PYG{n}{message}\PYG{p}{)}\PYG{p}{,}
                  \PYG{n}{solver\PYGZus{}iterations}\PYG{o}{=}\PYG{n+nb}{int}\PYG{p}{(}\PYG{n}{result}\PYG{o}{.}\PYG{n}{nit}\PYG{p}{)}\PYG{p}{,} \PYG{n}{support\PYGZus{}intervals}\PYG{o}{=}\PYG{n}{n}\PYG{p}{)}
    \PYG{c+c1}{\PYGZsh{} \CJKCode{只有明确的}LP\CJKCode{不可行状态才登记为不可行。}}
    \PYG{k}{if} \PYG{n}{result}\PYG{o}{.}\PYG{n}{status} \PYG{o}{==} \PYG{l+m+mi}{2}\PYG{p}{:}
        \PYG{k}{return} \PYG{p}{\PYGZob{}}\PYG{o}{*}\PYG{o}{*}\PYG{n}{report}\PYG{p}{,} \PYG{l+s+s2}{\PYGZdq{}}\PYG{l+s+s2}{status}\PYG{l+s+s2}{\PYGZdq{}}\PYG{p}{:} \PYG{l+s+s2}{\PYGZdq{}}\PYG{l+s+s2}{infeasible}\PYG{l+s+s2}{\PYGZdq{}}\PYG{p}{,} \PYG{l+s+s2}{\PYGZdq{}}\PYG{l+s+s2}{reason}\PYG{l+s+s2}{\PYGZdq{}}\PYG{p}{:} \PYG{l+s+s2}{\PYGZdq{}}\PYG{l+s+s2}{hard\PYGZus{}energy\PYGZus{}or\PYGZus{}day\PYGZus{}bounds\PYGZus{}infeasible}\PYG{l+s+s2}{\PYGZdq{}}\PYG{p}{,} \PYG{l+s+s2}{\PYGZdq{}}\PYG{l+s+s2}{weights}\PYG{l+s+s2}{\PYGZdq{}}\PYG{p}{:} \PYG{k+kc}{None}\PYG{p}{\PYGZcb{}}
    \PYG{k}{if} \PYG{o+ow}{not} \PYG{n}{result}\PYG{o}{.}\PYG{n}{success} \PYG{o+ow}{or} \PYG{n}{result}\PYG{o}{.}\PYG{n}{x} \PYG{o+ow}{is} \PYG{k+kc}{None}\PYG{p}{:}
        \PYG{k}{return} \PYG{p}{\PYGZob{}}\PYG{o}{*}\PYG{o}{*}\PYG{n}{report}\PYG{p}{,} \PYG{l+s+s2}{\PYGZdq{}}\PYG{l+s+s2}{status}\PYG{l+s+s2}{\PYGZdq{}}\PYG{p}{:} \PYG{l+s+s2}{\PYGZdq{}}\PYG{l+s+s2}{solver\PYGZus{}failure}\PYG{l+s+s2}{\PYGZdq{}}\PYG{p}{,} \PYG{l+s+s2}{\PYGZdq{}}\PYG{l+s+s2}{reason}\PYG{l+s+s2}{\PYGZdq{}}\PYG{p}{:} \PYG{l+s+s2}{\PYGZdq{}}\PYG{l+s+s2}{LP\PYGZus{}did\PYGZus{}not\PYGZus{}return\PYGZus{}optimum}\PYG{l+s+s2}{\PYGZdq{}}\PYG{p}{,} \PYG{l+s+s2}{\PYGZdq{}}\PYG{l+s+s2}{weights}\PYG{l+s+s2}{\PYGZdq{}}\PYG{p}{:} \PYG{k+kc}{None}\PYG{p}{\PYGZcb{}}
    \PYG{c+c1}{\PYGZsh{} \CJKCode{恢复原权重后按原量纲独立重算守恒与}CDF\CJKCode{目标。}}
    \PYG{n}{weights} \PYG{o}{=} \PYG{n}{lower} \PYG{o}{+} \PYG{n}{np}\PYG{o}{.}\PYG{n}{asarray}\PYG{p}{(}\PYG{n}{result}\PYG{o}{.}\PYG{n}{x}\PYG{p}{[}\PYG{p}{:}\PYG{n}{m}\PYG{p}{]}\PYG{p}{)}
    \PYG{n}{energy} \PYG{o}{=} \PYG{n+nb}{float}\PYG{p}{(}\PYG{n}{day\PYGZus{}energy} \PYG{o}{@} \PYG{n}{weights} \PYG{o}{\PYGZhy{}} \PYG{n}{target}\PYG{p}{)}
    \PYG{n}{conservation} \PYG{o}{=} \PYG{n+nb}{abs}\PYG{p}{(}\PYG{n+nb}{float}\PYG{p}{(}\PYG{n}{weights}\PYG{o}{.}\PYG{n}{sum}\PYG{p}{(}\PYG{p}{)} \PYG{o}{\PYGZhy{}} \PYG{n}{total}\PYG{p}{)}\PYG{p}{)}
    \PYG{n}{bound} \PYG{o}{=} \PYG{n+nb}{float}\PYG{p}{(}\PYG{n}{np}\PYG{o}{.}\PYG{n}{maximum}\PYG{p}{(}\PYG{n}{lower} \PYG{o}{\PYGZhy{}} \PYG{n}{weights}\PYG{p}{,} \PYG{l+m+mi}{0}\PYG{p}{)}\PYG{o}{.}\PYG{n}{max}\PYG{p}{(}\PYG{p}{)}\PYG{p}{)}
    \PYG{n}{normalized\PYGZus{}w1} \PYG{o}{=} \PYG{n+nb}{float}\PYG{p}{(}\PYG{n}{widths} \PYG{o}{@} \PYG{n}{np}\PYG{o}{.}\PYG{n}{abs}\PYG{p}{(}\PYG{n}{cdf} \PYG{o}{@} \PYG{n}{weights} \PYG{o}{\PYGZhy{}} \PYG{n}{observed\PYGZus{}cdf}\PYG{p}{)} \PYG{o}{/} \PYG{n}{mean}\PYG{p}{)}
    \PYG{n}{objective\PYGZus{}residual} \PYG{o}{=} \PYG{n+nb}{abs}\PYG{p}{(}\PYG{n}{normalized\PYGZus{}w1} \PYG{o}{\PYGZhy{}} \PYG{n+nb}{float}\PYG{p}{(}\PYG{n}{result}\PYG{o}{.}\PYG{n}{fun}\PYG{p}{)}\PYG{p}{)}
    \PYG{c+c1}{\PYGZsh{} \CJKCode{有限性和四类残差全部通过才交付权重。}}
    \PYG{n}{verified} \PYG{o}{=} \PYG{p}{(}\PYG{n}{np}\PYG{o}{.}\PYG{n}{isfinite}\PYG{p}{(}\PYG{n}{weights}\PYG{p}{)}\PYG{o}{.}\PYG{n}{all}\PYG{p}{(}\PYG{p}{)} \PYG{o+ow}{and} \PYG{n+nb}{abs}\PYG{p}{(}\PYG{n}{energy}\PYG{p}{)}\PYG{o}{/}\PYG{n+nb}{max}\PYG{p}{(}\PYG{n}{target}\PYG{p}{,} \PYG{l+m+mi}{1}\PYG{p}{)} \PYG{o}{\PYGZlt{}}\PYG{o}{=} \PYG{n}{tolerance}
                \PYG{o+ow}{and} \PYG{n}{conservation} \PYG{o}{\PYGZlt{}}\PYG{o}{=} \PYG{n}{tolerance}\PYG{o}{*}\PYG{n+nb}{max}\PYG{p}{(}\PYG{n}{total}\PYG{p}{,} \PYG{l+m+mi}{1}\PYG{p}{)} \PYG{o+ow}{and} \PYG{n}{bound} \PYG{o}{\PYGZlt{}}\PYG{o}{=} \PYG{n}{tolerance}
                \PYG{o+ow}{and} \PYG{n}{objective\PYGZus{}residual} \PYG{o}{\PYGZlt{}}\PYG{o}{=} \PYG{n}{tolerance}\PYG{o}{*}\PYG{n+nb}{max}\PYG{p}{(}\PYG{n}{normalized\PYGZus{}w1}\PYG{p}{,} \PYG{l+m+mi}{1}\PYG{p}{)}\PYG{p}{)}
    \PYG{k}{return} \PYG{p}{\PYGZob{}}\PYG{o}{*}\PYG{o}{*}\PYG{n}{report}\PYG{p}{,} \PYG{l+s+s2}{\PYGZdq{}}\PYG{l+s+s2}{status}\PYG{l+s+s2}{\PYGZdq{}}\PYG{p}{:} \PYG{l+s+s2}{\PYGZdq{}}\PYG{l+s+s2}{ok}\PYG{l+s+s2}{\PYGZdq{}} \PYG{k}{if} \PYG{n}{verified} \PYG{k}{else} \PYG{l+s+s2}{\PYGZdq{}}\PYG{l+s+s2}{solver\PYGZus{}failure}\PYG{l+s+s2}{\PYGZdq{}}\PYG{p}{,}
            \PYG{l+s+s2}{\PYGZdq{}}\PYG{l+s+s2}{reason}\PYG{l+s+s2}{\PYGZdq{}}\PYG{p}{:} \PYG{l+s+s2}{\PYGZdq{}}\PYG{l+s+s2}{independently\PYGZus{}verified\PYGZus{}hard\PYGZus{}energy\PYGZus{}and\PYGZus{}CDF\PYGZus{}objective}\PYG{l+s+s2}{\PYGZdq{}} \PYG{k}{if} \PYG{n}{verified} \PYG{k}{else} \PYG{l+s+s2}{\PYGZdq{}}\PYG{l+s+s2}{LP\PYGZus{}failed\PYGZus{}residual\PYGZus{}check}\PYG{l+s+s2}{\PYGZdq{}}\PYG{p}{,}
            \PYG{l+s+s2}{\PYGZdq{}}\PYG{l+s+s2}{weights}\PYG{l+s+s2}{\PYGZdq{}}\PYG{p}{:} \PYG{n}{weights}\PYG{o}{.}\PYG{n}{tolist}\PYG{p}{(}\PYG{p}{)} \PYG{k}{if} \PYG{n}{verified} \PYG{k}{else} \PYG{k+kc}{None}\PYG{p}{,}
            \PYG{l+s+s2}{\PYGZdq{}}\PYG{l+s+s2}{energy\PYGZus{}residual\PYGZus{}original\PYGZus{}units}\PYG{l+s+s2}{\PYGZdq{}}\PYG{p}{:} \PYG{n}{energy}\PYG{p}{,}
            \PYG{l+s+s2}{\PYGZdq{}}\PYG{l+s+s2}{relative\PYGZus{}energy\PYGZus{}residual}\PYG{l+s+s2}{\PYGZdq{}}\PYG{p}{:} \PYG{n}{energy}\PYG{o}{/}\PYG{n}{target}\PYG{p}{,}
            \PYG{l+s+s2}{\PYGZdq{}}\PYG{l+s+s2}{day\PYGZus{}conservation\PYGZus{}residual}\PYG{l+s+s2}{\PYGZdq{}}\PYG{p}{:} \PYG{n}{conservation}\PYG{p}{,}
            \PYG{l+s+s2}{\PYGZdq{}}\PYG{l+s+s2}{maximum\PYGZus{}lower\PYGZus{}bound\PYGZus{}violation}\PYG{l+s+s2}{\PYGZdq{}}\PYG{p}{:} \PYG{n}{bound}\PYG{p}{,}
            \PYG{l+s+s2}{\PYGZdq{}}\PYG{l+s+s2}{normalized\PYGZus{}W1}\PYG{l+s+s2}{\PYGZdq{}}\PYG{p}{:} \PYG{n}{normalized\PYGZus{}w1}\PYG{p}{,}
            \PYG{l+s+s2}{\PYGZdq{}}\PYG{l+s+s2}{W1\PYGZus{}original\PYGZus{}units}\PYG{l+s+s2}{\PYGZdq{}}\PYG{p}{:} \PYG{n}{normalized\PYGZus{}w1}\PYG{o}{*}\PYG{n}{mean}\PYG{p}{,}
            \PYG{l+s+s2}{\PYGZdq{}}\PYG{l+s+s2}{solver\PYGZus{}objective\PYGZus{}residual}\PYG{l+s+s2}{\PYGZdq{}}\PYG{p}{:} \PYG{n}{objective\PYGZus{}residual}\PYG{p}{,}
            \PYG{l+s+s2}{\PYGZdq{}}\PYG{l+s+s2}{caveat}\PYG{l+s+s2}{\PYGZdq{}}\PYG{p}{:} \PYG{l+s+s2}{\PYGZdq{}}\PYG{l+s+s2}{Minimum W1 on fixed selected profiles subject to these statistical constraints only; chronology and engineering operation unvalidated.}\PYG{l+s+s2}{\PYGZdq{}}\PYG{p}{\PYGZcb{}}
\end{Verbatim}

\endgroup

\noindent\textbf{结果核对。}先查\path{status=ok}及权重非空，再核对能量、天数、下界与目标残差均通过，并保存求解器状态和全部残差供后续复核。比较\path{normalized_W1}衡量分布拟合；权重成功仅是固定代表日的统计匹配，时序和工程运行证据须另行查验。

\Needspace{12\baselineskip}
\subsection*{17．热源启停三行紧凑凸包与容量关联爬坡}
连续展示可选机组组合的变量、开停出力边界、三行状态转换凸包和双向爬坡约束；开机状态为二元，启停事件保持连续而被相邻状态唯一确定。

\noindent\textbf{输入：}显式启用的机组组合配置、最大候选容量、既有容量和新增容量变量、最低出力比例、逐时热出力变量及已登记前驱。\par
\noindent\textbf{输出：}每时段一个二元开机状态、两个连续启停变量，以及出力、转换和爬坡线性约束；该段依赖build外层登记表与既有容量约束。\par

\noindent\textbf{实现思路。}以$u_t$为二元开机状态，$s_t,d_t$为连续启停事件，三行约束为$s_t-d_t=u_t-u_{t-1}$、$s_t+d_t\le u_t+u_{t-1}$及$s_t+d_t\le2-u_t-u_{t-1}$。相邻状态确定后，事件唯一取零或一；两条和约束与转换等式共同蕴含四条逐项上界，保持同一紧松弛而每时段少两行，并省去两个事件整数变量，从而减少模型规模。热出力边界与容量关联爬坡接续装配，启动或停机时通过最大候选容量留出转换余量。

{\small 来源：\path{model.py}，原第938--1018行；81行原代码、15行补充注释。}\par

\begingroup
\setstretch{1.0}
\begin{Verbatim}[commandchars=\\\{\},numbers=left,firstnumber=1,stepnumber=1,fontsize=\fontsize{9.5}{11.5}\selectfont,breaklines=true,breakanywhere=true,frame=single,framesep=3mm,numbersep=5pt,rulecolor=\color{black},xleftmargin=18pt]
\PYG{c+c1}{\PYGZsh{} \CJKCode{仅显式启用且候选容量为正时装配此块。}}
\PYG{k}{if} \PYG{n}{source}\PYG{o}{.}\PYG{n}{get}\PYG{p}{(}\PYG{l+s+s2}{\PYGZdq{}}\PYG{l+s+s2}{unit\PYGZus{}commitment}\PYG{l+s+s2}{\PYGZdq{}}\PYG{p}{,} \PYG{k+kc}{False}\PYG{p}{)} \PYG{o+ow}{and} \PYG{n}{maximum\PYGZus{}capacity} \PYG{o}{\PYGZgt{}} \PYG{l+m+mi}{0}\PYG{p}{:}
    \PYG{c+c1}{\PYGZsh{} \CJKCode{读入最低出力比例与上下爬坡比例。}}
    \PYG{n}{min\PYGZus{}fraction} \PYG{o}{=} \PYG{n+nb}{float}\PYG{p}{(}\PYG{n}{source}\PYG{o}{.}\PYG{n}{get}\PYG{p}{(}\PYG{l+s+s2}{\PYGZdq{}}\PYG{l+s+s2}{minimum\PYGZus{}output\PYGZus{}fraction}\PYG{l+s+s2}{\PYGZdq{}}\PYG{p}{,} \PYG{l+m+mf}{0.0}\PYG{p}{)}\PYG{p}{)}
    \PYG{n}{ramp\PYGZus{}up} \PYG{o}{=} \PYG{n+nb}{float}\PYG{p}{(}\PYG{n}{source}\PYG{o}{.}\PYG{n}{get}\PYG{p}{(}\PYG{l+s+s2}{\PYGZdq{}}\PYG{l+s+s2}{ramp\PYGZus{}up\PYGZus{}fraction\PYGZus{}per\PYGZus{}hour}\PYG{l+s+s2}{\PYGZdq{}}\PYG{p}{,} \PYG{l+m+mf}{1.0}\PYG{p}{)}\PYG{p}{)}
    \PYG{n}{ramp\PYGZus{}down} \PYG{o}{=} \PYG{n+nb}{float}\PYG{p}{(}\PYG{n}{source}\PYG{o}{.}\PYG{n}{get}\PYG{p}{(}\PYG{l+s+s2}{\PYGZdq{}}\PYG{l+s+s2}{ramp\PYGZus{}down\PYGZus{}fraction\PYGZus{}per\PYGZus{}hour}\PYG{l+s+s2}{\PYGZdq{}}\PYG{p}{,} \PYG{l+m+mf}{1.0}\PYG{p}{)}\PYG{p}{)}
    \PYG{n}{cap\PYGZus{}idx} \PYG{o}{=} \PYG{n+nb+bp}{self}\PYG{o}{.}\PYG{n}{variables}\PYG{p}{[}\PYG{l+s+s2}{\PYGZdq{}}\PYG{l+s+s2}{source\PYGZus{}capacity}\PYG{l+s+s2}{\PYGZdq{}}\PYG{p}{]}\PYG{o}{.}\PYG{n}{get}\PYG{p}{(}\PYG{n}{s\PYGZus{}name}\PYG{p}{)}
    \PYG{k}{for} \PYG{n}{t} \PYG{o+ow}{in} \PYG{n+nb+bp}{self}\PYG{o}{.}\PYG{n}{times}\PYG{p}{:}
        \PYG{c+c1}{\PYGZsh{} \CJKCode{每时段只给开机状态添加整数身份。}}
        \PYG{n}{on} \PYG{o}{=} \PYG{n}{p}\PYG{o}{.}\PYG{n}{add\PYGZus{}var}\PYG{p}{(}\PYG{l+s+sa}{f}\PYG{l+s+s2}{\PYGZdq{}}\PYG{l+s+s2}{source\PYGZus{}on[}\PYG{l+s+si}{\PYGZob{}}\PYG{n}{s\PYGZus{}name}\PYG{l+s+si}{\PYGZcb{}}\PYG{l+s+s2}{,}\PYG{l+s+si}{\PYGZob{}}\PYG{n}{t}\PYG{l+s+si}{\PYGZcb{}}\PYG{l+s+s2}{]}\PYG{l+s+s2}{\PYGZdq{}}\PYG{p}{,} \PYG{n}{ub}\PYG{o}{=}\PYG{l+m+mf}{1.0}\PYG{p}{,} \PYG{n}{integer}\PYG{o}{=}\PYG{k+kc}{True}\PYG{p}{)}
        \PYG{c+c1}{\PYGZsh{} Startup and shutdown are exact transition indicators once}
        \PYG{c+c1}{\PYGZsh{} the adjacent on/off states are binary.  Keeping these two}
        \PYG{c+c1}{\PYGZsh{} auxiliaries continuous removes two thirds of the unit\PYGZhy{}}
        \PYG{c+c1}{\PYGZsh{} commitment integer variables without relaxing the model;}
        \PYG{c+c1}{\PYGZsh{} the linking constraints below uniquely force them to 0/1.}
        \PYG{c+c1}{\PYGZsh{} \CJKCode{启停事件为连续变量，相邻二元状态将其唯一锁定。}}
        \PYG{n}{startup} \PYG{o}{=} \PYG{n}{p}\PYG{o}{.}\PYG{n}{add\PYGZus{}var}\PYG{p}{(}\PYG{l+s+sa}{f}\PYG{l+s+s2}{\PYGZdq{}}\PYG{l+s+s2}{source\PYGZus{}startup[}\PYG{l+s+si}{\PYGZob{}}\PYG{n}{s\PYGZus{}name}\PYG{l+s+si}{\PYGZcb{}}\PYG{l+s+s2}{,}\PYG{l+s+si}{\PYGZob{}}\PYG{n}{t}\PYG{l+s+si}{\PYGZcb{}}\PYG{l+s+s2}{]}\PYG{l+s+s2}{\PYGZdq{}}\PYG{p}{,} \PYG{n}{ub}\PYG{o}{=}\PYG{l+m+mf}{1.0}\PYG{p}{)}
        \PYG{n}{shutdown} \PYG{o}{=} \PYG{n}{p}\PYG{o}{.}\PYG{n}{add\PYGZus{}var}\PYG{p}{(}\PYG{l+s+sa}{f}\PYG{l+s+s2}{\PYGZdq{}}\PYG{l+s+s2}{source\PYGZus{}shutdown[}\PYG{l+s+si}{\PYGZob{}}\PYG{n}{s\PYGZus{}name}\PYG{l+s+si}{\PYGZcb{}}\PYG{l+s+s2}{,}\PYG{l+s+si}{\PYGZob{}}\PYG{n}{t}\PYG{l+s+si}{\PYGZcb{}}\PYG{l+s+s2}{]}\PYG{l+s+s2}{\PYGZdq{}}\PYG{p}{,} \PYG{n}{ub}\PYG{o}{=}\PYG{l+m+mf}{1.0}\PYG{p}{)}
        \PYG{n+nb+bp}{self}\PYG{o}{.}\PYG{n}{variables}\PYG{p}{[}\PYG{l+s+s2}{\PYGZdq{}}\PYG{l+s+s2}{source\PYGZus{}on}\PYG{l+s+s2}{\PYGZdq{}}\PYG{p}{]}\PYG{p}{[}\PYG{p}{(}\PYG{n}{s\PYGZus{}name}\PYG{p}{,} \PYG{n}{t}\PYG{p}{)}\PYG{p}{]} \PYG{o}{=} \PYG{n}{on}
        \PYG{n+nb+bp}{self}\PYG{o}{.}\PYG{n}{variables}\PYG{p}{[}\PYG{l+s+s2}{\PYGZdq{}}\PYG{l+s+s2}{source\PYGZus{}startup}\PYG{l+s+s2}{\PYGZdq{}}\PYG{p}{]}\PYG{p}{[}\PYG{p}{(}\PYG{n}{s\PYGZus{}name}\PYG{p}{,} \PYG{n}{t}\PYG{p}{)}\PYG{p}{]} \PYG{o}{=} \PYG{n}{startup}
        \PYG{n+nb+bp}{self}\PYG{o}{.}\PYG{n}{variables}\PYG{p}{[}\PYG{l+s+s2}{\PYGZdq{}}\PYG{l+s+s2}{source\PYGZus{}shutdown}\PYG{l+s+s2}{\PYGZdq{}}\PYG{p}{]}\PYG{p}{[}\PYG{p}{(}\PYG{n}{s\PYGZus{}name}\PYG{p}{,} \PYG{n}{t}\PYG{p}{)}\PYG{p}{]} \PYG{o}{=} \PYG{n}{shutdown}
        \PYG{c+c1}{\PYGZsh{} \CJKCode{登记的热电联产最小电出力随开机状态生效。}}
        \PYG{k}{if} \PYG{n}{t} \PYG{o+ow}{in} \PYG{n}{minimum\PYGZus{}electricity\PYGZus{}by\PYGZus{}time}\PYG{p}{:}
            \PYG{n}{p}\PYG{o}{.}\PYG{n}{add\PYGZus{}constraint}\PYG{p}{(}
                \PYG{p}{\PYGZob{}}
                    \PYG{n+nb+bp}{self}\PYG{o}{.}\PYG{n}{variables}\PYG{p}{[}\PYG{l+s+s2}{\PYGZdq{}}\PYG{l+s+s2}{chp\PYGZus{}electricity\PYGZus{}output}\PYG{l+s+s2}{\PYGZdq{}}\PYG{p}{]}\PYG{p}{[}\PYG{p}{(}\PYG{n}{s\PYGZus{}name}\PYG{p}{,} \PYG{n}{t}\PYG{p}{)}\PYG{p}{]}\PYG{p}{:} \PYG{l+m+mf}{1.0}\PYG{p}{,}
                    \PYG{n}{on}\PYG{p}{:} \PYG{o}{\PYGZhy{}}\PYG{n}{minimum\PYGZus{}electricity\PYGZus{}by\PYGZus{}time}\PYG{p}{[}\PYG{n}{t}\PYG{p}{]}\PYG{p}{,}
                \PYG{p}{\PYGZcb{}}\PYG{p}{,}
                \PYG{n}{lb}\PYG{o}{=}\PYG{l+m+mf}{0.0}\PYG{p}{,}
            \PYG{p}{)}
        \PYG{c+c1}{\PYGZsh{} \CJKCode{启停费用按事件口径入账，避免按普通功率计费。}}
        \PYG{n+nb+bp}{self}\PYG{o}{.}\PYG{n}{\PYGZus{}record\PYGZus{}time\PYGZus{}component}\PYG{p}{(}
            \PYG{l+s+s2}{\PYGZdq{}}\PYG{l+s+s2}{startup\PYGZus{}shutdown\PYGZus{}cost}\PYG{l+s+s2}{\PYGZdq{}}\PYG{p}{,} \PYG{n}{startup}\PYG{p}{,} \PYG{n+nb}{float}\PYG{p}{(}\PYG{n}{source}\PYG{o}{.}\PYG{n}{get}\PYG{p}{(}\PYG{l+s+s2}{\PYGZdq{}}\PYG{l+s+s2}{startup\PYGZus{}cost\PYGZus{}yuan}\PYG{l+s+s2}{\PYGZdq{}}\PYG{p}{,} \PYG{l+m+mf}{0.0}\PYG{p}{)}\PYG{p}{)}\PYG{p}{,} \PYG{n}{t}\PYG{p}{,} \PYG{n}{event}\PYG{o}{=}\PYG{k+kc}{True}
        \PYG{p}{)}
        \PYG{n+nb+bp}{self}\PYG{o}{.}\PYG{n}{\PYGZus{}record\PYGZus{}time\PYGZus{}component}\PYG{p}{(}
            \PYG{l+s+s2}{\PYGZdq{}}\PYG{l+s+s2}{startup\PYGZus{}shutdown\PYGZus{}cost}\PYG{l+s+s2}{\PYGZdq{}}\PYG{p}{,} \PYG{n}{shutdown}\PYG{p}{,} \PYG{n+nb}{float}\PYG{p}{(}\PYG{n}{source}\PYG{o}{.}\PYG{n}{get}\PYG{p}{(}\PYG{l+s+s2}{\PYGZdq{}}\PYG{l+s+s2}{shutdown\PYGZus{}cost\PYGZus{}yuan}\PYG{l+s+s2}{\PYGZdq{}}\PYG{p}{,} \PYG{l+m+mf}{0.0}\PYG{p}{)}\PYG{p}{)}\PYG{p}{,} \PYG{n}{t}\PYG{p}{,} \PYG{n}{event}\PYG{o}{=}\PYG{k+kc}{True}
        \PYG{p}{)}

        \PYG{c+c1}{\PYGZsh{} \CJKCode{各温位热出力合计受开机状态与最大候选容量限制。}}
        \PYG{n}{upper\PYGZus{}on} \PYG{o}{=} \PYG{p}{\PYGZob{}}\PYG{n}{idx}\PYG{p}{:} \PYG{l+m+mf}{1.0} \PYG{k}{for} \PYG{n}{idx} \PYG{o+ow}{in} \PYG{n}{heat\PYGZus{}by\PYGZus{}time}\PYG{p}{[}\PYG{n}{t}\PYG{p}{]}\PYG{p}{\PYGZcb{}}
        \PYG{n+nb+bp}{self}\PYG{o}{.}\PYG{n}{\PYGZus{}add}\PYG{p}{(}\PYG{n}{upper\PYGZus{}on}\PYG{p}{,} \PYG{n}{on}\PYG{p}{,} \PYG{o}{\PYGZhy{}}\PYG{n}{maximum\PYGZus{}capacity}\PYG{p}{)}
        \PYG{n}{p}\PYG{o}{.}\PYG{n}{add\PYGZus{}constraint}\PYG{p}{(}\PYG{n}{upper\PYGZus{}on}\PYG{p}{,} \PYG{n}{ub}\PYG{o}{=}\PYG{l+m+mf}{0.0}\PYG{p}{)}

        \PYG{c+c1}{\PYGZsh{} \CJKCode{启用最低出力时，把新增容量系数接入线性下界。}}
        \PYG{k}{if} \PYG{n}{min\PYGZus{}fraction} \PYG{o}{\PYGZgt{}} \PYG{l+m+mi}{0}\PYG{p}{:}
            \PYG{n}{lower\PYGZus{}on} \PYG{o}{=} \PYG{p}{\PYGZob{}}\PYG{n}{idx}\PYG{p}{:} \PYG{l+m+mf}{1.0} \PYG{k}{for} \PYG{n}{idx} \PYG{o+ow}{in} \PYG{n}{heat\PYGZus{}by\PYGZus{}time}\PYG{p}{[}\PYG{n}{t}\PYG{p}{]}\PYG{p}{\PYGZcb{}}
            \PYG{k}{if} \PYG{n}{cap\PYGZus{}idx} \PYG{o+ow}{is} \PYG{o+ow}{not} \PYG{k+kc}{None}\PYG{p}{:}
                \PYG{n+nb+bp}{self}\PYG{o}{.}\PYG{n}{\PYGZus{}add}\PYG{p}{(}\PYG{n}{lower\PYGZus{}on}\PYG{p}{,} \PYG{n}{cap\PYGZus{}idx}\PYG{p}{,} \PYG{o}{\PYGZhy{}}\PYG{n}{min\PYGZus{}fraction}\PYG{p}{)}
            \PYG{n+nb+bp}{self}\PYG{o}{.}\PYG{n}{\PYGZus{}add}\PYG{p}{(}\PYG{n}{lower\PYGZus{}on}\PYG{p}{,} \PYG{n}{on}\PYG{p}{,} \PYG{o}{\PYGZhy{}}\PYG{n}{maximum\PYGZus{}capacity}\PYG{p}{)}
            \PYG{n}{p}\PYG{o}{.}\PYG{n}{add\PYGZus{}constraint}\PYG{p}{(}\PYG{n}{lower\PYGZus{}on}\PYG{p}{,} \PYG{n}{lb}\PYG{o}{=}\PYG{n}{min\PYGZus{}fraction} \PYG{o}{*} \PYG{n}{existing\PYGZus{}capacity} \PYG{o}{\PYGZhy{}} \PYG{n}{maximum\PYGZus{}capacity}\PYG{p}{)}

    \PYG{k}{for} \PYG{n}{t} \PYG{o+ow}{in} \PYG{n+nb+bp}{self}\PYG{o}{.}\PYG{n}{times}\PYG{p}{:}
        \PYG{c+c1}{\PYGZsh{} \CJKCode{前驱来自已登记时序边界，本段据此连接相邻状态。}}
        \PYG{n}{prev} \PYG{o}{=} \PYG{n+nb+bp}{self}\PYG{o}{.}\PYG{n}{time\PYGZus{}predecessor}\PYG{p}{[}\PYG{n}{t}\PYG{p}{]}
        \PYG{n}{on} \PYG{o}{=} \PYG{n+nb+bp}{self}\PYG{o}{.}\PYG{n}{variables}\PYG{p}{[}\PYG{l+s+s2}{\PYGZdq{}}\PYG{l+s+s2}{source\PYGZus{}on}\PYG{l+s+s2}{\PYGZdq{}}\PYG{p}{]}\PYG{p}{[}\PYG{p}{(}\PYG{n}{s\PYGZus{}name}\PYG{p}{,} \PYG{n}{t}\PYG{p}{)}\PYG{p}{]}
        \PYG{n}{prev\PYGZus{}on} \PYG{o}{=} \PYG{n+nb+bp}{self}\PYG{o}{.}\PYG{n}{variables}\PYG{p}{[}\PYG{l+s+s2}{\PYGZdq{}}\PYG{l+s+s2}{source\PYGZus{}on}\PYG{l+s+s2}{\PYGZdq{}}\PYG{p}{]}\PYG{p}{[}\PYG{p}{(}\PYG{n}{s\PYGZus{}name}\PYG{p}{,} \PYG{n}{prev}\PYG{p}{)}\PYG{p}{]}
        \PYG{n}{startup} \PYG{o}{=} \PYG{n+nb+bp}{self}\PYG{o}{.}\PYG{n}{variables}\PYG{p}{[}\PYG{l+s+s2}{\PYGZdq{}}\PYG{l+s+s2}{source\PYGZus{}startup}\PYG{l+s+s2}{\PYGZdq{}}\PYG{p}{]}\PYG{p}{[}\PYG{p}{(}\PYG{n}{s\PYGZus{}name}\PYG{p}{,} \PYG{n}{t}\PYG{p}{)}\PYG{p}{]}
        \PYG{n}{shutdown} \PYG{o}{=} \PYG{n+nb+bp}{self}\PYG{o}{.}\PYG{n}{variables}\PYG{p}{[}\PYG{l+s+s2}{\PYGZdq{}}\PYG{l+s+s2}{source\PYGZus{}shutdown}\PYG{l+s+s2}{\PYGZdq{}}\PYG{p}{]}\PYG{p}{[}\PYG{p}{(}\PYG{n}{s\PYGZus{}name}\PYG{p}{,} \PYG{n}{t}\PYG{p}{)}\PYG{p}{]}
        \PYG{c+c1}{\PYGZsh{} Exact three\PYGZhy{}row state\PYGZhy{}transition hull.  The two sum}
        \PYG{c+c1}{\PYGZsh{} inequalities imply all four individual upper links}
        \PYG{c+c1}{\PYGZsh{} (startup \PYGZlt{}= on, startup \PYGZlt{}= 1\PYGZhy{}prev, shutdown \PYGZlt{}= prev,}
        \PYG{c+c1}{\PYGZsh{} shutdown \PYGZlt{}= 1\PYGZhy{}on) when combined with the equality.  This}
        \PYG{c+c1}{\PYGZsh{} is the same tight relaxation with two fewer rows per hour.}
        \PYG{c+c1}{\PYGZsh{} A one\PYGZhy{}row cyclic group has prev\PYGZus{}on == on. Accumulate}
        \PYG{c+c1}{\PYGZsh{} coefficients so its transition cancels exactly.}
        \PYG{c+c1}{\PYGZsh{} \CJKCode{转换等式为启动减停机等于当前开机减前驱开机。}}
        \PYG{n}{transition} \PYG{o}{=} \PYG{p}{\PYGZob{}}\PYG{n}{startup}\PYG{p}{:} \PYG{l+m+mf}{1.0}\PYG{p}{,} \PYG{n}{shutdown}\PYG{p}{:} \PYG{o}{\PYGZhy{}}\PYG{l+m+mf}{1.0}\PYG{p}{,} \PYG{n}{on}\PYG{p}{:} \PYG{o}{\PYGZhy{}}\PYG{l+m+mf}{1.0}\PYG{p}{\PYGZcb{}}
        \PYG{c+c1}{\PYGZsh{} \CJKCode{累加同一索引系数，单时段循环的状态项正确抵消。}}
        \PYG{n+nb+bp}{self}\PYG{o}{.}\PYG{n}{\PYGZus{}add}\PYG{p}{(}\PYG{n}{transition}\PYG{p}{,} \PYG{n}{prev\PYGZus{}on}\PYG{p}{,} \PYG{l+m+mf}{1.0}\PYG{p}{)}
        \PYG{n}{p}\PYG{o}{.}\PYG{n}{add\PYGZus{}constraint}\PYG{p}{(}\PYG{n}{transition}\PYG{p}{,} \PYG{n}{lb}\PYG{o}{=}\PYG{l+m+mf}{0.0}\PYG{p}{,} \PYG{n}{ub}\PYG{o}{=}\PYG{l+m+mf}{0.0}\PYG{p}{)}
        \PYG{c+c1}{\PYGZsh{} \CJKCode{第一条和约束限制事件之和不超过两期开机之和。}}
        \PYG{n}{active\PYGZus{}sum} \PYG{o}{=} \PYG{p}{\PYGZob{}}\PYG{n}{startup}\PYG{p}{:} \PYG{l+m+mf}{1.0}\PYG{p}{,} \PYG{n}{shutdown}\PYG{p}{:} \PYG{l+m+mf}{1.0}\PYG{p}{,} \PYG{n}{on}\PYG{p}{:} \PYG{o}{\PYGZhy{}}\PYG{l+m+mf}{1.0}\PYG{p}{\PYGZcb{}}
        \PYG{n+nb+bp}{self}\PYG{o}{.}\PYG{n}{\PYGZus{}add}\PYG{p}{(}\PYG{n}{active\PYGZus{}sum}\PYG{p}{,} \PYG{n}{prev\PYGZus{}on}\PYG{p}{,} \PYG{o}{\PYGZhy{}}\PYG{l+m+mf}{1.0}\PYG{p}{)}
        \PYG{n}{p}\PYG{o}{.}\PYG{n}{add\PYGZus{}constraint}\PYG{p}{(}\PYG{n}{active\PYGZus{}sum}\PYG{p}{,} \PYG{n}{ub}\PYG{o}{=}\PYG{l+m+mf}{0.0}\PYG{p}{)}
        \PYG{c+c1}{\PYGZsh{} \CJKCode{第二条和约束补上关机侧界，合计三行保留紧松弛。}}
        \PYG{n}{inactive\PYGZus{}sum} \PYG{o}{=} \PYG{p}{\PYGZob{}}\PYG{n}{startup}\PYG{p}{:} \PYG{l+m+mf}{1.0}\PYG{p}{,} \PYG{n}{shutdown}\PYG{p}{:} \PYG{l+m+mf}{1.0}\PYG{p}{,} \PYG{n}{on}\PYG{p}{:} \PYG{l+m+mf}{1.0}\PYG{p}{\PYGZcb{}}
        \PYG{n+nb+bp}{self}\PYG{o}{.}\PYG{n}{\PYGZus{}add}\PYG{p}{(}\PYG{n}{inactive\PYGZus{}sum}\PYG{p}{,} \PYG{n}{prev\PYGZus{}on}\PYG{p}{,} \PYG{l+m+mf}{1.0}\PYG{p}{)}
        \PYG{n}{p}\PYG{o}{.}\PYG{n}{add\PYGZus{}constraint}\PYG{p}{(}\PYG{n}{inactive\PYGZus{}sum}\PYG{p}{,} \PYG{n}{ub}\PYG{o}{=}\PYG{l+m+mf}{2.0}\PYG{p}{)}

        \PYG{c+c1}{\PYGZsh{} \CJKCode{上爬坡关联装机容量；启动事件提供候选最大容量余量。}}
        \PYG{n}{ramp\PYGZus{}up\PYGZus{}row} \PYG{o}{=} \PYG{p}{\PYGZob{}}\PYG{n}{idx}\PYG{p}{:} \PYG{l+m+mf}{1.0} \PYG{k}{for} \PYG{n}{idx} \PYG{o+ow}{in} \PYG{n}{heat\PYGZus{}by\PYGZus{}time}\PYG{p}{[}\PYG{n}{t}\PYG{p}{]}\PYG{p}{\PYGZcb{}}
        \PYG{k}{for} \PYG{n}{idx} \PYG{o+ow}{in} \PYG{n}{heat\PYGZus{}by\PYGZus{}time}\PYG{p}{[}\PYG{n}{prev}\PYG{p}{]}\PYG{p}{:}
            \PYG{n+nb+bp}{self}\PYG{o}{.}\PYG{n}{\PYGZus{}add}\PYG{p}{(}\PYG{n}{ramp\PYGZus{}up\PYGZus{}row}\PYG{p}{,} \PYG{n}{idx}\PYG{p}{,} \PYG{o}{\PYGZhy{}}\PYG{l+m+mf}{1.0}\PYG{p}{)}
        \PYG{k}{if} \PYG{n}{cap\PYGZus{}idx} \PYG{o+ow}{is} \PYG{o+ow}{not} \PYG{k+kc}{None}\PYG{p}{:}
            \PYG{n+nb+bp}{self}\PYG{o}{.}\PYG{n}{\PYGZus{}add}\PYG{p}{(}\PYG{n}{ramp\PYGZus{}up\PYGZus{}row}\PYG{p}{,} \PYG{n}{cap\PYGZus{}idx}\PYG{p}{,} \PYG{o}{\PYGZhy{}}\PYG{n}{ramp\PYGZus{}up}\PYG{p}{)}
        \PYG{n+nb+bp}{self}\PYG{o}{.}\PYG{n}{\PYGZus{}add}\PYG{p}{(}\PYG{n}{ramp\PYGZus{}up\PYGZus{}row}\PYG{p}{,} \PYG{n}{startup}\PYG{p}{,} \PYG{o}{\PYGZhy{}}\PYG{n}{maximum\PYGZus{}capacity}\PYG{p}{)}
        \PYG{n}{p}\PYG{o}{.}\PYG{n}{add\PYGZus{}constraint}\PYG{p}{(}\PYG{n}{ramp\PYGZus{}up\PYGZus{}row}\PYG{p}{,} \PYG{n}{ub}\PYG{o}{=}\PYG{n}{ramp\PYGZus{}up} \PYG{o}{*} \PYG{n}{existing\PYGZus{}capacity}\PYG{p}{)}

        \PYG{c+c1}{\PYGZsh{} \CJKCode{下爬坡对称处理，停机事件提供相应变化余量。}}
        \PYG{n}{ramp\PYGZus{}down\PYGZus{}row} \PYG{o}{=} \PYG{p}{\PYGZob{}}\PYG{n}{idx}\PYG{p}{:} \PYG{l+m+mf}{1.0} \PYG{k}{for} \PYG{n}{idx} \PYG{o+ow}{in} \PYG{n}{heat\PYGZus{}by\PYGZus{}time}\PYG{p}{[}\PYG{n}{prev}\PYG{p}{]}\PYG{p}{\PYGZcb{}}
        \PYG{k}{for} \PYG{n}{idx} \PYG{o+ow}{in} \PYG{n}{heat\PYGZus{}by\PYGZus{}time}\PYG{p}{[}\PYG{n}{t}\PYG{p}{]}\PYG{p}{:}
            \PYG{n+nb+bp}{self}\PYG{o}{.}\PYG{n}{\PYGZus{}add}\PYG{p}{(}\PYG{n}{ramp\PYGZus{}down\PYGZus{}row}\PYG{p}{,} \PYG{n}{idx}\PYG{p}{,} \PYG{o}{\PYGZhy{}}\PYG{l+m+mf}{1.0}\PYG{p}{)}
        \PYG{k}{if} \PYG{n}{cap\PYGZus{}idx} \PYG{o+ow}{is} \PYG{o+ow}{not} \PYG{k+kc}{None}\PYG{p}{:}
            \PYG{n+nb+bp}{self}\PYG{o}{.}\PYG{n}{\PYGZus{}add}\PYG{p}{(}\PYG{n}{ramp\PYGZus{}down\PYGZus{}row}\PYG{p}{,} \PYG{n}{cap\PYGZus{}idx}\PYG{p}{,} \PYG{o}{\PYGZhy{}}\PYG{n}{ramp\PYGZus{}down}\PYG{p}{)}
        \PYG{n+nb+bp}{self}\PYG{o}{.}\PYG{n}{\PYGZus{}add}\PYG{p}{(}\PYG{n}{ramp\PYGZus{}down\PYGZus{}row}\PYG{p}{,} \PYG{n}{shutdown}\PYG{p}{,} \PYG{o}{\PYGZhy{}}\PYG{n}{maximum\PYGZus{}capacity}\PYG{p}{)}
        \PYG{n}{p}\PYG{o}{.}\PYG{n}{add\PYGZus{}constraint}\PYG{p}{(}\PYG{n}{ramp\PYGZus{}down\PYGZus{}row}\PYG{p}{,} \PYG{n}{ub}\PYG{o}{=}\PYG{n}{ramp\PYGZus{}down} \PYG{o}{*} \PYG{n}{existing\PYGZus{}capacity}\PYG{p}{)}
\end{Verbatim}

\endgroup

\noindent\textbf{结果核对。}解后检查启停转移残差与事件小数偏差，均须符合数值容差，再结合全模型约束审计确认出力与爬坡。统计二元状态数、连续事件数及事件总量，可核查变量压缩；不能只凭费用下降判断机组组合成立。
